\documentclass[11pt]{article}
\usepackage[utf8]{inputenc}
\usepackage[T1]{fontenc}
\usepackage{amsmath,amssymb}
\usepackage{geometry}
\usepackage{float}
\usepackage{graphicx}
\usepackage[numbers]{natbib}
\usepackage[hidelinks]{hyperref}
\usepackage{parskip}
\usepackage{array}
\usepackage{longtable}
\usepackage[strings]{underscore}
\geometry{a4paper, margin=1in}
\setlength{\parindent}{1em}
\newcolumntype{L}[1]{>{\raggedright\arraybackslash}p{#1}}
\setcounter{secnumdepth}{0}
\makeatletter
\renewcommand{\@seccntformat}[1]{}
\makeatother

\author{Jacob Chinitz}
\title{Entropic Scalar EFT: From Entanglement Microstructure to Gravity and Cosmic Structure}
\date{\today}

\begin{document}
\maketitle

\begin{abstract}
\indent We propose that empty space is not a passive backdrop but a physical medium with a finite budget of quantum entanglement: the linking structure that allows parts of a quantum system to share state. Matter forms when some of that capacity becomes locked into stable, localized defects of the medium. A particle's mass measures how much entanglement is committed to such a defect. Gravity is the surrounding capacity-strain field: near matter, slightly less entanglement capacity is freely available, and in the weak-field limit the fractional shortfall gives the gravitational potential. The excess acceleration seen in galaxies, usually attributed to particle dark matter, is treated here as the large-scale continuation of the same capacity response rather than as a new unseen substance.

\indent The central result is that this picture is not free to be adjusted after the fact. Once one accepts the finite-capacity medium, the three founding postulates, and a specific minimal model for the smallest cell of space, the weak-field coefficients are fixed by counting the possible configurations of that cell. A single chain of calculation, with nothing left to tune galaxy by galaxy, then gives Newton's law, the acceleration scale $a_0$ at which galaxies begin to depart from Newtonian expectations, the tight observed relation between galaxy rotation and ordinary matter, and the leading no-slip lensing structure.

\indent The electron, the lightest charged particle, plays a double role. It fixes the exchange rate between committed entanglement and mass, and it also fixes the absolute size of the smallest cell. The second step uses Many-Pasts, one of the three founding postulates: the claim that the medium's present state is supported not by a single microscopic past but by many admissible histories, weighted by how well each leads to the present. In the operational branch used here this leaves all ordinary quantum-mechanical predictions intact --- Born-rule statistics and no-signaling are preserved --- while making the electron's microscopic dressing memoryless, and that memorylessness fixes the substrate length.

\indent The length fixed in this way --- from the electron Compton scale and the tetrahedral entropy, with no value of $G$ entering the chain --- implies a gravitational scale within about one percent of the measured Newton constant. We treat this as a calibrated coherence check, not as a clean independent prediction, because the construction's form was selected with that target partly in view. Appendix~L records that provenance and the associated fork accounting.

\indent The microscopic ensemble is also run. Coupled unmodified to a dynamical simplicial geometry, the admissibility weighting leaves the host geometry intact under matched-volume, placebo-controlled tests, while cells pinned into persistent closure failure (the theory's matter, inserted by hand) measurably strain the surrounding capacity state and deform the local geometry in a way a matched placebo does not. Five independent calculations exclude every mechanism the theory specifies by which the substrate could select a smooth four-dimensional vacuum on its own, forcing the reading the architecture already carries: the vacuum is supplied by the Many-Pasts conditioning on the present, and the medium supplies matter and its response. The background so selected is not inert --- the substrate writes a measured, predicted-in-advance contribution to its volume term --- and the same program derives the conservation law that must carry the matter-sourced strain to Newtonian range, whose direct measurement is in progress.

\indent Beyond the static weak-field sector, the framework extends to time-dependent transport, clusters, cosmology, the saturated early universe, black holes, and the charged-lepton spectrum, at varying and explicitly labeled levels of closure. The completed claim of the paper is the chain from the microscopic construction to ordinary weak gravity; the later sectors are presented as conditional extensions, frontier completions, or open tests.
\end{abstract}

\clearpage
\tableofcontents
\newpage

\section*{Part I. Physical Idea and Canonical Definitions}

\section{1. Introduction: The Physical Claim}

Space is not an empty container that matter sits inside. It is a finite medium of entanglement capacity, and the particles we call matter are stable defects of that medium rather than independent agents acting on it from outside. Gravity is then what the surrounding medium looks like once some of its capacity is locked up in such a defect. In brief:
\begin{itemize}
\item \textbf{Matter} is a localized capacity defect of the substrate.
\item \textbf{Mass} is the entanglement that defect commits, read in mass units.
\item \textbf{Gravity} is the extended capacity strain --- the fractional capacity deficit --- the medium carries around the defect.
\item \textbf{Dark-matter phenomenology} is not a new substance but the same medium read in two further regimes: the long-range reach of the capacity-strain field on galactic scales, and its saturated phase in the early universe.
\item \textbf{General relativity} is the low-energy geometry of this same capacity medium, not an external stage to which it is added.
\end{itemize}

The continuum statement is a scalar EFT for a vacuum-relative entanglement field $S_{\mathrm{ent}}(x)$ and its deficit $\delta S$ relative to the background capacity. The defect sector is written at continuum scale in ordinary stress-energy variables, but its ontology is unchanged; inertial mass enters through the mass-per-entropy map $\kappa_m$, and the weak-field potential is the fractional deficit $\delta S/S_\infty$.

The proposal replaces part of the usual dark-sector story rather than relabeling it. The standard picture keeps visible matter and Einstein gravity and adds dark components to supply the missing gravitational response; here the vacuum already carries a finite entanglement-capacity structure, and the same medium accounts for ordinary weak-field gravity, the galactic excess usually attributed to dark matter, and the homogeneous cosmological mode. It also goes beyond an ordinary scalar extension of Einstein gravity: GR is not a stage to which the entanglement field is appended but the low-energy capacity geometry of the same substrate, with the scalar sector tracking how that geometry is depleted and redistributed by localized defects.

The paper asks whether that ontology can be made quantitative, and the answer offered here is conditional. The finite-capacity substrate, the three postulates of Section~3, and the minimal tetrahedral boundary ensemble are theory-defining inputs; the paper does not derive them from a deeper Hamiltonian. Once they are fixed, the static weak-field coefficients are no longer phenomenological parameters but are determined by a finite ultraviolet counting problem: admissibility closure fixes the effective sharing entropy, edge transport and finite-loop dressing fix the stiffness, source projection fixes the coupling, and the weak-field bridge relates fractional capacity deficit to gravitational potential.

The primary closed branch is therefore the static weak-field chain
\[
\text{microstructure} \longrightarrow \text{coefficient chain} \longrightarrow \text{continuum EFT} \longrightarrow \{G,\ a_0,\ \mathrm{RAR},\ \mathrm{no\ slip}\}.
\]
It recovers the Newtonian point-source limit, the galactic acceleration scale, the radial-acceleration relation, and leading no-slip lensing without per-system tuning. The absolute substrate length is fixed through the electron anchor and the Many-Pasts memoryless dressing branch; because the construction was developed with Newton's constant partly in view, the one-percent agreement with $G$ is treated as a calibrated coherence check rather than a clean independent prediction, with the provenance recorded in Appendix~L.

The later parts extend the framework into regimes with lower closure status --- time-dependent transport, galaxy clusters, cosmology, the saturated early phase, strong fields, and the particle and gauge extensions --- with Part~VI recording the bookkeeping explicitly.

Many-Pasts belongs with the foundations. Its history weighting enters the scale-setting branch through memoryless electron dressing, and its operational consequences for quantum probability, branch realization, and the arrow of time are developed in Section~21 and Appendix~G.

\subsection{1.1 What Is Primitive, and What Is Closed}

The word ``closure'' is used here in a specific sense. The paper does not derive the existence of a finite entanglement substrate from no assumptions, nor does it derive the tetrahedral boundary ensemble from a deeper microscopic Hamiltonian in the main text; those are theory-defining ultraviolet inputs, and the closure claim begins only once they are fixed. There are five such inputs: finite local entanglement capacity; matter as localized defects of that capacity; mass--entropy equivalence, so that mass is committed defect entropy read in mass units; the Many-Pasts history weighting of Postulate~III; and the minimal tetrahedral boundary ensemble, in its physical realization as the full admissibility-closed sharing entropy, serving as the ultraviolet counting architecture.

Given those inputs, the static weak-field sector closes with no phenomenological freedom remaining. The finite boundary count fixes the admissibility weighting and the effective sharing entropy; a faithful sector-resolution principle fixes the substrate cell length $L_*$; edge kernel, finite-loop dressing, continuum stiffness, and source projection follow in turn; and the weak-field bridge $\delta S\leftrightarrow\Phi$ then yields the Newtonian limit, the galactic acceleration scale $a_0$, the radial-acceleration relation, and leading no-slip lensing, with no per-system tuning. This is the central result. Everything beyond the static branch extends the same ontology at a lower closure level and is labeled accordingly: time-dependent transport, clusters, the saturated-phase cosmology, the strong-field branch, and the particle and gauge extensions. Part~VI collects this bookkeeping in a closure-status table.

In compressed form, the central claim is
\[
\begin{aligned}
&\boxed{\text{primitive UV capacity hypothesis}}
\rightarrow
\boxed{\text{finite counting + admissibility}}
\rightarrow
\boxed{L_*}
\\
&\rightarrow
\boxed{\gamma,\ \kappa/\gamma}
\rightarrow
\boxed{\delta S \leftrightarrow \Phi}
\rightarrow
\boxed{G,\ a_0,\ \mathrm{RAR},\ \mathrm{no\ slip}}.
\end{aligned}
\]
The microstructure is therefore not the result being proven; it is the finite ultraviolet counting problem from which the weak-field sector is derived.

Only the absolute length calibration is deferred rather than defined here. The substrate cell length $L_*$ is fixed by identifying the electron --- the lightest clean charged defect --- with the maximally extended ground state of a history-space sector-resolution operator, supported over its reduced Compton scale. That principle is itself derived rather than assumed: its content is that each dressing channel carries the full admissibility entropy and the channels are mutually independent. The first holds because the dressing is memoryless in substrate time, which is the Many-Pasts content of Postulate~III; the second holds because inter-channel correlation would contract the defect and raise its mass, so the lightest charged defect carries none. The electron's two roles --- mass anchor and length anchor --- and the explicit formula are developed in Section~13 and derived technically in Appendices~D.4 and~H. The central technical claim is that, given the three postulates and the realized minimal ensemble, the static weak-field branch has no remaining phenomenological freedom.

The continuum action can superficially resemble an ordinary scalar-tensor theory of the Brans--Dicke lineage \citep{BransDicke1961}. In a conventional scalar-tensor model one asks whether variation of a covariant action containing a scalar can produce the desired metric phenomenology. Here the scalar is the continuum order parameter of an underlying capacity substrate, and the weak-field bridge is the emergent-geometry dictionary relating fractional capacity depletion to the metric lapse. The closure question is therefore not whether the tetrahedral microstructure has been derived from nothing, but whether the specified microscopic capacity ensemble fixes the weak-field gravitational response once it is adopted.

Several external tasks remain: derive the same microscopic ensemble from a deeper substrate Hamiltonian, independently audit the graph-level return operator, confront the resulting radial-acceleration law and no-slip prediction with data, derive the horizon boundary action from the graph ensemble, and extend the constrained-capacity black-hole branch into a fully dynamical collapse theory. These are reconstruction and validation tasks outside the already-closed static weak-field branch, not hidden fit parameters inside it.

A reader is therefore being asked to accept a small number of commitments, collected here in one place.

First, the vacuum has finite entanglement capacity: it is not empty background space but a medium with a bounded local capacity for entanglement.

Second, spacetime geometry and that capacity structure are the same substrate seen at different scales, so that in the weak field gravity is the fractional deficit of locally available capacity.

Third, matter is localized committed capacity: a particle is a stable defect of the medium, and its inertial mass is the entanglement content of that defect read in mass units.

Fourth, the entanglement network has many admissible microscopic pasts, and its present state is governed by a probability-weighted ensemble of them, with nearer histories weighted more heavily. This is the Many-Pasts postulate. Its operational branch preserves ordinary Born-rule and no-signaling laboratory statistics, and its memoryless history weighting is the ingredient used in the electron dressing that fixes the substrate length.

Fifth, the ultraviolet cell is the minimal tetrahedral boundary ensemble, and this is not a continuous tuning parameter. Given fermionic face data, maximum-capacity channel selection, injective four-face assignment, and minimality, the first admissible face alphabet has seven states, which gives the 1680-state ensemble (Section~5, Appendix~B).

The finite-capacity substrate is the primitive ultraviolet premise; the identification of geometry with capacity, the mass--entropy equivalence, and the Many-Pasts weighting are the three postulates of Section~3; and the tetrahedral ensemble is the minimal ultraviolet architecture. Once they are adopted, the static weak-field sector is claimed to close: the admissibility weighting, the effective sharing entropy $g_{\mathrm{share,eff}}$, the edge kernel, the loop dressing, the continuum stiffness, the source projection, the weak-field bridge, the Newtonian limit, the galactic acceleration scale, the radial-acceleration relation, and the leading no-slip lensing structure are all derived within that construction. The single number the construction produces, $g_{\mathrm{share,eff}}=7.4198$, is not a fitted parameter but a derived entropy value inside a selected construction: it has no adjustable dial and is reproducible from the stated rules alone, while the historical selection of the construction itself is audited and priced in Appendix~L. The sectors beyond the static branch are not hidden adjustments to that closed result; they are labeled separately as conditional extensions, frontier completions, empirical tests, or open tasks.

In one line, the architecture is three postulates, one finite-capacity substrate premise, one minimal ultraviolet ensemble, and a small number of explicitly labeled conditional readings --- and from those, a single derived entropy value carries the static weak-field chain.

\subsection{1.2 Physical Motivation for the Primitives}

These commitments are primitives of the framework, not theorems proved from deeper assumptions in this paper. They are nevertheless not arbitrary. Each one is motivated by a place where established physics already strains against its own foundations, and each earns its place by making a known difficulty look less mysterious once it is adopted. None of it is offered as proof; it is the reason the starting points are reasonable ones to adopt before the derivation begins.

Mainstream gravitational physics has been converging on a finite-capacity substrate for decades. A black hole's entropy scales with the area of its horizon rather than the volume it encloses, as though the contents of a region were written on its boundary; the Bekenstein bound limits the information a bounded region can hold; and holography and entanglement-based reconstructions of geometry tie the shape of spacetime directly to patterns of entanglement. Each of these is usually treated as a deep clue without a mechanism. The framework takes the clue literally: the vacuum is a medium with a finite local budget of entanglement, and geometry is the large-scale description of that budget. Several long-standing puzzles then become the ordinary behavior of a medium that can fill up. The area law is one. A black hole is a region whose capacity is exhausted, so there is nothing left in the interior to count, and the only live bookkeeping is at the surface where filled medium meets unfilled, which is why the entropy sits on the area and not in the volume. The black-hole information question softens for the same reason: if the interior holds no independent degrees of freedom, there is no interior store of information to lose, and what falls in is recorded in the boundary layer that later evaporation reads back out. And the old problem that a continuum theory of gravity develops infinities at short distances does not arise here, because below the smallest cell there is no continuum to diverge, and a finite state space has nothing to renormalize away.

The mass--entropy identification addresses a coincidence that general relativity encodes but does not explain. General relativity builds in the equality of inertial and gravitational mass geometrically, through the equivalence principle, but gives no microphysical account of why the mass that resists acceleration and the mass that sources attraction --- equal to fifteen decimal places --- should be one and the same. Here they are the same quantity seen from two sides: a defect's committed entanglement is at once its resistance to being accelerated and its pull on the surrounding medium, so the equality follows from the construction rather than being imposed by hand. The same identification addresses why gravity is so extraordinarily weak. Rather than inserting a tiny coupling by hand, the framework produces gravity's strength by counting: the induced gravitational scale carries the factor $e^{-14g_{\mathrm{share,eff}}}$, so a large entropy makes gravity exponentially feeble, and the famous gap between gravity and the other forces becomes a matter of arithmetic rather than fine-tuning. It also gives a concrete substrate reading of gravity: the extended capacity-strain field around a localized deficit. This makes its universality automatic (everything is made of substrate, so everything gravitates), its attractive-only character natural (a deficit draws the medium inward; it does not push), and the galactic acceleration scale a derived threshold rather than an unexplained constant. That galactic scale, $a_0\approx \epsilon\,cH_0$, ties the onset of anomalous rotation to the cosmic horizon, giving a concrete form to the old Machian suspicion that local inertia should somehow depend on the universe at large. And the tight, nearly scatter-free way a galaxy's rotation tracks its visible matter, the very feature that makes an independent dark particle look so strange, follows directly when a single medium is responding to the ordinary matter, rather than being a separate component that must be arranged to follow it.

A second consequence, in plain terms before its technical form: because the medium never holds still --- every instant it is re-drawn from all the ways it could be --- a particle is not an object sitting in space but a pattern the medium keeps re-forming against its own restlessness, and the entanglement it commits is a maintained quantity, re-established beat by beat. Two different questions can be asked about such a pattern, and their answers pull apart. One is how busy the medium is keeping the knot bound: it must hold its grip on every strand the knot is tied into at once, and separate holds add up, so this maintenance is large --- the medium works hard and continuously to keep the defect in being. The other is how often the knot fully re-completes, every strand falling into alignment at the same instant; because simultaneous coincidences multiply rather than add, that completion is exponentially rare. The particle's mass --- which is its rest energy --- follows the second question, not the first: it is fixed by how rarely the binding fully re-closes, not by how much work the holding takes. One depth of binding governs both, but holding adds while coinciding multiplies, so the maintenance is large and additive while the recurrence that sets the mass is exponentially small. This is why a particle can be vastly lighter than the natural substrate scale with no small number inserted anywhere: it is not weakly bound but deeply bound, and deep binding makes full re-completion exponentially rare. The electron is light not because its structure is weak but because that structure's coherent recurrence is exponentially delayed --- light because it is heavy to hold. Section~22 and Appendix~H give this its quantitative form, with binding additive across the sectors and recurrence multiplicative; here it is only the shape of the answer.

A memoryless update rule assumes the least: it is the rule one writes down when one declines to invent hidden internal machinery that carries information forward from tick to tick. It earns its place twice over. The construction that fixes the absolute size of the cell needs each dressing pass of the electron to sample the whole admissible ensemble afresh, and the allowed local single-label dynamics provably cannot do this: they freeze into disconnected sectors that never explore the full space (Appendix~D.4), so a genuine refresh is required, not optional. A single universal refresh rate also lets the theory carry a smallest length without running afoul of the experiments that ended earlier discrete-spacetime proposals: the granularity here lives in the capacity, not in a preferred spatial lattice, so gravitational waves and light travel at the same speed and there is no frame-dependent dispersion left to detect.

Many-Pasts holds that the present configuration of the entanglement network is supported not by one definite microscopic past but by a whole ensemble of admissible pasts, each weighted by how naturally it leads to the present. Its first payment is in the quantum puzzles that interpretations of quantum mechanics were invented to handle. The interference in a double-slit experiment is the persistence of the unrecorded alternative histories in that ensemble; making a which-path measurement conditions the ensemble on the new record, and the interference goes away. The correlations of an entangled pair come from weighting the histories of the whole joint system, which reproduces the nonclassical statistics without any signal passing between the two wings of the experiment. Measurement adds no separate collapse law; it lays down a durable record, after which the relevant histories are simply the ones compatible with it. In the branch used here all of this leaves ordinary laboratory predictions exactly intact (Born-rule statistics and no-signaling both hold), and the same weighting supplies the arrow of time. The postulate is more than interpretation because it is also load-bearing in a measured number: the same history weighting that satisfies the Born-rule consistency condition also makes the electron's dressing memoryless, and that memorylessness is one of the two conditions that fix the substrate length and hence the strength of gravity. Elsewhere in physics the interpretation of quantum mechanics is detachable philosophy; here it has a concrete job.

Tetrahedra are the standard building block in several independent approaches to quantum geometry, so the cell is a familiar object rather than one invented for this paper. What the framework adds is a set of its own restrictions that pick out one ensemble: the faces carry fermionic data, the maximum-capacity channel is selected, the four faces must be assigned distinct labels, orientations are counted with both parities, and the smallest admissible alphabet is taken. Those restrictions force a seven-label face structure and a boundary ensemble of exactly 1680 states. The same machinery that fixes the count also limits how many shell excitations the cell can carry before its closure structure degenerates, which terminates the charged-lepton ladder at three families, a candidate answer to the Standard Model's otherwise unexplained fact that matter comes in exactly three generations. How this particular ensemble came to be chosen, and how much freedom that choice had, is audited separately in Appendix~L. Within the stated construction, its entropy is computed from the rules with no adjustable dial.

Several mainstream results already point toward the finite-capacity reading of gravity: the derivation of Einstein's equations as a thermodynamic relation of state, the identification of entanglement entropy with horizon area, the reconstruction of spatial connectivity from entanglement, and the conjectured identity between entanglement and geometric bridges. This framework is a concrete, finite, and falsifiable instance of that program.

The primitives are not proved before the theory begins, because no theory proves its own starting point, and a list of mysteries explained is the easiest kind of case to assemble after the fact, since a proposal of this scope can almost always produce one. The discriminator is not how much the primitives explain, but whether the explanations were forced by the construction or fitted to the target (the question Appendix~L confronts directly), and whether the same construction then survives measurements it did not anticipate. This subsection claims only the weaker statement: the starting points are motivated by real and independent pressure points in gravity, quantum information, black-hole physics, and quantum measurement, and once adopted they constrain the weak-field sector tightly enough to be tested. The rest of the paper derives those constraints and confronts them with observation.

\section{2. Canonical Field Content and Definitions}

Before the symbols, four plain words recur throughout. \emph{Capacity} is the entanglement support locally available in the medium. A \emph{defect} is a stable, localized commitment of that capacity --- what we coarse-grain into a particle. A \emph{deficit} is capacity no longer freely available to the surrounding vacuum because a defect has committed it. \emph{Strain} is the extended profile of that deficit reaching out into the medium, whose fractional size the weak-field potential tracks. The field variables below are the precise versions of these words.

We define the fundamental continuum variable as the vacuum-relative coarse-grained entanglement assigned to a UV probe cell of size $L_*$ centered at $x$:
\[
S_{\mathrm{ent}}(x)\in \mathbb{R},
\]
measured in nats and therefore dimensionless. This is not a literal microscopic entropy density at a mathematical point. It is the leading scalar order parameter associated with a vacuum-relative entanglement defect after coarse-graining over a UV cell.

This definition keeps the microscopic and continuum pictures tied together. At continuum level, $S_{\mathrm{ent}}(x)$ is the field that appears in the action and field equations. At the microscopic level it is the coarse variable recording how much local entanglement capacity remains available in the underlying medium after averaging over a UV cell.

The asymptotic vacuum-capacity baseline is denoted $S_\infty$, and the deficit field is
\[
\delta S(x)\equiv S_\infty-S_{\mathrm{ent}}(x).
\]
Positive $\delta S$ denotes reduced available vacuum entanglement capacity in the neighborhood of a localized defect or defect distribution. It is the extended capacity-strain field sourced by the defect sector, not an independent medium acted on by matter from outside. For nonlinear work it is useful to define the bounded occupancy fraction
\[
q(x)\equiv \frac{S_{\mathrm{ent}}(x)}{S_\infty}=1-\frac{\delta S}{S_\infty}\in[0,1].
\]
The variables $S_{\mathrm{ent}}$, $\delta S$, and $q$ therefore describe the same local physics in three closely related ways: available capacity, missing capacity relative to vacuum, and surviving-capacity fraction. Each is used where it is most transparent: $\delta S$ for the weak-field theory, because it talks directly to the Newtonian potential; $q$ for the nonlinear and strong-field completion, because boundedness is built in from the start; and $S_{\mathrm{ent}}$ itself for the covariant EFT, because it is the field that appears in the action.
The operational meanings are:
\begin{itemize}
\item $q=1$: vacuum capacity fully available in the absence of local defect-induced capacity strain;
\item $0<q<1$: partial local capacity reduction around a defect configuration;
\item $q=0$: complete local exhaustion of available capacity on the physical branch.
\end{itemize}

\paragraph{Fixed-epoch normalization.}
The absolute normalization of $S_{\mathrm{ent}}$ and $S_\infty$ is a convention once an epoch and cell convention have been fixed. Under a constant rescaling
\[
S_{\mathrm{ent}}\mapsto K S_{\mathrm{ent}},
\qquad
S_\infty\mapsto K S_\infty,
\qquad
\delta S\mapsto K\delta S,
\]
the observable bridge
\[
\frac{\Phi}{c^2}=-\frac{\delta S}{2S_\infty}
\]
is unchanged. The source equation is invariant in the same sense: rescaling the entropy field rescales the source coefficient with it, so the observable Newtonian normalization depends on the gauge-invariant combination $\kappa/(\gamma S_\infty)$ rather than on $S_\infty$ alone. A cell-normalized description and a horizon-normalized description can therefore assign different numerical values to $S_\infty$ without changing $\Phi$, $G$, or the PPN limit. This is not a time-dependent gauge symmetry; it is a fixed-epoch entropy-unit convention. Gravity sees fractional capacity depletion.

\paragraph{Substrate length scale.}
The canonical UV cell length is not taken to be the conventional Planck length as an input. It is fixed later, through the electron anchor and the Many-Pasts memoryless dressing of Postulate~III (Section~13, Appendix~D.4); here we only record the canonical definition that the later derivation populates. The defining relation is the ground-state faithful sector-resolution principle
\[
\ln\!\left(\frac{\lambda_e}{L_*}\right)
=
7g_{\mathrm{share,eff}}-\ln\!\left(\frac32\right),
\qquad
\lambda_e=\frac{\hbar}{m_ec}.
\]
Equivalently,
\[
L_*=\frac32\,\lambda_e e^{-7g_{\mathrm{share,eff}}}
=
1.60771947\times10^{-35}\ {\rm m}.
\]
The corresponding induced gravitational scale is
\[
G_*:=\frac{c^3L_*^2}{\hbar}
=
\frac94\,\frac{\hbar c}{m_e^2}\,e^{-14g_{\mathrm{share,eff}}}
=
6.60399128\times10^{-11}\ {\rm m^3\,kg^{-1}\,s^{-2}}.
\]
The principle says that the lightest one-bit fermionic defect coherently resolves the seven face sectors exactly once and exports only the transverse $2/3$ share of that dressing block. The conventional Planck length $L_P=\sqrt{\hbar G/c^3}$ remains useful for comparison and for standard black-hole thermodynamic notation, but it is not the primitive scale-setting input here.

The principal coefficients and derived quantities used throughout are:
\begin{align}
\gamma &:\ \text{entanglement-field stiffness},\\
\kappa &:\ \text{defect--entropy coupling},\\
\kappa_m(\ell) &:\ \text{mass-per-entropy map at scale }\ell,\\
L_* &:\ \text{substrate cell length fixed by faithful sector resolution},\\
G_* &:\ \text{gravitational scale induced by }L_*,\\
g_{\mathrm{share,max}} &= \ln(1680),\\
g_{\mathrm{share,eff}} &:\ \text{admissibility-weighted effective sharing entropy},\\
J_{\mathrm{bare}},\,J_{\mathrm{eff}}^{\mathrm{tree}},\,J_{\mathrm{eff}}^{(\mathrm{ren})} &:\ \text{UV edge-kernel couplings},\\
a_0 &= \frac{cH_0g_{\mathrm{share,eff}}}{4\pi^2}.
\end{align}
The gravitational potentials are denoted $\Phi$ and $\Psi$, and the canonical weak-field bridge will be written as
\[
\frac{\Phi}{c^2}=-\frac{\delta S}{2S_\infty}.
\]
These same symbols reappear in the UV closure chain, in the continuum action, and in the phenomenology sections. From this point onward each one keeps the same meaning, so the later derivations can build on a single notation rather than shifting between parallel conventions.

\section{3. The Three Postulates}

The framework rests on exactly three postulates because they answer three different questions, and no fewer suffice to state the ontology. \emph{What is spacetime geometry?} --- Postulate~I: geometry is the long-wavelength expression of the capacity substrate itself. \emph{What is matter, and what is mass?} --- Postulate~II: matter is localized defect capacity, and mass is the defect's committed entropy in mass units. \emph{What is quantum probability, and what fixes the direction of history?} --- Postulate~III: quantum probability is a weighting over admissible past histories of the entanglement network. The first two define the gravitational ontology; the third defines its quantum and historical structure and, through memoryless dressing, also enters the scale-setting branch. The main text treats all three as theory-defining inputs, not derived outputs.

\subsection{3.1 Information--Geometry Equivalence}

The first postulate states that spacetime geometry is the continuum expression of the entanglement-capacity substrate. This is stronger than saying entanglement contributes an additional piece of stress-energy inside otherwise standard general relativity. The metric is not an independent background to which $S_{\mathrm{ent}}$ is appended; the geometric sector and the scalar capacity sector are two continuum projections of one finite-capacity medium. In the weak field, gravitational potential is the fractional deficit of available capacity.

Two consequences of this reading should be kept distinct from the start. First, absolute $S_{\mathrm{ent}}$ is not itself ``the gravitational potential''; the observable weak-field potential comes from the fractional deficit $\delta S/S_\infty$, which is why a fixed-epoch rescaling of the entropy units leaves gravity unchanged (Section~2). Second, because geometry and the scalar are projections of the same substrate, gradients and deficits of the entanglement field are not metaphorical contributions to gravity --- they belong in the gravitational bookkeeping as a matter of what the geometry \emph{is}, not as an extra source added to it. In the EFT this is realized concretely: the scalar enters a covariant action, contributes its own stress-energy, and couples to a trace-equivalent defect source written at continuum scale in the usual stress-energy variables.

\subsection{3.2 Mass--Entropy Equivalence}

The second postulate is that mass is not something added to the substrate from outside but the inertial reading of localized capacity commitment. At scale $\ell$,
\[
m(\ell)=\kappa_m(\ell)\,\Delta S.
\]
A particle is already a localized defect of the entanglement substrate, so $m=\kappa_m\Delta S$ does not assert an analogy between two independent things; it asserts that the inertial content of the defect \emph{is} its entanglement content, read in mass units.

For elementary fermionic sectors the canonical defect increment is
\[
\Delta S_f=\ln 2.
\]
The one bit here is not arbitrary. An elementary fermionic exclusion is binary --- the face is occupied or unoccupied --- and a binary distinction carries exactly $\ln 2$ of missing entanglement. This is the simplest possible defect increment, which is why the lightest such defect, the electron, becomes the cleanest anchor for the mass--entropy map (Section~13). For composite sectors the relevant quantity is the fully dressed bound-state entanglement budget, not a bare constituent count.

Two corollaries are used later. First, because mass and entanglement budget are two descriptions of the same defect, and the masses of separated defects add, capacity committed in service of one defect cannot simultaneously serve another: shared service would make the joint budget, and with it the joint mass, sub-additive. Commitment is therefore per-defect --- each committed unit carries the label of the defect it serves. Second, the same per-defect bookkeeping makes the saturated early phase countable: if mass is committed capacity, committed units cannot be double-counted across separated defects, and the abundance of committed capacity can be counted rather than fitted (Section~18.5).

\subsection{3.3 Many-Pasts Hypothesis}

The third postulate is one of the framework's three defining inputs. It concerns what the present entanglement network is supported by. The claim is that the present is not the endpoint of one microscopic past. For any present coarse configuration, many admissible microscopic histories of the network could have produced it, and the physical state is governed by a probability-weighted ensemble over those histories rather than by a single distinguished one.

Three objects make this precise:
\begin{itemize}
\item $H$: an admissible microscopic history of the entanglement network --- one self-consistent way the present could have been reached;
\item $P$: the present coarse configuration that all admissible histories must reproduce;
\item $D(H,P)$: a history-space distance measuring how far a given history lies from the present state.
\end{itemize}
In canonical closed form the operational history weight is
\[
P(H|P)\propto e^{-D(H,P)},
\]
so histories close to the present dominate and distant ones are exponentially suppressed. This is the branch $\alpha=1,\beta=0$ of a generalized weighting family.

This postulate occupies the conceptual role usually played by an interpretation of quantum mechanics. It is not many-worlds: reality does not branch forward into many co-real macroscopic outcomes. It is not a collapse theory: there is no added physical collapse event. It is not hidden-variable ignorance: probability is not ignorance over pre-existing classical values. It is a history-space ontology --- quantum probability read as a weighting over admissible past histories of the substrate. This is the framework's alternative to those accounts: ontologically central, though its laboratory predictions are deliberately conservative.

The operational branch is then fixed by two constraints. The Born rule is recovered: requiring exact Born statistics in the projective laboratory limit selects $\alpha=1$. No-signaling is preserved: forbidding any operational bias channel that could transmit information selects $\beta=0$. The substantive content is that the family admits a single branch satisfying both at once. Born recovery is therefore stated honestly as a consistency condition the weighting must satisfy, not as a fresh derivation of the rule from outside it; this differs from programs that obtain the rule from noncontextuality \citep{Gleason1957}, decision theory \citep{Deutsch1999}, or envariance \citep{Zurek2005}, with no-signaling \citep{PopescuRohrlich1994} the companion constraint. Beyond these, the same weight orients the arrow of time through conditional typicality and makes the dressing of an elementary defect memoryless in substrate time, each pass an independent draw from the admissibility ensemble.

That memorylessness is Postulate~III's entry into the central result. The coefficient chain that fixes the Newtonian limit, the galactic law, and lensing does not call on Many-Pasts. The substrate-scale branch does: the memoryless dressing is one of the two ingredients that fix the absolute substrate length in Appendix~D.4, the other being the mutual independence of the electron's seven dressing channels, which follows from its identity as the lightest charged defect (Appendix~H). Its operational consequences and the arrow of time are developed in Section~21 and Appendix~G.

\section{4. Relativistic Continuum Structure}

\subsection{4.1 Capacity budget and continuum symmetry}

In the present framework the continuum description is expected to be covariant not because a geometric axiom is added at the outset, but because the substrate itself is finite-capacity, isotropic, and relational.

The first ingredient is a finite maximal update rate, denoted by the same constant $c$ that later appears in the transport relation $D/\tau_0=c^2$. In the present interpretation, $c$ measures the largest rate at which the substrate can propagate and reorganize information. A defect at rest spends that budget entirely on local temporal evolution. A defect in motion must spend part of the same budget on spatial transport within the surrounding network. Because the substrate is isotropic, the cost of motion depends only on the rotational scalar $v^2$ at leading order, with the temporal rate maximal at $v=0$ and vanishing when the budget is exhausted at $v=c$. These endpoint conditions alone admit many interpolating functions and so do not fix the form of the time-dilation relation. The form is fixed once the finite update speed is treated as invariant across inertial coarse descriptions: homogeneity, isotropy, and the relativity principle then select the Lorentz group rather than the Galilean one, giving the invariant interval
\[
c^2 d\tau^2 = c^2 dt^2 - d\mathbf{x}^2,
\]
and hence
\[
\frac{d\tau}{dt}=\sqrt{1-\frac{v^2}{c^2}}.
\]
The capacity-budget picture supplies the substrate interpretation of this Lorentzian kinematics: motion allocates part of the finite update budget to spatial transport, leaving the remaining fraction as proper-time evolution.

The same capacity language also unifies motion-induced and gravity-induced clock slowing. In the nonlinear branch the surviving-capacity fraction is
\[
q=\frac{S_{\mathrm{ent}}}{S_\infty},
\]
so smaller $q$ means that less local update capacity remains available. Motion reduces the temporal share of the budget by consuming part of it in spatial transport; a nearby defect reduces the local budget by depleting available capacity. The two familiar time-dilation effects are therefore interpreted as two regimes of one mechanism.

The second ingredient is the relational character of the substrate. It is not embedded in a prior physical manifold whose coordinate labels carry independent meaning. The physical content is the pattern of local capacities, defects, and neighborhood relations within the network itself. Continuum coordinates are therefore descriptive labels imposed on that relational structure, not additional physical data. Smooth changes of coordinates relabel the same underlying configuration rather than altering the physics. In continuum language this is precisely why the low-energy description should be written in generally covariant form.

The upshot is that the metric sector of the EFT is not being introduced from outside. Lorentzian geometry is the natural coarse description of a finite-capacity, isotropic, relational substrate, and the Einstein sector is its lowest-order continuum gravitational expression. The entanglement scalar then tracks how that same capacity geometry is redistributed by localized defects. The resulting low-energy theory can therefore be written in the usual covariant language, but the intended logic runs from substrate properties to geometry, not the other way around. As with any discrete substrate, this is a continuum statement, and it faces a sharp known obstacle: a discrete structure that defines a preferred rest frame feeds dimension-four Lorentz-violating operators into the infrared with order-unity coefficients through loops \citep{Collins2004}, against laboratory bounds many orders of magnitude below unity. The protection here is structural. The tetrahedral ensemble is combinatorial and pre-geometric: it lives in the state counting from which the continuum is constructed, defines no embedding lattice in the emergent spacetime, and imprints on the EFT only through frame-independent scalars ($L_*$, $g_{\mathrm{share,eff}}$, $\eta_*$). Discreteness of this class is compatible with exact low-energy Lorentz symmetry, as causal-set sprinkling demonstrates by construction \citep{Bombelli1987,Dowker2004}. The cosmological bath does select a frame, in the same environmental sense the CMB does: a state rather than an operator, while the laboratory bounds constrain operators. The supporting calculation this argument calls for --- that substrate loop corrections generate no dimension-four Lorentz-violating operators --- is open and recorded as such in the closure table.

\subsection{4.2 Dependency Map of the Theory}

The logical flow of the theory begins with the three foundational postulates --- Information--Geometry, Mass--Entropy, and Many-Pasts --- on equal footing, and runs through the static weak-field chain before reaching the conditional sectors:
\[
\begin{gathered}
\{\text{Info--Geometry},\ \text{Mass--Entropy},\ \text{Many-Pasts}\}\\
\rightarrow\ \text{finite-capacity substrate ontology}
\ \rightarrow\ \text{tetrahedral boundary ensemble}\\
\rightarrow\ \text{admissibility closure}
\ \rightarrow\ \text{edge kernel / loop dressing / stiffness / source map}\\
\rightarrow\ \text{weak-field EFT}
\ \rightarrow\ \{\text{Newton},\ a_0,\ \text{RAR},\ \text{no-slip lensing}\}.
\end{gathered}
\]
Only then come the conditional and frontier sectors --- transport, clusters, cosmology, strong field, and the particle/gauge extensions --- each developed as a consequence or completion of the same framework.

Two features of this map matter. First, it is a dependency graph, not an epistemic-equality graph: the static weak-field sector, the UV coefficient chain that feeds it, and the operational Born-recovery branch are more tightly closed than the cosmological or strong-field sectors, and Part~VI makes that difference explicit in a closure-status table. Second, Many-Pasts appears among the foundational postulates rather than downstream, because its history weighting is already active in the scale-setting branch through memoryless dressing. The later Many-Pasts section (Section~21) develops the operational consequences of a postulate that has already done load-bearing work upstream.

\section*{Part II. UV Coefficient Chain}

Part~I fixed what the theory is about. The question now is whether the local capacity-sharing structure can actually be \emph{counted}. If the substrate has finite local capacity, the coefficients that appear in the continuum weak-field theory should not be free continuum parameters; they should descend from a finite local boundary problem. The next five sections follow that problem through: the smallest boundary cell that can carry capacity and close isotropically, the weighting that selects well-closed configurations, the cost of neighboring cells disagreeing, the local returns that dress that cost, and the continuum coefficient they leave behind. Nothing in this chain is matched to a gravitational observable; the observables enter only in Part~III.

Several of the ultraviolet choices below may look at first like independent tunings: the tetrahedral boundary cell, the seven face labels, the ordered injective assignment, the parity doubling, the admissibility kernel, the transverse ($2/3$) export, and the seven-sector electron anchor. None is a phenomenological knob in the closed branch. Section~5 and Appendix~B derive the minimal tetrahedral seven-state ensemble and its unique admissibility-closed entropy; Appendix~C fixes the transverse export from the tetrahedral bond-frame Green response; Appendices~D and~H derive the faithful sector-resolution length relation from the electron anchor and the memoryless dressing; and Appendix~L records the historical fork accounting and calibration status. The line to keep in view is between the primitive adoption of the finite-capacity tetrahedral substrate and the downstream coefficient chain, which contains no continuous fit once that substrate architecture is fixed. What the construction does not yet exhibit is a single microscopic action or path-integral measure whose continuum expansion generates these coefficients simultaneously; each is derived through its own local reduction, so the chain is a sequence of forced steps rather than one unified continuum limit, and consolidating it into a single measure is an open task.

\section{5. Why a Tetrahedral Boundary Ensemble}

The problem is to find the smallest discrete boundary cell that can carry finite channel entropy, close isotropically, and hand a single scalar response up to the continuum. The smallest structure that meets all three needs is a tetrahedral cell, fixed by four ingredients:
\begin{itemize}
\item a tetrahedral volumetric cell;
\item half-integer fermionic face data on each face;
\item injective face assignment;
\item binary orientation/parity.
\end{itemize}
This package is not presented as the only imaginable UV completion of emergent gravity. It is the minimal architecture used here to support the needed closure properties. The tetrahedron is the minimal volumetric simplex in $d=3$, injectivity preserves independent boundary information across the four faces, and parity doubling captures the two orientations of the cell. The face-state multiplicity is then not chosen from a menu. Postulate II identifies elementary defects as fermionic, so each face carries half-integer base spin
\[
j_0=\frac12,\frac32,\frac52,\ldots
\]
Two cells sharing a face therefore generate the effective boundary sector
\[
j_0\otimes j_0 = 0\oplus 1\oplus \cdots \oplus 2j_0.
\]
Postulate I selects the maximum-capacity boundary channel, so the effective face label is the top channel
\[
j_{\mathrm{eff}}=2j_0,
\]
with
\[
|M|=2j_{\mathrm{eff}}+1=4j_0+1
\]
distinguishable face states. Injectivity across four tetrahedral faces requires at least four distinct labels, so
\[
|M|\ge 4 \quad\Longrightarrow\quad 4j_0+1\ge 4.
\]
The only half-integer option below $j_0=3/2$ is $j_0=1/2$, which gives $j_{\mathrm{eff}}=1$ and $|M|=3$, so it fails the injectivity condition. The first fermionic choice that works is therefore
\[
j_0=\frac32,
\qquad
j_{\mathrm{eff}}=3,
\qquad
|M|=7.
\]
$j_0$ is sometimes mistaken for a free dial; it is not. Within this minimal construction $j_0$ is the smallest fermionic label that satisfies injectivity; a larger $j_0$ does not describe a fluctuation inside the same cell but a different, larger boundary ensemble. The minimal theory therefore has no $j$ to tune --- it has the first value that closes.

In that sense the seven-state face sector is derived from fermionic face data, maximum-capacity channel selection, tetrahedral injectivity, and minimality. The derivation chain itself makes no reference to the weak-field observables it later feeds; the order in which the construction was historically found, and where the evidential weight accordingly sits, are recorded in Appendix~L. The same face-level structure is also where the elementary matter sector enters: fermionic face exclusion creates the binary one-bit defect increment $\Delta S_f=\ln 2$ used later in the electron anchor.

The resulting combinatorial state count is
\[
\Omega_{\mathrm{tet}}=2\times P(7,4)=2\times 840=1680,
\]
so the combinatorial sharing ceiling is
\[
g_{\mathrm{share,max}}=\ln(1680)=7.42654907240.
\]
The exact $K^2$ spectrum and multiplicities are carried in the appendices. The $j$-labeled tetrahedron used here coincides with the quantum tetrahedron of simplicial spin networks \citep{Barbieri1998,BaezBarrett1999}, whose discrete geometric spectra \citep{RovelliSmolin1995} arise from the same $SU(2)$ representation theory; the present construction differs in weighting these states by admissibility closure rather than by a spin-foam amplitude, and in routing them to a capacity entropy rather than to area and volume operators. The UV theory thus begins with a finite microscopic counting problem rather than a free continuum ansatz.

\section{6. Admissibility Closure}

\subsection{6.1 Minimal isotropic kernel}

Not every boundary configuration should count equally. The raw combinatorial ensemble is too permissive to be the whole UV story: some configurations sit close to the regular closure pattern expected of a smooth local cell, while others are badly distorted. Admissibility closure is the statement, in its mildest form, that more poorly closed configurations contribute less to the coarse ensemble. The minimal rotationally invariant measure of that distortion is a single quadratic closure-defect scalar $K^2$, and the weighting it induces is
\[
p_\eta(b)\propto e^{-\eta K^2(b)}.
\]
This is not chosen because it works phenomenologically; it is the minimal isotropic maximum-entropy kernel under normalization and a fixed quadratic closure moment. Higher invariants such as $K^4$ carry additional UV information and so enter as subleading refinements, not as competing leading kernels.

\subsection{6.2 Closure condition and uniqueness}

The admissibility precision $\eta$ is not chosen externally; it is fixed by maximizing the normalized closure evidence. Tetrahedral closure is the vanishing of the three-component oriented-face sum, so the closure-defect space is three-dimensional, and the quadratic family on it carries a determinant weight $\eta^{3/2}$. The closure-evidence functional is therefore
\[
\mathcal F(\eta)=\ln Z(\eta)+\tfrac32\ln\eta,
\]
and its stationary point gives the closure condition
\[
\langle K^2\rangle_\eta = \frac{3}{2\eta},
\]
in which the factor $3/2$ is the determinant weight of the three independent closure components, while the discreteness and multiplicities of the spectrum stay inside the exact sum $Z(\eta)$. This is the stationary normalized-evidence point of the exact closure spectrum, and it is a maximum rather than a bare root (Appendix~B.2).

On that spectrum it is unique,
\[
\eta_* = 0.0298668443935.
\]
The closed branch is locally stiff: small fractional changes in $\eta$ produce only small fractional changes in the downstream effective sharing entropy.

\subsection{6.3 Effective sharing entropy}

The admissibility-weighted effective sharing entropy is
\[
g_{\mathrm{share,eff}} = 7.41980002357.
\]
The gap between $g_{\mathrm{share,max}}$ and $g_{\mathrm{share,eff}}$ is therefore not loss imposed by hand. It is the difference between the raw combinatorial ceiling and the admissibility-closed effective boundary entropy that actually propagates into observable couplings.

The continuum description does not inherit the naive channel-counting ceiling; it inherits the portion of the channel space that survives after closure is imposed. The downstream couplings should therefore be read as consequences of admissibility-closed sharing, not of raw combinatorics alone.

With $\eta_*$ fixed, the effective sharing entropy carries no remaining freedom; the exact spectrum, multiplicities, and uniqueness proof are given in Appendix~B.2.

\section{7. Edge Kernel and Tree-Level Coupling}

The admissibility closure says what one cell can carry. The edge kernel says how costly it is for two neighboring cells to carry different things, and that cost becomes stiffness in the continuum field: a medium whose cells resist disagreement strongly is one whose capacity deficits spread reluctantly. The same UV closure data fix this cost. The geometric bridge is the tetrahedral identity
\[
\sum_{i=1}^4 \hat n_i\hat n_i^{\mathsf T}=\frac{4}{3}I_3,
\]
which implies a channel-averaged transverse fraction of $2/3$ and gives the bare edge smoothness coupling
\[
J_{\mathrm{bare}}=\frac{2}{3}\eta_*.
\]
If adjacent cells disagree strongly the edge pays a larger penalty; if they agree, the penalty is small. The factor $2/3$ is the geometric fraction that survives after averaging the four tetrahedral channel directions into the isotropic continuum limit --- the part of the disagreement that the scalar sharing channel actually carries.

For a $z=4$ regular coarse adjacency graph, the tree-to-lattice reduction then yields
\[
J_{\mathrm{eff}}^{\mathrm{tree}}=\frac{J_{\mathrm{bare}}}{3}=\frac{2\eta_*}{9}.
\]
The division by $3$ comes from the branching geometry of the rooted $z=4$ graph. One neighboring link points back toward the source, while the remaining $z-1=3$ links carry the forward transport into the tree. Thus $J_{\mathrm{eff}}^{\mathrm{tree}}$ is not simply the microscopic edge penalty itself, but the part of that penalty that survives as net long-range transport after the local branching structure is taken into account.

\paragraph{Origin of the horizon target.}
The horizon target
\[
\sigma_*=\frac{\pi}{g_{\mathrm{share,eff}}}
\]
is the closure-consistency value required by the horizon-normalized field convention. In the admissibility-closed boundary ensemble, one active microscopic sharing unit carries effective entropy $g_{\mathrm{share,eff}}$. In the continuum normalization used for the weak-field scalar, the occupancy variable is normalized by
\[
S=\pi Q_{\mathrm{occ}},
\]
so a coarse horizon-normalized channel with occupancy $Q_{\mathrm{occ}}=1$ carries entropy $\pi$ in the $S$-field convention. If $\sigma_*$ denotes the asymptotic conditional-independence weight seen by the rooted shell hierarchy (Appendix~B.3), consistency between the boundary entropy count and the horizon-normalized continuum field requires
\[
\sigma_*\,g_{\mathrm{share,eff}}=\pi,
\]
and therefore
\[
\sigma_*=\frac{\pi}{g_{\mathrm{share,eff}}}=0.42340665\ldots .
\]
The role of $\sigma_*$ is to match the admissibility-closed microscopic entropy normalization to the horizon-normalized scalar-field convention; it is the closure target fixed by that choice of branch. The rooted shell observable converges to that target rapidly enough that the nonlocal correction is already strongly constrained by small shell depth. The tetrahedral identity keeps this bridge controlled: the four discrete channel directions average to the correct isotropic tensor structure in the continuum limit, so the same combinatorial data that fixed admissibility also fix the tree-level transport. The shell hierarchy and phase-selection checks are carried in Appendix~C.

\section{8. Finite-Loop Renormalization}

Tree level is not the whole UV story. The full lattice admits local closed-return motifs that recycle part of the transmitted information before it contributes to net coarse transport. The leading correction is organized as a local Dyson self-energy dressing,
\[
J_{\mathrm{eff}}^{(\mathrm{ren})}
=
\frac{J_{\mathrm{eff}}^{\mathrm{tree}}}{1+J_{\mathrm{eff}}^{\mathrm{tree}}\Sigma_{\mathrm{ret}}}.
\]

The need for this step is physically straightforward. A purely tree-like transmission rule would let the relevant amplitude move outward once and never locally return. A real coarse graph is not that simple. Some of the transmitted information cycles back through short closed motifs before contributing to long-distance transport. The renormalized coupling is therefore the true stiffness felt by the coarse field after these local returns have been resummed.

The structure of that self-energy is not a generic loop number. The returns split into seven sector-diagonal channels and one collective mode. The seven are the face-label channels, each returning independently without mixing. The one is the permutation-symmetric combination across channels, which returns as a shared closure-singlet rather than as a channel-specific loop, and it is weighted by the same transverse projection and branch-dilution factors that define the tree edge map,
\[
\left(\frac{2}{3}\right)\left(\frac{1}{3}\right)=\frac{2}{9}.
\]
The leading local self-energy is therefore
\[
\Sigma_{\mathrm{ret}}=7+\frac{2}{9}=\frac{65}{9}.
\]
Equivalently, on the seven-channel scalar return space,
\[
R_{\mathrm{ret}}=I_7+\frac{2}{9}P_{\mathrm{sing}},
\qquad
P_{\mathrm{sing}}=|u\rangle\langle u|,
\qquad
u=\frac{1}{\sqrt7}(1,\ldots,1),
\]
with $\Sigma_{\mathrm{ret}}=\mathrm{Tr}(R_{\mathrm{ret}})$. The orthogonal six-dimensional sum-zero sector carries no net scalar charge in the coarse branch and so adds no separate scalar return. The singlet weight is fixed by the tree map, not introduced here, so no new loop parameter appears.
\[
c_{\mathrm{loop}}^{(\mathrm{ren})}
\equiv
\frac{J_{\mathrm{eff}}^{(\mathrm{ren})}}{J_{\mathrm{eff}}^{\mathrm{tree}}}
=
\frac{1}{1+J_{\mathrm{eff}}^{\mathrm{tree}}\Sigma_{\mathrm{ret}}}
\approx 0.95426,
\]
and
\[
J_{\mathrm{eff}}^{(\mathrm{ren})}\approx 0.00633348.
\]
This reproduces the shell-target crossing near $J_{\mathrm{bare,cross}}\sim 0.019$ at the $0.05\%$ level.

The loop correction is no longer schematic: the finite renormalization is written as an explicit local self-energy. The remaining audit task is independent graph-level confirmation of the same scalar-return operator, not the introduction of any new loop parameter.

\section{9. Continuum Stiffness and SI Normalization}

The last UV step is not a thermodynamic one. The lattice quadratic form is interpreted as a Euclidean action weight,
\[
\frac{I_E}{\hbar}
=
\frac{J_{\mathrm{eff}}^{(\mathrm{ren})}}{2}
\sum_{a,i}(Q_a-Q_{a+L_*\hat n_i})^2,
\]
where the sum runs over all sites $a$ and all four outgoing nearest-neighbor directions $\hat n_i$, so each nearest-neighbor edge is counted twice. The microscopic four-cell has size
\[
\Delta V_4=\frac{L_*^4}{c}.
\]
Up to this point the derivation has determined a dimensionless lattice weighting. The continuum EFT, however, needs a dimensionful coefficient multiplying derivatives of a field in spacetime. The Euclidean-action interpretation upgrades the lattice closure data into a continuum action density with the right units and the right covariant target.

The same tetrahedral identity used in the edge-kernel reduction then yields the continuum coefficient for the occupancy field $Q_{\mathrm{occ}}$,
\[
\gamma_Q=\frac{4\hbar c}{3L_*^2}J_{\mathrm{eff}}^{(\mathrm{ren})}.
\]
Here $L_*$ is the canonical tetrahedral spacing, with one coarse cell carrying volume $L_*^3$ up to the fixed cell-shape convention, and $J_{\mathrm{eff}}^{(\mathrm{ren})}$ is the loop-dressed edge coupling. The numerical factor $4/3$ is the isotropic projection
\[
\sum_i \hat n_i\hat n_i^T=\frac43 I_3
\]
that turns the tetrahedral edge directions into the continuum gradient tensor.

The field normalization is fixed by horizon capacity:
\[
S=\pi Q_{\mathrm{occ}}.
\]
Therefore the canonical EFT coefficient in the $\frac{\gamma}{2}(\partial S)^2$ convention is
\[
\gamma=\frac{4\hbar c}{3\pi^2L_*^2}J_{\mathrm{eff}}^{(\mathrm{ren})}.
\]
Physically, $\gamma$ is the continuum stiffness of the entanglement-capacity field. A larger $\gamma$ makes spatial gradients more costly and suppresses the capacity-deficit response to a given source; a smaller $\gamma$ allows larger variations of the field.
The faithful sector-resolution principle fixes $L_*$ without using $G$. It is nevertheless useful to define the gravitational scale induced by this length,
\[
G_*:=\frac{c^3L_*^2}{\hbar}.
\]
Then the stiffness may be written in Einstein-normalized form as
\[
\gamma=\frac{4J_{\mathrm{eff}}^{(\mathrm{ren})}}{3\pi^2}\frac{c^4}{G_*}.
\]
This is the same algebra as the familiar Planck-cell rewrite, but read in the opposite direction: the substrate cell length induces the gravitational scale rather than being chosen by first inserting the measured value of $G$. Within the Euclidean-action convention already assumed by the EFT, the SI-normalized weak-field stiffness coefficient is fixed rather than left schematic.

This completes the micro-to-continuum coefficient chain. The tetrahedral ensemble determines the effective sharing entropy; the edge kernel and loop dressing turn that entropy into a discrete stiffness; and the Euclidean matching turns the discrete stiffness into the continuum coefficient $\gamma$ of the weak-field EFT.

\paragraph{Closed UV-to-IR chain.}
The UV coefficient chain can now be summarized as
\[
\{\Omega_{\mathrm{tet}},\,K^2,\,\eta_*,\,g_{\mathrm{share,eff}},\,L_*,\,J_{\mathrm{bare}},\,J_{\mathrm{eff}}^{\mathrm{tree}},\,\Sigma_{\mathrm{ret}},\,J_{\mathrm{eff}}^{(\mathrm{ren})},\,\gamma\}
\longrightarrow
\{\kappa,\,G,\,a_0,\,g_{\mathrm{obs}}(g_{\mathrm{bar}})\}.
\]
The first bracket is the micro-to-continuum closure chain; the second bracket collects the weak-field observables it feeds. The rest of the manuscript uses this chain rather than introducing independent weak-field coefficients.

The remaining microscopic question is independent confirmation of the same action-kernel interpretation from fuller inhomogeneous dynamics, not an unresolved normalization constant.
\section*{Part III. Weak-Field EFT and Static Phenomenology}

\section{10. Covariant Action}

With that continuum symmetry structure in place, the canonical weak-field EFT takes the covariant form
\[
I = \int d^4x\,\sqrt{-g}\left[
\frac{c^4}{16\pi G}R
-\frac{\gamma}{2}g^{\mu\nu}\partial_\mu S_{\mathrm{ent}}\partial_\nu S_{\mathrm{ent}}
-\lambda S_{\mathrm{ent}}
-\kappa\chi S_{\mathrm{ent}}
\right],
\]
with
\[
\chi(x)\equiv -\frac{T^\mu{}_\mu}{c^2}.
\]
The action should be read as the simplest weak-field continuum realization of the ontology already stated in Part I and the coefficient chain already derived in Part II. The metric sector remains the familiar Einstein one at low energy, but it is now interpreted as the continuum capacity geometry of the substrate rather than as an independent starting theory. It is coupled to a scalar field that tracks available entanglement capacity and to a source channel written in ordinary stress-energy notation while still being interpreted microscopically as the defect sector of the same medium.

At the EFT level $\chi$ is written in ordinary stress-energy language, but ontologically it is the coarse trace channel of the localized defect sector. Here $\gamma$ is the continuum stiffness fixed by the UV chain, while $\kappa$ is the defect--entropy coupling fixed by the canonical source map,
\[
\kappa=\frac{\Xi_\rho}{L_*^2\kappa_m(L_*)},
\]
and $\lambda$ controls the background branch. The source map is determined by matching the scalar mode on the lattice and in the continuum. Appendix~C fixes the rigid defect amplitude through the isotropic defect benchmark $\Delta S_{\mathrm{def}}=\ln(7/6)$ and the exact tetrahedral on-site Green constant $G_{\mathrm{tet}}(0)=0.448220394\ldots$. With defect-entropy density
\[
\sigma_{\mathrm{def}}=\frac{\rho}{\kappa_m(L_*)},
\]
the Green-matched continuum projection is
\[
\nabla^2\delta S
=
-\frac{3L_*}{4G_{\mathrm{tet}}(0)}\,\sigma_{\mathrm{def}},
\]
and comparison with $\nabla^2\delta S=-(\kappa/\gamma)\rho$ gives
\[
\frac{\kappa}{\gamma}
=
\frac{3L_*}{4G_{\mathrm{tet}}(0)\kappa_m(L_*)}.
\]
This is the source-side counterpart of the stiffness derivation: the same tetrahedral scalar mode that carries the edge stiffness is the mode sourced by localized mass defects. The fixed-epoch normalization against $S_\infty$ still enters the final gravity dictionary, but it is not an additional source freedom; $S_\infty$ and $\kappa$ rescale together under a change of entropy units, while $\kappa/(\gamma S_\infty)$ is invariant. Local weak-field dynamics are studied in the renormalized branch
\[
\lambda_{\mathrm{ren}}\equiv \lambda+\gamma\Box S_{\mathrm{bg}}=0,
\]
so that local perturbations are sourced only by the defect sector, written at continuum level in ordinary matter variables.

The Einstein--Hilbert coefficient is written here in the already-matched Einstein normalization of the weak-field EFT. This is the metric normalization in which weak-field gravity is observed, not a separate microscopic parameter added on top of the entanglement chain.

\paragraph{Three roles of \(G\) in the manuscript.}
There are three distinct uses of $G$ in what follows. First, the weak-field metric normalization gives the identity
\[
G=\frac{c^2\kappa}{8\pi\gamma S_\infty}.
\]
This is the gravitational scale implied by the scalar source coupling, scalar stiffness, and capacity-to-lapse bridge. The second use is the non-gravitational scale-setting route. The faithful sector-resolution principle fixes
\[
L_*=\frac32\,\lambda_e e^{-7g_{\mathrm{share,eff}}},
\qquad
\lambda_e=\frac{\hbar}{m_ec},
\]
and therefore induces
\[
G_*=\frac{c^3L_*^2}{\hbar}
=
\frac94\,\frac{\hbar c}{m_e^2}\,e^{-14g_{\mathrm{share,eff}}}.
\]
The electron anchor enters this route through the reduced Compton length and through the identification of the electron as the lightest one-bit fermionic defect, while the mass--entropy anchor still fixes $\kappa_m(\lambda_e)=m_e/\ln2$ for the source map. Third, formulas involving the conventional Planck length,
\[
L_P=\sqrt{\frac{\hbar G}{c^3}},
\]
are matched representations after the gravitational scale has been identified. They are useful for compactness, comparison, and black-hole thermodynamic notation, but the scale-setting argument runs through $L_*$, not through a Planck-cell rewriting.

\paragraph{Definition.} Faithful sector resolution fixes the substrate cell length $L_*$ without using $G$, and the induced gravitational scale $G_*=c^3L_*^2/\hbar$ follows from the electron Compton length and the tetrahedral entropy alone. The chain takes no value of $G$ as input.

\paragraph{Comparison.} Independently of that length, the scalar bridge and Green-matched source map fix the weak-field normalization through the invariant ratio $\kappa/(\gamma S_\infty)$,
\[
G=\frac{c^2}{8\pi}\frac{\kappa}{\gamma S_\infty},
\]
which lets $G$ be compared with $G_*$. Because both are built from the single substrate scale $L_*$, this is not a second measurement from independent inputs; it is a coherence check of whether that one scale propagates through the scalar stiffness, the source map, and the bridge to the observed normalization. The check has content because fixing $L_*$ does not by itself guarantee that $\kappa/\gamma$, $S_\infty$, and the bridge organize around it to give the right normalization.

\paragraph{Status.} The two agree to about one percent. We read this as a calibrated coherence check, not a clean independent prediction, because the construction was developed with the measured $G$ partly in view; Appendix~L records that provenance and the associated fork accounting. The dressing dynamics behind faithful sector resolution is derived in Section~22 and Appendix~H.

This is the simplest covariant realization of the closure chain: one metric, one scalar entanglement field, one trace-equivalent defect-source channel, and one renormalized background branch.

Once this weak-field covariant form is accepted as the correct low-energy language of the substrate, the earlier UV closure chain fixes the entanglement-side coefficients that appear in it; the terms in the action are then read off from that chain, not chosen term by term to fit phenomenology.

\section{11. Field Equations and Bridge Law}

Varying the action with respect to $S_{\mathrm{ent}}$ gives the sourced scalar equation
\[
\gamma\Box S_{\mathrm{ent}}=\lambda+\kappa\chi.
\]
Varying with respect to the metric yields
\[
G_{\mu\nu}=\frac{8\pi G}{c^4}
\left(
T_{\mu\nu}^{(\mathrm{matter})}
+T_{\mu\nu}^{(\mathrm{ent})}
\right),
\]
with the canonical scalar stress-energy induced by the entanglement sector.

These two equations separate the two jobs played by the scalar. The scalar equation tells us how the local entanglement-capacity variable responds to defect sources. The metric equation tells us how that scalar response then contributes back to spacetime curvature. The bridge law below turns those two statements into an ordinary weak-field gravitational potential.

The weak-field bridge law is not inserted as an arbitrary interpolation. It is the normalization map between fractional capacity depletion and the weak-field lapse. Let
\[
\epsilon(x)\equiv\frac{\delta S(x)}{S_\infty}.
\]
Vacuum normalization requires the lapse to approach its asymptotic value when $\epsilon=0$. Locality requires the leading metric response to depend on the local fractional deficit. Independent small deficits must add at leading order, while independent redshift factors multiply. These requirements force the lapse to take exponential form,
\[
N(\epsilon)=\exp[-C\epsilon],
\]
with a single constant $C$. The standard weak-field metric normalization,
\[
N=1+\frac{\Phi}{c^2}+O(\Phi^2/c^4),
\]
fixes $C=1/2$. Hence
\[
N=\exp\!\left[-\frac{\delta S}{2S_\infty}\right]
=1-\frac{\delta S}{2S_\infty}
 +O(\delta S^2/S_\infty^2),
\]
and therefore
\[
\frac{\Phi}{c^2}=-\frac{\delta S}{2S_\infty}.
\]
In the static weak-field branch the emergent Newton constant is therefore
\[
G=\frac{c^2\kappa}{8\pi\gamma S_\infty}.
\]

The bridge law connects the entropy description to observable gravity: the deficit field acquires a unique weak-field normalization in terms of the ordinary gravitational potential.

\paragraph{Matter enters once.}
The two equations are not two independent ways for ordinary matter to source the weak-field potential, and reading them as such would double-count it. Matter enters the weak-field branch through a single channel: the trace $\chi=-T^\mu{}_\mu/c^2$ sources the scalar equation, the scalar response $\delta S$ is converted to the lapse by the bridge law, and the observed Newtonian potential is that bridged lapse. The metric equation is the covariant image of the same response --- $T_{\mu\nu}^{(\mathrm{ent})}$ is the stress-energy of the scalar field that $\chi$ has already sourced, not a separate matter contribution --- so $G=c^2\kappa/(8\pi\gamma S_\infty)$ and the Newtonian limit of $G_{\mu\nu}=(8\pi G/c^4)(T^{(\mathrm{matter})}_{\mu\nu}+T^{(\mathrm{ent})}_{\mu\nu})$ are the same field written two ways. Ordinary matter is therefore not a second, independent source of $\Phi$ standing beside the bridge; it is the source of the scalar response that the bridge reads as the potential, and Bianchi consistency enforces matter conservation through that same trace channel rather than through a separate coupling. The Einstein equation above is thus written in the already-bridged physical-metric frame --- $T^{(\mathrm{ent})}_{\mu\nu}$ records the scalar response the bridge has converted into the lapse --- and not as an independent second scalar force acting on test bodies.

\section{12. Newtonian Gravity and the Point-Source Limit}

In the renormalized static weak-field sector the scalar equation reduces to
\[
\nabla^2\delta S=-\frac{\kappa}{\gamma}\rho.
\]
This is the point where the micro-to-macro chain becomes operationally familiar. Once the background is renormalized away and the source is nonrelativistic, the scalar sector obeys an ordinary Poisson equation for the deficit field. The unusual quantity is $\delta S$, but the mathematical structure is the same one that underlies standard weak-field gravity.

For a point source $M$,
\[
\delta S(r)=\frac{\kappa M}{4\pi\gamma r}.
\]
Using the bridge law,
\[
\frac{\Phi}{c^2}=-\frac{\delta S}{2S_\infty},
\]
the gravitational acceleration becomes
\[
g(r)=\frac{c^2\kappa}{8\pi\gamma S_\infty}\frac{M}{r^2}
=
\frac{GM}{r^2}.
\]
Thus Newtonian gravity is recovered as the weak-field response of the entanglement-capacity medium.

Nothing qualitatively exotic has to be inserted at the last step to recover ordinary gravity. The same sourced scalar equation and the same bridge law already imply the familiar point-mass force law. In that sense Newtonian gravity appears here not as a starting axiom but as the first infrared limit of the entanglement medium.

\paragraph{Interpretation.}
This is the first point where the capacity-strain picture becomes ordinary gravity. A point defect produces a $1/r$ capacity deficit; the bridge reads the gradient of that deficit as the Newtonian inverse-square force. Ordinary gravity is the small-deficit, weak-curvature limit of the extended capacity strain around localized defects, so the $1/r^2$ law is emergent, not fundamental. Section~23 exhibits the same structure on the discrete substrate: a capacity budget that merely re-equilibrates cell by cell screens at sub-cell range, while a conserved capacity current with maintenance sinks obeys the massless graph-Poisson equation and produces this $1/r$ deficit: the Newtonian form doubles as a diagnostic of the conservation law behind it.

\section{13. Electron Anchor: One-Bit Mass Scale and Seven-Sector Length Scale}

The electron does two separate jobs in this framework, and keeping them apart turns the weak-field normalization into a genuine consistency check rather than a single calibration. As the lightest one-bit fermionic defect, it is the cleanest \emph{mass} anchor for the mass--entropy map. As the lightest coherent defect whose dressing resolves the seven face sectors, its Compton scale is the \emph{length} anchor that fixes the substrate cell size. The two uses draw on the same particle but on different facts about it --- its mass on one side, its coherent dressing support on the other.

\subsection{13.1 Why the electron is the anchor}

The mass--entropy map needs a clean elementary anchor because, in this framework, the elementary matter sector \emph{is} the localized defect sector. The electron is the natural choice: it is the lightest simple charged fermionic defect, not a composite, and its mass is not obscured by hadronic or QCD dressing. A single fermionic face-exclusion defect carries the canonical increment
\[
\Delta S_f=\ln 2,
\]
one bit of missing entanglement, because an excluded face is a binary occupied/unoccupied defect of the local network.

\subsection{13.2 One-bit mass anchor}

At the electron Compton scale $\ell=\lambda_e$ the mass--entropy map reads
\[
\kappa_m(\lambda_e)=\frac{m_e}{\ln 2}.
\]
Dividing the electron mass by the fixed one-bit increment fixes the mass-per-entropy conversion at the electron's own scale. Run back to the cutoff cell, the conversion is
\[
\kappa_{m,\mathrm{UV}}=\frac{\hbar}{cL_*}\frac{1}{\ln 2},
\]
with canonical running law
\[
\kappa_m(\ell)=\kappa_{m,\mathrm{UV}}\left(\frac{L_*}{\ell}\right)^{1+\alpha_{\mathrm{cl}}},
\qquad
\alpha_{\mathrm{cl}}=0
\]
in the closed branch. One bit fixes the electron-scale conversion; the running law carries it to the UV scale; the same conversion then feeds the weak-field source map. This is the only point at which the mass anchor enters gravity.

\subsection{13.3 Seven-sector length anchor}

The same electron fixes the absolute cell length through a different fact about it. Its reduced Compton wavelength $\lambda_e$ is the coherent support of the lightest defect's ground-state dressing, and that dressing resolves the seven face sectors of the tetrahedral alphabet. Two properties make the resolution clean. The Many-Pasts weighting of Postulate~III makes each dressing pass memoryless in substrate time, so each pass carries the full admissibility entropy $g_{\mathrm{share,eff}}$; and the electron, as the lightest charged defect, carries no inter-channel correlation, so the seven sectors are resolved independently. Seven independent resolutions add, giving the support exponent $7g_{\mathrm{share,eff}}$ and hence
\[
\ln\!\left(\frac{\lambda_e}{L_*}\right)
=
7g_{\mathrm{share,eff}}-\ln\!\left(\frac32\right),
\]
with the $\ln(3/2)$ the transverse export factor of Appendix~C.5.

This does not mean an electron is a single tetrahedral cell carrying seven simultaneous labels. The local ensemble supplies the seven-sector alphabet; the electron is the lightest coherent one-bit defect whose ground-state dressing resolves those sectors once and exports the transverse share of the resulting support. The Compton scale calibrates the substrate length hierarchy; the one-bit mass calibrates the source map. The two uses impose a nontrivial joint requirement on the electron's role in the gravitational normalization without duplicating a single input.

\paragraph{The anchor is a gauge choice.}
Because counting fixes only dimensionless quantities, exactly one dimensionful measurement must be supplied, and which one is a convention. The physical content of the framework is its anchor-invariant statements: the lepton ratios $m_\mu/m_e$ and $m_\tau/m_e$ (Appendix~I.1), the hierarchy $m_e/m_P=\tfrac32e^{-7g_{\mathrm{share,eff}}}$, the horizon coefficient $1/4$ (Appendix~F.5), the abundance ratio $\Omega_c/\Omega_b$ (Section~18.5), and, within its branch assignment, $a_0/cH_0$ (Section~14). Anchoring on the electron reads the induced scale $G_*$ about one percent below the measured $G$ (Section~13.4); anchoring on the measured $G$ instead reads $m_e/m_P$ about half a percent below the measured ratio. These are the same residual in different units, and its status travels with the relation rather than with the anchor: from either end it remains the calibrated coherence check recorded in Appendix~L. Only dimensionful quantities connected to the counting by closed relations can serve as the anchor --- the electron mass, the gravitational scale, or the thermodynamic pair $(A,T_H)$ of a single horizon; the acceleration $a_0$ cannot, because its relation passes through the epoch-dependent $H(z)$. The electron serves as the anchor because it is the most precisely measured of these.

\subsection{13.4 Consistency checks}

Two cross-checks test the dual role rather than feed it. The first is the support exponent itself. The length relation reads the physical ratio as the \emph{first} power of the dressing support, $\lambda/L_*\propto e^{H}$ with $H=7g_{\mathrm{share,eff}}$ for the electron, and that linear power is not imposed. Letting it float as $\lambda/L_*\propto e^{\alpha H}$ and demanding the observed Newton constant from the electron branch requires $\alpha=0.9999$, while the same map applied to the charged-lepton ladder (Appendix~I.1) is independently solved by $\alpha_\mu=1.001$ and $\alpha_\tau=1.001$. A square-root map ($\alpha=\tfrac12$) would misplace $m_\mu/m_e$ near $14$ and a quadratic map ($\alpha=2$) near $4\times10^{4}$, so the spectrum excludes them by orders of magnitude. Three determinations across two sectors that share no gravitational input bracket $\alpha=1$ within $10^{-3}$. The dictionary itself is derived in Appendix~D.4 from the recurrence structure of the dressing: the waiting time for joint typicality of independent channels is the exponential of the summed entropies, so the linear power is the arithmetic of joint probability. The bracketing is the quantitative check of that derivation; identifying one phase cycle per dressing recurrence is the single assumption that remains beneath it.

The second cross-check is the induced gravitational scale. The length anchor fixes $L_*$, and hence $G_*=c^3L_*^2/\hbar$, from the electron Compton length and the tetrahedral entropy alone, with no gravitational input. Whether that one substrate scale also propagates through the scalar stiffness, the Green-matched source map, and the weak-field bridge to the observed normalization is a coherence test, not a second independent measurement (Section~10). The comparison agrees to about one percent, and Appendix~L records how the measured value of $G$ figured in selecting the construction.

The same comparison has a dimensionless form. Squaring the hierarchy gives the electron's gravitational coupling,
\[
\frac{Gm_e^2}{\hbar c}=\frac94\,e^{-14g_{\mathrm{share,eff}}},
\]
with the $9/4$ the square of the transverse export factor; it agrees with the measured coupling to the same one percent; this is the Appendix~L residual re-expressed. The relation also has an area reading. A horizon stores one bit in area $4\ln2\,L_P^2$, the saturated packing of Appendix~F.5, while the electron spreads its single bit $\Delta S_f=\ln2$ over the Compton area $\lambda_e^2$, so the ratio of the electron's bit density to the horizon's saturated density is $4\ln2\,\cdot Gm_e^2/\hbar c$. Gravity couples to the electron weakly for the same reason the electron is light: its one bit is diluted a factor $e^{14g_{\mathrm{share,eff}}}$ below capacity-saturated packing. The identity adds no new number --- it is the length anchor squared --- but it states the weakness of the electron's gravity and the smallness of its mass as one fact about packing density.

\subsection{13.5 Composite sectors}

For composite hadrons the claim is weaker and different in kind. The relevant quantity is the dressed, vacuum-subtracted bound-state entropy,
\[
m_{\mathrm{hadron}}=\kappa_m(\ell_H)\,S_{\mathrm{ent},H}^{\mathrm{dressed}},
\]
with the dressed budget generated by confinement, gluonic structure, trace-anomaly dynamics, and chiral vacuum reorganization. A finished lattice derivation of that dressed entropy is not yet available; what is claimed is structural compatibility between the mass--entropy map and the standard QCD mass budget. The elementary-fermion anchor is settled in the simple sectors; the hadronic sector remains structurally compatible but not yet coefficient-complete.

\section{14. Galactic Dynamics}

The static Newtonian branch already gives the ordinary baryonic force. The galactic excess appears when, at low acceleration, the transverse vacuum-capacity mode adds a bosonic emission factor on top of that force, and the scale at which this turns on is set by horizon thermality reduced by the microstructural sharing factor. The rest of the section is the precise version of those two sentences, and it separates cleanly into two questions: where the turn-on scale $a_0$ comes from, and why the resulting radial-acceleration relation has the shape it does.

The galactic sector is one of the main payoffs of the coefficient chain. Its characteristic acceleration scale, the radial-acceleration relation, and the deep-MOND limit all follow from one chain that uses no per-galaxy interpolation function. The chain has three structural inputs: a stable bosonic vacuum mode (Appendix~D.6), the Gibbons--Hawking thermal structure of that vacuum at the apparent horizon (Appendix~E), and the $1+2$ channel decomposition that places the cosmic horizon scale in the transverse sector.

\paragraph{Vacuum-state origin of the Bose--Einstein occupancy.}
The Bose--Einstein statistics in the galactic mode sector follow from the vacuum-state thermality of the entanglement field at the apparent horizon, with the cosmological boundary normalization fixed in Appendix~E. The mechanism is therefore equilibrium horizon thermality, not dynamical thermalization.

Appendix~D.6 establishes that $\delta S$ fluctuations around the on-shell background constitute a massless bosonic scalar at quadratic order with positive kinetic stiffness $\gamma>0$. For a massless bosonic scalar in a spacetime with apparent horizon $R_A$, the vacuum state restricted to the static patch is thermal at the de Sitter / apparent-horizon temperature
\[
k_B T_H=\frac{\hbar c}{2\pi R_A},
\]
with mode occupancy $n_B(\epsilon)=1/(e^{\epsilon/k_BT_H}-1)$. This is a Gibbons--Hawking statement of QFT in curved spacetime, fixed by the same horizon geometry that sets $S_\infty$.

It is useful to write a horizon temperature as an acceleration scale,
\[
a_T\equiv\frac{2\pi c}{\hbar}k_BT.
\]
For the apparent-horizon vacuum this gives
\[
a_T=\frac{c^2}{R_A}\equiv a_H.
\]
In the closed transport branch, $R_A^{-1}=H_0/c$ at the present epoch (from \S17 with $\tau_0^{-1}=H_0$), so
\[
a_H=\frac{c^2}{R_A}=cH_0.
\]
The acceleration scale defined by horizon thermality is therefore fixed by the same closure chain that fixes $S_\infty$. More generally $a_H(t)=cH(t)$ at any cosmological epoch, so the horizon thermal scale is epoch-dependent. The present-epoch value $cH_0$ is the one relevant to the local ($z\approx 0$) galactic RAR.

The MOND scale $a_0$ is related to $a_H$ by the microstructural sharing factor that connects $g_{\mathrm{share,eff}}$ to phase-space transverse-mode density in the $1+2$ channel decomposition:
\[
a_0=\frac{g_{\mathrm{share,eff}}}{4\pi^2}\,a_H=\frac{cH_0g_{\mathrm{share,eff}}}{4\pi^2}.
\]
Hence $a_0$ is the horizon thermal acceleration $a_H=cH_0$ reduced to the transverse scalar sector of the $1+2$ decomposition. The reduction factor $g_{\mathrm{share,eff}}/(4\pi^2)$ is fixed, not fitted. Its $(2\pi)^2$ is the canonical two-dimensional mode density $d^2k_\perp/(2\pi)^2$ --- the number of states per phase-space cell of a continuum two-plane, a quantization normalization rather than a Fourier convention --- and $g_{\mathrm{share,eff}}$ is the admissibility-sharing content carried per cell. The normalization is $(2\pi)^2$, not the spherical $4\pi$, because the weak-field response samples the local transverse two-plane orthogonal to the baryonic forcing, not the full horizon sphere. That the split exists at all is kinematic, not a modeling choice: the baryonic gradient picks the longitudinal direction $\hat n=\nabla\Phi_{\rm bar}/|\nabla\Phi_{\rm bar}|$, and its orthogonal complement in three-space is two-dimensional by dimension count, so the transverse two-plane and its $(2\pi)^2$ cell are forced. The one premise that is physical rather than geometric is the assignment of the low-acceleration response to the \emph{stationary} transverse vacuum branch rather than the telegrapher transport sector of Section~17 --- the step that ties $a_0$ to $cH_0$ rather than to a transport timescale, and that is testable through the epoch dependence $a_0(z)\propto H(z)$. Given it, the coefficient follows, and $a_0$ and $cH_0$ share one origin.

\paragraph{$1+2$ channel decomposition and the radial-acceleration law.}
The longitudinal acceleration scale is $g_{\mathrm{bar}}$ and the transverse vacuum reference scale is $a_0$. Each channel carries an acceleration frequency through the Unruh correspondence \citep{Unruh1976},
\[
\omega_\parallel\propto\frac{g_{\mathrm{bar}}}{c},
\qquad
\omega_\perp\propto\frac{a_0}{c},
\]
with a common prefactor that cancels in the occupancy argument below. The two channels mix, and a coupled pair of modes softens at the geometric mean of its two frequencies rather than their average; that softened mixed mode dominates the low-acceleration response. Precisely: the action is quadratic, so the two channels couple bilinearly, and the bilinear form fixes the cross-mode scale. For unit-normalized modes with stiffnesses $\omega_\parallel^2$ and $\omega_\perp^2$ and cross coupling $\lambda$, the determinant of the form is $\omega_\parallel^2\omega_\perp^2-\lambda^2$, so the mixed mode softens, and dominates the response, at the threshold scale
\[
\omega_\times=\sqrt{\omega_\parallel\omega_\perp}.
\]
The geometric mean is the spectral invariant of the bilinear cross term itself. Other exchange-symmetric combinations correspond to no invariant of the quadratic form; an arithmetic-mean scale would require a non-bilinear vertex the action does not contain. The cross-scale amplitude is therefore
\[
a_{\mathrm{eff}}=\sqrt{g_{\mathrm{bar}}\,a_0}.
\]
The bath temperature is set by the transverse vacuum scale, $k_BT_{\rm bath}=\hbar\,\omega_\perp$ in the normalization of Appendix~E, and the occupancy argument is the frequency ratio
\[
x=\frac{\hbar\,\omega_\times}{k_BT_{\rm bath}}=\frac{a_{\mathrm{eff}}}{a_0}=\sqrt{\frac{g_{\mathrm{bar}}}{a_0}},
\]
invariant under common rescaling of the frequencies and the temperature, so no Unruh normalization convention survives in the law.
When $g_{\mathrm{bar}}\gg a_0$ the system reduces to the ordinary Newtonian branch; when $g_{\mathrm{bar}}\ll a_0$ the response crosses into the low-acceleration completion.

The resulting radial-acceleration law is
\[
g_{\mathrm{obs}}
=
g_{\mathrm{bar}}\bigl(1+n_B(x)\bigr)
=
\frac{g_{\mathrm{bar}}}{1-\exp\!\left(-\sqrt{g_{\mathrm{bar}}/a_0}\right)},
\]
with the asymptotic limits
\begin{align}
g_{\mathrm{bar}}\gg a_0 &\Longrightarrow g_{\mathrm{obs}}\approx g_{\mathrm{bar}},\\
g_{\mathrm{bar}}\ll a_0 &\Longrightarrow g_{\mathrm{obs}}\approx \sqrt{a_0g_{\mathrm{bar}}}.
\end{align}
The two endpoint limits follow from the occupancy factor alone; the exact interpolation between them is conditional on kernel neutrality, $F(x)=1$, established structurally below and flagged there as the open gold-standard vertex computation.
The deep-MOND branch gives the baryonic Tully--Fisher law \citep{TullyFisher1977,McGaugh2000}
\[
v^4\approx a_0GM_b.
\]

\paragraph{Division of labor between the action and the state.}
The field equation of the weak-field sector is linear and remains linear here: it carries propagation and the Newtonian channel. The nonlinearity of the galactic law enters through the statistical state of the transverse vacuum, in the same way that Fermi--Dirac occupancy makes electrical conduction nonlinear in temperature while the Maxwell sector stays linear. The coefficient chain fixes the value of $a_0$; the bilinear form fixes the argument $x$; the occupancy enters as the bosonic emission factor
\[
g_{\rm obs}=g_{\rm bar}\bigl(1+n_B(x)\bigr).
\]
The factor is a rate statement, and the distinction carries the physics. Static linear response of a Gaussian theory is temperature-independent --- the retarded propagator does not depend on the state --- so no static susceptibility can bring the occupancy into the force law. The occupancy enters because the weak-field force here is the rate of a capacity-transfer process between the channels; the entropic ontology requires exactly this, since the force is a flux of capacity and fluxes are rates.

The rate factor then follows from first principles in three steps. First, for linear coupling to a bosonic mode the emission rate into occupation $n$ is exact,
\[
\Gamma=\Gamma_0\,\bigl|\langle n{+}1|a^\dagger|n\rangle\bigr|^2=\Gamma_0\,(1+n),
\]
which on the thermal cross mode reads $\Gamma_0(1+n_B(x))$. Second, the absorption channel that would restore detailed balance and cancel the stimulated term against $\Gamma_0\,n_B$ is structurally closed: the reverse process requires the defect to partially refill its capacity deficit, and that deficit is the quantized one-bit increment $\Delta S_f=\ln2$ of Section~5, which admits no partially refilled final state. The thermodynamic arrow of Part~V suppresses the remaining collective reverse fluctuations by the same conditional-typicality weighting that orients macroscopic entropy flow; the one-bit quantization that sets the substrate length thus also permits a low-acceleration excess to survive at all. Third, normalizing the bare rate on the Newtonian channel, $\Gamma_0\leftrightarrow g_{\rm bar}$, gives the law above with its two limits, since $1+n_B(x)=(1-e^{-x})^{-1}$. The identification beneath this derivation---that the static weak-field acceleration is the steady capacity-transfer rate---itself follows from the transport sector, in two parts.

\paragraph{The flux identification from transport.}
The static limit of the telegrapher dynamics of Section~17 is a conservation law: with capacity current $J=-D\nabla\,\delta S$ and a static source, continuity reduces to $\nabla\cdot J=\sigma$, so the field around a mass is a steady conserved capacity current with radial flux density $J(r)\propto M/r^2$. Because $\delta S$ is dimensionless capacity, $J$ is a \emph{bit-number} current, and its static density carries the Green-matched normalization of the weak-field chain with no new constant: the Newtonian field is a rate density, a statement the static transport limit derives rather than assumes. Delivery of this flux to a test defect is an \emph{inter-channel} process---the current resides in the longitudinal channel while the defect's dressing occupies the transverse sector through the $2/3$ export---and inter-channel transfer is generated solely by the bilinear cross term, so it proceeds at the cross-mode frequency, inside the thermally occupied transverse vacuum, at the field point. The stimulation is therefore local: the argument $x$ is evaluated where the transfer is delivered, the feature the wide-binary discriminator of Section~25.1 rests on.

\paragraph{Kernel neutrality.}
What remains to verify is that the golden-rule kernel contributes no dependence on $x$ beyond the occupancy, i.e.\ that $\Gamma=g_{\rm bar}\,F(x)\,(1+n_B(x))$ has $F\equiv1$. Four observations isolate the condition the kernel must meet. The factor $1+n$ is exact ladder algebra, $|\langle n{+}1|b^\dagger|n\rangle|^2$, common to all other kernel factors. A continuum density of states $\rho(\omega)\propto\omega^2$ would deform the law, but the transfer is into the per-cell transverse channel, whose state count per admissibility cell is fixed by the same $4\pi^2$ cell structure that fixes the coefficient of $a_0$; no frequency power enters. An energy-per-quantum conversion $\hbar\omega_\times$ would likewise deform it, but capacity is exchanged in the fixed one-bit unit $\Delta S_f=\ln2$, the same quantization that closes the absorption channel, so the transfer is number-counted. The number flux itself is the conserved bit current above, whose static density is $g_{\rm bar}$. Together these reduce the kernel to a single requirement: that the coarse inter-channel vertex couple to the admissibility-weighted ensemble \emph{mean} of the transition availability rather than to its spectral spread. Under that mean-coupling condition $F\equiv1$ and $g_{\rm obs}=g_{\rm bar}(1+n_B(x))$ exactly, and the structures enforcing it---the admissibility-cell channel and the one-bit quantization---are the same two facts that fix $a_0$ and forbid absorption, so the law's shape is tied to the framework's founding bookkeeping rather than to a free choice.

The main text establishes this mean-coupling condition structurally; its verification from the substrate Hamiltonian remains the audit task. A first ensemble-level computation shows why that check is nontrivial: one-step admissible transitions across the $1680$-state ensemble correlate with $K^2$ (Pearson $r\simeq0.5$, class means spanning $1$ to $6$ of $8$ moves), so the raw hopping kernel is not spectrally flat, and neutrality requires the coarse vertex to see only the ensemble mean of this availability---a single number---with the static field unable to re-weight the local ensemble and leak the spectral dependence into the kernel. Establishing that mean-coupling from the graph-level inter-channel vertex of Appendix~H is the remaining gold-standard computation. Empirically, any residual $F$ is bounded below the few-percent level across five decades of $x$ by solar-system precision and the measured flatness of the galactic and lensing relations, and the closure table records the status accordingly.

The law expresses vacuum-state mode occupancy at the apparent-horizon temperature. Departures from it would require either (a) an excitation mechanism beyond the vacuum, supplied by the causal nonequilibrium transport sector of Appendix~E for mergers and transients, or (b) a different transverse reference, which would require modifying the cosmological boundary normalization that fixes $a_H$. Neither route belongs to the closed stationary weak-field branch; both would be treated as extensions or departures from it.

The galactic branch is therefore fixed by the same channel-identification structure already used elsewhere in the weak-field EFT, with the radial-acceleration law read off from one chain rather than interpolated per galaxy.

\section{15. Lensing, PPN, and Weak-Field Consistency}

A modified-gravity account can reproduce galaxy rotation curves and still fail: if the field supplying the extra dynamics does not bend light the same way ordinary mass does, the two metric potentials slip apart and lensing comes out wrong. The weak-field branch here avoids that failure at leading order, and for a structural reason. Because the entanglement sector is scalar, it does not generate anisotropic stress at leading weak-field order. The scalar anisotropic stress is quadratic in gradients, schematically $\partial_iS\,\partial_jS=O(\Phi^2/c^4)$, so it enters beyond the linear branch. Hence
\[
\Phi=\Psi
\]
to the order treated in the present EFT. Light bending and dynamical mass estimates are therefore sourced by the same leading metric response. In effective-halo language, the entanglement response can be rewritten as
\[
\rho_{\mathrm{halo}}(r)=\frac{1}{4\pi Gr^2}\frac{d}{dr}\Big[r^2(g_{\mathrm{obs}}-g_{\mathrm{bar}})\Big],
\]
which yields the familiar $1/r^2$ outer-halo profile in the asymptotic branch.

The same response that governs the dynamics governs light deflection, so galactic support does not come at the cost of a leading weak-field inconsistency. Solar-system constraints such as Cassini require any PPN slip to remain at the $\sim10^{-5}$ level, consistent with the no-slip structure of the linear branch.

The same weak-field structure also returns the GR post-Newtonian values at the order treated. At leading post-Newtonian order,
\[
\gamma_{\mathrm{PPN}}=\beta_{\mathrm{PPN}}=1.
\]
Scalar-induced corrections to the slip $\Phi-\Psi$ enter only at quadratic weak-field order, schematically $O(\Phi^2/c^4)$. Thus the leading weak-field EFT reproduces galactic phenomenology without introducing gravitational slip or solar-system-scale pathologies.

The closed prediction here is the absence of leading gravitational slip, $\Phi=\Psi$ with $\gamma_{\rm PPN}=\beta_{\rm PPN}=1$. The corresponding open test is empirical: stacked weak-lensing measurements of the dynamical-to-lensing mass ratio, together with PPN and fifth-force searches, confront that prediction directly. At leading weak-field order the sector is closed; higher-order precision confrontation is an audit task, not an architectural gap.

\paragraph{The trace coupling is not a fifth force.}
The action's trace term $-\kappa\chi S_{\mathrm{ent}}$ could be read as a Brans--Dicke-like scalar charge that test bodies feel in addition to the metric, which solar-system fifth-force bounds would then constrain severely. That is not the operative reading. The trace channel sources the capacity deficit $\delta S$; the bridge law converts that deficit into the lapse $N=\exp[-\delta S/(2S_\infty)]$; and test bodies follow geodesics of the resulting physical metric. There is no second, independent scalar force beside geodesic motion --- the scalar's entire effect is already carried by the metric the bridge produces. The leading post-Newtonian values follow from that single bridged metric: writing $U=-\Phi>0$ and $N=e^{-U/c^2}$ gives $g_{00}=-N^2=-1+2U/c^2-2U^2/c^4+O(c^{-6})$, so $\beta_{\rm PPN}=1$, while the absence of leading scalar anisotropic stress gives $\gamma_{\rm PPN}=1$. The remaining formal audit is higher-order: if the canonical $T^{(\mathrm{ent})}_{\mu\nu}$ is retained as an independent Einstein-sector source, one should check explicitly that its quadratic pieces reproduce this bridge expansion rather than adding to it. That confrontation is the flagged audit task; the leading calculation returns the general-relativistic values.
\section*{Part IV. Time-Dependent, Transport, and Cosmological Sectors}

\section{16. Why Dynamics Requires Extension Beyond the Static Branch}

The static weak-field branch answers one question: what the capacity-deficit profile looks like once a source has settled. That is enough for ordinary galaxies, which are quasi-static. It is not enough for clusters and mergers, which ask a different question --- how that profile \emph{develops}: how it propagates when a source moves, how it lags behind a fast disturbance, how it saturates, and how it responds when matter separates into distinct dynamical phases. None of those is contained in a static Poisson law.

If the entanglement-capacity medium is physical, it must answer that second question too, with propagation and causal response built in. The time-dependent sector is therefore not an optional add-on but the natural dynamical completion of the medium that produces the static EFT, and the cluster and cosmological sectors draw on it. The next section gives its minimal causal form; Section~17.5 then uses it for clusters and mergers.

\section{17. Causal Transport and Telegrapher Dynamics}

The canonical time-dependent completion is most cleanly written relative to the substrate four-velocity $u^\mu$:
\[
\tau_0(u^\mu\nabla_\mu)^2\delta S
+u^\mu\nabla_\mu\delta S
=
D h^{\mu\nu}\nabla_\mu\nabla_\nu\delta S
+A\chi,
\qquad
h^{\mu\nu}=g^{\mu\nu}+u^\mu u^\nu .
\]
The vector $u^\mu$ is the local rest frame of the entanglement-capacity medium, not an additional ad hoc force carrier. In that frame, $u^\mu=(1,0,0,0)$, the equation reduces to the familiar telegrapher form
\[
\tau_0\partial_t^2\delta S+\partial_t\delta S
=
D\nabla^2\delta S+A\chi(x,t),
\]
with static-matching condition
\[
\frac{A}{D}=\frac{\kappa}{\gamma}.
\]
This equation is introduced because a physical medium should not respond instantaneously to changing sources. The static Poisson equation is appropriate when the source has already settled, but once sources evolve in time one needs both propagation and relaxation. The telegrapher form is the minimal causal extension that still reduces to the static branch when time dependence becomes negligible. The static capacity-strain field is not being replaced; the telegrapher equation describes how that same field propagates and settles when sources change. ``Relaxation'' in this section means genuine time-dependent settling, distinct from the static strain of the weak-field branch.

Causality requires
\[
\frac{D}{\tau_0}=c^2,
\]
so the transport sector propagates disturbances at finite speed. In the canonical no-new-IR-scale branch,
\[
\tau_0^{-1}=H_0,
\qquad
D=\frac{c^2}{H_0}.
\]
This transport equation separates two roles that must remain distinct. Ordinary galactic support belongs to the near-stationary static branch. The telegrapher sector governs how the same medium propagates, relaxes, and develops lag when sources evolve in time. It is therefore not used to generate ordinary static galactic support; it governs transport, lag, relaxation, and merger phenomenology around the near-stationary weak-field branch.

For galactic modes the Appendix~E analysis shows that the long relaxation time does not destroy the static limit. Galactic modes lie deep in the underdamped regime, so the static Poisson branch is recovered as the exact time average relevant to ordinary galactic dynamics. The assumption here is that the source is quasi-static on galactic timescales and supported on wavelengths far shorter than the critical scale $\lambda_c\sim 4\pi c/H_0$; under those conditions the oscillatory transient averages out instead of competing with the static branch.

The static branch still handles the ordinary galactic law. The transport equation describes what happens when the source history is no longer quasi-static: propagation delay, relaxation, and merger-era lag.

The transport branch is closed at the level of $D/\tau_0=c^2$ and the preferred choice $\tau_0^{-1}=H_0$. The cluster and merger phenomenology, which sets the source weights this transport then evolves, is developed in Section~17.5.

\section{17.5 Cluster Source Projection and the Diffuse--Decoupled Channel Split}

Clusters are where the simple galactic radial-acceleration relation stops being enough, and they fail it in two distinct ways. Relaxed clusters show a hook-shaped residual --- near unity in the stellar-dominated center, peaking at intermediate acceleration, converging back in the deep outskirts. Merging clusters show lensing peaks that stay with the collisionless galaxies while the dominant baryonic mass, the X-ray gas, is displaced. A viable account must produce both without granting galaxies hidden baryonic mass and without altering the galaxy law itself.

The mechanism proposed here is that the long-range entropic-excess channel does not couple equally to every baryonic phase. A diffuse, phase-averaged medium couples only through the transverse projection $\epsilon$ already fixed by the galactic sector; a dynamically decoupled collisionless component recovers the full projection; and a virialized coherent bath can be lifted above it. The rest of the section develops that one distinction and confronts it with data, and its status is kept explicit throughout: the projection coefficient is derived, the trend direction and hook shape are supported, the lift amplitude is not yet derived, and the decisive resolved-map test has not yet been performed.

The residual is not a single number but an acceleration-dependent profile: lensing and kinematic analyses of relaxed clusters \citep{Tian2020,Eckert2022,Li2023} find the ratio of observed to galaxy-RAR-predicted acceleration near unity in the stellar-dominated centers, rising to a peak of roughly $3$--$5$ at intermediate accelerations ($g_{\rm bar}\sim10^{-11}$--$10^{-10}\,{\rm m\,s^{-2}}$), and then apparently converging back toward the galaxy relation at the lowest probed accelerations, subject to gas-extrapolation caveats in the outskirts \citep{Li2023}. The older integrated value of $1.5$--$2$ from hydrostatic analyses at higher accelerations is the high-$g_{\rm bar}$ edge of this hook-shaped profile, not its amplitude. Merging clusters sharpen the problem. In systems of the Bullet type the dominant baryonic component---the intracluster plasma---is ram-pressure slowed and displaced from the collisionless galaxies, while the lensing peaks remain with the outgoing collisionless components rather than the X-ray gas. A viable cluster sector must explain both the relaxed normalization and the merger gas/lensing separation without assigning galaxies additional hidden baryonic mass and without altering the galactic mass anchor.

The mechanism is not transport lag and it is not extra galaxy mass; both are excluded. Transport lag is excluded because, in the canonical $\tau_0^{-1}=H_0$ branch, disturbances propagate at $c$ and communicate changes across a cluster-scale configuration on a light-crossing time of order $1\,\mathrm{Myr}$, three orders below the merger timescale, so the field tracks the moving source; an underdamped configuration retains memory of the \emph{pre-merger centroid}, not of the outgoing collisionless component, and a retarded forward solution sourced equally by all baryons therefore does not by itself keep the lensing peak on the outgoing collisionless component. Extra galaxy mass is excluded because the galactic acceleration scale and the gas-dominated-dwarf RAR fix the deficit per unit baryonic mass with no freedom to enhance it. The viable mechanism is instead a phase-dependent \emph{source projection} for the long-range entropic-excess channel, built from a distinction the framework already contains.

\paragraph{The projection coefficient as the galactic reduction factor.}
Section~14 fixes the galactic acceleration scale as the transverse reduction of the horizon thermal scale,
\[
a_0=\frac{g_{\mathrm{share,eff}}}{4\pi^2}\,a_H,
\qquad
a_H=cH_0 ,
\]
where the $(2\pi)^2$ is the two-transverse-mode phase-space cell of the $1+2$ split and $g_{\mathrm{share,eff}}$ is the admissibility-sharing content per cell. The same construction fixes the ratio between the transverse static projection and the full horizon projection,
\[
\epsilon
\equiv
\frac{a_0}{a_H}
=
\frac{g_{\mathrm{share,eff}}}{4\pi^2}
\simeq 0.188 .
\]
This is not a new coefficient. It is the already-fixed reduction factor of the galactic sector, read as a statement about source projection: a source that accesses only the transverse static two-plane couples to the long-range channel at strength $\epsilon$ of a source that accesses the full horizon projection.

\paragraph{The diffuse and decoupled source classes.}
The assignment follows from coherence under coarse-graining. Diffuse matter is a continuum of locally uncorrelated, thermalized source elements; under coarse-graining its off-diagonal source cross-terms average away and only the diagonal transverse static projection survives. It therefore couples at $\epsilon$. This includes shocked or unvirialized intracluster plasma, the warm--hot intergalactic medium, and cold but \emph{diffuse} galactic H\,\textsc{i}\ when the galaxy is treated as a single smooth source---which is why gas-dominated dwarfs and low-surface-brightness galaxies remain on the galaxy RAR. The suppressed projection is thus a coherence effect, not a temperature effect; a hot phase that has \emph{virialized} into a coherent bath is the exception, taken up below.

The unsuppressed projection is accessed by matter that is \emph{not} part of the phase-averaged continuum: a collisionless overdensity that is spatially separated from, and dynamically decoupled from, a surrounding diffuse medium. The criterion is relational, not intrinsic compactness. Compactness alone would misclassify: stars in ordinary galaxies, isolated ellipticals, and globular clusters are compact and bound yet must remain on the galaxy RAR, and they do, because none is a collisionless node decoupled from a distinct diffuse continuum. A galaxy in a cluster is different only because it is embedded in, and decoupled from, the intracluster medium. The suppression is the property of participating in the continuum; matter that has decoupled from the continuum escapes it. So decoupled collisionless matter sits at the baseline weight, and incoherent diffuse gas is suppressed to $\epsilon$, with $\epsilon=g_{\mathrm{share,eff}}/4\pi^2$.

A \emph{virialized} bath is the third state. Once the diffuse atmosphere relaxes into a coherent, extended phase it is no longer a sum of uncorrelated elements, and it opens a collective long-range response that lifts it \emph{above} the baseline, with a bath weight $W_{\rm bath}>1$ that grows with how completely the atmosphere has virialized, gated by $B_{\rm bath}\in[0,1]$ below. This super-baseline lift is a \emph{distinct} premise from the suppression rule above---suppression and its absence span only $[\epsilon,1]$. The lift carries a derived ceiling: the transverse channel holds $4\pi^2/g_{\mathrm{share,eff}}=1/\epsilon$ of capacity per admissibility cell relative to a single incoherent element, so no source weight can exceed $W_{\rm bath}=1/\epsilon\simeq5.32$---the reciprocal of the derived suppression coefficient, and not a new constant. A fully developed coherent bath saturates this bound, occupying every transverse cell of the channel it is otherwise suppressed within. The lift function therefore has fixed endpoints, $W_{\rm bath}=1$ at zero bath and $W_{\rm bath}=1/\epsilon$ at saturation; the minimal linear candidate $W_{\rm bath}=1+(1-\epsilon)B_{\rm bath}$ and the growth profile between the endpoints are evaluated against cluster data below.

The effective entropic-channel source $\chi_{\rm ent}$ is therefore regime-dependent. In a relaxed cluster the gas is a virialized bath,
\[
\boxed{\;
\chi_{\rm rel}=\rho_{\rm dec}+W_{\rm bath}\,\rho_{\rm bath}\;}
\]
with $\rho_{\rm dec}$ the decoupled collisionless substructure and $\rho_{\rm bath}$ the virialized continuum; in a non-equilibrium merger the central gas is shocked and incoherent while only a residual atmosphere stays virialized,
\[
\boxed{\;
\chi_{\rm merge}=\rho_{\rm dec}+\epsilon\,\rho_{\rm shock}+W_{\rm bath}\,\rho_{\rm vir}\;}
\]
which in the Bullet limit, where the displaced gas is shocked and little virialized bath remains on the cores, reduces to $\chi_{\rm Bullet}\simeq\rho_{\rm dec}+\epsilon\,\rho_{\rm shock}$. Ordinary matter continues to gravitate through the usual metric coupling; the projection rule concerns only the long-range entropic-excess channel.

\paragraph{The bath gate as a measured input.}
Decoupling requires something to decouple from. The gate is therefore tied to a measured property of the diffuse atmosphere: the fraction of the halo's cosmic baryon allotment that has become an extended virialized hot phase,
\[
\boxed{\;
B_{\rm bath}
=
\mathrm{clip}_{[0,1]}\!\left[
\frac{M_{\rm hot,vir}(<r_{500})}{f_{b,\rm cos}\,M_{500}}
\right]\;}
\qquad
f_{b,\rm cos}=\frac{\Omega_b}{\Omega_m}\simeq 0.156 ,
\]
with $f_{b,\rm cos}$ fixed by Planck values \citep{Planck2020} and $M_{\rm hot,vir}$, $M_{500}$ read from X-ray/SZ and total-mass estimates. The bookkeeping of the sector is then: $\epsilon$ is fixed and derived, $f_{b,\rm cos}$ is fixed cosmologically, and the per-system inputs ($f_\star$, $f_{\rm gas}$, $M_{\rm hot,vir}$, $M_{500}$) are all measured; the one undetermined object is the coherence-growth profile of the lift, $W_{\rm bath}(B_{\rm bath})\in[1,1/\epsilon]$, whose endpoints are fixed and whose interior form is not yet derived. In merging systems $M_{\rm hot,vir}$ is to be evaluated from the \emph{pre-merger} virialized atmosphere rather than the mid-collision gas state, since the shocked, displaced gas is not a relaxed continuum.

\paragraph{Relaxed-cluster residual.}
For a \emph{relaxed} cluster the diffuse gas is a virialized continuum, and the residual is carried by that continuum, not by the decoupled stellar component. Let
\[
f_{\rm cont}=\frac{M_{\rm hot,vir}}{M_{\rm baryon}}
\]
be the fraction of observed baryons in the virialized diffuse continuum. The residual relative to a galaxy-RAR extrapolation is then
\[
\boxed{\;
\mathcal R_{\rm rel}
=
1+\big(W_{\rm bath}-1\big)\,f_{\rm cont}\;}
\]
and the minimal linear candidate for the lift,
\[
W_{\rm bath}=1+(1-\epsilon)B_{\rm bath},
\qquad
1-\epsilon=1-\frac{g_{\mathrm{share,eff}}}{4\pi^2}\simeq 0.812 ,
\]
uses the amplitude $(1-\epsilon)$, the part of the full horizon channel that the suppressed transverse branch is missing---not a new coefficient. The physical reading is that as a virialized diffuse bath forms, the continuum itself opens a collective cluster response whose amplitude scales with how complete the bath is ($B_{\rm bath}$) and how much of the baryon budget sits in it ($f_{\rm cont}$).

This linear form has the qualitatively correct mass trend: both $B_{\rm bath}$ and $f_{\rm cont}$ \emph{increase} with halo mass---hot-gas fractions rise toward clusters \citep{Vikhlinin2006,Sanderson2003} and the atmosphere becomes more fully virialized---so their product rises monotonically from groups to massive clusters, reproducing the observed direction with no fitted parameter. Its amplitude, however, is excluded. For CLASH-scale clusters the gate gives $B_{\rm bath}\simeq0.83$, $f_{\rm cont}\simeq0.9$, hence $\mathcal R_{\rm rel}\simeq1.6$; but evaluating the source weight required to reproduce the published CLASH relation \citep{Tian2020} across its data-supported acceleration window gives $\mathcal R\simeq3.7$ at $g_{\rm bar}=10^{-10}\,{\rm m\,s^{-2}}$ rising to $\mathcal R\simeq7.8$ at $10^{-11}\,{\rm m\,s^{-2}}$. The linear lift $(1-\epsilon)B_{\rm bath}f_{\rm cont}$ is short of the observed peak residual by a factor of roughly $2.5$--$4$, and no escape through missing baryons (a multiple of the X-ray gas mass would be required), hydrostatic bias (the masses are lensing-based), or sample heterogeneity is available at that magnitude.

\paragraph{The linear candidate is excluded on amplitude.}
Because $B_{\rm bath}\le1$ and $f_{\rm cont}\le1$, the linear candidate bounds the relaxed residual by $\mathcal R_{\rm rel}<1+(1-\epsilon)\simeq1.81$. The measured peak residual of relaxed clusters is $3$--$5$ \citep{Tian2020,Eckert2022}, well above this bound, so the linear lift is excluded. The exclusion is specific to the lift function: the channel-split ontology itself makes a structural prediction about the residual's shape that the data support, taken up next.

\paragraph{The hook morphology.}
Independently of the lift amplitude, the channel split makes a radius-resolved prediction about the \emph{shape} of the cluster residual. In cluster centers the baryons are dominated by the BCG stellar component, which is decoupled and carries weight $1$, so the local residual starts near unity. Moving outward, the virialized gas continuum comes to dominate and the residual rises toward the bath-weighted value. At the lowest accelerations the deep branch compresses any bounded source weight $W$ toward $\sqrt W$ in acceleration terms, pulling the residual back down. The predicted profile is therefore a \emph{hook}: near unity in the stellar-dominated center, peaking where the bath dominates at intermediate acceleration, and converging back toward the galaxy relation in the deep outskirts. This is precisely the morphology class reported by current measurements \citep{Eckert2022,Li2023}---a shape that neither a total-baryon modified-gravity law (which predicts no relaxed-cluster excess at all) nor a constant offset produces.

The framework therefore gets the \emph{kind} of residual right and, with the linear candidate, its \emph{amplitude} wrong: the predicted peak is $\simeq1.5$ for CLASH-like parameters against the observed $3$--$5$. The constraint on the lift function is thus two-ended: $W_{\rm bath}$ must grow faster with bath development than the linear candidate, reaching peak residuals of $3$--$5$ at the developed-bath cluster scale, while remaining small enough at the group scale that X-ray-faint groups stay near the galaxy relation.

At the saturation weight $W_{\rm bath}=1/\epsilon$ the predicted mass residual for developed-bath parameters is $\mathcal R_{\rm sat}=f_{\rm dec}+f_{\rm cont}/\epsilon\simeq4.7$--$4.9$, inside the required window and at the upper edge of the measured peak band, and the radius-resolved profile then peaks at $\simeq3.9$ in acceleration terms for CLASH-like parameters against the observed $3$--$5$. The same weight applied uniformly at the group scale overshoots by a factor of $\sim2$, so saturation must be gated by coherence development: massive relaxed clusters sit at or near the bound, while X-ray-faint groups remain far below it, with the group-scale data requiring $W_{\rm bath}\lesssim1.4$ there.

One caveat accompanies the saturated profile: the predicted central residual depends on the stellar/gas decomposition and on excluding multiphase cool-core gas from the coherent bath. The deep-outskirt question---power-law continuation \citep{Tian2020} versus convergence toward the galaxy relation \citep{Li2023}---is adjudicated directly below.

\paragraph{The ceiling against cluster data.}
Inverting the X-COP hydrostatic measurements \citep{Eckert2019} into source-weight terms gives twenty-four direct tests of the bound: for each of the twelve clusters, at $R_{500}$ and at $R_{200}$, the weight $S$ solving $\nu(S\,g_{\rm bar}/a_0)\,S\,g_{\rm bar}=g_{\rm obs}$ with non-thermal-corrected masses, where $\nu(x)=[1-e^{-\sqrt{x}}]^{-1}$ is the response factor of the radial-acceleration law (Section~14). Every point respects $S\le1/\epsilon$; the sample maximum is $S=3.96$, at $R_{500}$ of the merging cluster A2319, the most model-dependent point, and the relaxed systems span $S\simeq1.8$--$2.7$ at $R_{500}$ and $1.4$--$2.2$ at $R_{200}$.

The same inversion adjudicates the deep end: the power-law continuation requires $S\simeq6$--$8$ at $g_{\rm bar}\simeq1$--$2\times10^{-11}\,{\rm m\,s^{-2}}$, precisely the accelerations of the $R_{500}$--$R_{200}$ points, which sit at $S\simeq1.5$--$2.7$; within this sample the deep end converges rather than continuing, and the strongest published challenge to the bound is not borne out.

The methodological caveat is that the continuation was fit to lensing-based masses of higher-redshift systems while the inversion here uses local hydrostatic masses; breaching the bound at $R_{500}$ would require the non-thermal-corrected masses to be low by a factor of $\simeq2.3$, well beyond any claimed hydrostatic bias. The measured radial run of the source weight---$S\simeq3.7$--$4.9$ in the hook-peak window, where the saturation band is reached, declining to $\simeq2.2$ at $R_{500}$ and $\simeq1.5$ at $R_{200}$---is the quantitative target the coherence-growth profile must reproduce.

Direct and fluctuation-based turbulence measurements find low non-thermal support in relaxed systems at all probed radii \citep{Dupourque2023,XRISM2025}, so the decoherence agent gating the lift cannot be the cluster-to-cluster turbulence level: it must grow with radius even in fully relaxed atmospheres. The coherence-growth profile between the fixed endpoints, so constrained, is the object the channel-selection theorem must deliver, with the group end requiring $W_{\rm bath}\lesssim1.4$.

\paragraph{The decoupled-fraction form is excluded by the mass trend.}
A second candidate class lets the residual ride on the decoupled stellar fraction, $\mathcal R=1+4.32\,B_{\rm bath}f_{\rm dec}$ with $f_{\rm dec}=f_\star/(f_\star+f_{\rm gas})$. This class is excluded by the mass trend rather than the amplitude. Because $f_{\rm dec}$ \emph{falls} with mass while $B_{\rm bath}$ rises, their product peaks at the group/poor-cluster scale and \emph{declines} toward massive clusters, predicting the largest anomalies at intermediate mass; the observed residual instead rises monotonically from groups to massive clusters, where it is most firmly established. A decoupled-fraction form passes only at the intermediate scale where its spurious hump crosses the observed band. The residual therefore rides on the continuum, not the decoupled fraction---the trend direction settles which component carries it, even while the lift amplitude remains underived. The broader lesson is that the relaxed-cluster amplitude and the Bullet morphology are distinct observables and must not be collapsed into a single scalar law.

\paragraph{Bullet-type mergers.}
The relaxed residual above and the Bullet morphology are governed by the same fixed coefficient $\epsilon$ but by \emph{different source expressions}, because the gas is in a different physical state. In a relaxed cluster the gas is a virialized continuum and drives the residual through the bath-lift term $(W_{\rm bath}-1)f_{\rm cont}$. In a Bullet-type merger the central gas is shocked, displaced, and out of equilibrium---not a relaxed bath---so it has small effective $B_{\rm bath}$ and couples to the long-range channel at the suppressed weight $\epsilon$, while the outgoing collisionless cores remain decoupled. The two regimes share one coefficient; they are not one scalar law, and we do not present them as such.

In the merger, then, the decoupled galaxies and subcluster cores retain the unsuppressed projection and the shocked diffuse gas couples at $\epsilon$. With a gas/galaxy baryon ratio near $5.7$, the gas contributes $\epsilon\times5.7\simeq1.07$ in the entropic channel against the galaxy contribution of $1.0$: the projection brings the two components to near-parity, removing the factor $\sim\!5.7$ by which the gas would otherwise dominate, but it does not by itself invert them. The inversion is completed by projected compactness. For two roughly symmetric outgoing components the ratio of one edge peak to the central gas contribution scales as
\[
\frac{\Sigma_{\rm edge}}{\Sigma_{\rm gas}}
\sim
\frac{f_{\rm dec}}{2\epsilon f_{\rm gas}}\,
\frac{A_{\rm gas}}{A_{\rm edge}},
\qquad
\frac{f_{\rm dec}}{2\epsilon f_{\rm gas}}\simeq0.47 ,
\]
so an edge peak dominates the projected map once the shocked gas is spread over more than about twice the projected area of a compact outgoing core---a condition the observed morphology satisfies by a wide margin. The projection rule and this geometry therefore produce the observed gas/lensing inversion---by projection and geometry together, not by projection alone and not by transport lag---as a spatial surface-density prediction rather than an integrated-mass argument. This is consistency, not yet a test of the coefficient. In standard flexible lens reconstructions the gas weight is degenerate with free halo and substructure components, so a model that fits comparably well with or without the fixed X-ray gas map constrains $\epsilon$ only weakly; Bullet-type mergers are thus consistent with the projection rule but do not yet measure $\epsilon$. Peak location alone is in any case insensitive to the coefficient, since sufficiently broad shocked gas yields clump-centered peaks across a wide range of gas weights; the coefficient is tested only by the resolved amplitude fit below.

\paragraph{Relation to the transport sector.}
The telegrapher sector of Section~17 is not the cluster mechanism; it governs how the field propagates and relaxes once the source weights are set. The source-projection rule supplies the static weights $\chi_{\rm ent}$; the transport sector then evolves them. This division avoids the failure mode of a transport-only account, in which a field sourced equally by all baryons cannot hold a lensing peak on the outgoing collisionless component. The full merger observable is obtained by evolving $\chi_{\rm ent}(x,t)=\rho_{\rm dec}(x,t)+[\,\cdots]$ through the causal equation with the observed geometry as input.

\paragraph{Falsifiers and open status.}
The rule makes sharp predictions beyond the relaxed normalization. \emph{(i)}~The linear candidate's ceiling $\mathcal R_{\rm rel}<1.81$ lies below the measured peak residual of relaxed clusters \citep{Tian2020,Eckert2022}, which excludes that form of the lift. The shape prediction of the channel split is the hook: residual near unity in BCG-dominated centers, a single peak at intermediate acceleration where the virialized bath dominates, and convergence back toward the galaxy relation in the deep outskirts. A relaxed cluster with a residual profile that is monotonic in acceleration, or that peaks in the stellar-dominated center, would falsify the channel split itself. The capacity bound supplies the sector's hard ceiling: $\mathcal R_{\rm rel}\le f_{\rm dec}+f_{\rm cont}/\epsilon\simeq5$ for developed-bath parameters, with developed baths reaching it in the hook-peak window and all twenty-four X-COP source-weight inversions respecting it; a robust relaxed-cluster mass residual well above this bound would falsify the capacity bound, the sector's central commitment.

\emph{(ii)}~At fixed mass, X-ray-bright bath-developed systems (larger $B_{\rm bath}$, larger $f_{\rm cont}$) should deviate more from the galaxy RAR than X-ray-faint systems; the residual turns on with the developed diffuse atmosphere, not with mass alone, so two systems of equal mass but different bath development should separate.

\emph{(iii)}~The decisive test is the resolved lensing map. Because the relaxed and merger regimes use different source expressions, the convergence is a \emph{three-component} channel-weighted map---a relaxed virialized-bath component, a shocked non-equilibrium continuum at the suppressed weight, and a decoupled collisionless component,
\[
\boxed{\;
\kappa_{\rm obs}(x,y)
=
A\Big[
W_{\rm bath}\,\Sigma_{\rm bath}(x,y)
+\epsilon\,\Sigma_{\rm shock}(x,y)
+\Sigma_{\rm dec}(x,y)
\Big]+b\;}
\]
with $\epsilon$ fixed at $g_{\mathrm{share,eff}}/4\pi^2\simeq0.188$ or fit freely, and $W_{\rm bath}$ fit per system within $[1,1/\epsilon]$, the saturated value $1/\epsilon\simeq5.3$ predicted for developed baths. Recovering both $\epsilon\simeq0.19$ on the shocked channel and saturation on the bath channel would be a two-coefficient parameter-free success. In the relaxed limit this reduces to the residual law $\mathcal R_{\rm rel}=1+(W_{\rm bath}-1)f_{\rm cont}$; in the Bullet limit the central gas is shocked ($\Sigma_{\rm shock}$ at weight $\epsilon$) and the peaks fall on $\Sigma_{\rm dec}$. A best-fit $\epsilon$ near $0.19$ across Bullet, MACS~J0025, the Musket Ball, a relaxed CLASH sample \citep{Tian2020}, and X-ray-bright versus X-ray-faint groups would be a parameter-free success that neither $\Lambda$CDM nor total-baryon modified gravity predicts, while a best fit near $1$ or near $0$ would falsify the sector. This map fit is the principal outstanding empirical task and has not yet been performed. It is decisive only if it removes the escape hatch that makes existing reconstructions non-discriminating: $\epsilon$ must be floated against the channel-weighted \emph{baryonic} maps ($\Sigma_{\rm bath}$, $\Sigma_{\rm shock}$, $\Sigma_{\rm dec}$) with free dark haloes disallowed, since free haloes absorb the gas weight and leave $\epsilon$ unconstrained---which is why current Bullet reconstructions neither recover nor exclude $\epsilon\simeq0.19$. The three-component decomposition makes the fit more demanding than a two-component one, but also more distinctive.

Three open items remain at the theory level. First, the relaxed residual and the Bullet morphology are governed by one fixed coefficient $\epsilon$ but by two regime-specific source expressions (virialized bath versus shocked non-equilibrium gas). This is justified by the differing physical state of the gas, but it is a weaker unification than a single scalar law, and the boundary between the regimes---how $B_{\rm bath}$ interpolates as a shocked atmosphere re-virializes---is not yet derived.

Second, the channel-selection rule rests on two distinct physical premises, not one. The first is a \emph{suppression} premise: participation in a phase-averaged diffuse continuum restricts a source to the transverse projection (weight $\epsilon$), so that a collisionless overdensity decoupled from that continuum recovers the full horizon projection (weight $1$). The second is a \emph{collective-lift} premise: once a diffuse continuum virializes into a coherent extended phase it opens a collective long-range response that raises it \emph{above} the decoupled baseline, to a weight $W_{\rm bath}>1$. These two are logically independent---escaping the suppression returns a source to weight $1$ but does not by itself carry it past unity---and it is the collective-lift premise, not the suppression premise, that carries the relaxed-cluster residual. The coefficient $\epsilon$ is derived; both premises are plausible within the ontology but neither is yet derived from the microscopic source map, and supplying those derivations is the channel-selection theorem. The collective-lift premise is the weaker of the two: the coherence/phase-averaging argument that motivates the suppression does not on its own generate super-baseline coupling, and the linear form of the lift is quantitatively excluded by the measured peak residual (above). The channel-selection theorem must therefore deliver the coherence-growth profile between the fixed endpoints, reproducing the measured radial decline of the source weight under uniformly low turbulence, constrained at the group end to $W_{\rm bath}\lesssim1.4$ and saturating at $1/\epsilon$ in the hook-peak window.

Third, the identification of the decoupled component with stellar/galaxy mass and of the continuum with the gas is a coarse split; intracluster light and tidally stripped stars blur it at a level the resolved map fit would expose. Abell~520, whose reported gas-coincident dark core is itself disputed, is a phase-state stress case rather than a clean test: a re-cohering or quasi-bound central component would raise its effective $B_{\rm bath}$ and return lensing toward the gas, while a confirmed \emph{young} merger with a robust gas-centered, galaxy-free lensing peak---where the re-coherence escape is unavailable on time grounds---would be a serious challenge.

This sector is therefore presented as a structured, falsifiable proposal, not as closed physics: the projection coefficient $\epsilon$ is derived; the bath gate is a measured input; the trend direction (residual riding on the continuum, rising from groups to clusters) and the hook morphology of the residual profile are supported by current data; the linear lift candidate is excluded on amplitude; the saturation amplitude is fixed by the capacity bound at the reciprocal of the derived projection coefficient, and the saturated residual lands inside the measured peak band; the coherence-growth profile, the resolved-map test, and the channel-selection theorem are the named open work.

\section{18. Cosmology and the Hubble-Tension Sector}

The scalar capacity field has two cosmological roles, and the sector works only if they stay separate. Its homogeneous mode $\overline S(t)$ affects the background expansion and the sound horizon; its inhomogeneous fluctuations $s(x,t)$ still govern local weak-field gravity. The cosmological sector is the homogeneous continuation of the same medium, not an unrelated dark-energy component appended to the weak-field theory: what changes is the kinematic regime, not the ontology, as the background mode becomes dynamically relevant on horizon scales while the local branch stays encoded in the fluctuations.

The cosmological sector uses the same field split,
\[
S(x,t)=\overline S(t)+s(x,t),
\]
where $\overline S(t)$ is the homogeneous mode and $s(x,t)$ the inhomogeneous sector responsible for local weak-field dynamics. The vacuum baseline is fixed by apparent-horizon capacity,
\[
S_\infty(t)=\pi\frac{R_A(t)^2}{L_*^2}.
\]

This is the horizon-normalized representation of the same entropy field used locally. It is compatible with the cell-normalized source theorem because local observables depend on $\delta S/S_\infty$ and $\kappa/(\gamma S_\infty)$ rather than on an absolute entropy unit.

Because the entanglement field couples to the trace of the stress-energy tensor, the homogeneous mode is largely dormant during radiation domination but becomes active near matter--radiation equality. This gives a transient early-energy component of the same general type used in early-dark-energy resolutions of the Hubble tension. In the closed cosmological branch treated here, the effect reduces the sound horizon and shifts the CMB-inferred Hubble constant upward from the high-$67$ range toward the high-$68$ to low-$69$ range.

The direction of the shift matters here, and so does the timing. A successful Hubble-tension mechanism must turn on near the right epoch, alter the sound horizon in the right direction, and then decouple cleanly enough from the local weak-field sector that the galactic branch is not spoiled. The entanglement medium has exactly that qualitative structure.

The local weak-field predictions are protected by the separation between $\overline S(t)$ and $s(x,t)$. This is the role of the shear-lock logic: changing the homogeneous background mode does not rewrite the local static Poisson branch that governs galactic dynamics and lensing.

The claim is therefore a mechanism with the right direction, timing, and qualitative separation of scales, not a finished precision cosmology package. What is shown is that the trace-coupled homogeneous mode turns on in the relevant epoch and pushes the sound horizon in the required direction; what remains open is the full perturbation propagation and likelihood-level confrontation. The homogeneous mode modifies the background history; the inhomogeneous branch continues to govern the local weak-field observables already fixed earlier in the derivation. That separation allows the cosmological extension to remain part of the same scalar medium rather than a re-tuning of the galactic sector.

This is a structurally supported and directionally successful extension, but it is not yet Boltzmann-closed.

\section*{Part V. Nonlinear, Interpretive, and Completion Sectors}

\section{18.5 The Saturated Phase and the Cosmic Microwave Background}

In galaxies the substrate has slack: a source's demand for capacity is met by a developing local response, and the radial-acceleration law of Section~14 follows. The early universe has no slack. The capacity bath is saturated across the linear perturbation regime: wherever a perturbation-scale source demands a response, the available transfer channel has already reached its ceiling, so there is no room for a developing response field. What gravitates is then not a field relaxing toward a source but a conserved count of \emph{committed} capacity channels: cells locked into transferring one bit of capacity per tick on behalf of a defect. That committed count is carried along with the expansion and dilutes as the volume grows, so it redshifts as $a^{-3}$ and gravitates as pressureless dust. This saturated phase plays, in the microwave-background sector, the role conventionally assigned to cold dark matter --- not a new particle species but a phase of the same medium.

\paragraph{Saturation is forced at recombination.}
At $z\simeq1100$ the accelerations of linear density perturbations satisfy $x=g_{\rm pert}/a_0(z)\sim10^{-3}$, where the response demand $\nu(x)\sim30$ far exceeds the capacity ceiling $1/\epsilon=5.32$. The bound derived for the cluster bath (Section~17.5) is therefore mandatory in the early universe: the response channel for the cosmological perturbation bath is saturated. The background itself sits permanently at the ceiling acceleration, since $a_0(z)=\epsilon\,cH(z)$ implies the invariant $cH(z)/a_0(z)=1/\epsilon$ at every redshift.

\paragraph{The pinned reading and the dust theorem.}
The \emph{pinned reading} is the statement that the saturated bath holds exactly at the ceiling and stays there. The transfer law quantizes channel transport at one bit per tick as an upper bound; saturation is the attainment of that bound, so on the pinned reading every committed channel advances by exactly one unit per tick, and while the phase remains pinned no channel can decommit, since there is no slack for it to relax into. The committed count is then conserved and dilutes only with the expanding volume. The coarse-grained field $\phi$ of the committed phase then carries a constraint rather than a generic kinetic term: $X=-\tfrac12 g^{\mu\nu}\partial_\mu\phi\,\partial_\nu\phi=\tfrac12$, with $\partial_\mu\phi$ a unit timelike covector. The leading action consistent with the constraint is
\[
S_{\rm comm}=\int d^4x\,\sqrt{-g}\;\lambda\left(X-\tfrac12\right),
\]
whose variation gives $\nabla_\mu(\lambda\,\partial^\mu\phi)=0$ and, on the constraint surface,
\[
T_{\mu\nu}=\rho\,u_\mu u_\nu,\qquad \rho=E_0\,n,\qquad p=0,\qquad c_s^2=0,\qquad \rho\propto a^{-3},
\]
with $u_\mu=\partial_\mu\phi$ and $n$ the conserved spatial density of committed cells. This is the constrained-scalar (mimetic) class \citep{ChamseddineMukhanov2013}: conditional on the pinned reading, the committed phase is pressureless dust, and the perturbation carrier is the conserved committed-capacity density itself. The gravitating component is therefore a phase of the medium, not an added species; throughout, ``dark matter'' names the conventional attribution this phase replaces. The known caveats of the class --- caustic formation and the need for a small-scale completion --- are inherited and recorded; they parallel the standard cold-dark-matter small-scale account. The continuous dynamical passage itself --- how the generic unconstrained kinetic term of the weak-field branch is driven onto the constraint surface $X=\tfrac12$ as the channel saturates, rather than the constrained action being posited at the ceiling --- is asserted here rather than derived; supplying it, together with the commit/decommit energy accounting recorded below, is an open task within this conditional sector. The constrained route also parallels the mechanism by which the one relativistic completion of Milgromian dynamics known to meet this test does so \citep{SkordisZlosnik2021}.

\paragraph{Uniqueness of the carrier.}
The constraint coupling is the unique survivor of the dynamics classes examined (Appendix~M). Relaxational response kernels are excluded twice over: a growth clock calibrated on the cluster radial decline misses cosmological development by a factor of order thirty, and the precision of the measured acoustic peaks bounds any oscillatory leakage of a lagged response below one part in five hundred, a rejection no causal filter achieves over the few oscillation periods available before recombination. Bound-type rail readings sit at the sound-speed pole and carry no perturbation; plateau approaches have $c_s^2=-1/(2n{+}1)<0$ and are gradient-unstable; the released branches have $c_s^2=\tfrac12$ and $1$ and free-stream. A general no-go for relaxation-plus-cap carriers is recorded in Appendix~M; the conserved committed density is precisely the additional integrating variable that the no-go requires and the constraint reading supplies.

\paragraph{Abundance.}
At full saturation the committed weight is the ceiling, giving
\[
\frac{\Omega_c}{\Omega_b}=\frac{1}{\epsilon}=\frac{4\pi^2}{g_{\mathrm{share,eff}}}=5.321
\qquad\text{against the measured } 5.364\pm0.065,
\]
a $0.7\sigma$ agreement with the abundance counted rather than fitted \citep{Planck2020}; equivalently $\Omega_m/\Omega_b=1+1/\epsilon=6.32$. A standard Einstein--Boltzmann computation \citep{Lewis2000} with the cold component tied to this value and the acoustic scale $\theta_*$ held fixed reproduces the Planck best-fit temperature spectrum with a root-mean-square residual of $0.15\%$ over $\ell=2$--$2500$ and $0.08\%$ in the third-peak region. At the likelihood level, the tied model carries one fewer free parameter than $\Lambda$CDM and sits at $\Delta\chi^2=+2.15$ against the six-parameter best fit on the compressed Planck TT,TE,EE likelihood \citep{PrinceDunkley2019} --- within noise of the full fit, and an upper bound on the constrained optimum since the remaining parameters were not re-optimized. The comparison must hold $\theta_*$ fixed: tying the densities at fixed $H_0$ instead breaks the acoustic scale and misreports the residual.

That the committed density attains the ceiling exactly, rather than a development-dependent fraction of it, follows from the recruitment structure of commitment. Commitment is source-driven: relaxation serves the demand of specific defects, and cells receiving no demand commit nothing. Each source's allocation is capped at $1/\epsilon$ per unit source mass, the same per-source bound that governs the cluster bath (Section~17.5). At recombination the demand side is fixed by the arithmetic above: perturbation-scale sources demand $\nu\sim30$, far above the cap, while the homogeneous background sits at the invariant acceleration $cH/a_0=1/\epsilon$, where the demand $\nu(1/\epsilon)=1.11$ lies below it. With supply abundant and commitment irreversible in the pinned regime, every source therefore recruits to its cap and the background recruits nothing. Allocations add, because committed units are per-defect: a unit serving two defects would make their joint entanglement budget, and hence their joint mass, sub-additive, contradicting the additivity of separated masses under Postulate~II (Section~3.2). The committed density is then
\[
\rho_c=\frac{1}{\epsilon}\,\rho_b
\]
pointwise at the commitment epoch, which delivers the abundance above and hands the perturbations adiabatic, baryon-tracking initial conditions, $\delta_c=\delta_b$ --- the conditions the Einstein--Boltzmann computation assumed. The derivation is conditional on the same pinned reading as the dust theorem and on the per-defect bookkeeping just invoked; the energy accounting of the committed budget remains the sector's open completion, and conservation of the perturbations is taken up directly below.

\paragraph{Conditional conservation: release at the caustic.}
The committed component must behave as conserved dust while the linear bath remains pinned, yet the galactic branch (Section~14) requires collapsed systems to follow the baryonic response law with no surviving collisionless halo, so the completion is a conditional-conservation law: $\nabla_\mu J^\mu_{\rm commit}=0$ in the pinned regime, with decommitment upon release. The constrained phase fixes the release mechanism itself. Because the committed flow is the gradient of a single clock field, $u_\mu=\partial_\mu\phi$, it is exactly irrotational and single-valued, and a multi-stream or vortical velocity field admits no such potential; gravitational collapse therefore terminates the phase at first shell-crossing. The caustic --- the known breakdown locus of the constrained-scalar class --- is here the decommitment event, converting committed capacity into the local response branch. Under the synchrony definition of commitment, one counter per cell ticking with the local phase, decommitment at first stream-crossing is forced rather than chosen: a single fixed cell cannot remain synchronized with two distinct phase branches, so conversion is event-like at the first caustic, a result holding at the same conditional grade as the pinned reading itself. Re-entry is forbidden on entropic grounds: recommitment of a virialized region would require shedding vorticity and re-synchronizing with the advanced global clock, a spontaneous recoherence, so decommitment is irreversible, with the door's orientation resting on the same history weighting that fixes the arrow of time (Section~21). The resulting partition is the one the data require: the linear cosmological field never shell-crosses and remains committed dust, while collapsed systems have shell-crossed and retain none. The alternative, commitment tracking the instantaneous local acceleration, is excluded directly by rotation-curve data: below $g_c$ the binned radial-acceleration residuals lie within $0.05$ dex of the response law \citep{McGaugh2016}, against the $+0.40$ dex excess that locally regrown dust at the cosmic ratio would produce. One completion remains open and is recorded, the energy bookkeeping of conversion at the caustic, together with the consequence that late nonlinear structure growth proceeds on the response branch. The sharp geometric prediction is registered as a falsifier: infalling material between turnaround and first shell-crossing is single-stream and still committed, so clusters should carry a surviving committed component on their infall streams, terminating at the splashback surface \citep{More2015}, with an excess gravitating rim in that shell, a sharp edge at the outermost caustic, and none inside it. The recruitment count fixes the rim's amplitude as well as its geometry: along a stream of which a fraction $f_{\rm rel}$ has already shell-crossed, the surviving committed density is $(1-f_{\rm rel})\,\rho_b/\epsilon$, so the prediction specifies location, edge, and magnitude together.

\paragraph{Domain structure.}
The release threshold $\nu(x_c)=1/\epsilon$ gives $x_c=0.0433$, an acceleration $g_c=x_c\,a_0\simeq5.2\times10^{-12}\,\mathrm{m\,s^{-2}}$ today. The linear cosmological field remains below threshold at all epochs ($x\lesssim2\times10^{-3}$ today, smaller as $\sqrt{a}$ earlier), so linear-regime observables --- baryon acoustic oscillations, the linear power spectrum, microwave-background lensing --- coincide with the $\Lambda$CDM form at the abundance above; in particular the framework inherits, and does not relieve, the $S_8$ tension. Departures are confined to released, collapsed structures, which is the galactic and cluster phenomenology of Sections~14 and~17.5.

\paragraph{Distinction from horizon saturation.}
The committed cosmological phase --- channels advancing at capacity, a running clock --- is distinct from the terminal saturation of the strong-field sector (Section~20), where capacity is exhausted and transactions cease. Whether the two termini are stages of one process is open and carries a stated consistency burden: the gravitating energy of committed capacity must coincide with the mass already accounted at infinity, with no double counting. This is recorded as an open question shared by the strong-field and cosmological sectors.

\section{19. Why These Sectors Belong}

The most directly constrained chain --- microstructure to static weak-field observables --- is now in hand, and the sectors on either side of it form one program rather than a list of add-ons. Each asks what the same finite-capacity substrate does once it leaves the static weak-field regime.

Transport asks how the capacity-strain field propagates and settles in time once a source moves (Section~17). Cosmology asks how the homogeneous capacity mode behaves on horizon scales (Section~18). The saturated phase asks what happens when the bath has no slack left and the carrier becomes a conserved committed density (Section~18.5). The strong-field branch asks the opposite-extreme question, what happens where local capacity falls to zero (Section~20). Many-Pasts asks what history and probability structure the substrate carries, and supplies the memoryless dressing the scale-setting already used (Section~21). The microstructure Hamiltonian asks what dynamics underlies the whole construction (Section~22).

These sectors share the static weak-field ontology but not its evidential status. The static chain is the closed result; the others are conditional extensions, frontier completions, or structural realizations, marked where each appears.

\section{20. Strong-Field Branch and Bounded Occupancy}

The weak-field theory used small capacity deficits. Black holes are the opposite regime: they ask what happens when the deficit is no longer small, and ultimately where the medium runs out of available capacity altogether. The branch therefore begins with the same variable as the weak-field bridge, written in the form that makes boundedness explicit:
\[
q(x)=\frac{S_{\mathrm{ent}}(x)}{S_\infty}\in[0,1].
\]
The physical meaning is direct: $q=1$ is undepleted vacuum capacity, $0<q<1$ is a partially depleted region, and $q=0$ is complete local exhaustion of continuum-defining capacity. The strong-field question is therefore not how to continue the weak-field scalar after it becomes large, but how to write the continuum theory on the domain where the capacity variable remains physically allowed.

The local lapse associated with the static asymptotic time must be a function of surviving capacity. If
\[
N^2=f(q),
\]
then vacuum normalization gives $f(1)=1$, horizon normalization gives $f(0)=0$, the weak-field bridge gives $f'(1)=1$, and serial composition of independent capacity-reduction layers gives
\[
f(q_1q_2)=f(q_1)f(q_2).
\]
The continuous solutions are $f(q)=q^\alpha$, and weak-field matching fixes $\alpha=1$. Thus
\[
N^2=q
\]
is the unique continuous multiplicative completion compatible with the weak-field branch. Once the capacity variable is adopted, the static lapse rule is fixed.

The constrained-capacity theory is then posed on the physical domain
\[
\mathcal{M}_q=\{x\mid q(x)>0\},
\]
with a Lagrange multiplier enforcing $N^2=q$. In bulk vacuum, away from the $q=0$ boundary and with no explicit matter source, the multiplier vanishes and the equations reduce to the ordinary vacuum Einstein equations on $\mathcal{M}_q$. The difference from GR is not in the exterior vacuum equations; it is in the physical domain on which those equations are allowed to live and in the boundary condition at capacity exhaustion.

Consequently the static spherical asymptotically flat vacuum branch is the Schwarzschild exterior,
\[
q(r)=N^2(r)=1-\frac{2GM}{c^2r},
\qquad
r>r_h=\frac{2GM}{c^2}.
\]
The continuum branch terminates at $q=0$. A real continuation of this static branch to $r<r_h$ would require $q<0$, which is outside the state space of the capacity variable. The classical Schwarzschild interior is therefore not interpreted as a physical continuation of the same entanglement-capacity EFT. It is a formal continuation of the GR manifold beyond the point where the substrate has no remaining local continuum channels.

This preserves the tested exterior physics. The near-horizon Euclidean regularity argument gives the usual Hawking temperature, and exterior perturbation theory gives the usual absorption and ringdown: horizon regularity at $q=0$ selects the ingoing mode at leading order, so the leading reflectivity vanishes. Echoes are correction-level rather than generic; they arise only if microscopic boundary channels carry finite UV reflectivity or if the effective reflection surface is displaced outward to a stretched layer $q=\epsilon>0$. The entropy result follows from substrate inputs already used elsewhere: Postulate~II identifies the per-channel cut entropy as $\ln 2$ via fermionic face exclusion, while the transverse bulk graph response fixes the channel density, and their product is the Bekenstein--Hawking $1/4$ as an exact identity. The geometric identification of $q$ with the areal-radius gradient invariant $|\nabla R|^2$ in spherical symmetry further shows that the substrate $q$ equals $1-2GM_{\rm MS}/(c^2R)$ on the exterior, so the $q=0$ saturation surface coincides with the apparent horizon and is inherited from ordinary trapped-surface formation in GR collapse rather than from a separate substrate-transport postulate.

The strong-field branch is therefore closed at the level of universal predictions. Static and stationary asymptotically flat vacuum exteriors reduce to Schwarzschild, Kerr, Reissner--Nordstr\"om, or Kerr--Newman through the standard vacuum Einstein and Einstein--Maxwell branches; the horizon location is fixed as the marginal-capacity surface; the area-law coefficient is closed under Postulate~II and the bulk graph response; and the leading reflectivity vanishes by horizon regularity. What remains is nonuniversal: the microscopic relaxation spectrum, possible stretched-layer corrections, transient response on the boundary, and the explicit graph-side consistency check on substrate transport during dynamical collapse.

\section{21. Many-Pasts: The History-Space Ontology}

Many-Pasts is the most conservative part of the theory to test and the most radical to adopt. In the selected operational branch, it changes nothing a laboratory measures: it reproduces Born-rule statistics and no-signaling exactly. Ontologically it replaces the entire interpretive layer of quantum mechanics --- branching worlds, collapse events, or hidden variables --- with a single object, a weighting over the admissible past histories of the entanglement substrate (Section~3.3).

\paragraph{The operational weight.}
The history weight $P(H|P)\propto e^{-D(H,P)}$ returns ordinary quantum statistics in the projective laboratory limit through the identity
\[
e^{-D(H,P)}=\mathrm{Tr}\!\left(\Pi_P\,\rho_{H\to\mathrm{now}}\right),
\]
the history-space distance reducing to the overlap of the present projector with the state evolved from $H$. Born recovery is a consistency condition on the weighting family, not a fresh derivation of the rule: requiring exact Born statistics selects $\alpha=1$, and forbidding any signaling-sensitive operational bias selects $\beta=0$. In that branch the rule and no-signaling both hold exactly, and nothing a laboratory measures distinguishes the framework from standard quantum mechanics. The conservatism is deliberate: the ontological content carries no operational cost.

\paragraph{Branch realization without many worlds.}
The weight also answers what it means for one outcome to be realized. There is no forward branching into co-real worlds and no collapse event. A definite present is a present with definite macroscopic records, and the histories the weight supports are exactly those compatible with those records; alternative outcomes correspond to alternative present records, not to coexisting branches. Probability is the measure this weighting assigns over the admissible pasts of the one realized present.

\paragraph{Familiar quantum examples.}
In the standard puzzles the weight changes the ontology, not the calculation. In the double-slit experiment the particle is not taken to have secretly travelled one slit while the other path was unreal: before any which-path record exists, the detection event is supported by the admissible histories through the apparatus, including the alternatives whose amplitudes interfere in ordinary quantum mechanics, and making a which-path record conditions that ensemble, removes the cross terms by decoherence, and erases the fringes. In an EPR or Bell experiment the present is a joint record of the pair and the detectors, and the histories are weighted as histories of the whole entangled system, which reproduces the nonclassical correlations while leaving the local Born marginals --- and so the impossibility of signaling --- untouched; the ``spooky'' element is the global conditioning of the history ensemble on a shared record, not a faster-than-light influence. Measurement is read the same way: it adds no collapse law but creates a durable record, and the operative ensemble is then the histories compatible with that record. The probabilities are the usual Born probabilities throughout; Many-Pasts supplies only the history-space reading of why those are the numbers being counted. The worked versions are in Appendix~G.10.

\paragraph{The arrow of time.}
The same weight orients time through conditional typicality. Among the histories consistent with the present macroscopic records, overwhelmingly many show entropy increasing toward the future direction those records define. This is not an added law of laboratory probability; it is a statement about which histories dominate the weight, and it places the substrate's account of temporal asymmetry on the same footing as its account of quantum statistics.

\paragraph{Where the weight is used.}
Many-Pasts enters the gravitational chain through the electron dressing. Its memorylessness in substrate time is one of the two conditions that fix the absolute substrate length: each dressing pass of the electron is an independent draw from the admissibility ensemble, so the seven channels carry the full sharing entropy and add (Sections~13, 22; Appendix~H). The same conditional-typicality weight that orients time also supplies the orientation of decommitment in the saturated phase, forbidding the spontaneous recoherence that re-entry into the committed phase would require (Section~18.5). The same weight therefore appears in four places: the scale-setting branch, the selection of the vacuum state by conditioning on the extended four-dimensional present (Section~23.4), the cosmological release law, and the quantum-foundational layer.

\paragraph{Status.}
Operational closure is exact: the laboratory sector is standard quantum mechanics, and the technical reduction is collected in Appendix~G. The ontology is central rather than secondary, but its distinctive claims --- branch realization and the arrow of time --- are structural and cosmological, not new laboratory predictions, and they should be read at that grade.

\section{22. Microstructure Hamiltonian and Underlying Dynamics}

The UV closure chain has a microscopic dynamical realization with two sides: the emergence of the continuum geometry from the substrate, and the defect dynamics that fixes the scale-setting principle of the weak-field chain. The first is the condensate-side compatibility check; the second closes the one microscopic principle the static branch was conditional on. Section~13 stated the electron anchor in physical terms; this section gives the microscopic dressing dynamics behind it, and Appendix~H carries the technical derivation.

On the geometry side, the realization is a GFT/condensate picture \citep{Oriti2016,Gielen2013} in which spacetime emerges from a condensate of discrete tetrahedral building blocks, while what is macroscopically read as matter appears as fermionic defects of that same substrate. In Madelung form,
\[
\sigma(x)=\sqrt{n(x)}e^{i\theta(x)},
\]
the condensate hydrodynamics generically generate a positive scalar stiffness for the logarithmic-density variable, providing the condensate-side origin of the EFT kinetic term. This is a structural compatibility check, not a replacement for the explicit coefficient closure in Appendix~C; it shows that the EFT kinetic term has a natural microscopic origin.

On the defect side, the same substrate supplies the dynamics behind the faithful sector-resolution principle that sets the absolute cell length $L_*$. The elementary fermionic defect is governed by a dressing Hamiltonian whose single-channel terms make the lightest charged defect bind all seven face sectors once --- the one-pass ground state that supplies the exponent $7$ --- and whose inter-channel term controls whether the seven-channel dressing cloud factorizes. Factorization turns the seven equal sector contributions into the additive $7g_{\mathrm{share,eff}}$, and needs two things: each channel must carry the full admissibility entropy, and the seven must be mutually independent. The first is memorylessness in substrate time --- the Many-Pasts content of Postulate~III, with the substrate relaxing to a fresh draw long before the electron is read out. The second follows from the electron's own lightness: inter-channel correlation would contract the defect and raise its mass, so the lightest charged defect carries none. Each channel at full entropy and the channels independent give $\dim_{\rm eff}=e^{7g_{\mathrm{share,eff}}}$ and $\lambda_e/L_*=\tfrac23 e^{7g_{\mathrm{share,eff}}}$, worked out in full in Appendix~H.

Reading the support ratio $\lambda_e/L_*$ as a physical Compton length used one further identification: that the defect advances one phase cycle per dressing recurrence. Appendix~H.6 gives that recurrence an operator form and removes the identification as a free assumption. The mass is the exported gap of the memoryless refresh super-operator, $m_ec^2=\tfrac32(\hbar/\tau_*)\,e^{-7g_{\mathrm{share,eff}}}$ (the $\tfrac32$ the transverse export of the length relation), so it is read as a recurrence rate, not a binding-energy sum, and the phase-cycle claim decomposes into the recurrence rate, a unit committed charge per pass from fermionic one-pass occupation, and the standard condensate phase relation, none with a free constant. What is not forced is the identification of the winding phase with the condensate $\mathrm{U}(1)$ (Fork~A); Appendix~H.7 derives the three dressing functionals from the mean-field condensate and shows that Fork~A and the rank-one inter-channel coupling reduce to a single premise --- a single scalar field on that condensate --- which is the founding premise of the framework and is tested in the charged-lepton ladder (Appendix~I.1). The hierarchy then rests on one asymmetry. Binding across the seven sectors is additive, but their coherent recurrence is multiplicative, so the continuous maintenance of the defect is large and linear in the commitment while the rest energy, fixed by the rare joint return, is exponentially small. The mass scale is the rare-return gap of the refresh generator, not the bare binding cost: the electron is light not because it is weakly bound but because its completed seven-sector dressing recurs rarely.

Together these two sides do more than check compatibility. The defect dynamics derives the faithful sector-resolution principle from the mass--entropy ontology and the electron anchor, so the substrate length and the induced gravitational scale follow from commitments the theory already makes. What remains microscopically open is twofold: the first-principles derivation of every inhomogeneous continuum coefficient from the full kernel, and --- conditional on Fork~A --- the GFT substrate actually producing the single-scalar mean-field condensate, with its emergent four-geometry, on which the defect spectrum and the geometry side both depend (Appendix~H.7).

The division of labor this section assigns (vacuum stiffness from the condensate branch, matter and its local response from the closure-and-refresh sector) is now supported from below. At cell order the static closure weight, the memoryless refresh, and closure-class history weightings all fall quantitatively short of supplying vacuum curvature stiffness, by more than an order of magnitude, and the conserved capacity field does no better once integrated out, free or budget-bounded (Section~23.3, Appendix~J.4); what the refresh does couple geometry to is the density of closure failure, the theory's matter. The lattice program of Section~23 measures that coupling dynamically, and the local geometric response it finds around pinned closure failures is the first measured instance of the mechanism this section describes. The exclusions also sharpen what the condensate branch owes: the microscopic form of the kinetic term, with the selection of the vacuum state supplied by the conditioning of Postulate~III (Section~23.4).
\section*{Part VI. Closure Status, Falsifiability, and Research Program}

\section{23. The Substrate on a Dynamical Lattice}

The coefficient chain of Part II is exact at the order of a single cell: the ensemble is enumerated, the closure point is solved, and every number that feeds the weak-field sector is evaluated in closed form. The postulates, however, also commit the theory to statements no single-cell calculation can reach, because they concern many cells in interaction: that the vacuum ensemble survives when its cells live on a fluctuating geometry; that a persistent closure failure (the theory's matter) strains the capacity state around it; and that the strain is carried outward far enough to become the $1/r$ deficit of Section~12. This section reports a numerical program built to test those statements directly, by placing the cell ensemble of Part II, unmodified, on a dynamical simplicial geometry and running the coupled system. The section carries the design logic and the results; ensemble definitions, control protocols, and the measured tables are collected in Appendix~J.

\subsection{23.1 The host geometry and the cell identification}

The host is a causal-dynamical-triangulations (CDT) ensemble: four-dimensional triangulations of $S^1\times S^3$ built from pentachora, weighted by the Regge action \citep{Regge1961}, and foliated into spatial slices that are closed simplicial $3$-manifolds of tetrahedra. CDT earns the role, because it is the discrete-gravity ensemble in which the dynamical emergence of an extended, four-dimensional, de~Sitter-like universe is an established result \citep{AmbjornJurkiewiczLoll2004,AmbjornGJL2012}: the machinery hosting the substrate is independently known to sustain the kind of state the theory's conditioning selects (Section~23.4).

The identification is immediate, and it respects a separation the theory requires. Each spatial tetrahedron of a slice hosts one cell's boundary data: its four triangles carry the labels $m=-3,\ldots,3$ of the $j_{\mathrm{eff}}=3$ sector of Section~5, each slice triangle is shared by exactly two cells, and the slice gluing supplies adjacency. Nothing else of the lattice enters. The theory's cell is its label configuration, a dimensionless unit of capacity; no size is ever assigned to it, and the granularity stays in the capacity, as Section~1.2 requires. The lattice simplices are regulator scaffolding, as they are throughout lattice gravity. The identification is nevertheless geometric in the sense that matters: the joint admissibility statistics of the labels depend on how the slice is glued, so the closure weighting can in principle tell one geometry from another, and that channel is the entire coupling between the theory and the host.

Two boundaries of the design are stated at the outset. First, the host supplies what the program needs from a substrate, a dynamical, curvature-carrying simplicial geometry with the theory's cell structure, and nothing more is asked of it. The program does not assume the CDT ensemble is the theory's own vacuum, and it withholds any claim about four-dimensional emergence, which belongs to the host's phase structure and is still being mapped at the couplings used here (Appendix~J.1). Second, the weighting coupled to the host is the paper's own: the admissibility energy of Sections~5--6 at the closed point, with $\eta=\eta_*$ and injectivity enforced. It is not a spin-foam vertex amplitude, and the distinction drawn in Section~27.8 was kept operational: a spin-foam amplitude was coupled first, as a neighboring-theory control, and the audit it forced supplied the control methodology used everywhere below (Appendix~J.3).

\subsection{23.2 The coupled ensemble and its controls}

The coupled system is the joint Gibbs measure
\[
\pi(g,m)\ \propto\ \exp\!\Big(-S_{\rm Regge}(g)-\varepsilon\,(N_{41}-\bar N)^2-\beta\sum_{\text{cells }c}\big[E_c(m)-\mu\big]\Big),
\]
\[
E_c(m)=\eta_*\,K^2(m_c)+\lambda\,n_{\rm coll}(m_c),
\]
where the sum runs over the slice cells, $K^2$ is the closure invariant of Appendix~B, $n_{\rm coll}$ counts label collisions (injectivity is imposed as a penalty whose hard limit is the constraint, with the residual collision fraction reported so the softness stays visible), $\varepsilon$ pins the slice volume, and $\mu$ is the per-cell label free energy, computed by thermodynamic integration so that the closure term cannot masquerade as a shift of the bare cosmological coupling. The labels carry a uniform base measure, so at $\beta=0$ the label entropy cancels exactly and the bare host is recovered identically. The theory point is $\beta=1$ at $\eta_*$: once the chain of Part II is accepted there is no coupling dial left free, and the intermediate $\beta$ values serve only as a diagnostic ramp.

Every run opens by recomputing the single-cell chain on its own tables and refuses to proceed unless $g_{\mathrm{share,eff}}=7.4198$ and $\langle K^2\rangle_{\eta_*}=3/(2\eta_*)=50.223$ are reproduced, so the object coupled to the lattice is verifiably the object counted in Part II. Results are read under a discipline fixed before the runs: matched total volume across arms; a global audit that every quoted ensemble's slices are closed $3$-manifolds correctly threaded by the foliation; and a placebo arm in which the $210$ label-orbit energies are shuffled (same value pool, same permutation symmetry, closure structure destroyed), so that only a real-versus-placebo difference is ever attributed to closure. Every claim in this section is a matched comparison at fixed regulator (real against placebo, coupled against bare, commitment level against commitment level), so regulator-scale artifacts cancel by construction. The intended reading of each possible outcome was recorded before the results existed (Appendix~J.3).

\subsection{23.3 What the closure sector cannot supply: vacuum stiffness}

Three calculations, each exact or numerically converged at cell order, bound what the closure sector can do for the vacuum before any coupled run is read. A transfer-matrix evaluation of the closure weighting's free energy on the rings of cells around a slice edge gives the induced coupling to edge coordination, the discrete curvature variable, as $c_0=+0.019$ per slice edge at the theory point, about one percent of the bare stiffness $k_0\simeq2.2$ that sustains the host's geometric phase, with a label correlation length of about half a lattice step. An exact solution of the refresh dynamics of Postulate~III at cell order (memoryless redraws from the admissibility ensemble, arriving cell by cell) gives the stationary state in closed form: the vacuum runs at admissibility $0.536$, carrying a $46\%$ density of \emph{transient} closure failure at any tick, and that stationary state is blind to the geometry it sits on to exponential accuracy, so the refresh induces no vacuum curvature action either. And the strongest vacuum role a Many-Pasts tilt over closure-class observables could play, a closure-maintaining history weight solved exactly by Doob transform on a ring, saturates at short range with a curvature signal roughly twenty times below the bare coupling even when amplified far beyond the theory point: no strength of closure preference reaches the requirement.

Two further calculations extend the exclusion beyond the label sector, to the conserved capacity field that Section~23.7 introduces as the carrier of the long-range sector. Even a field the labels cannot see might rank geometries on its own, because integrating out a conserved Gaussian field induces a purely geometric action: $\tfrac{1}{2}\ln\det'L(g)$ per channel, the spanning-tree entropy of the slice graph by the matrix-tree theorem. Computed on degree-matched proxies for smooth extended and for crumpled geometry, and on real engine slices, that entropy is nearly universal at fixed coordination: the per-cell differential is a few times $10^{-4}$, and with all seven channels the phase-tipping force is of order $10^{-3}$ against the bare $k_0\simeq2.2$. The theory's own non-Gaussianity closes the loophole tighter still. The finite budget of Postulate~I gives the field a saturation mass $m^2=2/g_{\mathrm{share,eff}}\simeq0.27$, and the massive determinant suppresses precisely the soft modes that carried the residual sensitivity, shrinking the differential by a further factor of three at the budget mass and toward zero beyond. What the field retains is a local renormalization of the host's couplings, of order $0.08$ per cell in the seven-channel count: enough to relocate the host's phase boundaries, and unable to create the phase.

Within everything the paper specifies, then --- static weight, refresh, closure-class history weighting, and the conserved field free or budget-bounded --- vacuum geometric stiffness is not available: five candidate mechanisms, five computed exclusions (Appendix~J.4). Two containers stay open and are named: long-range field kernels, and non-equilibrium geometry--field co-evolution, both unspecified by the theory; testing an invented version of either would test the invention. The same calculations return a positive result: the refresh couples geometry to the \emph{density of closure failure}, the theory's matter, at the order that leaves the vacuum sector quiet and the defect sector live. That coupling is the mechanism the defect experiment measures. The same pass also verified, by explicit construction, the fracture statement that makes Postulate~III load-bearing in the scale chain (Appendices~D.4, H.6): under injectivity-preserving unit shifts the $1680$ states fall into exactly $48$ components ($4!$ orderings times two orientations) of $35$ states each.

\subsection{23.4 The conditioned vacuum: why the background is the theory's own}

The exclusions settle where the vacuum comes from, and the answer was already written into the architecture. Two different things are called emergence of spacetime here, and they separate cleanly. The first is kinematic: the substrate \emph{encodes} geometry by construction, since the closure invariant $K^2$ measures the failure of the four oriented faces to close --- the closure condition of the quantum tetrahedron --- adjacency carries proximity, and the length anchor of Section~13 converts entanglement increments into meters. The encoding must be read as Section~1.2 states it: the cells are dimensionless units of capacity, the granularity lives in the capacity and not in a preferred spatial lattice, which evades the dispersion bounds that ended rigid-lattice discrete-spacetime proposals, and $L_*$ is a statistical conversion rate fixed by the electron's dressing --- a dictionary between counts and meters, not a grid pitch. No result in the program touches this. The second is dynamical selection: among the combinatorially overwhelming crumpled and branched configurations, what makes the realized vacuum a smooth, extended, four-dimensional one? Selection is a ranking job, and a ruler cannot rank. Section~23.3 is the computed statement that nothing the theory specifies for the substrate performs it.

The theory never asked the substrate to. Section~27.8 states the position directly: the framework does not sum over triangulations, and its history weighting is supplied by Many-Pasts rather than by a path integral over geometries. In the Many-Pasts weight $P(H|P)\propto e^{-D(H,P)}$ (Section~21), every history is conditioned on the present coarse configuration $P$, and the present is an observed, extended, four-dimensional universe. The vacuum therefore enters the theory by conditioning: the coefficients of Part~II are computed within the state the conditioning selects, and the dynamics of the medium act on that state rather than generating it. What Section~23.3 adds is necessity. With every specified substrate mechanism computed and excluded, conditioning is the only supplier of the vacuum the theory contains, so the division of labor --- conditioning holds the vacuum, the medium holds matter and its response --- is forced by calculation. Postulate~III thereby acquires a fourth job alongside the scale chain, the decommitment orientation, and the foundations: it holds the vacuum state itself. The condensate branch of Section~22 is correspondingly relieved of a burden it was never assigned; what it owes is the microscopic form of the continuum kinetic term, and the selection of the state that term acts on belongs to the conditioning.

The same reading legitimates the lattice host, and fixes which host it had to be. Regge calculus is the standard discrete spelling of an extended smooth four-geometry, so a Regge background implements the paper's own conditioning in the only form a Monte Carlo program can. Among Regge ensembles, CDT is the one in which the conditioned state is a demonstrated output: an extended four-dimensional de~Sitter-like universe emerges dynamically in its established phase \citep{AmbjornJurkiewiczLoll2004,AmbjornGJL2012}. The conditioning therefore runs on machinery that produces the very state it conditions on. And the substrate co-authors the background it is conditioned on. The closure sector's free energy per cell, computed analytically from $\eta_*$ and the injectivity penalty with no lattice input, is a contribution to the host action's volume sector with a magnitude predicted in advance, and the prediction has been measured: at demonstration volume the uncentered coupling must displace the equilibrium volume by $-62.2$ units relative to the centered arm, and the three-arm measurement returned $-62.8$, one percent from prediction (Appendix~J.6). The computed curvature-sector shifts ($0.019$ from the closure weight, $\simeq0.08$ per cell from the seven-channel field) are pre-registered as phase-boundary displacements at production statistics. The structural kinship runs to the level of what is being measured: the closure defect $K^2$ at cell order and the Regge deficit angle at hinge order are both failure-to-close functionals, and the smallness of the curvature-sector shifts is the measured statement of how weakly the two levels currently couple. The substrate does not derive the Regge action; the exclusions stand. It dresses it: a coefficient of the gravitational action, written by the ultraviolet counting, at the predicted size.

\subsection{23.5 Compatibility: the weighting on dynamical geometry}

At $N_{41}=20{,}000$ (about $47{,}000$ pentachora in all, hosting $10{,}000$ cells across eighty slices), arms at $\beta=0$, $0.3$, $1$, and the $\beta=1$ placebo were run to matched volume with clean foliation audits (Appendix~J.5). The label sector orders exactly as the closure weight demands: the collision fraction falls from $0.648$ at $\beta=0$ (against the uniform-measure prediction $1-840/2401=0.650$, confirming that the control arm samples the base measure) to $0.082$ at the theory point, while the placebo, carrying the same energy values without the closure structure, stays at $0.73$. The geometry does not move: the Hausdorff dimension and the slice-volume profile are statistically indistinguishable across all four arms at single-seed resolution. That outcome was the quantitative prediction of Section~23.3, since an induced coupling of order one percent of the bare stiffness cannot steer these observables at this precision. The reading is compatibility: switching the theory's weighting on at full strength does not destabilize the host geometry, and the placebo separation in the label sector shows the coupling was live while the geometry held still.

\subsection{23.6 The defect experiment: geometry responds to the theory's mass}

The theory's matter is persistent closure failure: committed capacity, maintained against the refresh (Sections~3.2, 22). The experiment inserts it by hand. In a thermalized coupled ensemble at the theory point, one hundred well-separated cells were pinned into maximal closure failure (all four faces at $m=0$: six collisions, $K^2=48$) and held for three thousand sweeps. A placebo arm pinned one hundred cells into the best-closed injective configuration ($m=\{0,1,2,3\}$, the minimum $K^2=40.67$) under identical anchoring and identical protection from the geometry moves, so the real-minus-placebo comparison subtracts everything the pinning machinery does; freshly sampled unpinned cells far from any pin supply the vacuum reference inside both arms. Around each pin, four observables were accumulated shell by shell in the slice adjacency, for roughly $10^5$ defect-environment samples per arm: the mean closure energy and the collision fraction of the shell, which read the capacity state, and the edge coordination and cell count of the shell, which read the local geometry.

The environment responds, and the placebo comparison ties the response to the failure content. At the first shell the capacity state separates: mean closure energy $1.6782\pm0.0010$ around the failure pins against $1.7291\pm0.0011$ around the closed pins and $1.7306\pm0.0011$ in the vacuum, with the collision fraction displaced the same way ($0.0690$ against $0.0794$ and $0.0797$). By the second shell both label observables have returned to their vacuum values: the strain profile dies within one lattice step, the range the correlation length of Section~23.3 requires. The local geometry separates as well, and reaches further. At matched anchoring, the failure pins carry higher shell coordination than the closed pins ($5.2997\pm0.0007$ against $5.2513\pm0.0008$ at the first shell) and systematically more cells per shell, a local volume excess still present at the third shell: $11.296\pm0.006$ cells against the placebo's $10.905\pm0.007$ and the vacuum's $10.909\pm0.012$.

The licensed statement is local and specific: persistent closure failure sources a measurable strain in the surrounding capacity state and deforms the local geometry around it, and it does so through its failure content, since the placebo carries the anchoring without the failure. This is the microscopic seed of the source-to-strain mechanism the weak-field sector runs on, demonstrated dynamically. The volumes license nothing long-range: no $1/r$, no $G$, no curvature in the continuum sense; the measured reach is three lattice steps. The result is single-seed, its quoted errors are statistical and not yet thinned for autocorrelation, and the host's phase location is still being mapped; a detected response survives those caveats, which govern the precision of the numbers rather than the existence of the effect.

\subsection{23.7 Reaching Newtonian range: the conservation requirement}

The strain of Section~23.6 dies within a lattice step, as the sub-cell correlation length says it must. The weak-field sector needs it carried to macroscopic range, and the lattice makes sharp what that requires. Suppose free capacity were only a per-cell budget, re-equilibrating locally by maximum entropy around each commitment. The linear response then solves $(2z-L)\,\delta\mu=-m/\chi$ on the slice adjacency ($z=4$ faces per cell, $L$ the slice Laplacian), and that operator has no small-momentum pole: on a real slice the response falls eight orders of magnitude within thirteen steps. A budget that merely re-balances screens gravity at the substrate scale. Suppose instead that free capacity obeys a genuine continuity law (transported through shared faces, exchanged only with commitment, globally conserved) and that a defect draws the steady maintenance flux Section~22 assigns it, linear in its commitment. The steady state then solves the massless graph-Poisson equation, massless because conservation forbids any local restoring term, and its Green function on a three-dimensional slice falls as $1/r$: the deficit field is $\delta f\propto m/r$, the $M$-linearity inherited from the maintenance postulate and the $1/r$ from the conservation law. On the lattice, Newton's form of Section~12 is the signature of conserved capacity with maintenance sinks, and nothing weaker produces it (Appendix~J.8).

The continuum sector never had to say what is conserved; the lattice does. Two carriers are available: a capacity stock transported between cells, which would require a conserved quantity the label sector demonstrably does not possess, or the refresh bandwidth itself, each cell's one redraw per $\tau_*$, exactly conserved tick by tick, with a defect's maintenance consuming the refresh coherence of its neighborhood, and the deficit field the steady allocation of that bandwidth. The paper's own commitments lean to the second: mass is already a recurrence rate (Section~22, Appendix~H.6), $a_0\propto cH(z)$ is a rate, and maintenance linear in commitment is the natural law for a bandwidth share.

The measurement instrument is built and gated in advance (Appendix~J.8). A conserved field on the slice cells, initialized at the vacuum anchor $7.4198$ and transported by antisymmetric face fluxes with global conservation asserted at machine precision every sweep, is coupled to pinned defects through the microphysics alone: cells absorb capacity in proportion to their instantaneous closure failure, so the sink strength a defect exerts is an output of the label dynamics, with nothing in the transport law specifying it. Absorption reads the failure a cell carries above the vacuum's own standing level, so the vacuum is absorption-free and the transported field massless by construction; without that, the vacuum's own absorption generates a screening length of its own, which the first full-volume run measured at the predicted scale (Appendix~J.8). A ladder of pinned commitments at four mass levels then measures the two quantities the theory does not preordain: whether the emergent maintenance flux is linear in the commitment, which is the maintenance postulate tested dynamically, and the deficit per unit flux at the source junction, the lattice form of the statement that fixes $G$. Two profile gates are fixed in advance: the deficit must follow the graph-Coulomb form, and a fitted screening mass must be consistent with zero, since any leakage in the conservation law would appear as Yukawa screening of the far field. The production measurement, at the volumes the far-field gates require, is in progress; its readings are committed in advance and will be reported at whatever grade they earn.

\subsection{23.8 What the program establishes}

The division of labor the results draw matches the paper's architecture. The closure sector owns what mass is: the counting of Part II, reproduced as a fidelity gate in every run. The refresh owns how mass acts locally: a persistent failure strains its neighborhood's capacity state and deforms its local geometry, measured against a placebo that isolates the failure content. Conservation owns how the action travels: the long-range sector requires a conserved capacity current with maintenance sinks, and given one, Newton's form is structural. And the vacuum is owned by none of them: five computed exclusions leave the conditioning of Postulate~III as the only supplier of the smooth four-dimensional state the theory contains, with the Regge background implementing the conditioning and the substrate's measured contribution to the volume term showing that the relation runs in both directions (Section~23.4). Section~24 carries these entries at their measured grade: a five-branch exclusion with its forced architectural reading, a measured contribution to the host action, one compatibility result, one dynamical mechanism result at single-seed grade, and one derived requirement with its measurement in progress.

\section{24. Closure-Status Table}

The closure bookkeeping is concentrated here in one place so the rest of the text can simply derive, state, and move on.

The status column uses a fixed vocabulary. \emph{Closed} means derived within the stated postulates and ensemble. \emph{Fixed} means no phenomenological freedom remains once the branch is adopted. \emph{Conditional} means the result follows if a named reading or completion holds. \emph{Frontier} or \emph{open} means the piece is structured but not yet complete. \emph{Empirical support} means a comparison with data rather than a derivation, and \emph{audit task} means an independent derivation or check is still required. Rows follow the order of the text, grouped by sector from the UV counting through the foundational, weak-field, galactic, transport, cosmological, strong-field, and quantum-foundational sectors, with the lattice-dynamics entries of Section~23 ahead of the validation layer.

\begingroup
\scriptsize
\setlength{\tabcolsep}{3pt}
\renewcommand{\arraystretch}{1.12}
\begin{longtable}{|L{0.18\textwidth}|L{0.13\textwidth}|L{0.16\textwidth}|L{0.33\textwidth}|L{0.13\textwidth}|}
\hline
\textbf{Quantity / Claim} & \textbf{Sector} & \textbf{Status} & \textbf{Type of Support} & \textbf{Where Established} \\
\hline
\endfirsthead
\hline
\textbf{Quantity / Claim} & \textbf{Sector} & \textbf{Status} & \textbf{Type of Support} & \textbf{Where Established} \\
\hline
\endhead
$\Omega_{\mathrm{tet}}$, $g_{\mathrm{share,max}}$ & UV counting & Closed & exact combinatorics & Part II, App. B \\
\hline
$\eta_*$ & admissibility closure & Closed & stationary normalized closure evidence $\ln Z+\tfrac32\ln\eta$ maximized on the exact $K^2$ spectrum; the $3/2$ is the determinant weight of the three closure-defect components & Part II, App. B \\
\hline
$g_{\mathrm{share,eff}}$ & UV entropy & Closed & exact weighted evaluation & Part II, App. B \\
\hline
Substrate length $L_*$ (faithful sector resolution / support-to-length) & scale setting & Derived from dressing recurrence; identification discharged under Fork~A & $\lambda_e/L_*=\tfrac23 e^{7g_{\mathrm{share,eff}}}$ from the recurrence time of the memoryless seven-channel dressing (waiting time for joint typicality); support exponent bracketed at unity within $10^{-3}$ by $G_*$ and the lepton ladder; phase-cycle identification discharged in App.~H ($\Gamma\oplus\Delta Q{=}1\oplus$ Josephson), residual is Fork~A (winding phase $=$ condensate $\mathrm{U}(1)$) and the unit-winding integer, both tested by the lepton ladder & Part I, App. D, H \\
\hline
$J_{\mathrm{bare}}$, $J_{\mathrm{eff}}^{\mathrm{tree}}$ & UV edge kernel & Closed & tetrahedral isotropy identity & Part II, App. C \\
\hline
$\Sigma_{\mathrm{ret}}=65/9$ & finite-loop UV & Fixed in the minimal UV return sector & explicit seven-channel plus singlet return count & Part II, App. C \\
\hline
$J_{\mathrm{eff}}^{(\mathrm{ren})}$ & finite-loop UV & Fixed by the UV return resummation & derived from $\Sigma_{\mathrm{ret}}$ & Part II, App. C \\
\hline
$\gamma$ & continuum stiffness & Closed in canonical EFT convention & Euclidean-action normalization & Part II, App. C \\
\hline
Green-matched source projection & UV source map & Closed in the canonical weak-field branch & exact defect counting + tetrahedral on-site Green function & App. C \\
\hline
$\kappa/\gamma$ & source-to-stiffness ratio & Closed in the canonical weak-field branch & $\sigma_{\mathrm{def}}=\rho/\kappa_m(L_*)$ plus $G_{\mathrm{tet}}(0)$ and tetrahedral $4/3$ projection & Part III, App. C--D \\
\hline
Weak-field bridge law & EFT / gravity & Closed in canonical weak-field branch & uniqueness from multiplicative composition & Part III, App. D \\
\hline
$G_*$ & gravitational scale & Conditional derivation at percent level; calibration, with provenance recorded & $G_*=(9/4)(\hbar c/m_e^2)e^{-14g_{\mathrm{share,eff}}}$ & Part I, App. D, L \\
\hline
Matched $G$ & weak-field gravity & Closed in the matched weak-field branch & bridge law + Green-matched $\kappa/\gamma$ + fixed-epoch normalization of $\kappa/(\gamma S_\infty)$; comparison with $G_*$ tests coherent propagation of the same $L_*$, not a disjoint-input second derivation of $G$ & Part III, App. C--D \\
\hline
Electron anchor & mass / length sector & Fixed empirical elementary anchor & one-bit fermionic defect supplies $\lambda_e$ for length setting and $m_e/\ln2$ for mass--entropy map & Part III, App. D \\
\hline
Seven-sector additivity & UV scale theorem & Closed conditional on the support-to-length identification & $H_{\rm seven}=7g_{\mathrm{share,eff}}$ from subadditivity plus mass minimization: full per-channel entropy from the memoryless dressing, channel independence forced because correlation would raise the defect mass ($\Delta=0$ for the lightest one-bit defect) & App. D, H \\
\hline
$a_0$ & galactic EFT & Fixed in the closed weak-field realization & UV entropy $\times$ horizon scale $cH_0$, with $4\pi^2=(2\pi)^2$ the canonical two-transverse-mode phase-space cell of the $1+2$ split & Part III, App. C \\
\hline
RAR law & galactic EFT & Closed at the rate-bookkeeping level & $a_0$ fixed by the coefficient chain; $x=\sqrt{g_{\mathrm{bar}}/a_0}$ forced by the spectral invariant of the bilinear form; $1+n_B(x)$ from exact bosonic emission statistics, with absorption closed by the quantized one-bit deficit; flux identification derived from the static transport limit (conserved bit current, Green-matched normalization), with inter-channel delivery fixing the locality of $x$; kernel neutrality $F\equiv1$ from per-cell channels and bit quantization; graph-level substrate computation of the inter-channel vertex is the remaining gold-standard check & Part III, App. C \\
\hline
No slip / lensing consistency & weak-field metric & Closed at leading weak-field order & scalar-stress structure & Part III, App. D \\
\hline
PPN leading values & weak-field metric & Structurally supported & weak-field expansion & Part III, App. F \\
\hline
Telegrapher relation $D/\tau_0=c^2$ & transport & Closed in canonical transport branch & causal closure & Part IV, App. E \\
\hline
Canonical $\tau_0^{-1}=H_0$ branch & transport & Fixed in the minimal transport closure & no-new-IR-scale choice & Part IV, App. E \\
\hline
Hubble-tension mechanism & cosmology & Structurally supported extension & homogeneous trace-coupled mode & Part IV, App. E \\
\hline
Saturated-phase committed dust & cosmology & Conditional on the pinned transfer-law reading & saturation forced at recombination; constrained-scalar dust theorem ($p=0$, $c_s^2=0$, $\rho\propto a^{-3}$); relaxational carriers excluded & Part IV, \S18.5, App.~M \\
\hline
Committed-capacity abundance $\Omega_c/\Omega_b=1/\epsilon$ & cosmology & Derived from forced recruitment, conditional on the pinned reading and per-defect bookkeeping & over-demanded sources recruit to the per-source cap, the background sits below it, and per-defect allocations add (Postulate~II corollary), giving $\rho_c=\rho_b/\epsilon$ with adiabatic seeding $\delta_c=\delta_b$; temperature-spectrum check at fixed abundance; counted, not fitted & Part IV, \S18.5 \\
\hline
Caustic release / conditional conservation & cosmology & Conditional --- forced under the synchrony reading & committed dust decommits at first shell-crossing of the irrotational clock flow; splashback-rim falsifier; conversion bookkeeping open & Part IV, \S18.5 \\
\hline
Diffuse source projection $\epsilon=g_{\mathrm{share,eff}}/4\pi^2$ & cluster sector & Derived (no new coefficient) & same transverse reduction factor that fixes $a_0/a_H$ (Section~14) & Part IV, \S17.5 \\
\hline
Bath gate $B_{\rm bath}$ & cluster sector & Operationally closed; microscopic origin open & $M_{\rm hot,vir}/(f_{b,\rm cos}M_{500})$ from X-ray/SZ; no per-system fitted knob & Part IV, \S17.5 \\
\hline
Cluster residual $\mathcal R_{\rm rel}=1+(W_{\rm bath}-1)f_{\rm cont}$ & cluster sector & Hook morphology and trend direction supported; linear lift candidate excluded; lift function open & linear candidate's ceiling $1.81$ lies below measured peak residual $3$--$5$ \citep{Tian2020,Eckert2022,Li2023}; capacity bound fixes saturation at $1/\epsilon$, giving $\mathcal R_{\rm sat}\simeq4.7$--$4.9$ and a radius-resolved peak $\simeq3.9$ inside the measured band; the bound holds against all twenty-four X-COP source-weight inversions, which also exclude the deep power-law continuation within that sample; coherence-growth profile open & Part IV, \S17.5 \\
\hline
Relaxed vs.\ merger source expressions & cluster sector & Conditional --- two regimes, one coefficient & residual rides on virialized bath; Bullet on shocked gas $+$ decoupled clumps; regime boundary not yet derived & Part IV, \S17.5 \\
\hline
Channel-selection rule (suppression $+$ collective lift) & cluster sector & Conditional --- suppression open; linear lift excluded & transverse-suppression premise (participation $\Rightarrow$ weight $\epsilon$) not yet derived from the source map; linear form of the lift excluded by the measured peak residual; saturation endpoint fixed at $1/\epsilon$ by the channel capacity bound; coherence-growth profile between the endpoints not yet derived & Part IV, \S17.5 \\
\hline
Bullet gas/lensing inversion & cluster sector & Structurally supported; consistent but untested ($\epsilon$ not yet measured) & $\epsilon$ brings gas/galaxy to near-parity, compactness completes the inversion; existing flexible reconstructions non-discriminating because gas weight is halo-degenerate & Part IV, \S17.5 \\
\hline
Resolved cluster lensing-map test & cluster sector & Open --- principal empirical task & three-component $\kappa\propto(1+(1-\epsilon)B_{\rm bath})\Sigma_{\rm bath}+\epsilon\Sigma_{\rm shock}+\Sigma_{\rm dec}$; predicted best-fit $\epsilon\simeq0.19$, floated on baryonic maps with free dark haloes disallowed; not yet performed & Part IV, \S17.5 \\
\hline
Bounded capacity $q$, $N^2=q$ & strong field & Closed as constrained-capacity rule & uniqueness from multiplicative composition and weak-field matching & Part V, App. F \\
\hline
Constrained action on $\mathcal{M}_q=\{q>0\}$ & strong field & Closed within the EFT & ADM action plus lapse-capacity constraint & Part V, App. F \\
\hline
Static spherical exterior & strong field & Closed within the constrained-capacity branch & bulk vacuum reduces to Einstein vacuum on $\mathcal{M}_q$ & Part V, App. F \\
\hline
Capacity-exhaustion horizon $q=0$ & strong field & Closed as static-domain boundary & Schwarzschild exterior terminates where $q$ leaves the physical state space & Part V, App. F \\
\hline
Hawking temperature and exterior ringdown & strong field & Closed at leading EFT order & unchanged Schwarzschild exterior; absorbing boundary at $q=0$ fixed by horizon regularity; echoes only as UV / stretched-layer corrections & App. F \\
\hline
Bekenstein--Hawking area-law bridge & strong field & Closed within the constrained-capacity EFT under Postulate~II and the transverse graph response & per-channel cut entropy $\ln 2$ from fermionic face exclusion + channel density from $G_\perp=(2/3)G_{\mathrm{tet}}(0)$; product $n_{\rm hor}S_\infty^{\rm cell}=1/4$ as exact identity & App. C, F \\
\hline
Horizon formation and boundary microphysics & strong field & Formation closed geometrically; nonuniversal spectroscopy open & $q=|\nabla R|^2=1-2GM_{\rm MS}/(c^2R)$ in spherical symmetry, with $q=0$ at the marginal-trapped-surface inherited from standard apparent-horizon formation; remaining work covers relaxation spectra, stretched-layer corrections, and transient response & Part V, App. F \\
\hline
Rotating / charged stationary exteriors & strong field & Closed within the constrained-capacity EFT & vacuum Einstein / Einstein--Maxwell reduction on $\mathcal{M}_q$ gives Kerr / Reissner--Nordstr\"om / Kerr--Newman exteriors; $q=0$ identified with outer Killing horizon; inner Cauchy horizons are non-physical analytic continuations into $q<0$ & App. F \\
\hline
Many-Pasts Born compatibility & quantum foundations & Consistency condition & exact Born statistics selects $\alpha=1$ within the weighting family; the content is the existence of a branch meeting Born and no-signaling simultaneously & Part V, App. G \\
\hline
No-signaling in operational branch & quantum foundations & Consistency condition & forbidding signaling-sensitive bias selects $\beta=0$ & Part V, App. G \\
\hline
Arrow-of-time account & quantum foundations & Coherent extension & conditional typicality / counting & Part V, App. G \\
\hline
Microstructure Hamiltonian & UV realization & Defect side closes scale-setting; geometry side coherent & dressing Hamiltonian with one-pass ground state, memoryless dressing, and minimal-mass channel independence (scale-setting); mass as the refresh-operator gap with the phase-cycle identification discharged into $\Gamma\oplus\Delta Q{=}1\oplus$ Josephson (App.~H.6) and $\varepsilon_0,F_m,V$ from the mean-field condensate with rank-one $V$ as one scalar (App.~H.7); GFT/condensate origin of the kinetic term (geometry) & Part V, App. H \\
\hline
Charged-lepton spectrum & particle-sector extension & Derived; coefficients forced, one map-form input & three-generation termination from $K^2$ permutation-invariance at $N=3$; $m_N/m_e=720^N(2/7)^{N^2}$ with $720=6!$, $2/7$ the singlet projection, $N^2$ the ordered-pair count; matches PDG at $0.5\text{--}0.65\%$ with no fitted parameters & App. I \\
\hline
Gauge-redundancy extension & gauge sector & Coherent extension & baseline-redundancy construction with Maxwell/Yang--Mills form & App. I \\
\hline
Vacuum stiffness from the substrate & lattice dynamics & Excluded, five branches & induced curvature coupling $c_0=+0.019$ against bare $k_0\simeq2.2$; refresh stationary state geometry-blind to exponential accuracy; history-tilt ceiling twenty times below requirement; free conserved field induces near-universal spanning-tree entropy; budget saturation mass suppresses the residual soft modes & \S23.3, App.~J.4 \\
\hline
Vacuum selection by conditioning & lattice dynamics & Forced division of labor & with all five specified substrate mechanisms excluded, the Many-Pasts conditioning on the extended four-dimensional present is the only vacuum supplier the theory contains; the Regge background implements that conditioning & \S23.4, App.~J.4 \\
\hline
Ensemble fracture under local dynamics & lattice dynamics & Verified & explicit construction: $48$ components ($4!\times2$) of $35$ states each under injectivity-preserving unit shifts, conserved (ordering, parity) charge & \S23.3, App.~D.4, H.6, J.4 \\
\hline
Substrate contribution to the host volume term & lattice dynamics & Measured, demonstration volume & closure free energy per cell, computed from $\eta_*$ alone, displaces the equilibrium volume by the predicted amount: $-62.2$ predicted, $-62.8$ measured; curvature-sector shifts pre-registered & \S23.4, App.~J.6 \\
\hline
Admissibility weighting on dynamical geometry & lattice dynamics & Compatibility demonstrated & matched-volume coupled ensemble at the theory point: label sector orders (collision fraction $0.648\to0.082$; placebo $0.73$), geometric observables statistically unchanged, as the one-percent induced coupling predicts & \S23.5, App.~J.5 \\
\hline
Local response of geometry to persistent closure failure & lattice dynamics & Demonstrated dynamically; single-seed grade & pinned maximal-failure cells against matched closed-pin placebo: first-shell capacity strain, coordination and shell-volume deformation through three shells, separating from placebo well outside quoted (unthinned) errors & \S23.6, App.~J.7 \\
\hline
Long-range propagation of the strain field & lattice dynamics & Conservation law derived; measurement in progress & budget-only response screens at sub-cell range; a conserved free-capacity current with maintenance sinks gives the massless graph-Poisson form, $\delta f\propto m/r$; transport instrument built with pre-registered linearity and screening-mass gates & \S23.7, App.~J.8 \\
\hline
Numerical robustness checks & validation layer & Supportive audit layer & cross-sector consistency tests & App. K \\
\hline
EFT consistency checklist & field-theory audit & Supportive audit layer & no-ghost / no-tachyon / causal-propagation checklist with explicit vacuum dispersion stability & App. D \\
\hline
\end{longtable}
\endgroup

This table is the epistemic map used for the rest of the discussion.

\paragraph{Cosmological row.}
The saturated phase enters the ledger as follows: the dust form of the committed phase is a theorem of the constrained-scalar class, conditional on the pinned reading of the transfer law; the abundance $1/\epsilon=5.321$ stands against the measured $5.364\pm0.065$, derived from forced recruitment conditional on the same pinned reading and the per-defect bookkeeping of Postulate~II; the energy accounting of the committed budget is open; the linear regime is closed by the domain structure.

\section{25. Falsifiability and Observational Tests}

\subsection{25.1 Static weak-field falsifiers}

The static weak-field sector stands or falls on a small number of concrete checks. The most direct are the shape and tightness of the galaxy RAR transition \citep{McGaugh2016,Lelli2016}, the baryonic Tully--Fisher scaling in systems where the EFT should apply, and the weak-field lensing sector, where stacked galaxy--galaxy lensing now probes the relation two decades below the rotation-curve regime \citep{Brouwer2021}. A persistent need for gravitational slip where the scalar stress predicts none would be especially damaging, because it would break the same no-slip structure used to keep dynamics and lensing aligned. Solar-system bounds, especially Cassini-class PPN tests \citep{Bertotti2003,Will2014}, also require the leading no-slip branch to survive at high precision, and laboratory fifth-force searches \citep{Adelberger2003} bound any residual static scalar channel in the same regime.

Wide binaries supply an independent solar-neighborhood discriminator. The occupancy argument of Section~14 is local, a feature its transport derivation fixes rather than assumes: the acceleration entering $x$ is the total local baryonic field, so a binary embedded in the Galactic field $g_{\rm ext}\simeq1.4$--$1.9\times10^{-10}\,{\rm m\,s^{-2}}$ has its internal dynamics amplified once the pair's own field falls below the Galactic term. The predicted internal-force boost rises from unity at small separations to $1+n_B(\sqrt{g_{\rm ext}/a_0})\simeq1.4$--$1.5$ beyond a few thousand au. This lands close to the AQUAL external-field value through a different mechanism --- locality of the occupancy argument in the total baryonic field, with no nonlinear field equation --- so the wide-binary regime separates the framework from Newtonian gravity while leaving it degenerate with AQUAL-class theories. Current analyses divide: a low-acceleration boost of about $1.4$ is reported in Gaia DR3 samples \citep{Chae2023}, and consistency with Newtonian gravity is reported in an independent analysis of the same mission's data \citep{Banik2024}; the disagreement is methodological. The framework's commitment is fixed in advance: a settled Newtonian verdict falsifies the locality of the occupancy rule.

\subsection{25.2 Dynamical falsifiers}

The dynamical extension is more vulnerable, and its failure modes are correspondingly sharper. Time-dependent halo lag, cluster-scale acceleration relations, merger offsets, or relaxation signatures that cannot be reconciled with the telegrapher relation $D/\tau_0=c^2$ would indicate that the causal completion has the wrong propagation structure even if the static branch survives. Cluster data accordingly split into two tests of distinct parts of the framework: \emph{source-projection} tests, which ask whether the relaxed-cluster residual profile follows the predicted hook morphology---near unity in BCG-dominated centers, peaked where the virialized bath dominates, converging in the deep outskirts \citep{Tian2020,Eckert2022,Li2023}---and whether resolved merger maps prefer the projection coefficient $\epsilon\simeq0.19$ (Section~17.5); the linear lift candidate is excluded on amplitude (Section~17.5); and \emph{transport} tests, which ask whether the causal propagation law $D/\tau_0=c^2$ correctly evolves those source weights through a merger. Systems such as the Bullet Cluster \citep{Clowe2006} probe both, and should not be treated merely as larger versions of the static galaxy problem.

\subsection{25.3 Cosmological falsifiers}

Cosmology presents a different kind of test. The question there is whether a full Boltzmann treatment allows the trace-coupled homogeneous mode to reduce the sound horizon without spoiling the CMB or structure-growth observables. If it cannot, the cosmological extension fails on its own terms. The empirical target is set by the current measurement spread: early-universe inferences near $67.4$ \citep{Planck2020}, distance-ladder determinations ranging from $\simeq70$ \citep{Freedman2025} to $73$ \citep{Riess2022}, and a tension whose proposed resolutions are reviewed in \citet{DiValentino2021}.

The epoch dependence $a_0(z)=cH(z)\,g_{\mathrm{share,eff}}/4\pi^2$ is the framework's most exposed kinematic prediction, and high-redshift rotation curves already bear on it. At $z\simeq2.3$ the predicted scale is $3.5\,a_0$; a massive star-forming disk with $g_{\rm bar}\simeq2\times10^{-10}\,{\rm m\,s^{-2}}$ at $8\,$kpc is then predicted to show $g_{\rm obs}/g_{\rm bar}\simeq2.0$, against $1.36$ for an epoch-independent scale. Observed outer rotation curves of such disks decline and indicate strong baryon dominance within the disk scale \citep{Genzel2017,Lang2017}, the direction opposite to the predicted enhancement. Pressure-support corrections at these redshifts are large and the stacked-profile normalization is disputed, so the confrontation is unsettled. The framework requires that massive $z\simeq2$ disks, once those corrections are controlled, show stronger low-acceleration boosts than matched $z=0$ systems; the opposite outcome falsifies the identification of $a_0$ with $cH(z)$ while leaving the static $z=0$ sector intact.

Direct measurements at intermediate redshift have since entered. A MUSE sample of $79$ star-forming galaxies at $0.33<z<1.44$ finds the radial-acceleration relation persisting with a characteristic acceleration that rises systematically with redshift \citep{Ciocan2026}, and a resolved low-redshift H\,\textsc{i} sample independently reports tentative evolution in the same direction \citep{Varasteanu2025}. The sign is the framework's: an epoch-independent $a_0$ predicts no evolution at all. The magnitude is not yet a test, since absolute normalizations at these redshifts carry mass-to-light and pressure-support systematics of the same order as the effect; the clean confrontation is the survey-internal ratio $a_0(z_{\rm high})/a_0(z_{\rm low})$ against $E(z_{\rm high})/E(z_{\rm low})$, a comparison with no free parameter that the prediction must pass.



\paragraph{Saturated-phase falsifiers.}
Three tests bind the saturated phase: a full Einstein--Boltzmann implementation of the constrained committed fluid must reproduce the measured temperature and polarization spectra with the abundance fixed at $1/\epsilon$; the recruitment derivation must survive scrutiny of its two named conditions, the pinned reading and the per-defect bookkeeping; and departures from the $\Lambda$CDM form must be confined to accelerations above $g_c\simeq5.2\times10^{-12}\,\mathrm{m\,s^{-2}}$, so canonical behavior failing below that boundary, or non-canonical behavior appearing in the linear regime, would falsify the domain assignment. The release law adds a fourth: a surviving committed component on cluster infall streams terminating at the splashback surface; its robust absence, or a halo-like committed component persisting inside collapsed systems, would falsify the caustic release mechanism.

\subsection{25.4 Correlated-constant falsifiers}

One of the more distinctive signatures of the framework is that the same microstructural chain feeds the substrate scale, the weak-field gravitational normalization, and the galactic acceleration scale. A precision program that could test the inferred $L_*$ and $G_*$ scale against matched $G$, $a_0$, and the RAR normalization would probe the theory more sharply than isolated single-observable fits, because it would confront the shared coefficient origin directly.

\subsection{25.5 Many-Pasts status}

The Many-Pasts sector is not likely to be challenged first by ordinary laboratory deviations from quantum mechanics, because it is built to reproduce the usual operational structure there. Its more immediate points of failure are internal ones: failure of exact Born recovery, failure of no-signaling, or incompatibility with the thermodynamic arrow structure it is supposed to illuminate.
\section{26. What the Theory Would Have to Get Wrong to Fail}

The failure modes are not all equally severe, and it helps to order them by how much each would bring down.

\paragraph{Kills the core ontology.} A demonstrated failure of mass--entropy equivalence, an internal incoherence in the Many-Pasts weighting, or evidence that geometry cannot be read as the long-wavelength form of an entanglement-capacity substrate would remove the foundations on which everything else rests.

\paragraph{Kills the static weak-field closure.} A persistent gravitational slip where the scalar stress predicts none; an RAR transition shape that systematically departs from the derived bosonic occupancy law in systems well described by the static branch; a source projection contradicted by the data; or a weak-field UV coefficient chain that cannot be reconciled with an independently validated microscopic derivation --- any of these breaks the central completed result while leaving the ontology in principle intact.

\paragraph{Kills an extension only.} If the cosmological trace-coupled homogeneous mode cannot survive a full Boltzmann likelihood confrontation, the cosmology sector fails while the static weak-field branch stands. If the saturated phase fails any of its commitments --- the dark-to-baryonic abundance departing from the capacity ceiling, linear-regime observables departing from the inherited form below the release threshold, or no committed component surviving on cluster infall streams out to the splashback surface --- the committed-dust reading fails in the same contained way. If the constrained-capacity branch cannot retain the tested Schwarzschild exterior, standard absorbing-boundary ringdown, or the required horizon thermodynamics once its boundary action is derived, the strong-field completion fails. None of these touches the weak-field core.

\paragraph{Requires modification, not death.} Some results could be wrong without bringing down a sector: an adjustment to the finite-loop self-energy from a fuller graph-level computation, a revision of the strong-field boundary microphysics, or a failure of the charged-lepton ladder, which is a cross-check on the mass--entropy ontology rather than part of the closed gravitational chain.

One scale-setting commitment cuts across these tiers and deserves its own line. The percent-level $G_*$ rests on the seven-fold additivity, which rests in turn on the electron being the lightest one-bit charged defect, since inter-channel correlation in its dressing would raise the defect mass. If that identification or the mass--entropy ontology failed, the scale-setting derivation of Appendix~H would not go through.
\section{27. Comparison with Other Approaches}

Because the framework reattributes the dark sector (the galactic excess to medium response, the cosmological abundance to a saturated phase of the same medium) and partially reorganizes the usual dark-energy story, the sections below set out how its logic differs from nearby alternatives.

\subsection{27.1 Relative to \texorpdfstring{$\Lambda$}{Lambda}CDM}

The contrast with $\Lambda$CDM begins at the level of ontology. Standard cosmology explains the relevant phenomenology by adding dark matter and an independent cosmological constant or dark-energy sector to otherwise standard gravity. Here the visible matter sector is retained, but it is interpreted as the macroscopic description of localized defects in a vacuum-capacity medium whose weak-field response supplies the effective extra gravitating component. The same closure chain is then asked to feed $G$, $a_0$, the RAR law, weak-field lensing consistency, and the homogeneous cosmological mode.

\subsection{27.2 Relative to MOND-like interpolation programs}

MOND-like programs \citep{Milgrom1983,SandersMcGaugh2002,FamaeyMcGaugh2012} usually begin from an acceleration law or interpolation function and ask how much galaxy phenomenology it can explain. The present logic runs the other way. The interpolation law is not taken as primary; it is downstream of the UV entropy, the $1+2$ channel geometry, and the bosonic occupancy branch. The galactic law is thus treated as an output of the same micro-to-IR closure chain rather than as the phenomenological starting point.

\subsection{27.3 Relative to Verlinde-style emergent gravity}

Verlinde-style emergent-gravity programs share the broad intuition that gravity may be entropic \citep{Jacobson1995,Padmanabhan2010,Verlinde2011,Verlinde2017}, but they are usually formulated at the level of thermodynamic reasoning or horizon-inspired force laws. The present framework is narrower and more explicit: finite tetrahedral boundary counting, admissibility closure, edge coupling, finite renormalization, Euclidean-action normalization, and only then a continuum scalar EFT. Whether that chain is ultimately correct is an empirical matter, but it is a different kind of proposal from a purely macroscopic entropic argument.

\subsection{27.4 Relative to TeVeS and other multi-field modified gravities}

Multi-field relativistic MOND completions such as TeVeS \citep{Bekenstein2004} typically introduce additional vector or tensor sectors in order to repair lensing or cosmological problems. The present weak-field construction instead keeps a single scalar entanglement field within a low-energy Einstein continuum sector that is itself interpreted as emergent from the substrate, and relies on the no-slip structure $\Phi=\Psi$ at leading order to keep lensing and dynamics aligned. That economy is attractive if the branch survives confrontation with data, and immediately vulnerable if future observations demand persistent slip or extra weak-field structure.

\subsection{27.5 Relative to AeST}

The closest modern comparator is the AeST theory of Skordis and Z\l{}o\'snik \citep{Skordis2021}, which reproduces MOND-scale galaxy phenomenology and fits the CMB power spectrum using a vector field and two scalar degrees of freedom. AeST establishes that galaxy-scale modified dynamics and a viable linear cosmology can coexist in a single relativistic theory, and any framework in this space is measured against it. The present construction differs in starting point and in risk profile: the acceleration scale, the interpolation argument, and the source projections descend from a finite UV counting problem with coefficients fixed in advance, so nothing can be adjusted where AeST retains functional freedom, and the cosmological sector is correspondingly less developed --- the full Boltzmann treatment AeST has passed is named open work here (Section~25.3). The two programs fail differently: AeST through accumulating parameter pressure, the present framework through any single fixed coefficient meeting a contrary measurement.

\subsection{27.6 Relative to scalar-tensor gravity}

The action of Section~10 can be mistaken for a Brans--Dicke-type scalar-tensor theory \citep{BransDicke1961}, since it carries a metric, a scalar, and a coupling between them. The difference is ontological, and it removes the usual scalar-tensor freedom. In a scalar-tensor model the scalar is an independent field added to general relativity, and its kinetic and coupling functions are free to be chosen or fit to data. Here the metric and the scalar are two continuum expressions of one finite-capacity substrate: the scalar is not added to gravity but the same medium read in a different variable, and its stiffness, source coupling, and bridge normalization are fixed by the UV counting chain rather than left as functions. The resemblance is at the level of the written action; the content is that the coefficients are not adjustable.

\subsection{27.7 Relative to quantum-mechanical interpretations}

Because Many-Pasts occupies the role of an interpretation of quantum mechanics (Sections~3.3, 21), it should be placed against the standard options. Unlike many-worlds, it posits no forward branching into co-real macroscopic outcomes; the multiplicity is over admissible pasts of a single realized present, not over futures. Unlike collapse theories, it adds no physical collapse event and no modification of unitary evolution. Unlike hidden-variable accounts, its probabilities are not ignorance over pre-existing classical values. It is closest in spirit to consistent-histories and decoherence-based reasoning, but it differs in carrying an explicit weight $e^{-D(H,P)}$ over histories and in using that weight to do physical work outside the foundations --- the memoryless dressing that sets the substrate length, and the orientation of cosmological decommitment. It does not claim to out-derive the Born rule; it recovers Born statistics and no-signaling as consistency conditions on the weight (Section~21).

\subsection{27.8 Relative to CDT, spin foams, and group field theory}

The construction uses discrete tetrahedral boundary data, and the $j=3/2$ tetrahedron coincides with the quantum tetrahedron of simplicial spin networks \citep{Barbieri1998,BaezBarrett1999,RovelliSmolin1995}, so it sits near the causal-dynamical-triangulation, spin-foam, and group-field-theory programs. The way it uses those structures differs on three points. It does not sum over arbitrary triangulations or evaluate a spin-foam vertex amplitude as the primary weighting; it solves a finite local boundary-counting problem and weights configurations by admissibility closure, routing the result to a capacity entropy rather than to area and volume operators. Matter is not added on top of the geometry but is a fermionic capacity defect of the same substrate. And the history weighting is supplied by Many-Pasts rather than by a path integral over geometries. The GFT/condensate picture enters only as a compatibility check on the continuum kinetic term (Section~22), not as the engine of the derivation. The relation is therefore one of shared combinatorial vocabulary with a different selection principle and a different target. Since Section~23 that relation is also a working interface: the admissibility weighting has been coupled, unmodified, to a causal-dynamical-triangulations ensemble \citep{AmbjornJurkiewiczLoll2004,AmbjornGJL2012}, where it leaves the host geometry intact at matched volume and supports the pinned-defect and capacity-transport measurements reported there. CDT supplies the dynamical substrate --- in the paper's terms, the conditioned vacuum in discrete form (Section~23.4) --- and the selection principle under test is this paper's. The debated points of the CDT program itself, chiefly the fixed foliation, do not transfer: the lattice serves here as conditioning apparatus and regulator, so those criticisms reach the scaffolding and stop there; the substrate's granularity lives in the capacity (Section~1.2).

\section{28. Conclusion}

One physical picture runs through the paper. The vacuum is a finite medium of entanglement capacity; matter is that capacity tied up in stable, localized defects; a particle's mass measures the entanglement its defect commits; and gravity is the capacity strain the surrounding medium carries once that commitment is made. General relativity is not assumed underneath the picture but recovered as its low-energy geometry, and the extra gravity seen around galaxies, usually read as dark matter, is the long-range reach of the same medium rather than a separate substance.

The central result is that this picture is not free to be adjusted. Once the finite-capacity substrate and the three postulates are granted, the static weak-field sector follows from a finite counting problem with nothing left to tune. The minimal tetrahedral ensemble fixes the sharing entropy; the electron anchor fixes the substrate length; the edge kernel, loop dressing, continuum matching, and Green-matched source map fix the coefficients of the scalar EFT; and the weak-field bridge turns the capacity deficit into the gravitational potential. The rest of the paper rests on that chain, the only part offered as closed.

Several quantities then fall out of the one chain without separate tuning. Newtonian gravity is its point-source limit. The galactic acceleration scale is the horizon thermal scale reduced by the sharing factor, and the radial-acceleration relation follows from it with no per-galaxy fitting. Lensing and dynamics share a single potential at leading order, so galactic support does not come at the cost of a lensing inconsistency. The same substrate length implies a value for Newton's constant within about one percent of the measured one; we read this as a check that a single length propagates correctly through the whole chain, not as a second independent measurement, and Appendix~L records its provenance. The charged-lepton mass ratios come out to better than a percent from the same shell algebra, as a cross-check on the mass--entropy ontology rather than as part of the gravitational chain.

The medium reaches further, into territory the paper holds more tentatively. Transport gives the field a finite propagation speed and the lag that clusters and mergers need. The cluster sector reads the lensing anomalies as a phase-dependent projection of an already-fixed coefficient. The cosmological mode pushes the sound horizon in the direction the Hubble tension wants. In the saturated early universe the medium becomes a conserved committed density that gravitates as pressureless dust, with an abundance that is counted rather than fitted and lands on the measured dark-to-baryonic ratio. The strong-field branch ends the continuum where capacity runs out, recovering the Schwarzschild and Kerr exteriors and the horizon thermodynamics. These are extensions and frontier completions, not closed results, and the closure table marks the distinction.

Part of the microscopic picture has now been run as well as written. Coupled at full strength to a dynamical simplicial geometry, the admissibility ensemble orders its label sector as the closure weight demands and leaves the host geometry intact; cells pinned into persistent closure failure dent the capacity state and the geometry of their immediate neighborhood where a matched placebo does not; and the long-range sector is shown to require a conserved capacity current with maintenance sinks, from which Newton's form follows structurally; the direct lattice measurement of that current is in progress (Section~23). The same calculations settle where the vacuum comes from: nothing the theory specifies lets the substrate rank geometries, so the smooth four-dimensional state is held by the conditioning of Postulate~III, and the background it selects carries a measured, predicted contribution from the substrate's own free energy. These first dynamical results are graded plainly: single-seed, short-range, on a host whose phase structure is still being mapped. The ledger of Section~24 carries them at that grade.

Many-Pasts carries more weight here than an interpretation of quantum mechanics usually does. It is the claim that the present state of the medium is supported by many admissible microscopic pasts, weighted by how well each leads to the present. In the operational branch used throughout, it changes nothing a laboratory measures: Born-rule statistics and no-signaling both hold, and it supplies the account of the arrow of time. But the same weighting also makes the electron's dressing memoryless, and that memorylessness is one of the two conditions that fix the substrate length. The postulate therefore does real work well beyond the foundations of quantum theory: it fixes the scale that sets gravity, and, with the lattice exclusions of Section~23 in hand, it is the only supplier of the smooth four-dimensional vacuum the theory contains.

The remaining work is specific and limited. The finite-loop self-energy should be confirmed by an independent graph-level calculation. The cosmological sector should be carried through a full Boltzmann likelihood. The cluster lift should have its coherence-growth profile derived and its lensing-map test performed. The strong-field boundary should have its non-universal spectroscopy worked out. The lattice program should complete its propagation measurement --- the emergent maintenance flux linear in commitment, the deficit profile on the graph-Coulomb form, a screening mass consistent with zero --- and replicate the defect response across seeds and across the host's phase diagram. And the substrate length's recurrence derivation, which rested on identifying one phase cycle per dressing recurrence, is now discharged in Appendix~H into a recurrence rate, a unit committed charge per pass, and a condensate phase relation; the residual that remains is the sharper Fork~A identification --- that the winding phase is the condensate's --- together with the single integer fixing the electron as the unit-winding defect, both exposed to the charged-lepton ladder. The paper states these plainly, because its value, if it has one, lies in making them sharp enough to settle.

The case for the picture is its economy: one finite-capacity medium, counted once in the ultraviolet, fixes the coefficients in advance rather than fitting them one at a time. That same economy is its exposure --- with nothing left to adjust in the closed sector, one clean contrary measurement would settle it. The aim of the paper has been to make that confrontation possible.

\clearpage
\section*{Appendix A: Symbol Dictionary and Canonical Conventions}

Appendix A gathers the conventions used throughout the technical material that follows, fixing the units, field definitions, and couplings in one place before the denser calculations begin.

\paragraph{Plain-language terms.}
Several physical words recur throughout the paper and are collected here in plain form before the symbols:
\begin{itemize}
\item \emph{Capacity} --- the entanglement support locally available in the vacuum medium.
\item \emph{Defect} --- a stable, localized commitment of that capacity; coarse-grained, a particle.
\item \emph{Deficit} --- capacity no longer freely available near a defect, $\delta S=S_\infty-S_{\mathrm{ent}}$.
\item \emph{Capacity strain} --- the extended deficit profile whose fractional value gives the weak-field potential and whose gradient gives the gravitational field.
\item \emph{Committed capacity} --- capacity locked into transferring on behalf of a defect; the conserved carrier of the saturated phase (Section~18.5).
\item \emph{Dressing} --- the ground-state cloud an elementary defect builds by resolving the boundary sectors; its coherent support fixes the defect's length scale.
\item \emph{Saturation} --- the regime in which the capacity bath has no slack left, with the available transfer channel at its ceiling.
\item \emph{Pinned reading} --- the statement that the saturated phase holds exactly at that ceiling and stays there.
\item \emph{Admissibility} --- the weighting that favors boundary configurations close to a regular, isotropic local cell.
\item \emph{Closure} --- the condition that the four oriented face data sum to zero, so the cell closes into a regular volume; $K^2$ measures the failure of closure.
\item \emph{Many-Pasts} --- the postulate that the present is supported by a probability-weighted ensemble of admissible microscopic pasts, with weight $e^{-D(H,P)}$.
\end{itemize}

\subsection*{A.1 Units, signature, and entropy normalization}

All dimensional quantities are expressed in SI units unless noted otherwise. The metric signature is $(-,+,+,+)$. Entropies are measured in nats, so Boltzmann's constant is absorbed into the entropy normalization. The canonical UV cell has spatial scale $L_*$ and volume $V_*=L_*^3$. In the canonical branch $L_*$ is fixed by faithful sector resolution,
\[
L_*=\frac32\,\lambda_e e^{-7g_{\mathrm{share,eff}}},
\qquad
\lambda_e=\frac{\hbar}{m_ec}.
\]
The conventional Planck length $L_P=\sqrt{\hbar G/c^3}$ is used only as a comparison scale or in standard gravitational thermodynamic expressions after the gravitational scale has been identified.

These conventions matter because the argument repeatedly moves between a dimensionless UV counting problem and a dimensionful continuum EFT. The units and signature make those two descriptions genuinely comparable.

\subsection*{A.2 Core scalar variables}

The canonical continuum variable is the vacuum-relative coarse-grained entanglement field
\[
S_{\mathrm{ent}}(x),
\]
with vacuum baseline $S_\infty$ and deficit
\[
\delta S(x)=S_\infty-S_{\mathrm{ent}}(x).
\]
For nonlinear work the bounded occupancy fraction is
\[
q(x)=\frac{S_{\mathrm{ent}}(x)}{S_\infty}=1-\frac{\delta S}{S_\infty}\in[0,1].
\]
The absolute entropy unit is fixed only after choosing a cell or horizon normalization. Under a constant rescaling of $S_{\mathrm{ent}}$, the quantities $S_\infty$ and $\kappa/\gamma$ rescale together, leaving $\delta S/S_\infty$ and $\kappa/(\gamma S_\infty)$ invariant.
The source channel is
\[
\chi(x)=-\frac{T^\mu{}_\mu}{c^2},
\]
which is the continuum trace channel of the localized defect sector and reduces to the ordinary mass density $\rho$ in the nonrelativistic static limit.

\subsection*{A.3 Couplings and derived observables}

The main-text conventions are
\begin{align}
\gamma &:\ \text{entanglement-field stiffness},\\
\kappa &:\ \text{continuum defect--entropy coupling},\\
\kappa_m(\ell) &:\ \text{mass-per-entropy map at scale }\ell,\\
L_* &= \frac32\,\lambda_e e^{-7g_{\mathrm{share,eff}}},\\
G_* &= \frac{c^3L_*^2}{\hbar}
=\frac94\,\frac{\hbar c}{m_e^2}e^{-14g_{\mathrm{share,eff}}},\\
G_{\mathrm{tet}}(0) &:\ \text{tetrahedral on-site Green constant},\\
g_{\mathrm{share,max}} &= \ln(1680),\\
g_{\mathrm{share,eff}} &:\ \text{admissibility-weighted sharing entropy},\\
J_{\mathrm{bare}},\,J_{\mathrm{eff}}^{\mathrm{tree}},\,J_{\mathrm{eff}}^{(\mathrm{ren})} &:\ \text{UV edge couplings},\\
a_0 &= \frac{cH_0 g_{\mathrm{share,eff}}}{4\pi^2}.
\end{align}
The canonical weak-field bridge and Newton closure are
\[
\begin{aligned}
\frac{\Phi}{c^2}=-\frac{\delta S}{2S_\infty},
&\qquad
\frac{\kappa}{\gamma}
=
\frac{3L_*}{4G_{\mathrm{tet}}(0)\kappa_m(L_*)},\\
G&=\frac{c^2\kappa}{8\pi\gamma S_\infty}.
\end{aligned}
\]

Collected in one place, these formulas also make clear which quantities are downstream of the closure chain. The UV data determine the stiffness and source-to-stiffness ratio first; the observable weak-field constants appear after the bridge and fixed $S_\infty$ normalization are applied.

\subsection*{A.4 Notation map}

One notation set is used throughout. The effective sharing entropy is denoted $g_{\mathrm{share,eff}}$, the scalar variable is always the vacuum-relative field $S_{\mathrm{ent}}$ or its deficit $\delta S$, and the weak-field bridge is used in the single form stated above.

\section*{Appendix B: UV Boundary Ensemble and Admissibility Closure}

Appendix~B records the finite ultraviolet counting problem in its explicit form. It supports Sections~5--6 by showing that the seven-state tetrahedral ensemble and the admissibility weighting are finite and auditable rather than phenomenological dials: the theory begins from a discrete boundary ensemble and ends with a unique admissibility-closed entropy, not with an unconstrained continuum ansatz.

\subsection*{B.1 Minimal tetrahedral package}

The canonical UV cell is a tetrahedron with four structural ingredients:
\begin{itemize}
\item a tetrahedral volumetric cell;
\item half-integer fermionic face data on each face;
\item injective face assignment across the four faces;
\item binary orientation/parity.
\end{itemize}
Postulate II identifies the elementary defect sector as fermionic, so each face carries half-integer base spin $j_0$. For a shared face the effective boundary sector is
\[
j_0\otimes j_0 = 0\oplus 1\oplus \cdots \oplus 2j_0.
\]
Postulate I selects the maximum-capacity channel, hence $j_{\mathrm{eff}}=2j_0$ with $|M|=2j_{\mathrm{eff}}+1=4j_0+1$ distinguishable face states. Injectivity across four faces requires $|M|\ge 4$. The $j_0=1/2$ option fails because it gives $j_{\mathrm{eff}}=1$ and $|M|=3$. The first half-integer choice that works is therefore $j_0=3/2$, giving $j_{\mathrm{eff}}=3$ and the canonical seven-state face sector.
The resulting state count is
\[
\Omega_{\mathrm{tet}}=2\times P(7,4)=1680,
\qquad
g_{\mathrm{share,max}}=\ln(1680)=7.42654907240.
\]
This is the minimal discrete package used in the framework to obtain a finite, isotropic, auditable boundary-channel structure.

The important feature is that the counting closes for structural reasons. Fermionic face data, injectivity, and maximum-capacity channel selection together force the seven-state face sector instead of leaving it as a tunable menu choice.

The minimality statement can also be written as a short proof. A volumetric cell in $d=3$ needs at least four faces, so a tetrahedron is the first admissible simplex. The closure surrogate is three-component, so the face sector must be rich enough to support a nontrivial quadratic spectrum in $d=3$ rather than a degenerate one-dimensional label count. Postulate II makes the face data fermionic, hence half-integer. Maximum-capacity channel selection then gives
\[
j_{\mathrm{eff}}=2j_0,
\qquad
|M|=2j_{\mathrm{eff}}+1=4j_0+1.
\]
Injectivity across four faces requires $|M|\ge 4$. The only half-integer option below $j_0=3/2$ is $j_0=1/2$, which gives $j_{\mathrm{eff}}=1$ and $|M|=3$, so it fails. The first admissible fermionic choice is therefore $j_0=3/2$, giving $j_{\mathrm{eff}}=3$ and the canonical seven-state face sector. In that precise sense, the $(4\text{-face},7\text{-state})$ tetrahedral package is the minimal architecture compatible with a three-component isotropic closure mode, injective boundary information, and finite volumetric counting.

This also settles the status of $j_0$ as a parameter: it has none to tune. Within the minimal construction $j_0=3/2$ is forced as the smallest fermionic label meeting injectivity. A larger $j_0$ does not describe a fluctuation inside the same cell; it specifies a different, larger boundary ensemble, with a different state count $\Omega_{\mathrm{tet}}$ and a different closure spectrum. The minimal theory therefore fixes $j_0$ rather than leaving it open, and the seven-state sector is not one option among a family but the first that closes.

\subsection*{B.2 Closure invariant, kernel, and unique fixed point}

The canonical scalar closure invariant is
\[
K^2(b)=48-\frac{1}{3}\bigl(S^2-\Sigma^2\bigr),
\qquad
S=\sum_{i=1}^4 m_i,
\qquad
\Sigma^2=\sum_{i=1}^4 m_i^2.
\]
The admissibility family is
\[
p_\eta(b)=\frac{1}{Z(\eta)}e^{-\eta K^2(b)},
\qquad
Z(\eta)=\sum_{b\in B}e^{-\eta K^2(b)}.
\]
The admissibility precision $\eta$ is fixed by stationary normalized closure evidence. The closure constraint is the vanishing of the three-component oriented-face sum $\mathbf c(b)=\sum_i\hat n_i m_i\in\mathbb R^3$, of which $K^2$ is the quadratic invariant; marginalizing the three continuous closure components against the family at precision $\eta$ contributes a determinant factor $\eta^{3/2}$, so the normalized closure evidence is $\eta^{3/2}Z(\eta)$ and its logarithm is
\[
\mathcal F(\eta)=\ln Z(\eta)+\tfrac32\ln\eta.
\]
With $\partial_\eta\ln Z=-\langle K^2\rangle_\eta$, the stationary condition $\mathcal F'(\eta)=0$ is the closure relation
\[
\langle K^2\rangle_\eta=\frac{3}{2\eta},
\]
so the factor $3/2$ is the determinant weight of the three independent closure components, not an equipartition rule imported onto the bounded discrete spectrum; the discreteness, multiplicities, and positive floor of the spectrum remain inside the exact finite sum $Z(\eta)$. The second derivative is $\mathcal F''(\eta)=\mathrm{Var}_\eta(K^2)-3/(2\eta^2)$, and on the exact spectrum $\mathcal F'$ has a single interior zero,
\[
\eta_*=0.0298668443935,
\]
at which $\mathrm{Var}_{\eta_*}(K^2)=15.69$ is dwarfed by $3/(2\eta_*^2)=1681.6$, giving $\mathcal F''(\eta_*)=-1665.9<0$: $\eta_*$ is the unique local maximum of the normalized closure evidence.
Because the parity-symmetric ensemble is finite, both $Z(\eta)$ and $\langle K^2\rangle_\eta$ are exact finite sums over the spectrum. The distinct closure-defect values and their degeneracies are
\[
\begin{array}{c|cccccc}
K^2 &
\frac{122}{3} & \frac{134}{3} & \frac{142}{3} & \frac{146}{3} & \frac{152}{3} & \frac{154}{3} \\ \hline
\mathrm{mult} &
96 & 96 & 96 & 288 & 192 & 144
\end{array}
\]
\[
\begin{array}{c|ccccc}
K^2 &
\frac{158}{3} & 54 & \frac{164}{3} & \frac{166}{3} & \frac{170}{3} \\ \hline
\mathrm{mult} &
384 & 192 & 48 & 96 & 48
\end{array}
\]
with total multiplicity $1680$ as required. In particular,
\[
Z(\eta)=\sum_a n_a e^{-\eta K_a^2},
\qquad
\langle K^2\rangle_\eta=\frac{\sum_a n_a K_a^2 e^{-\eta K_a^2}}{\sum_a n_a e^{-\eta K_a^2}},
\]
where $(K_a^2,n_a)$ run over the table above. The closed-branch value $\eta_*$ is therefore the unique interior maximum of an exact finite-spectrum functional, not an unseen numerical fit.
The corresponding effective sharing entropy is
\[
g_{\mathrm{share,eff}}
=
-\sum_{b\in B}p_{\eta_*}(b)\ln p_{\eta_*}(b)
=
7.41980002357.
\]
The closed-branch moments used in the UV stiffness discussion are
\[
\langle K^2\rangle_{\eta_*}=50.2229154254,
\qquad
\mathrm{Var}_{\eta_*}(K^2)=15.6889750078,
\qquad
a_{\mathrm{UV}}=0.0637390269.
\]
These values quantify the local stiffness of the canonical closure point rather than a tunable phenomenological uncertainty.
The closure-saturation product is
\[
C_{\rm cl}:=\eta_*\langle K^2\rangle_{\eta_*}=\frac32,
\qquad
C_{\rm cl}^{-1}=\frac23.
\]
This same reciprocal $2/3$ reappears as the tetrahedral transverse export fraction in Appendix~C and as the global correction in the faithful sector-resolution scale setting.

The admissibility parameter stops being free here. The kernel introduces $\eta$, and the closure condition removes its arbitrariness again by demanding that the fluctuation scale produced by the weighting agree with the weighting itself.

\subsection*{B.3 Rooted reduction and local benchmarks}

Rooting on the shared face reduces the exact parity-symmetric ensemble to $140$ rooted microstates and $69$ rooted closure classes. The rooted classes can be labeled by $\alpha=(m_\bullet,K^2)$, so the same reduced state space already supports the local evaluation, the cavity benchmark, and the later shell propagation. The local information observable
\[
\sigma_{\mathrm{ind}}^{(r)}=\frac{H(X\mid Y_r)}{H(X)}
\]
has the principal pre-nonlocal benchmarks
\begin{align}
\sigma_{\mathrm{ind}}^{\mathrm{toy}} &= 0.44997,\\
\sigma_{\mathrm{ind}}^{\mathrm{loc}} &= 0.44708,\\
\sigma_{\mathrm{ind}}^{\mathrm{Bethe}}(J=0) &= 0.44749.
\end{align}
Here the Bethe value is the homogeneous cavity evaluation on the $69\times 69$ rooted-class interaction graph at zero transport coupling,
\[
\mu_\alpha
\propto
w_\alpha
\left(\sum_\beta U_{\alpha\beta}(0)\mu_\beta\right)^{z-1},
\qquad
\sum_\alpha \mu_\alpha=1,
\]
with $z=4$ and $U_{\alpha\beta}(0)$ the rooted shared-face compatibility matrix before shell transport is turned on. In other words, $\sigma_{\mathrm{ind}}^{\mathrm{Bethe}}(J=0)$ is the cavity-theory benchmark of the same explicit rooted ensemble, not a disconnected numerical insert.
The horizon target implied by the effective sharing entropy is
\[
\sigma_*=\frac{\pi}{g_{\mathrm{share,eff}}}=0.42340665.
\]
The gap between the local benchmarks and $\sigma_*$ is therefore a genuinely shell / loop problem rather than a failure of the local admissibility closure.

That separation matters for the later UV story. It means the remaining work is not to repair the local closure ensemble, but to propagate it more accurately through transport and return structure.

\subsection*{B.4 What is fixed at this stage}

By the end of the admissibility calculation, the framework has already fixed the microscopic counting ceiling, the unique closure point, the effective sharing entropy, and the local stiffness moments. What remains for the next appendix is not another entropy choice, but the propagation of those quantities into edge transport, finite renormalization, and continuum normalization.

\section*{Appendix C: Edge Kernel, Finite Renormalization, and Continuum Matching}

Appendix~C carries the local boundary ensemble of Appendix~B into edge transport, loop dressing, and the continuum stiffness coefficient of the weak-field EFT.

\subsection*{C.1 Channel-averaged isotropy identity and tree coupling}

Let $\hat n_i$ be the four face normals of a regular tetrahedron. The exact identity
\[
\sum_{i=1}^4 \hat n_i\hat n_i^{\mathsf T}=\frac{4}{3}I_3
\]
implies a channel-averaged transverse fraction of $2/3$. The bare edge stiffness is therefore
\[
J_{\mathrm{bare}}=\frac{2}{3}\eta_*=0.0199112296.
\]
For a rooted $z=4$ coarse adjacency graph, the tree-to-lattice map gives
\[
J_{\mathrm{eff}}^{\mathrm{tree}}=\frac{J_{\mathrm{bare}}}{z-1}=\frac{2\eta_*}{9}=0.0066370765.
\]

This is the first place where local closure data become a transport law. The tetrahedral identity fixes the isotropic projection, and the rooted branching structure determines how much of the microscopic edge penalty survives as net outward propagation on the coarse graph.

\subsection*{C.2 Horizon target and shell convergence}

The horizon-capacity target is
\[
\sigma_*=\frac{\pi}{g_{\mathrm{share,eff}}}=0.42340665.
\]
At the derived coupling the explicit shell values are
\[
\sigma_{\mathrm{ind}}^{(2)}=0.42143,
\qquad
\sigma_{\mathrm{ind}}^{(3)}=0.42166,
\qquad
\Delta_{2\to 3}=0.00023.
\]
The residual shift from the target is already small and stable by shell depth $r=2$, isolating the remaining correction to the loopy local-return sector rather than a broad nonlocal ambiguity.

So the shell calculation narrows the open problem substantially. The tree branch already lands very near the target, and the residual discrepancy can be assigned specifically to local returns rather than to an uncontrolled long-range correction.

\subsection*{C.3 Finite-loop self-energy closure}

The leading loopy correction is organized as a local Dyson dressing:
\[
J_{\mathrm{eff}}^{(\mathrm{ren})}
=
\frac{J_{\mathrm{eff}}^{\mathrm{tree}}}{1+J_{\mathrm{eff}}^{\mathrm{tree}}\Sigma_{\mathrm{ret}}}.
\]
The structural decomposition is
\[
\Sigma_{\mathrm{ret}}=7+\frac{2}{9}=\frac{65}{9},
\]
and each term has a concrete return-channel origin. A short return motif leaves a shared face, explores a local closed loop, and re-enters the same coarse edge before contributing to net long-range transport. In the canonical label basis $m=-3,-2,\ldots,3$, there are exactly seven ways to do this without changing sector. These are the seven sector-diagonal returns, one for each face-label channel, and together they contribute
\[
\mathrm{Tr}(I_7)=7.
\]
In addition to these label-preserving loops, permutation symmetry allows one collective mode shared across all channels. Writing
\[
P_{\mathrm{sing}}=|u\rangle\langle u|,
\qquad
u=\frac{1}{\sqrt 7}(1,1,\ldots,1),
\]
this shared return is rank one. Any additional off-diagonal return sector would break the permutation symmetry of the canonical local ensemble, so there is no second independent collective channel to count. Only the transverse scalar branch feeds back into the coarse transport law, so the singlet first acquires the same $2/3$ projection factor that appeared in the tree coupling. It is then diluted by the rooted branching factor $1/(z-1)=1/3$ on the $z=4$ graph, because only one of the three outward branches returns to the original edge. The collective contribution is therefore
\[
\mathrm{Tr}\!\left(\frac{2}{3}\frac{1}{3}P_{\mathrm{sing}}\right)=\frac{2}{9},
\]
since $\mathrm{Tr}(P_{\mathrm{sing}})=1$. Equivalently,
\[
R_{\mathrm{ret}}=I_7+\frac{2}{9}P_{\mathrm{sing}},
\qquad
\Sigma_{\mathrm{ret}}=\mathrm{Tr}(R_{\mathrm{ret}})=7+\frac{2}{9}.
\]
This is the sense in which the finite-loop coefficient is counted rather than guessed: seven independent label-preserving returns plus one shared singlet return with exactly the same projection and branching weights already fixed in the tree map. Hence
\[
c_{\mathrm{loop}}^{(\mathrm{ren})}
\equiv
\frac{J_{\mathrm{eff}}^{(\mathrm{ren})}}{J_{\mathrm{eff}}^{\mathrm{tree}}}
=
\frac{1}{1+J_{\mathrm{eff}}^{\mathrm{tree}}\Sigma_{\mathrm{ret}}}
\approx 0.95426,
\]
and
\[
J_{\mathrm{eff}}^{(\mathrm{ren})}\approx 0.00633348.
\]
This reproduces the shell-target crossing near $J_{\mathrm{bare,cross}}\sim 0.019$ at the stated level of agreement.

The local Dyson dressing is therefore doing one precise job: it corrects the tree branch by accounting for the short motifs that recycle amplitude before it contributes to true coarse transport. The renormalized coupling is not a new parameter, but the tree coupling after local returns have been summed.

\subsection*{C.4 Euclidean-action normalization and continuum stiffness}

The lattice quadratic form is interpreted canonically as a Euclidean action weight,
\[
\frac{I_E}{\hbar}
=
\frac{J_{\mathrm{eff}}^{(\mathrm{ren})}}{2}
\sum_{a,i}(Q_a-Q_{a+L_*\hat n_i})^2,
\]
where the sum runs over all sites $a$ and all four outgoing nearest-neighbor directions $\hat n_i$ (each nearest-neighbor edge counted twice). The microscopic four-cell is
\[
\Delta V_4=\frac{L_*^4}{c}.
\]
The same tetrahedral identity then yields
\[
\gamma_Q=\frac{4\hbar c}{3L_*^2}J_{\mathrm{eff}}^{(\mathrm{ren})}
\]
for the occupancy field $Q_{\mathrm{occ}}$. With the horizon-capacity normalization
\[
S=\pi Q_{\mathrm{occ}},
\]
the canonical convention $\frac{\gamma}{2}(\partial S)^2$ gives
\[
\gamma
=
\frac{4\hbar c}{3\pi^2L_*^2}J_{\mathrm{eff}}^{(\mathrm{ren})}
=
\frac{4\hbar c}{3\pi^2L_*^2}
\frac{2\eta_*/9}{1+(2\eta_*/9)(65/9)}.
\]
Using the substrate-induced scale
\[
G_*:=\frac{c^3L_*^2}{\hbar},
\]
this is
\[
\gamma
=
\frac{4J_{\mathrm{eff}}^{(\mathrm{ren})}}{3\pi^2}\frac{c^4}{G_*}
\approx
8.556\times10^{-4}\frac{c^4}{G_*}.
\]

This is the decisive stiffness-side matching step. Up to here the derivation has produced a dimensionless lattice weighting; after Euclidean normalization and faithful sector-resolution scale setting, that same weighting becomes the dimensionful continuum stiffness that appears in the weak-field action.

\subsection*{C.5 Local defect insertion and the source-side lattice constant}

The stiffness-side matching is not the only UV quantity that can be closed locally. For the canonical rigid defect insertion, excluding one of the seven admissible face labels from one face removes exactly one-seventh of the isotropically averaged local partition weight. Therefore the logarithm of the isotropically averaged partition ratio is exactly
\[
\Delta S_{\mathrm{def}}
:=
-\ln\!\Bigl\langle \frac{Z_{\mathrm{def}}}{Z_{\mathrm{vac}}}\Bigr\rangle_{\mathrm{iso}}
=
\ln\!\frac{7}{6}.
\]
This is the exact isotropic source benchmark in the canonical seven-label ensemble. The isotropically averaged defect free-energy cost differs from it only at $O(10^{-5})$ because the admissibility weighting breaks label symmetry only weakly.

The local source benchmark $\ln(7/6)$ should not be confused with the elementary fermionic one-bit anchor $\ln2$. The former is the isotropically averaged partition-ratio shift produced by removing one admissible label from the seven-label boundary ensemble. The latter is the intrinsic binary entropy of an elementary occupied/unoccupied fermionic face-exclusion defect. The source theorem uses $\ln(7/6)$ to normalize the local scalar insertion into the lattice response, while the electron anchor uses $\ln2$ to fix the mass--entropy unit of the elementary fermionic defect.

To propagate that local defect into the lattice field equation one needs the on-site Green function of the tetrahedral/diamond nearest-neighbor Laplacian. The dual graph of face-sharing tetrahedral cells is the four-valent diamond graph. Eliminating the two-sublattice structure gives the standard Brillouin-zone representation for the diamond lattice Green function \citep{Guttmann2010}
\[
G_{\mathrm{tet}}(0)
=
\frac{1}{(2\pi)^3}
\int_{[-\pi,\pi]^3}
\frac{4\,d^3k}
{16-\left|1+e^{ik_1}+e^{ik_2}+e^{ik_3}\right|^2}.
\]
This integral is the reproducible source-side lattice constant: it is the self-energy of a unit point insertion for the same scalar mode whose long-wavelength stiffness was matched in Appendix~C.4. Joyce's exact evaluation of the diamond-lattice Green function, in this normalization, gives \citep{Joyce1994,Guttmann2010}
\[
G_{\mathrm{tet}}(0)
=
\frac{3\Gamma(1/3)^6}{2^{14/3}\pi^4}
=
0.4482203943883814\ldots .
\]
Thus the source-side graph constant is an exact lattice invariant rather than a fitted numerical coefficient. Direct quadrature with endpoint extrapolation reproduces the same value.

\paragraph{Green-tensor transverse response.}
The same graph response also fixes the transverse export weight $w_\perp$, the single graph quantity that propagates downstream into the horizon channel count of Appendix~F.5 and into the global $\ln(3/2)$ closure-saturation factor of the scale-setting relation. We compute it directly from the four nearest-neighbor bond frame, with no fitting freedom.

Let $d_i$, $i=1,\ldots,4$, be the four diamond nearest-neighbor bond directions,
\[
d_1=\frac{(1,1,1)}{\sqrt3},\quad
d_2=\frac{(1,-1,-1)}{\sqrt3},\quad
d_3=\frac{(-1,1,-1)}{\sqrt3},\quad
d_4=\frac{(-1,-1,1)}{\sqrt3}.
\]
They obey
\[
\sum_{i=1}^4 d_i^a d_i^b=\frac43\delta^{ab}.
\]
Distributing the scalar on-site Green response over the tetrahedral bond frame gives
\[
{\cal G}_{\rm loc}^{ab}
=
G_{\mathrm{tet}}(0)\frac14\sum_i d_i^a d_i^b
=
\frac{G_{\mathrm{tet}}(0)}{3}\delta^{ab}.
\]
For a local horizon normal $\hat r$,
\[
P_\perp^{ab}=\delta^{ab}-\hat r^a\hat r^b,
\]
so
\[
G_\perp=P_\perp^{ab}{\cal G}_{\rm loc}^{ab}
=
\frac23G_{\mathrm{tet}}(0).
\]
Thus the transverse export weight $w_\perp=2/3$ is the transverse part of the exact local graph response.

Using the field normalization $S=\pi Q_{\mathrm{occ}}$, the rigid local defect shift is
\[
\delta Q_{\mathrm{def}}=\frac{\Delta S_{\mathrm{def}}}{\pi}
=
\frac{\ln(7/6)}{\pi}.
\]
The corresponding local source amplitude in lattice units is therefore
\[
s_{\mathrm{def}}
=
J_{\mathrm{eff}}^{(\mathrm{ren})}
\frac{\delta Q_{\mathrm{def}}}{G_{\mathrm{tet}}(0)},
\]
so that
\[
\frac{s_{\mathrm{def}}}{J_{\mathrm{eff}}^{(\mathrm{ren})}}
=
\frac{\ln(7/6)}{\pi G_{\mathrm{tet}}(0)}
=
0.109472228\ldots
\]
is a pure number fixed by the same UV lattice geometry.

The Green-function constant turns the local insertion into a continuum source theorem. Defining the defect-entropy density by
\[
\sigma_{\mathrm{def}}=\frac{\rho}{\kappa_m(L_*)},
\]
the tetrahedral projection used in the stiffness mapping gives
\[
\nabla^2\delta S
=
-\frac{3L_*}{4G_{\mathrm{tet}}(0)}\,\sigma_{\mathrm{def}}.
\]
Equating this with the weak-field source equation
\[
\nabla^2\delta S=-\frac{\kappa}{\gamma}\rho
\]
closes the canonical source-to-stiffness ratio:
\[
\boxed{
\frac{\kappa}{\gamma}
=
\frac{3L_*}{4G_{\mathrm{tet}}(0)\kappa_m(L_*)}
}.
\]
This is the source-side counterpart of the stiffness derivation. The edge-kernel calculation fixes how the scalar capacity mode resists gradients; the Green-matched defect calculation fixes how localized matter defects source that same mode.

\paragraph{The three cancellations, made explicit.}
The passage from the dimensionless lattice insertion $s_{\mathrm{def}}/J_{\mathrm{eff}}^{(\mathrm{ren})}=\ln(7/6)/(\pi G_{\mathrm{tet}}(0))$ to the source theorem turns on three reductions, none of which leaves a free constant.
\emph{(i) The factor $\pi$ cancels against the field normalization.} The insertion is written for the occupancy field, where one defect shifts $\delta Q_{\mathrm{def}}=\ln(7/6)/\pi$. Converting to the entropy field through $S=\pi Q_{\mathrm{occ}}$ multiplies by $\pi$, so the $S$-field shift is $\delta S_{\mathrm{def}}=\pi\,\delta Q_{\mathrm{def}}=\ln(7/6)$ and the explicit $\pi$ does not survive into the source theorem.
\emph{(ii) $\ln(7/6)$ is carried by the entropy density, not the coefficient.} A defect contributes entropy $\ln(7/6)$ at the isotropic benchmark, so a defect number density sources an entropy density proportional to it; writing that source as $\sigma_{\mathrm{def}}=\rho/\kappa_m(L_*)$ folds the benchmark and the mass--entropy unit into the single density $\sigma_{\mathrm{def}}$. The geometric coefficient $3L_*/(4G_{\mathrm{tet}}(0))$ that multiplies it is fixed by the tetrahedral $4/3$ projection and the cell length alone and contains no $\ln(7/6)$.
\emph{(iii) $J_{\mathrm{eff}}^{(\mathrm{ren})}$ cancels in the ratio.} The source amplitude $s_{\mathrm{def}}$ and the stiffness $\gamma$ each carry one power of the renormalized edge coupling, so it cancels in the source-to-stiffness ratio $\kappa/\gamma$, leaving the lattice-geometric constant $3L_*/(4G_{\mathrm{tet}}(0)\kappa_m(L_*))$. The renormalized loop coupling therefore drops out of the observable normalization.

Equivalently, the same source closure fixes the continuum quantity $\Xi_\rho$ appearing in
\[
\kappa=\frac{\Xi_\rho}{L_*^2\kappa_m(L_*)}.
\]
The absolute scale of $S_\infty$ remains a fixed-epoch normalization convention in the bridge law, not a residual freedom in the source projection. In the cell-normalized gauge natural to the local source theorem, one may write
\[
S_{\infty}^{\rm cell}
=
\frac{3\ln2}{32\pi G_{\mathrm{tet}}(0)}
=
0.0461482516\ldots .
\]
In a horizon-normalized gauge, $S_\infty$ instead carries the much larger apparent-horizon capacity. These are not two physical constants. A constant rescaling of the entropy field rescales $S_\infty$ and $\kappa/\gamma$ together, leaving $\kappa/(\gamma S_\infty)$ and hence $G$ unchanged. Thus the weak-field source map is closed in the canonical branch once the mass--entropy map $\kappa_m(L_*)$, the tetrahedral Green constant, and the standard cell convention are specified.

\subsection*{C.6 Local susceptibility cross-check}

The exact local moments of the admissibility-closed ensemble supply an independent non-degeneracy check on the source-side result. From the variance in Appendix~B,
\[
a_{\mathrm{UV}}
:=
\frac{1}{\mathrm{Var}_{\eta_*}(K^2)}
=
0.0637390269,
\]
which is the local zero-mode inverse susceptibility of the closure scalar. Two features of this number matter for the source theorem in C.5. First, $a_{\mathrm{UV}}$ is finite and positive, confirming that the closed branch at $\eta_*$ has a non-degenerate zero-mode response rather than a critical singularity that would invalidate the linear Green-function matching used to derive $\kappa/\gamma$. Second, the same local susceptibility controls the stability of the closure mode that is being propagated into the shell and loop calculations, so the residual fractional discrepancy between the local benchmark and $\sigma_*$ in C.2 is genuinely loop-sector rather than a local-closure failure.

The actual closure of $\kappa/\gamma$ in the canonical branch is the Green-matched theorem in Appendix~C.5. The role of $a_{\mathrm{UV}}$ here is restricted to confirming non-degeneracy of the local mode being matched.

\subsection*{C.7 UV-to-IR payoff}

At this stage the weak-field UV coefficient chain is explicit:
\[
\Omega_{\mathrm{tet}}
\rightarrow
g_{\mathrm{share,eff}}
\rightarrow
L_*
\rightarrow
J_{\mathrm{bare}}
\rightarrow
J_{\mathrm{eff}}^{\mathrm{tree}}
\rightarrow
\Sigma_{\mathrm{ret}}
\rightarrow
J_{\mathrm{eff}}^{(\mathrm{ren})}
\rightarrow
\gamma.
\]
The same chain feeds
\[
a_0=\frac{cH_0g_{\mathrm{share,eff}}}{4\pi^2},
\]
and faithful sector resolution gives the substrate-induced scale
\[
G_*=\frac{c^3L_*^2}{\hbar}
=
\frac94\,\frac{\hbar c}{m_e^2}e^{-14g_{\mathrm{share,eff}}}.
\]
and the Green-matched source theorem fixes
\[
\frac{\kappa}{\gamma}
=
\frac{3L_*}{4G_{\mathrm{tet}}(0)\kappa_m(L_*)}.
\]
The weak-field bridge then converts this source-to-stiffness ratio into the observed Newtonian normalization through the invariant combination $\kappa/(\gamma S_\infty)$. The comparison with the same $G_*$ fixed by faithful sector resolution tests that this normalization propagates the substrate scale coherently; it is not a disjoint-input second derivation of $G$, since the matched route carries $L_*$ within it.

The coefficients this appendix supplies to the main weak-field chain are $J_{\mathrm{bare}}$, $J_{\mathrm{eff}}^{\mathrm{tree}}$, $\Sigma_{\mathrm{ret}}$, $J_{\mathrm{eff}}^{(\mathrm{ren})}$, $\gamma$, and the Green-matched source projection $\kappa/\gamma$. The remaining uses of $S_\infty$ belong to the fixed normalization of the weak-field bridge, not to the source sector itself.

\section*{Appendix D: Weak-Field Technical Derivations, Electron Anchor, and EFT Consistency}

Appendix~D collects the weak-field derivations that are central but too dense for the main line: the bridge law, Newtonian recovery, the electron anchor, and the EFT consistency checks.

\subsection*{D.1 Bridge-law uniqueness}

The weak-field bridge is derived once and then used throughout. Let the lapse be written as
\[
N=e^{-F(\delta S/S_\infty)}.
\]
Additivity of independent deficits requires $F(x+y)=F(x)+F(y)$, so continuity implies $F(x)=cx$. Standard weak-field metric normalization fixes $c=1/2$, giving
\[
N=e^{-\delta S/(2S_\infty)}
\]
and therefore
\[
\frac{\Phi}{c^2}=-\frac{\delta S}{2S_\infty}
\]
to leading order. This is the unique weak-field bridge compatible with locality, additive independent deficits, and multiplicative redshift composition.

Writing the argument this explicitly removes one of the most common ambiguities in modified-gravity proposals. The bridge from entropy deficit to gravitational potential is not being chosen phenomenologically after the fact; it is fixed by the structural requirements of the weak-field limit itself.

\subsection*{D.2 Point source, Newton limit, and lensing}

In the renormalized static branch,
\[
\nabla^2\delta S=-\frac{\kappa}{\gamma}\rho.
\]
For a point source $M$,
\[
\delta S(r)=\frac{\kappa M}{4\pi\gamma r},
\qquad
g(r)=\frac{c^2\kappa}{8\pi\gamma S_\infty}\frac{M}{r^2}
=
\frac{GM}{r^2}.
\]
Because the leading entanglement stress carries no anisotropic stress,
\[
\Phi=\Psi
\]
at the order treated. The effective-halo rewrite is
\[
\rho_{\mathrm{halo}}(r)=\frac{1}{4\pi Gr^2}\frac{d}{dr}\Big[r^2(g_{\mathrm{obs}}-g_{\mathrm{bar}})\Big].
\]
Thus the same deficit field controls both orbital dynamics and light bending in the leading weak-field regime.

That shared control is the key weak-field consistency test. A viable branch must not reproduce galactic support only by sacrificing lensing, and the scalar deficit sector avoids that failure at the order treated.

\subsection*{D.3 Electron anchor and composite matter}

The canonical fermionic entropy increment is
\[
\Delta S_f=\ln 2.
\]
The UV mass normalization is
\[
\kappa_{m,\mathrm{UV}}=\frac{\hbar}{cL_*}\frac{1}{\ln 2},
\]
and the running law in the closed branch is
\[
\kappa_m(\ell)=\kappa_{m,\mathrm{UV}}\left(\frac{L_*}{\ell}\right)^{1+\alpha_{\mathrm{cl}}},
\qquad
\alpha_{\mathrm{cl}}=0.
\]
At the electron Compton scale $\ell=\lambda_e$ this gives
\[
\kappa_m(\lambda_e)=\frac{m_e}{\ln 2},
\]
which is the clean elementary anchor used here. Composite hadrons are not reduced to a bare constituent count. Their mass budget is assigned to a dressed bound-state entropy
\[
m_{\mathrm{hadron}}=\kappa_m(\ell_H)S_{\mathrm{ent},H}^{\mathrm{dressed}},
\]
whose microscopic decomposition must include confinement, gluonic structure, trace-anomaly contributions, and chiral vacuum reorganization.

The contrast between the two sectors is deliberate. The electron is a clean one-bit defect anchor; hadrons are not. Their inertial content must therefore be assigned to a dressed entropy budget rather than to a naive constituent count.

\subsection*{D.4 Faithful sector resolution and induced \texorpdfstring{$G_*$}{G*}}

The electron anchor enters the gravitational normalization through the absolute substrate length, and this subsection makes the link explicit. The task is to turn a dimensionless count --- the effective number of dressing configurations the electron's seven channels resolve --- into a physical length ratio $\lambda_e/L_*$, and hence into the induced gravitational scale $G_*$. The argument has three parts: a recurrence derivation that fixes the linear, first-power map from the seven-channel sharing entropy to the length; a demonstration that the allowed local single-label dynamics cannot reach the full ensemble, so the memoryless refresh of Postulate~III is required; and an explicit transfer-operator realization that makes the count concrete. One identification sits beneath the map --- one phase cycle per dressing recurrence --- which Appendix~H.6 gives an operator form and discharges into a recurrence rate, a unit charge per pass, and a condensate phase relation, conditional on the Fork~A identification settled in H.7; the linear power itself is then cross-checked three independent ways.

\paragraph{The recurrence derivation of the dictionary.}
The support-to-length identification follows from the recurrence structure of the dressing. By Postulate~III the dressing is memoryless in substrate time: each tick, each of the seven channels draws independently from the admissibility ensemble. By the asymptotic-equipartition property --- used here only in its elementary form, that a distribution with entropy $g$ has an effective typical support of size $e^{g}$, so a specified typical configuration appears with probability $e^{-g}$ per independent draw --- a single channel realizes its typical dressing configuration with probability $e^{-g_{\mathrm{share,eff}}}$ per draw, and the seven mutually independent channels (Appendix~H) realize the complete dressing jointly with probability $e^{-7g_{\mathrm{share,eff}}}$. The defect therefore re-completes its dressing once per $e^{7g_{\mathrm{share,eff}}}$ substrate ticks. Identifying the defect's internal phase clock, $\tau_e=\hbar/m_ec^2$, with this recurrence period---one phase cycle per complete re-dressing---and applying the transverse export factor of Appendix~C.5 gives
\[
\frac{\tau_e}{\tau_*}=\frac{2}{3}\,e^{7g_{\mathrm{share,eff}}},\qquad \tau_*=\frac{L_*}{c},
\]
the length-setting relation below in temporal form. The linearity of the support exponent is forced rather than chosen: probabilities of independent rare events multiply, so their exponents add, and the waiting time for the joint event is the exponential of the sum. The three-way bracketing of the exponent at unity within $10^{-3}$, summarized in Section~13 and tested again in Appendix~I.1, is the quantitative check of this derivation. One identification sits beneath it---that the defect's phase advances one cycle per dressing recurrence---which Appendix~H.6 decomposes into the recurrence rate, a unit committed charge per pass, and the condensate phase relation, relocating the residual to the Fork~A identification (the winding phase is the condensate's) and the single integer the lepton ladder tests.

\paragraph{The transfer-support dual.}
The same relation follows from a spatial reading. The entropy of the one-sector transfer distribution fixes an effective support size $N^{(1)}_{\rm eff}=e^{g_{\mathrm{share,eff}}}$, and on the lightest charged branch the seven sectors carry no mutual information (Appendix~H), so the joint support is $N^{(7)}_{\rm eff}=e^{7g_{\mathrm{share,eff}}}$. A coherent transfer orbit over $N$ cells of step length $a$ has first phase mode of wavelength $\lambda_1=aN$; with the transverse export fixing the step at $a=\tfrac23 L_*$ (Appendix~C.5), the first mode over the typical support is $\lambda=\tfrac23L_*e^{7g_{\mathrm{share,eff}}}$, the recurrence relation above in spatial form. The two readings are reciprocal faces of one asymptotic-equipartition fact---a designated configuration within a typical set of size $e^{S}$ is reached with probability $e^{-S}$ per draw, and traversed at unit step in $e^{S}$ cells---and they share a single residual identification: the temporal form assumes one phase cycle per dressing recurrence, the spatial form that the Compton mode is the first phase mode over the typical support. These are the same statement in two grammars, so the branch carries a single residual identification rather than two --- the one Appendix~H.6 discharges into the recurrence rate, a unit committed charge per pass, and the condensate phase relation, conditional on Fork~A.

\paragraph{Locality cannot replace the refresh kernel.}
The ensemble entering both readings is reached only by the memoryless refresh of Postulate~III, and the requirement is topological. Under any local dynamics---single-label shifts preserving admissibility---the state graph of one parity copy fractures into $4!=24$ disconnected ordering sectors of $\binom{7}{4}=35$ states each, because a unit shift can never reorder two slots without passing through the collision injectivity forbids. Slot ordering is a frozen charge of local motion: no local transfer operator has $p_{\eta_*}$ as its equilibrium, the strength-weighted walk measure saturates at entropy $7.374$ against $g_{\mathrm{share,eff}}=7.4198$, and the dominant transfer mode localizes within a single sector. The full count is available only to a kernel that redraws from the ensemble, which is the content of Postulate~III; the postulate is therefore load-bearing in the scale chain by necessity.

The reduced electron Compton wavelength is
\[
\lambda_e=\frac{\hbar}{m_ec},
\]
and the admissibility-closed sharing entropy is
\[
g_{\mathrm{share,eff}}=7.41980002357.
\]
The closure fixed point also gives
\[
\langle K^2\rangle_{\eta_*}=\frac{3}{2\eta_*},
\qquad
\eta_*=0.0298668443935,
\]
so the closure-saturation factor is
\[
C_{\rm cl}:=\eta_*\langle K^2\rangle_{\eta_*}=\frac32,
\qquad
C_{\rm cl}^{-1}=\frac23.
\]

The faithful sector-resolution principle is
\[
\boxed{
\ln\!\left(\frac{\lambda_e}{L_*}\right)
=
7g_{\mathrm{share,eff}}-\ln C_{\rm cl}
=
7g_{\mathrm{share,eff}}-\ln\!\left(\frac32\right)
}.
\]
Equivalently,
\[
\boxed{
L_*=\frac32\,\lambda_e e^{-7g_{\mathrm{share,eff}}}
}.
\]
This is the named sector-resolution principle of the canonical branch. Once it is adopted, the cell length is fixed from non-gravitational input: the electron Compton scale and the already-derived substrate combinatorics. It gives
\[
L_*=1.60771947\times10^{-35}\ {\rm m},
\]
which is about $0.528\%$ below the conventional CODATA Planck length.

The induced gravitational scale follows from the same algebra used in the stiffness matching:
\[
G_*:=\frac{c^3L_*^2}{\hbar}.
\]
Substituting the sector-resolution relation gives the closed form
\[
\boxed{
G_*=
\frac{c^3}{\hbar}
\left(\frac32\,\lambda_e e^{-7g_{\mathrm{share,eff}}}\right)^2
=
\frac94\,\frac{\hbar c}{m_e^2}\,e^{-14g_{\mathrm{share,eff}}}
}.
\]
Numerically,
\[
G_*=6.60399128\times10^{-11}\ {\rm m^3\,kg^{-1}\,s^{-2}},
\]
about $1.053\%$ below the CODATA value. The uncertainty imported from $\lambda_e$ is many orders of magnitude smaller than this residual, so the difference is theoretical rather than metrological. The residual carries no statistical interpretation within the chain: $g_{\mathrm{share,eff}}$ is an exact weighted evaluation, $G_*\propto e^{-14g_{\mathrm{share,eff}}}$, and matching the CODATA central value would require $\delta g/g\simeq7.5\times10^{-4}$, a shift the $1680$-state combinatorics does not admit. The residual therefore measures a physical correction omitted at this order. The charged-lepton ladder residuals of $-0.51\%$ and $-0.65\%$ (Section~13, Appendix~I) cannot share an algebraic origin with it, because the ladder contains no factor of $g_{\mathrm{share,eff}}$; the framework carries two independent systematic residuals, each assigned to dressing physics beyond leading order. All three residuals are negative, a regularity that constrains candidate corrections without establishing a common origin. A candidate correction enters the closure ledger only by predicting the sign and magnitude of the residual it addresses in advance of any fit. That standard is met at the level of sign by the framework's most natural correction channel. Once the dressing is read as a recurrence of independent draws (Appendix~D.4), the dressing exponent fluctuates about its ensemble mean, and convexity fixes the direction of every induced shift: for fluctuating exponent $h$, $\langle e^{h}\rangle\ge e^{\langle h\rangle}$ and $\langle e^{-h}\rangle\ge e^{-\langle h\rangle}$, so the combinatorial mass ratios shift upward and, through the lengthening of $L_*$, the induced $G_*$ shifts upward as well. The channel admits no opposite-signed member, and all three measured values exceed their combinatorial predictions: three deficits, one forced sign, three matches. The magnitudes are not yet derived; the branch variances are computable from the ensemble statistics with no free parameters, and the resulting ratio predictions $\delta_\tau/\delta_\mu$ and $\delta_{\ln G}/\delta_\mu$ are the most direct test this sector can state in advance of measurement.

The factor $3/2$ is fixed structurally by the tetrahedral transverse-export geometry. The tetrahedral identity gives the transverse export fraction
\[
\frac14\sum_{i=1}^4(\hat n_i\cdot\hat u)^2=\frac13,
\qquad
f_\perp=1-\frac13=\frac23,
\]
so $f_\perp^{-1}=3/2$. The admissibility fixed point returns the same value,
\[
C_{\rm cl}^{-1}
=
\left(\eta_*\langle K^2\rangle_{\eta_*}\right)^{-1}
=
\frac23,
\]
but it is built in. The condition defining $\eta_*$ is $\langle K^2\rangle_\eta=3/(2\eta)$, i.e. $\eta\langle K^2\rangle=3/2$, so $C_{\rm cl}=3/2$ holds by construction --- the fixed point returns the value the condition put in. The $3/2$ is fixed once, by the transverse-export geometry, and the closure saturation only reproduces it. In the sector-resolution relation it enters once as a global subtraction $\ln(3/2)$, not seven times, because the closure condition is one scalar fixed-point condition on the admissibility family rather than seven independent sector constraints.

A short look-elsewhere check makes the role of this correction more precise. Write a possible log-gap correction as
\[
L_*=\lambda_e\exp[-7g_{\mathrm{share,eff}}+\Delta].
\]
Several framework numbers lie near the required correction, but they do not play the same structural role:
\begin{center}
\small
\begin{tabular}{p{0.22\textwidth}p{0.14\textwidth}p{0.16\textwidth}p{0.34\textwidth}}
\hline
Correction $\Delta$ & Value & Induced $G$ residual & Structural status \\
\hline
$\ln(3/2)$ & $0.40546511$ & $-1.05\%$ & fixed by transverse tetrahedral geometry; mirrored by the closure-saturation normalization \\
$g_{\mathrm{share,eff}}-7$ & $0.41980002$ & $+1.82\%$ & nearby, but not independently selected \\
$G_{\mathrm{tet}}(0)$ & $0.44822039$ & $+7.78\%$ & local Green constant, not global export factor \\
$7[-\ln c_{\rm loop}^{(\mathrm{ren})}]$ & $0.32774720$ & $-15.30\%$ & transport renormalization, wrong role \\
\hline
\end{tabular}
\normalsize
\end{center}
The preferred correction is the one already singled out by both the tetrahedral export geometry and the admissibility fixed-point saturation; the other nearby numbers play different structural roles.

The electron anchor carries three related roles. Its reduced Compton wavelength $\lambda_e$ is the non-gravitational length used in faithful sector resolution. Its status as the lightest clean one-bit fermionic defect identifies which elementary excitation calibrates the seven-channel dressing block. Its mass fixes the mass--entropy map through
\[
\kappa_m(\lambda_e)=\frac{m_e}{\ln2}.
\]
The same elementary defect enters the two normalization channels in distinct roles: its Compton length sets the UV cell scale, while its mass per one-bit defect entropy sets the source normalization.

The seven-sector scale-setting relation is not a counting heuristic; it has a concrete finite-dimensional realization on a transfer-operator history space, which we now construct. The face label algebra is
\[
\mathcal A_7=\bigoplus_{m=-3}^{3}\mathbb C E_m,
\qquad
E_mE_n=\delta_{mn}E_m,
\qquad
\sum_{m=-3}^{3}E_m=I_7.
\]
The admissibility-closed tetrahedral state space has $1680$ oriented injective states, with stationary weight
\[
p_{\eta_*}(b)=Z^{-1}e^{-\eta_*K^2(b)}.
\]
For a single sector, the refresh kernel
\[
P_{\eta_*}(b,b')=p_{\eta_*}(b')
\]
has Perron stationary entropy $g_{\mathrm{share,eff}}$. The electron does not simultaneously occupy seven mutually exclusive face labels in one tetrahedron; the labels are sector channels in the dressing history of the one-bit defect. The appropriate support object is therefore a history space, with one admissibility-closed sector layer for each $m=-3,\ldots,3$.

Let
\[
\mathcal H_{\rm hist}
=
\bigotimes_{m=-3}^{3}\mathcal H_B^{(m)},
\qquad
\dim \mathcal H_B=1680,
\]
and let $\mathcal T_m$ act as $P_{\eta_*}$ on the $m$th factor and as the identity on the others. A representative one-pass dressing operator is
\[
\mathcal D_e^{(0)}
=
\mathcal T_{-3}\mathcal T_{-2}\cdots\mathcal T_{3}.
\]
The displayed order is only a representative of the symmetrized one-pass class. It does not introduce a physical ordering of the face sectors, nor does it multiply the support by an extra factor of $7!$. What matters is that the ground state of the elementary fermionic defect resolves each simple face sector once: fewer sectors leave unresolved label memory, while repeated sectors describe excited or additionally dressed states rather than the elementary anchor.

A sector-resolution step should not be identified with conditioning the
1680-state ensemble on boundary states containing a given label $m$.  Such
conditioning changes the entropy and does not reproduce
$g_{\mathrm{share,eff}}$.  The sector label instead specifies which simple
face-algebra channel is being resolved while the admissibility cloud sampled
in that step remains the full closed boundary ensemble.

With the effective dimension defined by the stationary Shannon entropy of the positive history kernel, each sector contributes $g_{\mathrm{share,eff}}$, so
\[
\dim_{\rm eff}(\mathcal D_e^{(0)})
=
\exp(7g_{\mathrm{share,eff}})
=
3.60286052\times10^{22}.
\]
Only the transverse part of this support is exported into the weak-field scalar channel. The normalized export weight is
\[
w_\perp=\frac23,
\]
the same fraction fixed by the tetrahedral transverse projection and by the reciprocal closure-saturation factor. Hence
\[
\dim_{\rm eff}(\mathcal D_{e,\perp})
=
w_\perp\,\dim_{\rm eff}(\mathcal D_e^{(0)})
=
\frac23\,e^{7g_{\mathrm{share,eff}}}
=
2.40190701\times10^{22}.
\]
The faithful sector-resolution principle is then the support-scale relation
\[
\boxed{
\frac{\lambda_e}{L_*}
=
\dim_{\rm eff}(\mathcal D_{e,\perp})
=
\frac23\,e^{7g_{\mathrm{share,eff}}}
}.
\]
The justification of this relation is given in Appendix~H. An explicit dressing Hamiltonian (Appendix~H.2) makes the lightest charged fermionic defect the one-pass ground state that resolves each of the seven sectors once. The refresh kernel $P_{\eta_*}$ used above is the memoryless dressing of Postulate~III: with the substrate relaxed long before readout, each pass is an independent draw from the admissibility ensemble, so every channel carries the full $g_{\mathrm{share,eff}}$. The seven-fold additivity then needs the channels to be mutually independent, and they are, for a reason internal to the anchor: inter-channel correlation would contract the defect and raise its mass, and the electron is the lightest one-bit charged defect, so its dressing is the product and correlated dressings describe heavier excitations (Appendix~H.5). With the $2/3$ factor the tetrahedral transverse projection of Appendix~C.5, the faithful sector-resolution principle is a consequence of the mass--entropy ontology and the electron anchor, and nothing in the chain remains free.

The status of the relation should be named precisely. The effective dimension $e^{H}$ is the perplexity of the coherent dressing --- the effective number of mutually resolvable support positions in the defect's history, not a count of cells in a spatial volume. Reading it as a physical length ratio is the recurrence identification derived above: the dressing re-completes once per $e^{7g_{\mathrm{share,eff}}}$ ticks, and one phase cycle per recurrence converts that waiting time into the Compton scale, so the linear, first-power map is the arithmetic of joint probability rather than a chosen dictionary. The map is cross-validated by the induced gravitational scale, by the charged-lepton ladder, and by the mass-minimization argument above, which together bracket the support exponent at unity within $10^{-3}$. One identification remains beneath the derivation --- one phase cycle per dressing recurrence, equivalently that the Compton mode is the first phase mode over the typical support --- and a one-defect transfer operator $T$ whose coherence length obeys $\lambda=C_{\rm cl}L_*\,e^{H[p_T]}$ is the object whose spectrum would certify it; exhibiting that operator is the branch's remaining microscopic completion.

This is the microscopic content of the length formula above: the elementary one-bit fermionic defect is supported over one complete transverse-exported dressing block. The matched weak-field gravitational constant then follows from the gauge-invariant bridge
\[
G=\frac{c^2}{8\pi}\frac{\kappa}{\gamma S_\infty},
\]
with the entropy-unit normalization handled by the rescaling convention described in Appendix~C.5.

\subsection*{D.5 EFT consistency checklist}

The weak-field EFT does not rely only on successful phenomenology; it also passes a standard consistency checklist at the level claimed here.
\begin{itemize}
\item `No ghost`: the scalar kinetic term carries positive sign because $\gamma>0$.
\item `No tachyon`: the quadratic fluctuation operator contains no mass term at this order.
\item `Correct-sign sourcing`: the defect-source coupling lowers the available entanglement capacity around positive-mass defect configurations rather than generating repulsive static behavior in the weak-field branch.
\item `Causal propagation`: the transport completion satisfies $D/\tau_0=c^2$, so the time-dependent sector propagates at finite signal speed.
\item `Weak-field unitarity below cutoff`: once the scalar sector is quantized around the weak-field branch, the absence of ghost or tachyonic modes leaves an ordinary sub-cutoff scalar EFT rather than an obviously pathological one.
\item `Energy-condition role`: the scalar gradient sector contributes positive local stiffness energy, while cosmological acceleration enters through the background branch rather than through a ghost-like local degree of freedom.
\end{itemize}
These statements are made at the EFT level claimed here. They do not replace the need for a fuller UV derivation, but they do show that the weak-field scalar sector is not buying phenomenology by obvious field-theoretic pathology.

The checklist is intentionally modest. Its role is not to prove ultraviolet completion of the full framework, but to show that the low-energy scalar sector used in the weak-field branch passes the standard first tests of EFT health.

The one place where an explicit formula is worth recording is linear vacuum stability in the time-dependent sector. Writing a small perturbation $\delta s$ about the vacuum branch, the linearized telegrapher equation is
\[
\tau_0\ddot{\delta s}+\dot{\delta s}-D\nabla^2\delta s=0.
\]
For a plane-wave mode $e^{-i\omega t+i\mathbf{k}\cdot\mathbf{x}}$, this gives the dispersion relation
\[
\tau_0\omega^2+i\omega-Dk^2=0.
\]
With $\tau_0>0$ and $D>0$, the corresponding mode frequencies have non-growing time dependence, so the vacuum is linearly stable. The same sign structure underlies the earlier no-ghost and no-tachyon statements: positive kinetic stiffness, positive transport coefficients, and no negative mass-squared term in the linearized sector.

\subsection*{D.6 Quadratic fluctuations and weak-field stability}

Expanding the action about an on-shell background yields the quadratic fluctuation operator
\[
I^{(2)}[\delta S]
=
-\int d^4x\,\sqrt{-g}\,
\frac{\gamma}{2}\,
g^{\mu\nu}\partial_\mu\delta S\,\partial_\nu\delta S.
\]
There is no quadratic mass term at this order, so the low-energy scalar sector contains one massless bosonic mode. Stability requires $\gamma>0$, which is reinforced in the microscopic realization appendix by condensate hydrodynamics.

This is also the local EFT reason the bosonic occupancy language in the galactic section is natural rather than decorative. The weak-field branch genuinely contains a stable massless scalar mode whose occupation can be discussed meaningfully.

Appendix~D provides the technical support layer for the weak-field bridge, Newton limit, electron anchor, substrate length branch, and EFT consistency audit.

\section*{Appendix E: Transport, Cosmology, and Hubble-Tension Implementation}

Appendix~E collects the time-dependent and homogeneous extensions of the static branch: the same scalar medium propagating causally, relaxing toward its static limit, and supporting a cosmological background mode without losing contact with the weak-field structure already derived. The status differs across these pieces: the transport relation is closed at the preferred-branch level ($D/\tau_0=c^2$ with $\tau_0^{-1}=H_0$); the homogeneous Hubble-tension mechanism is directionally supported but not Boltzmann-closed; and the full perturbation likelihood treatment remains open.

\subsection*{E.1 Telegrapher equation and causal closure}

The time-dependent deficit field obeys
\[
\tau_0\partial_t^2\delta S+\partial_t\delta S=D\nabla^2\delta S+A\chi,
\qquad
\frac{A}{D}=\frac{\kappa}{\gamma}.
\]
Causality requires
\[
\frac{D}{\tau_0}=c^2.
\]
In the canonical no-new-IR-scale branch,
\[
\tau_0^{-1}=H_0,
\qquad
D=\frac{c^2}{H_0}.
\]

This is the minimal causal completion of the static Poisson sector. The telegrapher form supplies propagation and relaxation, but it is chosen so that the static weak-field law remains the exact late-time limit rather than being replaced by a new phenomenological rule.

\subsection*{E.2 Static-limit recovery for galaxies}

For a Fourier mode $k$, the telegrapher characteristic equation
\[
\tau_0 s^2+s+Dk^2=0
\]
has the roots
\[
s=-\frac{1}{2\tau_0}\pm i\omega_k,
\qquad
\omega_k\simeq ck
\]
whenever $4\tau_0Dk^2\gg 1$. Galactic wavelengths are far below the critical scale
\[
\lambda_c=\frac{4\pi c}{H_0}\approx 54\ \mathrm{Gpc},
\]
so galactic modes are deeply underdamped. Time-averaging the sourced solution over intervals large compared with $2\pi/\omega_k$ returns the static Poisson branch exactly, and the residual ponderomotive correction scales parametrically as
\[
\frac{\delta F_{\mathrm{pond}}}{F_{\mathrm{static}}}
\sim
e^{-T/(2\tau_0)}
\left(\frac{\omega_{\mathrm{orb}}}{\omega_k}\right)^2
\sim 10^{-8}.
\]

That estimate is why the transport sector does not undercut the static galactic results. The oscillatory contribution is present, but it is parametrically too small to compete with the near-stationary weak-field branch in ordinary galactic systems.

\subsection*{E.3 Homogeneous mode and cosmological sourcing}

The cosmological split is
\[
S(x,t)=\overline S(t)+s(x,t),
\]
with $\overline S(t)$ the homogeneous mode and $s(x,t)$ the inhomogeneous weak-field sector. The background capacity is normalized by the apparent horizon,
\[
S_\infty(t)=\pi\frac{R_A(t)^2}{L_*^2},
\qquad
R_A(t)=\frac{c}{\sqrt{H^2+k c^2/a^2}}.
\]
Because the field couples to the trace of the stress-energy tensor, the homogeneous mode is suppressed during radiation domination and turns on near matter--radiation equality.

This timing is the central cosmological virtue of the mechanism. The homogeneous mode is quiet when it must be quiet, then becomes relevant close to the epoch where a sound-horizon shift is most useful.

\subsection*{E.4 Sound-horizon shift and shear lock}

In the closed cosmological branch, the trace-sourced homogeneous mode acts as a transient early-energy contribution. The qualitative payoff is a smaller sound horizon and an upward shift of the CMB-inferred Hubble constant toward the upper-$68$ / low-$69$ km s$^{-1}$ Mpc$^{-1}$ range. Local weak-field predictions are protected by the separation between $\overline S(t)$ and $s(x,t)$: the homogeneous mode changes the background branch without rewriting the local static Poisson law.

The result is qualitative but substantial. The homogeneous mode can matter cosmologically without forcing a re-tuning of the local weak-field sector that already fixed the galactic branch.

The transport relation and preferred branch are closed; the cosmological sector remains structurally supported but not yet Boltzmann-closed.

\section*{Appendix F: Constrained-Capacity Strong-Field Branch}

Appendix~F records the strong-field black-hole branch in the form used by the main text. The result is not a full microscopic theory of collapse, but it is stronger than a mere horizon mnemonic. Once the bounded capacity variable is adopted and the lapse is tied to surviving capacity, the static spherical exterior is fixed: the continuum EFT lives on the domain $q>0$, reduces to vacuum Einstein dynamics there, and terminates at the $q=0$ surface.

\subsection*{F.1 Bounded capacity and the unique lapse map}

The strong-field order parameter is the surviving-capacity fraction
\[
q(x)=\frac{S_{\mathrm{ent}}(x)}{S_\infty}\in[0,1].
\]
This bound follows directly from finite local channel capacity. If the vacuum channel count is finite and $S_{\mathrm{ent}}$ is the logarithmic coarse entropy of the surviving local ensemble, then no physical branch can have either negative capacity or more than the asymptotic vacuum capacity.

In a static exterior, the lapse associated with the asymptotic Killing time is determined by the local surviving capacity. Let
\[
N=f(q).
\]
The conditions are:
\[
f(1)=1,\qquad \lim_{q\to0^+}f(q)=0.
\]
The substrate-level composition axiom is that independent serial capacity losses compose multiplicatively on the lapse:
\[
f(q_1q_2)=f(q_1)f(q_2).
\]
This is the assumption that extends the linear weak-field match to a nonlinear lapse map. With continuity, the positive solutions on $(0,1]$ are $f(q)=q^\alpha$. Expanding near $q=1-\epsilon$ gives
\[
N=q^\alpha=1-\alpha\epsilon+O(\epsilon^2).
\]
The weak-field bridge gives
\[
N=1-\frac{\epsilon}{2}+O(\epsilon^2),
\]
so $\alpha=1/2$ and therefore
\[
N=\sqrt q,
\qquad
N^2=q.
\]
Equivalently, if one writes $N^2=F(q)$, the unique continuous multiplicative completion is $F(q)=q$. The nonlinear lapse rule is therefore fixed by capacity composition and weak-field matching; it is not a freely chosen black-hole ansatz.

\subsection*{F.2 Constrained-capacity action on the physical domain}

The physical strong-field domain is
\[
\mathcal{M}_q=\{(t,x)\mid q(t,x)>0\}.
\]
The minimal constrained-capacity action is the ADM gravitational action on $\mathcal{M}_q$, with a multiplier enforcing the lapse-capacity relation. In canonical form,
\[
I_{\rm strong}
=
\int_{\mathcal{M}_q}dt\,d^3x
\left[
\pi^{ij}\dot h_{ij}
-N\mathcal H_{\rm GR}
-N^i\mathcal H^{\rm GR}_i
+\sqrt h\,\lambda\,(N^2-q)
\right]
+
I_{\rm matter}
+
I_\infty
+
I_{\partial\mathcal{M}_q}.
\]
The standard ADM constraint densities are
\[
\mathcal H_{\rm GR}
=
\frac{16\pi G}{c^4\sqrt h}
\left(\pi_{ij}\pi^{ij}-\frac12\pi^2\right)
-
\frac{c^4}{16\pi G}\sqrt h\,{}^{(3)}R,
\qquad
\mathcal H^{\rm GR}_i=-2D_j\pi^j{}_i.
\]
The same content can be written in Lagrangian form as
\[
I_{\rm strong}
=
\frac{c^4}{16\pi G}
\int_{\mathcal{M}_q}dt\,d^3x\,
N\sqrt h\left({}^{(3)}R+K_{ij}K^{ij}-K^2\right)
\]
\[
\quad
+
\int_{\mathcal{M}_q}dt\,d^3x\,\sqrt h\,\lambda\,(N^2-q)
+
I_{\rm matter}
+
I_\infty
+
I_{\partial\mathcal{M}_q}.
\]
Here $h_{ij}$ is the spatial metric, $N$ and $N^i$ are the lapse and shift, and
\[
K_{ij}=\frac{1}{2N}
\left(\dot h_{ij}-D_iN_j-D_jN_i\right)
\]
is the extrinsic curvature. The term $I_\infty$ is the usual asymptotic boundary term required for a well-posed variational principle and finite ADM energy, while $I_{\partial\mathcal{M}_q}$ is the effective action carried by the $q=0$ capacity-exhaustion surface.

Varying $\lambda$ gives
\[
N^2=q.
\]
The $q$ variation gives, schematically,
\[
-\sqrt h\,\lambda
+
\frac{\delta I_{\rm matter}}{\delta q}
+
\frac{\delta I_{\partial\mathcal{M}_q}}{\delta q}
=0.
\]
Away from the $q=0$ boundary, and in bulk vacuum with no explicit matter coupling directly to $q$, this gives $\lambda=0$. The lapse variation gives
\[
\mathcal H_{\rm GR}+\mathcal H_{\rm matter}
=
2N\sqrt h\,\lambda,
\]
so in bulk vacuum
\[
\mathcal H_{\rm GR}=0.
\]
The shift variation gives the ordinary momentum constraint,
\[
\mathcal H^{\rm GR}_i+\mathcal H^{\rm matter}_i=0.
\]
The metric evolution equations then reduce to the vacuum Einstein equations on $\mathcal{M}_q$ \citep{ADM1962}. The constrained branch is thus not a new scalar-hair exterior. It is Einstein vacuum evolution on the surviving-capacity domain, together with the algebraic statement that the lapse is the square root of local capacity.

\subsection*{F.3 Horizon boundary action and boundary conditions}

The only new strong-field ingredient not present in ordinary exterior GR is the inner boundary action on the capacity-exhaustion surface. At the EFT level its variation may be written schematically as
\[
\delta I_{\partial\mathcal{M}_q}
=
\int_{\partial\mathcal{M}_q}d^3y\,\sqrt\sigma
\left(
\frac12 s^{AB}\delta\sigma_{AB}
+\chi_{\partial}\,\delta q
+\cdots
\right),
\]
where $\sigma_{AB}$ is the induced metric on the boundary world tube, $s^{AB}$ is the effective boundary stress response, and $\chi_{\partial}$ is the response conjugate to the capacity variable. The universal boundary condition is not optional:
\[
q\big|_{\partial\mathcal{M}_q}=0.
\]
The further constitutive information--whether the boundary is absorbing or partially reflecting, how quickly its channels relax, and what microscopic degeneracy they carry--is encoded in $I_{\partial\mathcal{M}_q}$.

The static exterior can therefore close before the full microscopic black-hole sector is finished. The exterior bulk equations only require $\lambda=0$ on $\mathcal{M}_q$ and the algebraic condition $N^2=q$. The detailed horizon-channel physics lives in the boundary action. A shape variation of the moving boundary would give a force-balance condition of the schematic form
\[
\delta_\xi I_{\rm strong}
=
\int_{\partial\mathcal{M}_q}d^3y\,\sqrt\sigma\,
\xi_n\left(\mathcal P_{\rm bulk}+\mathcal P_{\partial q}\right)=0,
\]
where $\xi_n$ is the normal displacement and the two pressures encode the limiting bulk stress and the boundary-channel response. Deriving $\mathcal P_{\partial q}$ from the tetrahedral graph ensemble is one of the remaining strong-field closure tasks.

\subsection*{F.4 Static spherical vacuum exterior}

In static spherical vacuum, the result follows immediately. On $\mathcal{M}_q$, the bulk equations are the vacuum Einstein equations. The unique asymptotically flat static spherical solution is the Schwarzschild exterior,
\[
ds^2
=
-\left(1-\frac{2GM}{c^2r}\right)c^2dt^2
+
\left(1-\frac{2GM}{c^2r}\right)^{-1}dr^2
+
r^2d\Omega^2,
\]
so the capacity variable is
\[
q(r)=N^2(r)=1-\frac{2GM}{c^2r},
\qquad
r>r_h,
\]
with
\[
r_h=\frac{2GM}{c^2}.
\]
This also agrees with the weak-field capacity deficit:
\[
\frac{\delta S(r)}{S_\infty}
=
\frac{2GM}{c^2r},
\qquad
q(r)=1-\frac{\delta S(r)}{S_\infty}.
\]

The domain carries the physical content. The exterior $r>r_h$ is exactly the standard Schwarzschild exterior. At $r=r_h$, $q=0$. Continuing the same real capacity branch to $r<r_h$ would require $q<0$, which is not in the state space. In this framework the classical Schwarzschild interior is therefore not the physical continuation of the same continuum EFT. It is the GR manifold analytically continued beyond the surface where the entanglement-capacity substrate has no remaining local continuum channels.

This exterior-domain result does not by itself decide what an infalling observer experiences at the $q=0$ surface. Whether the capacity-exhaustion boundary is locally smooth, dissipative, or anomalous is a question about the boundary microphysics encoded in $I_{\partial\mathcal{M}_q}$, not a conclusion fixed by the exterior Schwarzschild closure alone.

\paragraph{Geometric identification of the capacity variable.}
The static rule $N^2=q$ together with the bulk reduction to vacuum Einstein on $\mathcal{M}_q$ forces a geometric reading of the substrate variable. In a spherically symmetric foliation with areal radius $R$, the normalization-independent gradient invariant on the orbit space is
\[
q=h^{\mu\nu}\partial_\mu R\,\partial_\nu R=|\nabla R|^2,
\]
which in Schwarzschild gives
\[
q=1-\frac{2GM}{c^2r}=N^2,
\]
matching the substrate definition $q=S_{\rm ent}/S_\infty$ on the exterior. In spherical dynamical collapse with Misner--Sharp mass $M_{\rm MS}(R,t)$,
\[
q(R,t)=1-\frac{2GM_{\rm MS}(R,t)}{c^2R}.
\]
The substrate variable $q$ and the geometric areal-radius gradient invariant therefore agree as a theorem of the framework's structural commitments, not as an additional axiom. The marginal-trapped-surface condition $\theta_+\theta_-=0$ on a closed two-surface coincides with $q=0$ independently of the null normalization used to define $\theta_\pm$: $q>0$ is the untrapped exterior, $q=0$ is the marginal / apparent horizon, and $q<0$ would be the classical trapped interior. The capacity domain $\mathcal{M}_q=\{q>0\}$ is the untrapped exterior, and the saturated boundary $\partial\mathcal{M}_q=\{q=0\}$ is the apparent-horizon two-surface. Standard apparent-horizon formation theorems for gravitational collapse therefore produce the $q=0$ saturation surface inside the framework, without an additional substrate-transport postulate; the bounded transport law of Appendix~F.7 is a substrate-side consistency check on this geometric formation rather than its primary mechanism.

\paragraph{Rotating and charged stationary exteriors.}
The same vacuum-domain reduction extends to the rotating and charged cases. On $\mathcal{M}_q$ the bulk equations reduce to vacuum Einstein, or to Einstein--Maxwell when a conserved gauge sector hosted by the substrate (Appendix~I.2) is present. The standard no-hair theorems then identify the stationary asymptotically flat exterior branches as Kerr, Reissner--Nordstr\"om, or Kerr--Newman, and the saturation surface $q=0$ coincides with the outer Killing horizon. The Bekenstein--Hawking entropy
\[
S=\frac{k_BA}{4L_*^2}
\]
retains its form with $A$ the appropriate horizon area, and the Hawking temperature follows from the surface gravity as usual,
\[
T_H=\frac{\hbar \kappa_{\rm sg}}{2\pi k_Bc}.
\]
The inner Cauchy horizons of Reissner--Nordstr\"om and Kerr are not physical substrate interiors; they are analytic continuations into the $q<0$ region the substrate does not enter.

\subsection*{F.5 Horizon thermodynamics and boundary capacity}

Because the exterior geometry is unchanged, semiclassical quantities depending only on the exterior near-horizon saddle are unchanged. The Euclidean continuation is used here as an exterior-saddle calculation: the resulting periodicity depends on regularity of the near-horizon exterior geometry, not on a physical continuation of the substrate past $q=0$. The Euclidean regularity argument therefore gives the standard Hawking temperature \citep{Hawking1975,GibbonsHawking1977},
\[
T_H=\frac{\hbar c^3}{8\pi GM k_B}.
\]
The same exterior Euclidean saddle gives the Bekenstein--Hawking area law \citep{Bekenstein1973,Hawking1975},
\[
S_{\rm BH}=\frac{k_BA}{4L_P^2}.
\]

We now derive the $1/4$ area coefficient from substrate inputs already in the framework, rather than from a stipulated channel-counting rule. The derivation uses two independent ingredients, each of which is independently closed within the bulk EFT before any horizon machinery is introduced.

\paragraph{Per-channel cut entropy from fermionic face exclusion.}
At a saturation surface $q=0$, the physical continuum domain terminates at $\partial\mathcal{M}_q$. A tetrahedral cell on the exterior side whose outward face would, in the unsaturated bulk, be paired with a neighboring cell across that face instead has an unfilled pairing slot, because the would-be partner lies outside $\mathcal{M}_q$. By Postulate~II the elementary face slot is fermionic and admits only the occupied (paired) state or the excluded (unpaired) state. A horizon cut face is therefore an instance of the same elementary face-exclusion defect that anchors the one-bit fermionic sector in the bulk. The seven-state $j_{\rm eff}=3$ channel belongs to the paired bulk link $j_0\otimes j_0\to j_{\rm eff}$, not to the cut face itself: with the partner cell absent, the paired representation is never formed. The seven-state structure enters the boundary count through the bulk graph response and hence through the channel density below, not through the per-channel entropy. The entropy carried by an elementary cut defect is consequently the same primitive fermionic increment that anchors the electron at $\kappa_m(\lambda_e)=m_e/\ln 2$,
\[
\Delta S_f=\ln 2.
\]

\paragraph{Channel density from the transverse bulk graph response.}
Because the local graph Green tensor is isotropic,
\[
{\cal G}_{\rm loc}^{ab}=\frac{G_{\mathrm{tet}}(0)}{3}\delta^{ab},
\]
a codimension-one horizon cut with local normal $\hat n$ exports only the transverse two-plane component,
\[
G_\perp=(\delta^{ab}-\hat n^a\hat n^b){\cal G}_{\rm loc}^{ab}=\frac23G_{\mathrm{tet}}(0).
\]
The horizon channel density is therefore not a new boundary response coefficient; it is the transverse projection of the same bulk Green response already fixed in Appendix~C. The number of active channels per outward angular direction is $G_\perp/\ln 2$, and integrating over the horizon two-surface (the spherical angular measure in Schwarzschild, the smooth axisymmetric horizon for Kerr, with isotropic transverse response in either case) gives
\[
n_{\rm hor}=4\pi\frac{G_\perp}{\ln 2}=\frac{8\pi G_{\mathrm{tet}}(0)}{3\ln 2}.
\]

\paragraph{Closure as an exact identity.}
Using the cell-normalized capacity baseline from Appendix~C,
\[
S_\infty^{\rm cell}=\frac{3\ln 2}{32\pi G_{\mathrm{tet}}(0)},
\]
the product of the two substrate inputs is
\[
n_{\rm hor}S_\infty^{\rm cell}=\frac14,
\]
an exact identity in which the Joyce diamond-lattice constant $G_{\mathrm{tet}}(0)$ and the fermionic increment $\ln 2$ cancel between the two factors. A horizon of area $A$ therefore carries the dimensionless entropy
\[
\frac{S_{\rm hor}}{k_B}=\frac{A}{4L_*^2},
\]
and the Bekenstein--Hawking coefficient is recovered after the gravitational scale is matched so that $L_P(G_*)=L_*$. The area coefficient is closed within the constrained-capacity EFT under Postulate~II and the transverse bulk graph response, both of which are independent of the horizon construction. The remaining strong-field task concerns nonuniversal boundary spectroscopy --- the explicit microscopic Hamiltonian behind the relaxation spectrum, stretched-layer corrections, and transient response --- rather than the universal area coefficient.

\paragraph{Composition of the coefficient.}
Written as a product of its surviving factors, the coefficient is
\[
\frac14=\frac23\times\frac38,
\]
and the $2/3$ is the transverse projection $G_\perp/G_{\mathrm{tet}}(0)$ --- the same export factor that sets the electron's length relation $\lambda_e/L_*=\tfrac23e^{7g_{\mathrm{share,eff}}}$ (Section~13.3) and enters the loop self-energy through the edge map $(2/3)(1/3)=2/9$ (Section~8). One microscopic projection is load-bearing in the electron mass, the bulk stiffness, and the horizon coefficient at once, so these sectors cannot be adjusted independently: a change to the transverse export that repaired any one of them would displace the other two.

\paragraph{Coherence with the bulk normalization.}
The identity also closes a thermodynamic loop. With the exterior-saddle temperature above, the Clausius relation $\delta Q=T\,dS$ applied to local horizon patches at entropy density $1/(4L_*^2)$ reproduces the field equations with induced constant $c^3L_*^2/\hbar=G_*$ --- the Jacobson construction \citep{Jacobson1995} run inside the framework --- so $1/4$ is exactly the coefficient at which the horizon first law returns the same induced scale the bulk transport chain produces. The horizon route inherits its normalization from the same conventions that close the source theorem (Appendix~C.5), so the agreement is a coherence closure within those conventions. Its content is that the substrate's horizon thermodynamics and its bulk transport meet at one Newton constant rather than two.

\subsection*{F.6 Absorption, ringdown, and echoes}

Exterior wave propagation is unchanged. With $q=0$ identified as the marginal-capacity / apparent-horizon surface (Appendix~F.4), the standard near-horizon tortoise coordinate $r_*\sim r_h\ln q$ pushes the surface to $r_*\to-\infty$, and the wave equation reduces to $(\partial_t^2-\partial_{r_*}^2)\psi\simeq 0$ near the boundary. Horizon regularity at the future horizon then selects the purely ingoing mode
\[
\psi\sim e^{-i\omega(t+r_*)},
\]
so the leading reflectivity is
\[
\mathcal R=0.
\]
Perturbations outside the capacity-exhaustion surface obey the usual Schwarzschild Regge--Wheeler/Zerilli scattering problem \citep{ReggeWheeler1957,Zerilli1970}, with the absorbing boundary fixed by regularity rather than stipulated. The greybody factors, absorption cross sections, and quasinormal ringdown agree with the standard Schwarzschild exterior calculation.

Echoes are therefore correction-level rather than generic. They arise only if microscopic boundary channels carry finite UV reflectivity or if the effective reflection surface is displaced outward to a stretched layer
\[
q=\epsilon>0,
\]
in which case the region between the exterior potential barrier and the partially reflecting layer behaves as a cavity, and a typical echo delay scales as
\[
\Delta t_{\rm echo}
\sim
\frac{2r_h}{c}|\ln\epsilon|+\tau_{\rm ch},
\]
where $\tau_{\rm ch}$ is a boundary-channel relaxation time. Absence of echoes is the default absorbing-boundary prediction; echoes, if detected, would probe boundary microphysics rather than follow automatically from the absence of a classical interior.

\subsection*{F.7 Dynamical formation as a free-boundary problem}

The geometric identification of Appendix~F.4 supplies the primary horizon-formation story: standard apparent-horizon formation in gravitational collapse produces the $q=0$ saturation surface as the first marginally outer trapped surface, with no additional substrate-transport postulate required. What remains is a substrate-side dynamical consistency check --- writing a bounded causal transport law for $q$ and verifying that its evolution agrees with the geometric formation, while bounding $q$ in its physical range $[0,1]$ throughout the collapse. The natural completion is a bounded causal transport system for $q$, for example
\[
\partial_t q+D_iJ^i=-\Gamma(q)\Sigma[T_{\mu\nu}],
\]
\[
\tau_J(\partial_t+\mathcal{L}_v)J^i+J^i=-D(q)D^i q.
\]
Here $J^i$ is the capacity flux, $\Sigma[T_{\mu\nu}]$ is a positive depletion source built from the collapsing stress-energy, $D(q)$ is a bounded mobility, $\Gamma(q)$ is a bounded depletion rate, and $\tau_J>0$ is a relaxation time. Eliminating $J^i$ gives a telegrapher-type equation with finite characteristic speed
\[
v_{\rm cap}\sim \sqrt{\frac{D_0}{\tau_J}}.
\]
Choosing constitutive functions such as
\[
D(q)=D_0q(1-q),
\qquad
\Gamma(q)=\Gamma_0q
\]
makes $q=0$ and $q=1$ invariant sets, preventing the evolution from overshooting into $q<0$.

When $q$ first reaches zero on a two-surface, that surface becomes a moving boundary
\[
\partial\mathcal{M}_q(t)=\{x\mid q(t,x)=0\}.
\]
The level-set kinematics are fixed by differentiating $q(t,X(t))=0$ along the moving surface:
\[
V_n=-\frac{\partial_t q}{|\nabla q|}\bigg|_{q\to0^+}.
\]
In spherical symmetry this becomes
\[
\frac{dr_f}{dt}
=
-\frac{\partial_t q}{\partial_r q}\bigg|_{r=r_f(t)}.
\]
During continued infall the exterior should be Vaidya-like with a slowly varying mass parameter, settling to the Schwarzschild exterior after the front stabilizes. This is a well-posed program, but not yet a closed derivation: the transport coefficients, boundary action, and channel relaxation spectrum must be computed from the graph ensemble or constrained by simulation.

The dynamical system is presumed to preserve the usual covariant conservation of the combined matter-plus-capacity stress-energy, with any local matter depletion balanced by flux, boundary work, or capacity-sector stress. Showing that this conservation structure follows from a graph-derived transport action, rather than imposing it as a constitutive condition, is part of the dynamical closure work.

The concrete closure tests are correspondingly specific. A spherical collapse simulation should show formation of the first $q=0$ surface without overshoot into $q<0$. A coupled matter-plus-capacity run should approach a Vaidya exterior during accretion and a Schwarzschild exterior after settling while satisfying the combined conservation law. Exterior perturbation simulations with an absorbing boundary should reproduce standard Schwarzschild greybody factors and ringdown, while partial-reflectivity runs should produce controlled echo delays. Finally, a microscopic boundary-action calculation or a graph-ensemble Monte Carlo of saturated boundary channels should reproduce the channel-counting rule that yields $n_{\rm hor}S_\infty^{\rm cell}=1/4$. These are not new fit knobs; they are the numerical and microscopic tests that would close the dynamical and boundary sectors.

\subsection*{F.8 PPN boundary and weak-field breakdown}

In the weak-field Solar-System regime, the scalar sector yields, at leading post-Newtonian order,
\[
\gamma_{\mathrm{PPN}}=\beta_{\mathrm{PPN}}=1,
\]
with scalar-induced corrections to $\Phi-\Psi$ appearing only at quadratic weak-field order, schematically $O(\Phi^2/c^4)$, and with the remaining PPN coefficients vanishing in the canonical covariant branch. Weak-field truncations fail only when
\[
\frac{|\Phi|}{c^2}=O(1),
\qquad
\frac{\delta S}{S_\infty}=O(1),
\]
which is exactly the regime where $q$ rather than $\delta S$ is the correct variable.

Appendix~F therefore separates the strong-field claims cleanly. The bounded capacity variable, the lapse rule $N^2=q$, the constrained-capacity exterior action, and the static Schwarzschild exterior on $\mathcal{M}_q$ are closed within the EFT. The geometric identification of $q$ with the marginal-trapped-surface function fixes horizon location and inherits apparent-horizon formation from standard GR collapse. The area coefficient closes under Postulate~II and the transverse bulk graph response, and the leading reflectivity vanishes by horizon regularity. Stationary rotating and charged exteriors reduce to Kerr, Reissner--Nordstr\"om, and Kerr--Newman through the same vacuum-domain bulk reduction. What remains open is nonuniversal boundary spectroscopy: the explicit graph-derived microscopic boundary Hamiltonian behind the relaxation spectrum, stretched-layer corrections, and the substrate-side transport check on the dynamical evolution of $q$ during collapse.

\section*{Appendix G: Many-Pasts, Operational Closure, Branch Realization, and the Arrow of Time}

Appendix~G is the technical support layer for the Many-Pasts postulate (Sections~3.3, 21). It defines the three objects the postulate weights --- a history, the present, and the distance between them --- then derives the operational consequences (Born recovery, no-signaling), explains branch realization and the arrow of time, connects the same weight to the scale-setting branch through memoryless dressing, and locates the proposal among the standard interpretations. The operational core is conservative: it reproduces laboratory quantum mechanics exactly. The ambition is in the ontology built on that core.

\subsection*{G.1 What is a history of the entanglement network?}

A history $H$ is a coarse-grained trajectory of the substrate's microstate: a sequence of admissibility-closed configurations of the entanglement network, one per substrate tick, each obtained from the previous by an allowed update. ``Admissible'' carries the same meaning as in the UV ensemble of Appendix~B --- a configuration close to a regular, isotropic local cell. A history is not a path through an external spacetime, since there is none at this level; it is a path through the network's own configuration space. The set of histories compatible with a fixed present is large, and the postulate is a rule for weighting them.

\subsection*{G.2 What is the present coarse configuration $P$?}

The present $P$ is the macroscopic, record-bearing coarse-graining of the network now. It is the quantity held fixed: every admissible history must terminate in a microstate consistent with $P$. Operationally $P$ is represented by a projector $\Pi_P$ onto the subspace of microstates carrying the present records. The coarse-graining is the ordinary one of statistical mechanics --- macroscopically distinguishable configurations --- not a new structure introduced for this purpose.

\subsection*{G.3 What is the distance $D(H,P)$?}

The history-space distance is
\[
D(H,P)=-\ln\operatorname{Tr}\!\bigl(\Pi_P\,\rho_{H\to\mathrm{now}}\bigr),
\]
where $\rho_{H\to\mathrm{now}}$ is the state obtained by evolving the history $H$ forward to the present. It is the negative logarithm of the overlap between the present projector and the state that history would produce. A history leading naturally to the present records sits close, with small $D$ and large weight; a history that would require an improbable conspiracy to reproduce them sits far, with large $D$ and suppressed weight. The weight $P(H|P)\propto e^{-D(H,P)}$ is then exactly that overlap, normalized over the admissible histories.

\subsection*{G.4 Born recovery in the projective laboratory limit}

In the ordinary projective laboratory limit the evolved state is the prepared state and the present projector is the measured outcome, so the distance collapses to a single overlap,
\[
e^{-D(H,P)}=\operatorname{Tr}\!\bigl(\Pi_P\,\rho\bigr),
\]
which is the Born probability of that outcome. This is the branch $\alpha=1$ of the generalized weighting family. The recovery is a consistency condition rather than a fresh derivation: requiring that the weight reproduce exact Born statistics selects $\alpha=1$, and the substantive content is that the family contains a branch for which this holds at the same time as no-signaling. Nothing in the laboratory is changed; the operational theory is standard quantum mechanics.

\subsection*{G.5 No-signaling}

The companion condition forbids any operational bias channel whose statistics could depend on a spacelike-separated setting. Imposing it selects $\beta=0$. With $\beta=0$ the weight depends only on the local present records, so a choice made at one location does not alter the marginal statistics at another, and no-signaling holds exactly in the operational branch. Born recovery and no-signaling are therefore the two consistency conditions $\alpha=1$, $\beta=0$; the postulate's nontrivial claim is that they are simultaneously satisfiable, not that either is derived from outside.

\subsection*{G.6 Branch realization}

The postulate makes one realized present, not many. There is no forward branching into co-real macroscopic worlds, and no collapse event selecting among them. What is real is the present with its records; the multiplicity the weight ranges over is the admissible \emph{pasts} of that single present, not parallel futures. ``Which outcome occurred'' is therefore not a question about which branch the world fell into but a statement about which present records obtain, with the Born weight giving their statistics. Branch realization is thereby answered without the ontological cost of many worlds or the dynamical cost of collapse.

\subsection*{G.7 Arrow of time from conditional typicality}

Let $h=\{M_t\}_{t<t_0}$ be a macrohistory conditioned on present records $M_{t_0}$. If the count of compatible microhistories is $N_h$, then
\[
P(h\mid M_{t_0})\propto N_h.
\]
Under coarse-grained factorization,
\[
\ln P(h\mid M_{t_0})
\approx
\sum_{t<t_0}S(M_t)
\;+\;
\sum_{t<t_0}\ln T(M_{t+\Delta t}\mid M_t)
\;+\;
\mathrm{const},
\]
so the dominant record-compatible histories are those with entropy increasing toward the present. The arrow of time appears as a counting-dominance effect among record-compatible histories, not as a new dynamical coupling or a new law of laboratory probability. It is conditional typicality: given the records, overwhelmingly many admissible pasts show entropy rising toward them.

\subsection*{G.8 Memoryless dressing and the connection to $L_*$}

The same weight does work outside the foundations. Because $e^{-D}$ weights whole histories with no memory term, the dressing of an elementary defect is memoryless in substrate time: between one readout and the next the substrate relaxes to a fresh draw from the admissibility ensemble, so each dressing pass is statistically independent of the last. That independence is exactly the input the scale-setting recurrence of Appendix~D.4 requires: it lets the seven channels each carry the full sharing entropy $g_{\mathrm{share,eff}}$ and multiply. The mixing-time hierarchy that justifies treating each pass as a fresh draw is the $\sim\!10^{22}$ ratio of the electron Compton time to the substrate tick (Appendix~H.4). Postulate~III therefore enters the gravitational chain directly: its memoryless dressing is load-bearing in the length that sets $G_*$.

\subsection*{G.9 Relation to many-worlds, collapse, hidden variables, and decoherence}

Placed among the standard interpretations, Many-Pasts is a form of history-space realism. Unlike many-worlds, it posits no co-real forward branches; the only realized macrostate is the present. Unlike objective-collapse models, it adds no stochastic collapse term and leaves unitary evolution intact. Unlike hidden-variable theories, its probabilities are not ignorance over pre-existing sharp values. It is nearest to the consistent-histories framework and to decoherence-based accounts of classicality, but it differs from them in two ways: it carries an explicit normalizable weight $e^{-D(H,P)}$ over admissible histories, and it puts that weight to physical use outside the measurement problem, in the memoryless dressing of G.8. The proposal does not claim to out-derive the Born rule from more primitive axioms; it recovers Born statistics and no-signaling as the consistency conditions $\alpha=1$ and $\beta=0$ on the weight (G.4, G.5), and offers in addition a history-space account of branch realization and the arrow of time.


\subsection*{G.10 Familiar quantum examples: double-slit, EPR/Bell, and measurement}

In the familiar quantum examples, Many-Pasts changes the ontology rather than the laboratory calculation. In a double-slit experiment the particle is not treated as having secretly chosen one classical path while the other was unreal. Before a which-path record exists, the present detection event is supported by an ensemble of admissible microscopic histories through the apparatus, and the histories compatible with the final record include the alternatives whose amplitudes interfere in ordinary quantum mechanics. When a which-path record is made, the admissible history ensemble is conditioned on that record, the cross terms are removed by decoherence, and the interference pattern disappears. The probabilities are the usual Born probabilities; Many-Pasts supplies the history-space reading of why those are the probabilities being counted.

The same point applies to entangled pairs. In an EPR or Bell experiment the two detector outcomes are not produced by a signal sent from one wing of the experiment to the other. The present record is a joint record of the pair and the apparatus, and the admissible microscopic histories are weighted as histories of that whole entangled system. This gives the standard nonclassical correlations while preserving the local Born marginals, and because the local marginal statistics are unchanged, no observer can use the correlations to send a signal. The ``spooky'' feature is therefore not a superluminal force in spacetime; it is the global conditioning of the history ensemble on the shared entangled record.

Measurement is read the same way. A measurement does not add a new physical collapse law; it creates a durable record, and the relevant history ensemble is then the ensemble compatible with that record. Operationally this is ordinary quantum mechanics. Ontologically, the realized branch is not selected by a forward-propagating collapse event but by conditioning the present on its admissible microscopic pasts.

None of this is a new derivation of the quantum phenomena. In the operational branch the calculations and the probabilities are exactly those of standard quantum mechanics; what Many-Pasts adds is the ontology under that branch --- a history-space account of what is being counted and why --- with the double slit reading as unrecorded alternatives that remain in the admissible ensemble until a record conditions it, Bell correlations as the global weighting of a shared entangled record rather than a faster-than-light signal, and measurement as record-conditioned history selection rather than new collapse physics.

\section*{Appendix H: Microscopic Realization and Coarse-Graining}

Appendix~H addresses a different question from the weak-field appendices. Instead of asking whether the coefficient chain is internally closed, it asks whether a plausible microscopic realization exists in which the same scalar stiffness and defect ontology arise naturally.

\subsection*{H.1 GFT condensate realization and coarse-graining}

The candidate microscopic realization is a GFT/condensate picture with bosonic tetrahedral quanta $\phi(g_1,\ldots,g_4)$ and fermionic defects $\psi$. In the condensate regime, the coarse field may be written as
\[
\sigma(x)=\sqrt{n(x)}\,e^{i\theta(x)}.
\]
The hydrodynamic identity
\[
|\nabla_\mu \sigma|^2
=
\frac{(\nabla_\mu n)^2}{4n}
+n(\nabla_\mu\theta)^2
\]
shows that if
\[
S_{\mathrm{ent}}(x)=S_0+\alpha \ln\!\frac{n(x)}{n_{\mathrm{bg}}},
\]
then the coarse action contains a positive scalar stiffness
\[
\gamma \sim \frac{Z_\sigma n_{\mathrm{bg}}}{2\alpha^2}>0.
\]
The coarse source channel arises from fermionic face exclusion: what is macroscopically read as matter is a localized defect of the condensate, and the surrounding reduction of available occupancy is the long-wavelength field captured by the EFT. The microscopic appendix therefore supports the EFT without replacing it: the continuum kinetic term and source channel can arise from a concrete substrate realization, while a full first-principles derivation of every inhomogeneous continuum coefficient from the underlying kernel remains to be done.

It does not replace the explicit coefficient derivation given earlier, but it shows that the ontology and sign choices of the EFT are compatible with a concrete microscopic picture.

The condensate picture addresses the emergence of the continuum geometry. The complementary microscopic question --- the defect dynamics that fixes the faithful sector-resolution principle of Appendix~D.4 --- is taken up in the remainder of this appendix, where the principle is shown to follow from the physically realized admissibility ensemble of Part~II together with the Many-Pasts postulate.

\subsection*{H.2 The dressing Hamiltonian and the one-pass ground state}

The faithful sector-resolution relation of Appendix~D.4 was written there as a support-scale identity on the history space $\mathcal H_{\rm hist}=\bigotimes_{m=-3}^{3}\mathcal H_B^{(m)}$, using the refresh kernel $P_{\eta_*}(b,b')=p_{\eta_*}(b')$ on each sector layer. Its dynamical content is carried by an explicit defect Hamiltonian, which we now write down so that the choice of that kernel can be derived.

The seven face-label channels of the admissibility-closed ensemble define the orthogonal projectors $E_m$, $m=-3,\ldots,3$, of the face algebra $\mathcal A_7$. By Postulate~II the elementary defect is fermionic, so each channel carries an occupation number $n_m\in\{0,1\}$ with fermionic creation and annihilation operators $c_m^\dagger,c_m$ and $n_m=c_m^\dagger c_m$. When channel $m$ is occupied it is dressed by a cloud state described by a density operator $\rho_m$ on the boundary ensemble $\mathcal H_B$. The defect Hamiltonian is
\[
H
=
\underbrace{\sum_{m}\bigl(\varepsilon_0\,n_m-n_m\,F_m(\rho_m)\bigr)}_{\text{single-channel terms}}
+
\underbrace{\sum_{m<m'}V_{mm'}(\rho_m,\rho_{m'})}_{\text{inter-channel coupling}},
\]
with $\varepsilon_0$ the bare cost to occupy a channel, $F_m$ the free energy released by dressing channel $m$ with its cloud, and $V_{mm'}$ the residual coupling between the clouds of distinct channels. Two properties of the ground state of $H$ supply the two ingredients of the sector-resolution relation.

\paragraph{One-pass occupation and the exponent seven.}
A channel is occupied in the ground state whenever its dressing free energy beats the bare cost, $F_m(\rho_m)>\varepsilon_0$. The lightest charged defect is the configuration that minimizes $\varepsilon_0$. Because the seven channels are symmetry-equivalent, each binding with the same free-energy gain, that configuration is exactly the one for which all seven bind: any unbound sector is an excitation above the ground state. ``Lightest'' and ``resolves each of the seven sectors once'' are then the same statement: the elementary fermionic defect fills each face sector exactly once, with fewer sectors leaving unresolved label memory and repeated sectors describing excited or further-dressed states. This one-pass occupation supplies the factor seven in $7g_{\mathrm{share,eff}}$, the seven being the channel count $|M|=7$ fixed in Part~II.

\paragraph{Factorized cloud and the additive support.}
The dressing dimension multiplies across channels --- so that the seven equal contributions $g_{\mathrm{share,eff}}$ add rather than merge --- precisely when the joint cloud state is a product, $\rho=\bigotimes_m\rho_m$. Whether the ground state of $H$ has this product form is controlled by the inter-channel term $V_{mm'}$. Beyond the one-pass count, then, faithful sector resolution turns on the seven-channel cloud being a product of full-entropy channels: each channel must sample its whole ensemble, and $V_{mm'}$ must leave the channels uncorrelated.

The dressing Hamiltonian thus reduces the sector-resolution principle to that product structure in its ground state. The next three subsections establish it.

\subsection*{H.3 Slot coupling, layer factorization, and the additivity theorem}

Two distinct correlation structures appear in the closure data, and separating them is essential: conflating them leads to the false conclusion that the cloud cannot factorize.

\paragraph{The closure invariant is a pure pair coupling.}
Expanding $S^2=\bigl(\sum_i m_i\bigr)^2=\Sigma^2+2\sum_{i<j}m_im_j$ in $K^2(b)=48-\tfrac13(S^2-\Sigma^2)$ cancels the self-terms exactly and leaves the identity
\[
K^2(b)=48-\frac23\sum_{i<j}m_im_j.
\]
The admissibility weight $e^{-\eta_*K^2(b)}$ therefore contains only cross-terms between distinct face slots of a single boundary state; it does not factorize over those four slots. This is the exact origin of the residual correlation between face slots on the closed ensemble,
\[
I(\text{slot}_0;\text{slot}_1)=0.1545~\text{nats},
\qquad
\frac{I}{H(\text{slot})}=0.079,
\]
so the four faces of one tetrahedron are about eight percent correlated, a structural feature of the closure invariant.

\paragraph{Slot coupling does not obstruct layer factorization.}
The factorization the support relation requires is over the seven \emph{sector layers} $b^{(-3)},\ldots,b^{(3)}$ of $\mathcal H_{\rm hist}$, each layer being a full boundary state drawn from the entire $1680$-state ensemble. The slot coupling $-\tfrac23 m_im_j$ lives \emph{inside} a single layer's boundary state: it relates the four faces of that one tetrahedron and never couples layer $m$ to layer $m'$. The two structures act on different objects:
\begin{center}
\small
\begin{tabular}{p{0.26\textwidth}p{0.30\textwidth}p{0.30\textwidth}}
\hline
Structure & What it couples & Role \\
\hline
slot coupling $-\tfrac23 m_im_j$ & the four faces \emph{within} one boundary state $b$ & the eight-percent mutual information; lives inside each $\mathcal H_B^{(m)}$ \\
layer coupling $V_{mm'}$ & the seven sector layers $b^{(m)}$ of $\mathcal H_{\rm hist}$ & controls factorization; acts \emph{between} the factors \\
\hline
\end{tabular}
\normalsize
\end{center}
The substrate's intrinsic correlation, the most natural candidate obstruction to factorization, therefore acts at the wrong level to obstruct it: it is internal to a layer, not between layers.

\paragraph{Additivity of the history distance.}
The Many-Pasts history weight of Appendix~G.1 is a log-overlap, $D(H,P)=-\ln\operatorname{Tr}(\Pi_P\rho_{H\to\mathrm{now}})$. If the dressing state and the resolution projector factorize over channels, $\rho=\bigotimes_m\rho_m$ and $\Pi_P=\bigotimes_m\Pi_m$, the trace factorizes and the logarithm converts the product into a sum,
\[
\operatorname{Tr}(\Pi_P\rho)=\prod_m\operatorname{Tr}(\Pi_m\rho_m)
\quad\Longrightarrow\quad
D=\sum_m\bigl[-\ln\operatorname{Tr}(\Pi_m\rho_m)\bigr]=\sum_m D_m .
\]
Additivity of $D$ over the seven channels --- and hence the multiplicativity $\dim_{\rm eff}=\prod_m e^{g_{\mathrm{share,eff}}}=e^{7g_{\mathrm{share,eff}}}$ used in the length relation --- is therefore a theorem once the product structure holds. The implication needed downstream runs in one direction: a product cloud gives additive $D$ and the full $7g_{\mathrm{share,eff}}$. Whether the electron's dressing is that product is decided by two conditions, established in the next two subsections.

\paragraph{The ceiling and the two conditions.}
Subadditivity bounds the support from above. With each channel marginal fixed to the admissibility weight, $H(B_m)=g_{\mathrm{share,eff}}$, the joint entropy of a single readout obeys
\[
H\bigl(B_{-3},\ldots,B_3\bigr)\le\sum_{m=-3}^{3}H(B_m)=7g_{\mathrm{share,eff}},
\]
and the deficit
\[
\Delta=\sum_{m}H(B_m)-H\bigl(B_{-3},\ldots,B_3\bigr)\ge 0
\]
measures the total correlation among the seven channels. The full support $\dim_{\rm eff}=e^{7g_{\mathrm{share,eff}}}$ is reached precisely when each channel carries its full entropy $g_{\mathrm{share,eff}}$ and the channels are mutually independent, $\Delta=0$. The first is a condition on each layer in substrate time; the second is a condition on the seven layers at a single readout. The next two subsections establish them.

\subsection*{H.4 Each channel at full entropy: memorylessness in substrate time}

A channel reaches its full entropy $g_{\mathrm{share,eff}}$ only if its dressing keeps no memory of its own past. To see what memory would cost, consider the one-parameter family of per-channel kernels
\[
K_a(b,b')=a\,\delta(b,b')+(1-a)\,p_{\eta_*}(b'),
\qquad a\in[0,1],
\]
which repeats the current state with probability $a$ and otherwise redraws from the stationary weight. Every member has $p_{\eta_*}$ as its stationary distribution and the same single-time marginal, so the equilibrium ensemble cannot say which one governs the dressing. The per-channel conditional entropy $H(b'\mid b)=\sum_b p_{\eta_*}(b)\,H\bigl(K_a(b,\cdot)\bigr)$, evaluated on the exact $1680$-state ensemble, nonetheless slides with the memory $a$:
\begin{center}
\small
\begin{tabular}{ccc}
\hline
$a$ (memory) & per-channel $H(b'\mid b)$ & status \\
\hline
$0.0$ (refresh) & $7.41980=g_{\mathrm{share,eff}}$ & full entropy \\
$0.1$ & $6.99953$ & reduced \\
$0.3$ & $5.80153$ & reduced \\
$0.5$ & $4.40051$ & reduced \\
$0.9$ & $1.06642$ & reduced \\
\hline
\end{tabular}
\normalsize
\end{center}
Only the memoryless endpoint $a=0$ --- the refresh kernel $P_{\eta_*}(b,b')=p_{\eta_*}(b')$ of Appendix~D.4 --- returns the full $g_{\mathrm{share,eff}}$. The kernel carries more than its stationary marginal, so this first condition is settled by the dressing dynamics, not by the equilibrium ensemble. The requirement is in fact topological: single-label local moves conserve slot ordering, fracturing each parity copy into twenty-four sectors of thirty-five states, so no local kernel reaches the full ensemble at any parameter value (Appendix~D.4).

The dynamics settle it cleanly. Take any ergodic generator $Q$ with $p_{\eta_*}$ as its detailed-balance stationary weight --- a continuous-time single-label-swap generator built from the same admissibility weights serves, with no memorylessness assumed. The kernel over substrate-time $\tau$ is $e^{Q\tau}$, and the mutual information between its endpoints decays to zero,
\[
I(\tau{=}0.1)\approx 3.4~\text{nats},
\qquad
I(\tau{=}1)\approx 0.030,
\qquad
I(\tau{=}2)\approx 1\times10^{-4},
\]
so $e^{Q\tau}$ relaxes to the refresh kernel. The readout sits far inside that limit: the substrate time is $\tau_*=L_*/c\approx5.4\times10^{-44}$~s, the electron's Compton time is $\tau_e=\hbar/(m_ec^2)\approx1.3\times10^{-21}$~s, and their ratio $\approx2.4\times10^{22}$ is the same hierarchy $\lambda_e/L_*$ the length relation expresses. The ground state is read out some $10^{22}$ mixing times after relaxation, so each channel draws afresh and carries its full $g_{\mathrm{share,eff}}$. This is the Many-Pasts weight $P(H|P)\propto e^{-D(H,P)}$ seen on the dressing: a sum over pasts that keeps no trace of the particular one realized.

\subsection*{H.5 Independent channels: the electron as the lightest defect}

The second condition is that the seven channels are mutually independent at one readout, so that the deficit $\Delta$ of H.3 vanishes. This needs no separate dynamical postulate; it follows from what the electron is.

Keep the deficit explicit in the support relation. With each channel at its full entropy, the joint entropy is $7g_{\mathrm{share,eff}}-\Delta$, and the support reads
\[
\frac{\lambda}{L_*}=\frac23\,e^{\,7g_{\mathrm{share,eff}}-\Delta}.
\]
At a fixed substrate scale $L_*$, a one-bit charged defect whose dressing carries correlation $\Delta$ has spatial support $\lambda\propto e^{-\Delta}$, and through $m=\hbar/(c\lambda)$ a mass
\[
m\propto e^{\Delta}.
\]
Correlation among the dressing channels contracts the defect and raises its mass. Because $\Delta\ge 0$, the lightest one-bit charged defect is the one with $\Delta=0$: the product dressing. The electron is that defect, so its seven channels are independent and its support is the full $\dim_{\rm eff}=e^{7g_{\mathrm{share,eff}}}$. Correlated dressings are still allowed; they describe heavier, additionally dressed excitations, not the elementary anchor.

The selection is physical, not an appeal to ignorance. Independence is not read off from the absence of an inter-channel constraint in the ensemble; it is forced because correlation carries a mass, and the electron is defined as the lightest charged one-bit defect. The argument uses only the mass--entropy identification at the centre of the theory: the same relation $m=\kappa_m\,\Delta S$ that anchors the electron makes inter-channel correlation an excess of defect entropy, and so of mass.

With both conditions met, the scale-setting branch reads end to end without a free step. One-pass occupation fixes the seven channels (H.2); memorylessness in substrate time gives each its full entropy $g_{\mathrm{share,eff}}$ (H.4); the electron's lightness forces the channels independent (this subsection). Hence
\[
\frac{\lambda_e}{L_*}=\frac23\,e^{7g_{\mathrm{share,eff}}},
\qquad
L_*=\frac32\,\lambda_e e^{-7g_{\mathrm{share,eff}}},
\qquad
G_*=\frac{c^3L_*^2}{\hbar},
\]
with no gravitational quantity used as input.

The condensate realization gives the EFT kinetic term a natural microscopic origin, and the dressing analysis derives the faithful sector-resolution principle from the mass--entropy ontology, the one-bit electron anchor, and the memoryless dressing, so the substrate length and the induced scale follow from commitments the theory already makes. The first-principles derivation of every inhomogeneous continuum coefficient from the full kernel remains open.

\subsection*{H.6 Mass as a rate: the dressing operators}

The scale-setting branch above fixes the support exponent $7g_{\mathrm{share,eff}}$ and hence the ratio $\lambda_e/L_*$. Reading that ratio as a physical length used one identification: that the electron advances one Compton phase cycle per complete dressing recurrence. Memoryless sampling delivers a recurrence \emph{rate}, not by itself a Compton pole, so this subsection gives the recurrence an explicit operator form and asks what the phase clock must be. The identification does not close to zero residual; it decomposes into pieces that are either forced or falsifiable, with one identification (Fork~A below) remaining.

\paragraph{The rest mass is read as a recurrence rate, not a binding-energy sum.}
Writing the length relation as a mass,
\[
m_e c^2=\frac{\hbar c}{\lambda_e}=\frac{3}{2}\,\frac{\hbar}{\tau_*}\,e^{-7g_{\mathrm{share,eff}}},\qquad \tau_*=L_*/c,\qquad e^{-7g_{\mathrm{share,eff}}}=2.78\times10^{-23},
\]
shows the electron is lighter than the substrate quantum $E_*=\hbar c/L_*$ not by a small factor but by $e^{-7g_{\mathrm{share,eff}}}$. That suppression is a waiting time: seven independent full-entropy channels jointly reaching their resolved configuration, probability $e^{-7g_{\mathrm{share,eff}}}$ per substrate tick. This rules out a binding-energy reading. Seven channels each releasing free energy of order $g_{\mathrm{share,eff}}E_*$ would give a mass linear in $g_{\mathrm{share,eff}}$, of order $50\,E_*$ --- heavy, and wrong by twenty-three orders. The object that produces an exponentially light particle is the spectral gap of a stochastic generator whose slowest mode relaxes at $e^{-7g_{\mathrm{share,eff}}}$ per tick, not the ground-state energy of a Hamiltonian. This rate reading sets the scale, not the content of the mass--entropy map: the electron's committed entanglement is the one bit $\Delta S_f=\ln2$ of Section~13.2, while the recurrence rate fixes the exchange rate $\kappa_m$ at which that bit is read as a mass (the length anchor of Section~13.3). Mass measures committed entanglement; the rate fixes how small that measure is.

\paragraph{The killed-refresh return rate.}
By Postulate~III the dressing is memoryless: the refresh kernel $P(b,b')=p_{\eta_*}(b')$ redraws each tick from the admissibility Gibbs state. On the seven-layer history space the marked configuration $\star$ --- the resolved dressing, each channel at its designated typical configuration --- has stationary weight $\prod_m p_{\eta_*}(\star_m)=e^{-7g_{\mathrm{share,eff}}}$. The dynamics of ``at $\star$ versus not'' is a two-state renewal chain with per-tick return probability $e^{-7g_{\mathrm{share,eff}}}$; the killed (absorbing-at-$\star$) operator has leading eigenvalue $1-e^{-7g_{\mathrm{share,eff}}}$, so its gap is the raw return rate
\[
\Gamma_{\rm raw}=\frac{e^{-7g_{\mathrm{share,eff}}}}{\tau_*},\qquad \langle T_{\rm rec}\rangle=e^{+7g_{\mathrm{share,eff}}}=3.6\times10^{22}\ \text{ticks}.
\]
The physical electron Compton rate carries the transverse export factor $2/3$ of the length relation, $\Gamma_e=\tfrac32\Gamma_{\rm raw}$, so $m_ec^2=\hbar\Gamma_e=\tfrac32(\hbar/\tau_*)e^{-7g_{\mathrm{share,eff}}}$, consistent with $\lambda_e/L_*=\tfrac23 e^{7g_{\mathrm{share,eff}}}$. This is the recurrence of the scale-setting branch re-expressed as a transfer-operator gap; the rate scale was never the contested part.

\paragraph{The locality obstruction.}
The generator that carries this recurrence cannot be local. Define a local move as shifting one face label by $\pm1$ while preserving injectivity. Under these moves the $1680$ states fracture into
\[
48\ \text{components}=4!\ \text{orderings}\times2\ \text{orientations},\qquad \text{each of size}\ C(7,4)=35,
\]
verified directly on the ensemble. Slot ordering is a conserved charge of any local single-label dynamics, so the closed seven-sector loop that winds the phase once is not reachable by a local Hamiltonian; the generator that closes it is the non-local refresh kernel of Postulate~III. This is an honest cost of the construction: a referee asking for a local ``Hamiltonian spectrum'' is owed the statement that the operative generator is the non-local refresh super-operator. It is the content of memorylessness rather than a device introduced to evade locality.

\paragraph{The phase clock, by elimination.}
The mass is $\hbar$ times a phase-turn rate, so closing ``one turn per recurrence'' requires the turning clock and the refresh counter to be the same object, not two rates set equal. Two candidates exist: (B) a clock private to the defect, or (A) the global $\mathrm{U}(1)$ phase $\theta$ of the condensate order parameter $\sigma=\sqrt n\,e^{i\theta}$. A private clock is not forced to complete exactly one turn per recurrence; pinning it to the refresh would require declaring the rates equal, the identification one is trying to remove. The condensate phase is forced by counting: one complete refresh commits exactly one defect (the lightest charged defect binds the seven sectors once --- one fermion, Postulate~II), one committed defect is one unit of the condensate $\mathrm{U}(1)$ charge (that charge being the count of committed cells), and one unit of charge on a $\mathrm{U}(1)$ phase is one full turn. Hence
\[
\text{one refresh}=\text{one defect}=\text{one unit charge}=\text{one }2\pi\text{ turn},
\]
each link a count or a definition, with no rate chosen. A lone defect has no conserved tally for the refresh to increment, so the counting argument is unavailable to it: the clock is the medium's.

\paragraph{The charge-tilted generator.}
Tilting the renewal generator by a counting field $\phi$ conjugate to committed charge gives the reduced tilted transfer matrix $\left[\begin{smallmatrix}1-r & re^{i\phi}\\ 1 & 0\end{smallmatrix}\right]$ with $r=e^{-7g_{\mathrm{share,eff}}}$. Its leading eigenvalue is exactly $2\pi$-periodic in $\phi$, so committed charge is integer --- the fermionic $n_m\in\{0,1\}$ appearing as periodicity, not a tuned input --- and a complete seven-sector pass is one closed loop on the label cycle, winding number one. The identification ``one phase cycle per dressing recurrence'' therefore decomposes into three pieces, none carrying a knob:
\[
\underbrace{\Gamma=e^{-7g_{\mathrm{share,eff}}}/\tau_*}_{\text{memorylessness}+\text{entropy}}
\ \oplus\
\underbrace{\Delta Q=1\ \text{per recurrence}}_{\text{fermionic one-pass}}
\ \oplus\
\underbrace{\Delta\theta=2\pi\,\Delta Q}_{\text{Josephson}}.
\]
The Compton clock is now a derived winding rate rather than an asserted one.

\paragraph{What remains: Fork~A.}
Two things do not close, and are stated as such. First, the identification of the winding $\mathrm{U}(1)$ with the condensate phase, and its $2\pi$ with the Compton cycle (Fork~A), is consistent with these operators but not forced by them; its support is the elimination argument above, not a computation, and it is settled in H.7, where the mean-field condensate is shown to carry exactly one phase. Second, the residual is now a single integer --- that the electron is the $\Delta Q=1$ (one-loop) recurrence rather than a higher winding. Unlike the original identification this is falsifiable: the winding-$N$ recurrences are the muon and tau, and their rates fix $m_\mu/m_e$ and $m_\tau/m_e$ with no new input (Appendix~I.1). The phase-cycle identification has thus been replaced by counting plus a condensate relation, with the residual reduced to one integer that the lepton spectrum constrains.

\subsection*{H.7 The dressing functionals from the mean-field condensate}

The dressing Hamiltonian named in Section~22,
\[
H=\sum_m\big(\varepsilon_0\,n_m-n_m F_m(\rho_m)\big)+\sum_{m<m'}V_{mm'}(\rho_m,\rho_{m'}),
\]
carries three named functionals $\varepsilon_0,F_m,V$. This subsection derives them from the mean-field condensate. The result is not that everything closes: the two residuals of H.6 collapse into a single premise, and that premise is the one the geometry side has not yet established.

\paragraph{Mean-field reduction.}
In the GFT condensate the single-tetrahedron expectation is the order parameter over boundary data, $\langle\hat\varphi(b)\rangle=\sigma(b)=\sqrt{n(b)}\,e^{i\theta}$, with one global phase $\theta$ --- the Goldstone of the broken GFT number $\mathrm{U}(1)$. Extremizing the condensate free energy under the closure constraint gives the admissibility Gibbs state $|\sigma(b)|^2=n\,p_{\eta_*}(b)$, so the closure weight is a Boltzmann factor at inverse temperature $\eta_*$ with $K^2$ as energy. The substrate temperature and closure condition are then thermodynamic,
\[
T_*=\frac{1}{\eta_*}=33.48,\qquad \langle K^2\rangle_{\eta_*}=\tfrac32 T_*=50.22,
\]
the $3/2$ being equipartition over the three quadratic closure components, each at $\tfrac12 T_*$; this is the identity $\eta_*\langle K^2\rangle=3/2$ read as a temperature.

\paragraph{Two generators: $H$ selects, $L$ rates.}
The mass scale does not live in the spectrum of $H$. $H$ is combinatorial: its ground state fixes \emph{which} configuration the lightest defect is. The rate is the gap of the dissipative refresh super-operator $L$ whose stationary state is $p_{\eta_*}$ and whose memorylessness is Postulate~III, and the physical mass is that gap carrying the transverse export $2/3$ of the length relation, so
\[
m_e c^2=\tfrac32\,\hbar\,\mathrm{gap}(L)=\frac{3}{2}\frac{\hbar}{\tau_*}\,e^{-7g_{\mathrm{share,eff}}},\qquad \mathrm{gap}(L)=\frac{e^{-7g_{\mathrm{share,eff}}}}{\tau_*},\quad e^{-7g_{\mathrm{share,eff}}}=|\sigma(\star)|^2=\prod_m p_{\eta_*}(\star_m),
\]
the exponential smallness being the condensate density at the resolved configuration, not a fine-tuned eigenvalue. This is the operator content of H.6.

\paragraph{$\varepsilon_0$, the bare cost.}
$\varepsilon_0$ is the condensate chemical potential $\mu_0=\partial E/\partial N$, the energy to commit one cell against the mean field, of order $E_*=\hbar c/L_*$. Deriving it from the substrate names it as $\mu_0$: it is the single irreducible dimensionful input, relocated rather than removed, since a mass cannot come from combinatorics alone. Energies in $H$ are measured in units of $E_*$, so $\varepsilon_0$ is the one place the physical scale is attached, and the substrate temperature $T_*=1/\eta_*$ entering $F_m$ below is the dimensionless closure temperature ($K^2$ as energy); the physical mass is set by the rate $L$, not by the spectrum of $H$.

\paragraph{$F_m$, the dressing free energy.}
A bound channel carries a cloud $\rho_m$, a distribution over the boundary ensemble; the unbound channel carries none. The released free energy is the cloud entropy at the substrate temperature, $F_m(\rho_m)=T_*\,S[\rho_m]$. Extremizing it at fixed mean closure returns the Gibbs state $\rho_m=p_{\eta_*}$ with $S[\rho_m]=g_{\mathrm{share,eff}}$: the equilibrium cloud is the condensate density itself. The binding test $F_m>\varepsilon_0$ is met for all seven channels, and the lightest defect binds each sector once with a full-entropy cloud --- one-pass occupation, total bound entropy $7g_{\mathrm{share,eff}}$. This last equality bridges the two generators: the cloud entropy that $H$ maximizes in binding is exactly the rarity exponent that $L$ carries in its gap, since the joint recurrence needs all seven full-entropy clouds at once, probability $\prod_m e^{-S[\rho_m]}=e^{-7g_{\mathrm{share,eff}}}$. Binding and lightness are one extremum: the most bound defect, carrying maximal cloud entropy, is the rarest and hence the lightest. The $F_m$ step is otherwise a consistency loop --- its form is entropy, the condensate density is $p_{\eta_*}$, and the maximum-entropy condition hands $p_{\eta_*}$ back --- so it confirms the one-pass structure rather than independently generating it.

\paragraph{Two energies, kept distinct.}
Two quantities both invite the name ``energy of the defect,'' and conflating them is what would reintroduce the hierarchy the rate reading removes. The first is the \emph{separation cost}: the entropic free energy needed to unbind the defect from the network. In a medium whose only currency is shared entanglement there is no stored potential energy in the spring sense, so the cost of any operation is the free-energy cost of changing the admissible-arrangement count, and the natural energy of a defect is what it would take to dissolve it back into the substrate. That cost is large and linear in $g_{\mathrm{share,eff}}$ --- of order the $\sim50\,E_*$ that the binding-energy reading of H.6 produces, a number that is correct as a separation cost and wrong only as a mass. The second is the rest energy $m_ec^2=\hbar\Gamma_e$, which is the recurrence rate and is tiny, exponential in $-7g_{\mathrm{share,eff}}$. Binding-equals-lightness is the statement that these are one bound entropy $g_{\mathrm{share,eff}}$ read two ways: the deeper the binding, hence the larger the separation cost, the rarer the full re-completion, hence the smaller the rate. They are proportional through $g_{\mathrm{share,eff}}$ but not equal, and the distinction is load-bearing --- reading the separation cost as the rest mass would set the electron at tens of $E_*$ and undo the exponential suppression. The cost to separate and the rate of recurrence are two faces of one entanglement, not the same number.

\paragraph{$V$, the inter-channel coupling.}
This is the one piece with new content. Channels couple by exchanging the field they share, and the framework has exactly one infrared field: the scalar capacity fluctuation $\phi=\delta S_{\mathrm{ent}}$, coupled to matter through the trace channel $-\kappa\chi S_{\mathrm{ent}}$ of Section~10. That scalar is the sector singlet $|s\rangle=\tfrac{1}{\sqrt7}\sum_m|m\rangle$, so each channel couples to it through $\langle m|s\rangle=1/\sqrt7$, the same for every $m$. Integrating out the single shared scalar gives an outer product,
\[
V_{mm'}=v\,\langle m|s\rangle\langle s|m'\rangle=\frac{v}{7}\quad\text{for all }m,m',
\]
with $v$ the zero-momentum scalar-exchange amplitude: the uniform $7\times7$ matrix is rank one. The per-incidence overlap that enters the ladder is $\langle m|s\rangle\langle s|m'\rangle=1/7$, doubled by orientation to the $2/7$ of Appendix~I.1, while the overall coupling $v$ sets the mass scale and not the ratios. One exchanged mode is one outer product, which is rank one. A second light field would add a second outer product, raising the rank to two and splitting the uniform entries, which breaks the single $2/7$ and the $(2/7)^{N^2}$ ladder. So $\mathrm{rank}(V)=1$ is not an assumption layered onto the construction; it is the statement ``one scalar field,'' the same premise that gives no anisotropic stress and $\Phi=\Psi$ in the PPN sector (Section~15). A higher-rank $V$ is a second light field, which that sector already excludes. This also supplies the minimum-mass content used in H.5: rank-one is the minimal inter-channel correlation, so the lightest charged defect is forced to it, and ``the coupling is rank one'' and ``this is the lightest defect'' are the same statement.

\paragraph{The collapse, and the wall.}
The two residuals of H.6 --- (A) the winding phase is the condensate $\mathrm{U}(1)$ and its $2\pi$ is the Compton cycle, and (B) the lightest defect's coupling is the rank-one singlet --- reduce to one fact. The mean-field condensate has exactly one global phase, so the mode that mediates $V$ and the mode that winds the Compton clock cannot be different: there is no second $\mathrm{U}(1)$ to be the wrong one, which settles the underdetermination of Fork~A. And rank-one $V$ is ``one scalar.'' A and B reduce to the same single-scalar condensate assumption --- a single scalar field on the mean-field condensate --- which is the founding premise of the framework, tested and passing in the charged-lepton sector to sub-percent. The cost is that this premise is now load-bearing twice. Every step here is conditional on the GFT substrate supplying that condensate, with density $|\sigma|^2=p_{\eta_*}$, one $\mathrm{U}(1)$, one scalar infrared mode, and the emergent four-geometry to go with it --- the geometry-side compatibility problem left open in Section~22 and H.1. The defect sector and the geometry sector are two outputs of the same $\sigma$: the sub-percent lepton ladder cannot be banked as evidence independent of the geometry problem, because it is the same condensate graded on its other output. The construction falls out cleanly because it rests on the single-scalar condensate, and that condensate is the unproven object; the two are one coin. Writing a fully explicit $H$ past this point would return $\varepsilon_0,F_m,V$ as they were put in, and the rank-one $V$ is already as derived as it gets without the condensate itself.


\paragraph{Reproduction.}
The ensemble invariants ($\eta_*$, $\langle K^2\rangle_{\eta_*}$, $g_{\mathrm{share,eff}}$, and the identity $\eta_*\langle K^2\rangle=3/2$), the killed-refresh return rate, the $48\times35$ locality fracture, the $2\pi$-periodicity of the charge-tilted generator, the rank-one spectrum $\{1,0^6\}$ of $P_{\rm sing}$, and the charged-lepton ratios of Appendix~I.1 are reproduced end to end from the bare admissibility ensemble by a standalone script provided as supplementary material (requiring only \texttt{numpy} and \texttt{networkx}).

\section*{Appendix I: Mass and Gauge Extensions}

Appendix~I collects sectors that are structurally connected to the same entanglement logic but are not part of the closed weak-field core. They are kept here because they show how the framework may extend, not because the main derivation depends on them.

\subsection*{I.1 Charged-lepton spectrum from the shell algebra}

Before the algebra, the physical picture. The electron is the ground-state one-bit charged defect --- the $N=0$ configuration whose seven-channel dressing sets the substrate length (Appendix~D.4). The muon and tau are not new kinds of object but heavier shell excitations of the same defect: each adds one radial entanglement shell, realized microscopically as the exclusion of one face label from the seven-state alphabet of Appendix~B. The number of generations is not put in by hand; it is limited by the closure machinery, which stops having anything to weight after the third admissible shell, so there are three charged leptons and not four. The mass of each lepton is the exponentiated dressing free energy of its shell source, in the same log-overlap form that carries the gravitational dressing (Appendix~H.3). This sector is a structural extension of the mass--entropy ontology, outside the closed weak-field chain. With the dressing operators of Appendix~H.6 in place it is operator-backed: the muon and tau are the winding-$N$ recurrences of the same generator, so the ratios below follow with nothing fitted to lepton data, conditional on the Fork~A identification of Appendix~H.7. Because those ratios carry no dimensionful input, the sector is a parameter-free, falsifiable test of the deeper construction rather than only a cross-check.

With that picture fixed, the spectrum is written as a shell expansion,
\[
\log\frac{m_N}{m_e}=B_0 N+A_0 N^2,
\qquad
N=0,1,2,
\]
with the electron the $N=0$ ground state and the muon and tau the $N=1,2$ radial entanglement-shell excitations of the same core structure. Each shell excitation is the exclusion of one face label from the seven-state alphabet of the boundary ensemble in Appendix~B.

\paragraph{Three-generation termination from closure-spectrum collapse.}
Combinatorial injectivity alone allows up to $N=3$ shells, since the reduced alphabet must retain at least four distinct labels to support an injective face assignment ($7-N\geq 4$). The sharper structural bound comes from the admissibility-closure machinery. For the reduced ensemble at shell $N$, the closure invariant $K^2$ inherits a discrete spectrum whose dispersion gives the admissibility kernel $e^{-\eta K^2}$ real weight. For $N=0,1,2$ the $K^2$ spectrum has multiple distinct values, and the closure equation $\langle K^2\rangle_\eta=3/(2\eta)$ has a non-degenerate solution.

At $N=3$ the spectrum collapses. With only four labels left, every injective assignment to the four faces uses all of them, and $K^2$ depends on the labels only through the permutation-invariant sums $S$ and $\Sigma^2$, so every microstate of the reduced four-label ensemble carries the same value of $K^2$. With a single spectral value $K_0^2$ the closure equation $\langle K^2\rangle_\eta=3/(2\eta)$ still fixes $\eta=3/(2K_0^2)$ formally; what fails at $N=3$ is selection, not solvability. The kernel $e^{-\eta K^2}$ acts trivially when every microstate carries the same $K^2$, so it weights nothing and leaves no admissibility dispersion to define a branch. Shell $N=3$ therefore does not define a non-degenerate admissibility-closed branch in the same sense as the earlier shells: the cutoff rests on the collapse of the $K^2$ dispersion, not on the closure equation admitting arbitrary $\eta$.

The number of charged-lepton generations is consequently the number of shells with a non-degenerate closure branch,
\[
N\in\{0,1,2\}\qquad\Longrightarrow\qquad\text{three generations,}
\]
a structural bound sharper than the combinatorial one. It is this mechanism, rather than injectivity alone, that fixes the generation count at three rather than four.

\paragraph{Linear coefficient from the first-shell reduced-alphabet entropy.}
The linear coefficient is the entropy cost of adding one shell excitation. With one face label excluded, the reduced ensemble has
\[
\Omega_1=2\,P(6,4)=720=6!,
\]
i.e., the full symmetric group $S_6$ with orientation degeneracy. Direct admissibility evaluation on the reduced ensemble gives $g_1=\ln\Omega_1=\ln 720$ to within $0.1\%$, since the admissibility correction is small on the reduced ensemble where the alphabet symmetry is nearly intact. Hence
\[
B_0=\ln\Omega_1=\ln 720=6.579\ldots,
\]
compared with the empirical fit $B_0=6.586$ (agreement at $0.1\%$).

\paragraph{Why the same $720$ at every shell.}
The per-shell factor is the \emph{same} $\Omega_1=720$ for each added shell, not a decreasing sequence $720\cdot240\cdots$ in which each shell excludes a further label. The two readings diverge sharply at the tau: cumulative exclusion would give $m_\tau/m_e\simeq1151$ ($-67\%$), whereas the repeated factor gives $3454.6$ ($-0.65\%$), so the data require the same reduction at each shell. The structural reason is the rank of the inter-shell coupling. The shells couple only through the permutation-invariant singlet projector $P_{\rm sing}$, whose spectrum is $\{1,0,0,0,0,0,0\}$: rank one. A rank-one projector pins exactly one shared direction, and it is the same direction however many shells are already stacked, so each added shell surrenders only that one shared singlet direction and retains the other six --- the same direction at every shell, not a fresh label removed cumulatively. The injective count realizes this exactly: $P(6,4)/P(7,4)=360/840=3/7$, so $1680\to720$ once and again at each subsequent shell. This also separates the two reduced-alphabet notions that otherwise look inconsistent. The \emph{generation cutoff} above counts $7-N$ surviving labels because it asks whether a non-degenerate closure \emph{branch} still exists; the \emph{mass factor} here is the rank-one singlet pinning one shared direction per shell. They are different operations, which is why the cutoff scales with $7-N$ while the ladder carries a fixed $720$.

\paragraph{Quadratic coefficient from singlet-projection shell algebra.}
The quadratic coefficient follows from three premises already in the framework: (i) mass arises from scalar capacity response (Postulate II); (ii) the scalar EFT response is quadratic in the total source; (iii) the coarse scalar branch projects onto the permutation-invariant singlet of the seven-state face alphabet.

Let
\[
\mathcal H_7=\mathrm{span}\{|m\rangle:m=-3,\ldots,3\},
\qquad
|u\rangle=\frac{1}{\sqrt 7}\sum_{m=-3}^{3}|m\rangle,
\qquad
P_{\rm sing}=|u\rangle\langle u|.
\]
A shell excitation that marks the face state $|a_r\rangle$ contributes to a coherent shell source
\[
|J_N\rangle=\sum_{r=1}^{N}|a_r\rangle.
\]
The coarse scalar field $\delta S$ is permutation-symmetric on face labels and therefore couples only through $P_{\rm sing}$. The scalar-channel response is quadratic in the total source, so the relevant object is the source norm
\[
\langle J_N|P_{\rm sing}|J_N\rangle
=
\sum_{r,s=1}^{N}\langle a_r|P_{\rm sing}|a_s\rangle.
\]
Because $\langle a|P_{\rm sing}|b\rangle=1/7$ for any $a,b$, this reduces to $N^2/7$ independent of which specific face labels are excluded. Including the binary orientation degeneracy of the boundary ensemble, the orientation-summed scalar transfer weight per ordered shell incidence is
\[
\mathcal M_{rs}=2\,\langle a_r|P_{\rm sing}|a_s\rangle=\frac{2}{7}.
\]
The ordered-pair count decomposes as $N^2=N+2\binom{N}{2}$: $N$ self-incidences along the diagonal and $2\binom{N}{2}$ directed cross-incidences. Two distinct steps enter here, and only the first is the quadratic EFT response. The quadratic singlet response fixes that there are $N^2$ ordered incidences, each carrying the per-incidence overlap $\mathcal M_{rs}=2/7$; it does not by itself make the mass multiplicative. The passage from these $N^2$ additive incidences to a multiplicative factor is the log-overlap mass map: the mass is the exponential of a dressing free energy whose pairwise term is $\sum_{r,s}\ln\mathcal M_{rs}=N^2\ln(2/7)$, the same log-overlap form that carries the gravitational dressing (Appendix~H.3) and stated explicitly below. With that map, the $N$-shell dressing is
\[
\prod_{r,s=1}^{N}\mathcal M_{rs}=\left(\frac{2}{7}\right)^{N^2},
\]
so
\[
A_0=\ln\frac{2}{7}=-1.253\ldots,
\]
compared with the empirical fit $A_0=-1.255$ (agreement at $0.15\%$).

The ordered-bilinear structure is not a free assumption. It is the unique scalar-channel response to a multi-source state under the three premises just stated: the source norm $\langle J_N|P_{\rm sing}|J_N\rangle$ runs over both indices independently because the integrated-out scalar contribution is quadratic in $J_N$.

\paragraph{Mass ladder and numerical accuracy.}
Combining the two coefficients,
\[
\boxed{
\frac{m_N}{m_e}=720^N\left(\frac{2}{7}\right)^{N^2},
\qquad N=0,1,2.
}
\]
Equivalently,
\[
\log\frac{m_N}{m_{N-1}}=\ln 720+(2N-1)\ln\frac{2}{7},
\]
so each added shell contributes $\ln 720$ plus an odd number of new scalar-overlap incidences ($1,3,5,\ldots$, summing to $N^2$). The second difference is fixed at
\[
\Delta^2\log m_N=2\ln\frac{2}{7}=-2.506\ldots,
\]
compared with the empirical $-2.509$ ($0.15\%$). With no parameters fitted to lepton data,
\[
\frac{m_\mu}{m_e}=\frac{720\cdot 2}{7}=\frac{1440}{7}=205.71,\qquad
\text{PDG: }206.77\qquad(0.5\%);
\]
\[
\frac{m_\tau}{m_e}=720^2\left(\frac{2}{7}\right)^4=3454.56,\qquad
\text{PDG: }3477.23\qquad(0.65\%).
\]
\paragraph{The dressing map, explicitly.}
The three factors follow from one construction. Write the $N$-th charged lepton as a coherent source raised on the electron ground dressing by exciting $N$ shells, $\mathcal S_N=\sum_{k=1}^{N}s_k$, where each $s_k$ excludes one face label from the seven-state alphabet. Its dressing free energy is a log-overlap functional of the kind used for the gravitational dressing in Appendix~H.3, and it separates into a per-shell term and a pairwise term. Each shell carries the reduced-alphabet entropy $\ln\Omega_1=\ln 6!=\ln 720$, the permutation count of the six surviving labels, which builds the linear factor $720^N$. Each ordered pair of shells contributes a singlet log-overlap: a single label projected on the closure singlet $u=\tfrac1{\sqrt7}(1,\dots,1)$ of Section~8 has weight $|\langle u|m\rangle|^2=1/7$, and orientation summation doubles it to $2/7$. The $N$ shells form $N^2$ ordered pairs, so the pairwise term is $N^2\ln(2/7)$, and exponentiating the free energy returns
\[
\frac{m_N}{m_e}=\Omega_1^{\,N}\left(\frac{2}{7}\right)^{N^2}=720^N\left(\frac{2}{7}\right)^{N^2}.
\]
The coefficients are fixed, not chosen: $720=6!$ and the $N^2$ pair count are combinatorial, and $2/7$ is the orientation-summed singlet projection behind $\Omega_{\rm tet}=2\,P(7,4)=1680$. What remains an input is the form of the map --- that the lepton mass is the exponentiated log-overlap free energy of the coherent multi-shell source --- a log-overlap structure of the kind that carries the gravitational dressing, taken here as the form of the mass map.

\paragraph{The support exponent, tested against the spectrum.}
The ladder also tests the exponent in the support-to-length identification of Appendix~D.4, which reads the physical ratio as the first power of the entropy support, $\lambda/L_*\propto e^{H}$. Letting that power float, $\lambda/L_*\propto e^{\alpha H}$, turns each lepton ratio into a measurement of $\alpha$, since $m_N/m_e=\bigl[720^N(2/7)^{N^2}\bigr]^{\alpha}$:
\[
\alpha_N=\frac{\ln(m_N/m_e)_{\rm PDG}}{\ln\bigl[720^N(2/7)^{N^2}\bigr]},
\qquad
\alpha_\mu=1.0010,
\qquad
\alpha_\tau=1.0008.
\]
Read the same way, the electron-to-$G_*$ branch needs $\alpha=0.9999$ to reproduce the observed Newton constant, so three determinations across two sectors land within $10^{-3}$ of unity and bracket it. What the leptons settle is the power, not its final digit: a square-root map ($\alpha=\tfrac12$) would put $m_\mu/m_e$ near $14$ and a quadratic map ($\alpha=2$) near $4\times10^{4}$, and only the linear rule reproduces the spectrum with the natural shell count $720=6!$. The identification that sets the substrate scale is therefore not a single assumption tuned to gravity; the same linear exponent is preferred in a sector that never enters the $G_*$ chain.

\paragraph{Status and remaining audit.}
The three-generation termination, the linear coefficient, and the quadratic coefficient are thus fixed from machinery already native to the framework, with no parameters fitted to lepton data: the collapse at $N=3$ is the permutation-invariance of $K^2$, and the $720$, $2/7$, and $N^2$ are the reduced-alphabet, singlet-projection, and ordered-pair counts of the map above. The form of the map is not a separate posit. Mass is $\hbar/(c\lambda)$ with $\lambda/L_*=e^{H}$, the same support-to-length identification that fixes $L_*$ in Appendix~D.4; the ladder is its $N>0$ continuation, with the electron the $N=0$ entry, and the equality of the support exponent across the two sectors is the $\alpha=1$ check above. The per-shell term is the log-overlap additivity theorem of Appendix~H.3, and the pairwise term is the singlet overlap of shells sharing the seven-state alphabet. Beyond that one identification, the only lepton-specific input is the structure of the excited source --- a product across the six spectator channels, coherent in the shared singlet --- which produces the $2/7$ pairings.

The electron anchor remains the clean weak-field entry point; the heavier charged leptons are conditionally derived outputs of the same shell algebra, not closure-defining ingredients of the gravitational chain. Composite hadrons remain part of the dressed bound-state entropy program rather than a completed output of the present construction. The neutrino and quark sectors are not treated here. The same closure-spectrum-collapse argument should apply universally if the framework is right --- the Standard Model has three generations across all fermion sectors --- but verifying that requires extending the shell construction to the relevant defect classes.

\subsection*{I.2 Gauge-structure extension and Standard Model ontology}

The framework's primary target is the gravitational and dark sector. Standard Model gauge structure is hosted on the substrate as external content rather than derived from it. This subsection specifies which gauge components have substrate-natural analogues, which are substrate-compatible but external, and how the baseline-redundancy template accommodates the full Standard Model gauge group $SU(3)_c\times SU(2)_L\times U(1)_Y$.

\paragraph{Baseline-redundancy template.}
The same logic that organizes the gravity sector extends to any conserved-charge sector. For a conserved charge $Q$, introduce an entropy-like potential $S_Q(x)$ and require that physical observables depend only on differences of that potential, not on its absolute baseline. Localizing that redundancy requires a compensating connection. In the Abelian case,
\[
D_\mu S_Q=\partial_\mu S_Q-qA_\mu,
\qquad
S_Q\to S_Q+\alpha(x),
\qquad
A_\mu\to A_\mu+\frac{1}{q}\partial_\mu\alpha,
\]
giving Maxwell-type dynamics for $A_\mu$. The non-Abelian generalization, for multiplet-valued entropic potentials transforming under a compact Lie group $G$, gives Yang--Mills covariant derivatives and field strengths in the standard form.

The template tells us what form a gauge interaction takes once a conserved-charge sector with baseline redundancy is present. It does not select which specific groups are physically realized.

\paragraph{Substrate-natural components.}
Two ingredients of Standard Model gauge structure are naturally accommodated by the substrate.

\emph{$U(1)_Y$.} The baseline-redundancy template naturally accommodates an Abelian conserved-charge sector of the hypercharge type. The associated entropy potential is conserved-charge in character, and localizing its baseline redundancy produces the corresponding Maxwell-type dynamics. The argument is independent of the gravity-sector derivation, applying the same template to a separate conserved-charge potential.

\emph{$SU(2)_L$.} The substrate naturally carries an $SU(2)$-representation structure through its half-integer face data. The spin-$3/2$ commitment forced by maximum-capacity channel selection and tetrahedral injectivity (Section~5) carries an $SU(2)$ action on face data through the standard double-cover relation. The static $K^2$ ensemble breaks this $SU(2)$ to its $U(1)$ Cartan via the magnetic-quantum-number projection, but the full $SU(2)$ is recoverable at the level of the quantum face data the projection discards. Identifying this representation-theoretic $SU(2)$ with the internal weak group $SU(2)_L$, including chirality and doublet assignments, remains external to the present derivation.

\paragraph{Substrate-compatible but external: $SU(3)_c$.}
Color $SU(3)$ does not emerge from the substrate by any of the standard mechanisms (direct decomposition of the seven-state face algebra, Seiberg duality, preon compositeness, topological soliton, monopole condensation). The obstructions are structural: $SU(3)$ is not a spin group, so it does not inherit from rotational covariance; its fundamental representation requires three equivalent (non-hierarchical) states, while the substrate's natural three-fold structures (shell index, $K^2$ classes, residual face positions) are all hierarchical; its rank-2 Lie algebra requires two independent quantum-number axes, while the substrate provides only rank-1 axis structures; and its non-Abelian commutation relations cannot be reproduced from the substrate's scalar edge couplings.

A positive ontological reading is nonetheless available within the framework. Color is the substrate's name for an internal three-valued label on fermionic defects that lives in the non-singlet complement of the face-state algebra. The scalar gravitational branch couples to defect sources only through the permutation-invariant singlet (Appendix~D), so the macroscopic scalar field is blind to non-singlet structure by construction. Color labels and color dynamics therefore live in a sector that gravity cannot resolve: the gravitational sector and the color sector occupy orthogonal subspaces of the substrate's algebra. The strong interaction can be represented, within this hosting picture, as the local gauge dynamics of the color label under the general baseline-redundancy template. The choice of $SU(3)$ specifically is empirical input; it is the realized non-Abelian structure that fits the template.

This account explains several structural features of color without claiming to derive them. Gravity's blindness to color follows from the scalar branch being the singlet projection of the substrate. Color confinement itself is not derived here; the substrate account provides that the macroscopic scalar-gravity branch resolves only singlet structure and is therefore blind to color non-singlets, which is consistent with the observed fact that hadronic macroscopic signatures are color-singlet but is not a replacement for the QCD confinement mechanism. The independence of color from generation, charge, and spin reflects the independence of the corresponding substrate sectors: color in the non-singlet complement of the face algebra, generation in the shell-excitation index (Appendix~I.1), electromagnetic charge in the baseline redundancy of a separate conserved potential, and spin in the fermionic face data.

\paragraph{A sharper reading: the transverse complement as a color arena.}
The obstructions above isolate to a single contingent ingredient, which turns ``$SU(3)$ is external'' into a definite derivation target. Write the face-state space as
\[
\mathcal H_7=\langle s\rangle\oplus\mathcal H_\perp,\qquad |s\rangle=\tfrac{1}{\sqrt7}(1,1,1,1,1,1,1),\qquad \dim_{\mathbb R}\mathcal H_\perp=6,
\]
the same singlet/non-singlet split that carries the scalar return operator and the lepton ladder. Suppose the six real transverse directions carry a complex structure: a preferred pairing into three complex modes, $\mathcal H_\perp\simeq\mathbb C^3$. The local frame redundancy of a three-complex-dimensional space is $U(3)$, and its overall phase is the condensate $U(1)$ already in use, so the traceless remainder is $SU(3)$. The rank-two algebra and eight generators come from the frame redundancy of the complex space itself, not from substrate quantum-number axes, so this route is not blocked by the rank obstruction above. The representation content then follows with no further input: a single defect is a vector in the $\mathbf 3$, hence non-singlet, and the scalar branch resolves only singlets, so lone defects are not asymptotic; $\mathbf 3\otimes\bar{\mathbf 3}=\mathbf 1\oplus\mathbf 8$ supplies the meson's singlet, and the antisymmetric $\epsilon_{abc}$ in $\mathbf 3\otimes\mathbf 3\otimes\mathbf 3$ supplies the baryon's. Three color directions and their conjugates, eight gluons, the absence of free single defects from the singlet spectrum, and gravity's blindness to color all read off one split.

\paragraph{The missing ingredient, and a route ruled out.}
The skeleton rests entirely on the complex structure $6_{\mathbb R}\to3_{\mathbb C}$, which the construction does not supply. Without a preferred pairing of the six transverse directions into three complex ones, the symmetry of the complement is $SO(6)$, not $SU(3)$, and the pairing must in addition be degenerate, the three directions equivalent --- which is the original reason the substrate's hierarchical three-fold structures (shell index, $K^2$ classes, residual face positions) do not furnish color. The most concrete candidate is a cyclic readout $T:|m\rangle\mapsto|m{+}1 \bmod 7\rangle$, whose Fourier modes would pair as conjugates $(k,7{-}k)$ into $\mathbf1\oplus3_{\mathbb C}$; it does not survive contact with what the labels are. They are not residues modulo seven but the weights $m=-3,\dots,3$ of the spin-$3$ multiplet ($j_{\mathrm{eff}}=3$, Section~5), a real and irreducible representation of the rotation group: it carries no invariant complex structure and does not split off the singlet under rotations at all. The closure invariant uses the integer weights, $K^2=48-\tfrac13(S^2-\Sigma^2)$ with $S=\sum_i m_i$ and $\Sigma^2=\sum_i m_i^2$, so the cyclic shift is not a symmetry of the dynamics: shifting $(-3,-2,-1,0)$ by one carries $(S,\Sigma^2)=(-6,14)$ to $(-2,6)$. A cyclic $\mathbb Z_7$ would deliver $\mathbf1\oplus\mathbf3\oplus\bar{\mathbf3}$, but it is not the structure the seven labels carry. The status is therefore the one above and no stronger: color has a natural location in the non-singlet complement the scalar branch cannot resolve, but the operator that would make that complement $\mathbb C^3$ is absent, and the cyclic route to it is excluded by the spin-weight structure. A derivation of $SU(3)_c$ would need an additional readout symmetry independent of the rotational dynamics, not presently exhibited; the skeleton reaches none of quark flavor, fractional charge, the QCD scale, confinement, or any hadron mass.

\paragraph{A better arena: tetrahedral routing.}
A more natural three is available, not from the seven labels but from the cell. Four boundary faces admit exactly three pairings,
\[
A=(12\,|\,34),\qquad B=(13\,|\,24),\qquad C=(14\,|\,23),
\]
the opposite-edge routings of the tetrahedron, and the tetrahedral symmetry group permutes them as the full $S_3$, so the three are symmetry-equivalent with no channel singled out --- the property the cyclic route could not get from the ordered spin labels. Read as quantum amplitudes the three routings form a complex vector $q\in\mathbb C^3$; if the local routing basis is a gauge choice, the comparison between neighbouring cells is a link $U_{yx}\in U(3)$, and removing the common phase already carried by the charged/condensate $\mathrm U(1)$ leaves $SU(3)$ with eight connection modes. Quarks are open routing vectors $q^a$, antiquarks the conjugate, and singlet closure is the standard count --- $q^a\bar q_a$ for the meson, $\epsilon_{abc}q^aq^bq^c$ for the baryon, the latter using the oriented three-form the cell's orientation/parity primitive already supplies.

\paragraph{What the closed ensemble gives, and why.}
The decisive question is whether the admissibility kernel treats the three routings as equivalent and independent, $K_{\mathcal R}\propto I_3$. It does not. For the natural pair observable $R_A=m_1m_2+m_3m_4$ and its images $R_B,R_C$, the kernel over the admissibility-weighted $1680$-state ensemble is symmetric with nonzero off-diagonal entries, so its spectrum is one isolated value and a degenerate pair: the routing triplet splits as $\mathbf3=\mathbf1\oplus\mathbf2$, not an irreducible color triplet. The reason is exact,
\[
R_A+R_B+R_C=\sum_{i<j}m_im_j=\tfrac12\bigl(S^2-\Sigma^2\bigr)=72-\tfrac32 K^2,
\]
so the symmetric routing mode is an affine function of the closure invariant $K^2$: it is the scalar sector itself, in routing variables. The closed cell offers only two genuinely transverse directions in the three-dimensional routing space (distinct from the six-dimensional label complement above), the third routing direction already committed to gravity, and a color $\mathbb C^3$ cannot be a closed-cell observable.

\paragraph{The open-flux direction.}
The split locates the missing principle rather than closing the route. The closure tie forcing $R_A+R_B+R_C\propto K^2$ is a property of closed scalar observables; an open routing imbalance carries no such constraint, and its three routes are independent local amplitudes. The natural object is then a lattice gauge action,
\[
S_{\mathrm{open}}=\sum_{\langle xy\rangle}\bigl\|q_y-U_{yx}q_x\bigr\|^2,\qquad q_x\in\mathbb C^3,\ \ U_{yx}\in SU(3),
\]
which exists precisely when open non-singlet routing flux is parallel-transported between cells rather than memorylessly refreshed --- the same line that separates the closed singlet sector (refreshed and scalar-projected, the gravity and lepton branches) from a non-singlet flux that cannot be locally redrawn without breaking continuity. This is one added dynamical principle, not a consequence of the three postulates, and it is what would release the three routings from the closure tie and make them a local $SU(3)$ frame. The status is therefore sharp, and is offered as a research direction rather than a result: the tetrahedral cell supplies three equivalent routing channels and an orientation form, the closed dynamics give $\mathbf1\oplus\mathbf2$ with the singlet equal to $K^2$, and color would live in an open-flux gauge sector governed by a transport principle the present construction does not contain.

\paragraph{Hadronic matter and the mass--entropy bridge.}
The mass--entropy bridge $m=\kappa_m(\ell)\,\Delta S$ applies to fermionic defects regardless of their color label. For charged leptons, the defect is color-singlet and the dressed entropy is well-approximated by the shell-excitation budget of Appendix~I.1, with mass ratios matching observation at the sub-percent level. For hadrons, the dressed entropy is dominated by the QCD-internal contributions of confinement-scale gluonic flux, trace-anomaly structure, chiral vacuum reorganization, and quark binding. The mass--entropy bridge then organizes the full dressed bound-state entropy budget,
\[
S_{\mathrm{ent,H}}^{\mathrm{dressed}}=S_{\mathrm{defect}}+S_{\mathrm{bind}}+S_{\mathrm{conf}}+S_{\chi\mathrm{SB}},
\]
without requiring the framework to re-derive QCD. The structural claim is compatibility: the dressed entropy budget is the quantity QCD calculates, and the mass--entropy bridge converts it to inertial mass through the same running $\kappa_m(\ell)$ used for elementary sectors.

\paragraph{Scope summary.}
The Standard Model's full gauge group $SU(3)_c\times SU(2)_L\times U(1)_Y$ is hosted on the substrate at three distinct levels of derivational support: $U(1)_Y$ from baseline redundancy (substrate-natural via independent argument); $SU(2)_L$ from spin-$3/2$ fermionic face representation theory (substrate-natural with chirality identification external); $SU(3)_c$ as a non-singlet internal label (substrate-compatible but not substrate-derived). The framework's target scope remains the gravitational and dark sector, and the gauge sector is hosted accordingly.

The sectors in this appendix are structurally linked to the same entanglement logic but are not part of the closed static weak-field chain.

\section*{Appendix J: The Lattice Microstructure Program}

Appendix~J records the technical content of Section~23: the host ensemble and its integrity gates, the coupled measure with its declared modeling forks, the control methodology, the exact cell-order calculations, and the measured tables of the compatibility and defect experiments, together with the specification of the capacity-transport instrument. Where Appendix~K recomputes the coefficient chain from its definitions, this appendix couples that chain to a dynamical substrate and reports what the coupled system does. The engine, run configurations, audit scripts, and raw measurement tables are maintained with the manuscript materials at \texttt{jaigp.org}.

\subsection*{J.1 Host ensemble and engine discipline}

The host is a standard causal-dynamical-triangulations ensemble \citep{AmbjornJurkiewiczLoll2004,AmbjornGJL2012}: triangulations of $S^1\times S^3$ with eighty spatial slices, sampled by local Pachner-class moves under the Regge action with bare couplings $k_0=2.2$ and $k_4$ tuned near criticality, and the $(4,1)$-simplex number pinned by a quadratic penalty $\varepsilon(N_{41}-\bar N)^2$ with $\varepsilon=0.01$. The action is tracked incrementally and audited against full recomputation throughout every run, with drifts held at the $10^{-11}$ level, and every quoted ensemble carries a terminal verification of the manifold conditions (gluing, links, Euler characteristic, connectivity).

The foliation is enforced exactly, and the enforcement has a history worth recording. Validating the cell identification of Section~23.1 exposed a defect in the move set as first written: it preserved the four-dimensional manifold conditions while permitting moves that fold a spatial slice, leaving a few percent of slice triangles in violation of the closed-$3$-manifold condition the identification requires. The audit layer caught the defect; the move set now rejects any proposal that would break the foliation, through a local test that is complete (an affected triangle's link necessarily passes through a newly created pentachoron) and preserves detailed balance in the same way as the other locality filters; and every ensemble quoted in this paper was generated with the corrected engine and passes a global foliation audit: every slice triangle in exactly two spatial tetrahedra, each pentachoron correctly threaded between its slices.

At these couplings and volumes the bare host has Hausdorff dimension rising from $2.7$ to $2.9$ to $3.1$ across $N_{41}=10{,}000$, $20{,}000$, $40{,}000$, with the slice-volume-profile score falling as the volume grows. The host's position relative to the de~Sitter phase established for this ensemble class \citep{AmbjornJurkiewiczLoll2004} is an open scan at this coupling point, and no reading below leans on the host being in that phase: every claim is comparative, coupled against bare and real against placebo, at matched volume.

\subsection*{J.2 The coupled measure and its fidelity anchors}

The coupled measure is the joint Gibbs distribution of Section~23.2, with the per-face label sum normalized so that each label carries base measure $1/7$. Under that base measure two properties hold exactly: at $\beta=0$ the label entropy cancels and the geometry marginal is identically the bare host, so the control arm is faithful by construction; and uniform label assignment at geometry moves is exactly detailed-balanced, so no proposal correction is needed for label births. The centering constant $\mu$ is the per-cell label free energy,
\[
\mu(\beta)=-\frac{1}{\beta N_{\rm cells}}\,\ln\,\mathbb{E}_{m\sim{\rm unif}}\!\left[e^{-\beta\sum_c E_c}\right]
=\frac{1}{\beta}\int_0^\beta\!\langle \bar E\rangle_s\,ds,
\]
computed by thermodynamic integration on the thermalized base configuration and logged with each arm, so the extensive part of the closure term is volume-neutral and cannot pose as a shift of $k_4$; matched total volume across arms is nevertheless retained as an independent gate rather than assumed from the centering.

Every run opens by rebuilding the single-cell tables and checking the Part~II anchors ($\Omega_{\rm tet}=1680$, $\eta_*=0.0298668$, $\langle K^2\rangle_{\eta_*}=50.223$, $g_{\mathrm{share,eff}}=7.4198$), and refuses to sample if any fails.

Three modeling forks are declared. Injectivity is imposed as a penalty $\lambda$ per colliding pair ($\lambda=3$ in the quoted runs) because a hard seven-color constraint at the maximal conflict degree does not guarantee an ergodic single-face heat bath; the hard constraint is the $\lambda\to\infty$ limit, the residual collision fraction is reported in every table so the softness stays visible, and the theory point is read at the ordered end. The closure sum is evaluated in the unoriented reading, in which the two cells sharing a face see the same label; the oriented alternative, in which they see opposite signs, requires oriented-slice bookkeeping through all moves and is recorded as an open fork. Parity contributes exactly $\ln 2$ per cell, extensive at pinned volume, and drops out of the coupled dynamics.

\subsection*{J.3 Controls and the reading discipline}

Four gates precede any reading: the manifold and foliation audits of J.1 on every arm; thermalization of the action and volume over the measured window; matched total simplex number across arms, since volume differences masquerade as dimension differences in the observables of interest; and a live chain, since a frozen chain is a failed coupling and never a result. The interpretation of each possible outcome, including which outcomes would license no claim at all, was written and committed before the corresponding results existed.

The load-bearing control is the placebo arm. The closure energy takes $210$ distinct values on the sorted-label orbits of the cell; the placebo shuffles those values across orbits, preserving the value pool and the permutation symmetry while destroying the closure structure, and runs the full pipeline at the top coupling. Only a real-versus-placebo difference at matched volume is ever attributed to the closure weighting; anything the placebo reproduces is machinery.

That discipline was learned on a neighboring theory, and the episode is part of the record. Before the admissibility coupling was built, an EPRL vertex amplitude \citep{EPRL2008} at frozen boundary spin $j=3$, computed with the \texttt{sl2cfoam-next} library \citep{Gozzini2021}, was coupled to the same host. A pre-registered audit of that coupling found five distinct misspecifications: a centering constant calibrated off-equilibrium and acting as an uncontrolled shift of the bare cosmological coupling; a label heat bath running at the wrong exponent; label births outside the acceptance ratio; an order-unity dependence on an arbitrary slot convention of the vertex tensor; and the positivization that any sign-alternating amplitude must undergo to be sampled at all. Each vanishes identically at $\beta=0$ and grows with $\beta$, so each masquerades as the signal the sweep was designed to detect while remaining invisible to the control arm. The decisive instrument was an entry-shuffled copy of the vertex tensor: it reproduced the apparent geometric steer at matched coupling, identifying the visible effect as sampler machinery. The sampler was rebuilt as an exact joint Gibbs measure (the fixes are inherited by the closure coupling, as the construction of J.2 shows), and the shuffled-placebo requirement became standing policy. The EPRL exercise is reported here as methodology: its physics scope is in any case bounded by the frozen-spin truncation and the positivization, and the object it samples is a neighboring theory's weighting, whose seven-dimensional intertwiner space is a different object from the seven face labels of this paper's ensemble (Section~27.8).

\subsection*{J.4 Exact results at cell order}

Six calculations, independent of any Monte Carlo run, fix what the substrate can and cannot do for the vacuum before the coupled measurements are read.

\paragraph{Induced curvature coupling.} The geometry marginal of the coupled measure weights a slice by the free energy of its labels, and the leading geometric dependence enters through the coordination $q$ of slice edges, the discrete curvature variable. A transfer-matrix evaluation of the label free energy on the ring of cells around an edge gives, at the theory point, an induced coupling $c_0=+0.019$ per slice edge, against the bare $k_0\simeq2.2$: about one percent. The label correlation length is about half a lattice step, and the beyond-quadratic dependence on $q$ is negligible at physical coordinations. The closure weighting alone therefore cannot supply the host's geometric stiffness, and its predicted effect on the coupled sweep is a steer too small to see at single-seed resolution, a prediction the sweep of J.5 confirms.

\paragraph{The refresh at cell order.} Modeling Postulate~III's refresh as memoryless whole-boundary redraws from the admissibility ensemble, arriving at unit rate per cell with shared faces taking the most recent writer's value, the single-cell stationary state is exact: the probability that $k$ of a cell's four faces retain their last self-drawn values is $k!\,(4-k)!/120$, the vacuum admissibility is $0.536$ (a $46\%$ standing density of transient closure failure), and $\langle K^2\rangle=48.74$ against the ensemble value $50.22$. The stationary state is geometry-blind to exponential accuracy: the keep-set law is uniform in the neighbors, overwrites arrive independently, and transmission around an edge ring decays as a power of the per-face keep probability. Unconditioned refresh therefore induces no vacuum curvature action either. What it does induce, for any geometry dynamics gated on closure, is a coupling of geometry to the local density of closure failure, the channel the defect experiment measures. The alternative refresh reading, redraws conditioned on the current neighborhood, is a Gibbs block sampler of the static weight and coarse-grains back to the transfer-matrix result above; the fork changes nothing downstream.

\paragraph{The history-weighting ceiling.} Two computations bound the strongest vacuum role a Many-Pasts tilt over closure-class observables could play. Amplifying the closure link interaction to $128$ times the theory point saturates the link-agreement ratio near $0.55$ with correlation length $1.7$ links and a curvature-sector signal near $0.1$: a structural ceiling, set by the ground-state degeneracy of the closure-plus-injectivity interaction, twenty times below the bare coupling. Independently, the exact Doob-transform solution of a closure-maintaining history tilt (paths weighted by time-integrated closure quality) on a ring of six seven-state faces drives the nearest-neighbor label correlation from $-0.12$ through zero to only $+0.13$ across the full tilt range, with longer-range correlations near $0.03$ throughout: no long-range order at any tilt strength. Together with the two paragraphs above this closes the label sector: static weight, refresh, and closure-class history weighting each fail, quantitatively, to supply vacuum geometric stiffness at cell order.

\paragraph{The free conserved field, integrated out.} A conserved Gaussian field back-coupled to the geometry induces, on integration, the geometric action $\tfrac{1}{2}\ln\det'L(g)$ per channel and slice: the spanning-tree entropy of the slice graph, by the matrix-tree theorem. Computed on degree-matched graph classes, the entropy per cell is $0.607$--$0.613$ on diamond-cubic lattices (the smooth extended proxy), $0.609$ on random $4$-regular graphs (the crumpled proxy), and $0.56$--$0.57$ on real engine slices: nearly universal at fixed coordination, with the genuinely nonlocal soft-mode part diluting as $\ln N/N$ per cell. The per-cell differential between extended and crumpled geometry is a few times $10^{-4}$; with all seven channels the phase-tipping force is of order $10^{-3}$ against $k_0\simeq2.2$.

\paragraph{The budget-bounded field.} The finite budget of Postulate~I makes the field non-Gaussian, with saturation curvature $m^2=2/g_{\mathrm{share,eff}}\simeq0.27$, and the massive determinant differential between the same proxies shrinks from $1.2\times10^{-3}$ per cell at zero mass to $4\times10^{-4}$ at the budget mass and toward zero beyond: the theory's own non-Gaussianity suppresses the modes where the sensitivity lived. What survives, measured on eight engine slices, is a local coupling renormalization of scale $\simeq0.08$ per cell in the seven-channel count, absorbed into the host's bare-coupling scan; it moves phase boundaries and supplies no phase of its own. The two containers no branch covers are named in Section~23.3: long-range field kernels and non-equilibrium geometry--field co-evolution, both unspecified by the theory.

\paragraph{The fracture construction.} The locality obstruction of Appendices~D.4 and~H.6 was verified by explicit construction of the move graph on the $1680$-state ensemble: under unit label shifts preserving injectivity the graph has exactly $48$ connected components (the $4!$ slot orderings times two orientations) of exactly $35$ states each, with the pair (ordering, parity) conserved. The claim that local dynamics cannot equilibrate the admissibility ensemble, which makes the non-local refresh of Postulate~III load-bearing in the scale chain, holds as stated.

\subsection*{J.5 Compatibility at matched volume}

The coupled sweep at $N_{41}=20{,}000$ ran four arms from a common thermalized bare checkpoint: the $\beta=0$ control, an intermediate ramp point at $\beta=0.3$, the theory point $\beta=1$, and the shuffled placebo at $\beta=1$. All arms pass the manifold, foliation, thermalization, and matched-volume gates; total simplex numbers agree within $1.6\%$.

\begin{center}
{\small
\begin{tabular}{lcccc}
\hline
arm & $N_4$ & $\langle E\rangle$ per cell & collision fraction & $d_H$ \\
\hline
$\beta=0$ (control) & $47{,}034$ & $3.99$ & $0.648$ & $2.89$ \\
$\beta=0.3$ & $46{,}966$ & $2.78$ & $0.403$ & $2.87$ \\
$\beta=1$ (theory point) & $46{,}286$ & $1.74$ & $0.082$ & $2.85$ \\
$\beta=1$ placebo & $46{,}406$ & $1.95$ & $0.730$ & $2.98$ \\
\hline
\end{tabular}
}
\end{center}

The label sector behaves as designed: the control arm's collision fraction $0.648$ sits on the uniform-measure prediction $1-840/2401=0.650$, the theory point drives it to $0.082$, and the placebo, carrying the same energy values with no closure structure, stays at $0.730$. The geometric observables (Hausdorff dimension above; the slice-volume profile behaves the same way) are statistically indistinguishable across arms at single-seed resolution, as the one-percent induced coupling of J.4 predicts. The run is single-seed; the reading is the comparative one stated in Section~23.5, and the residual few-percent spreads are within the drift of the base ensemble.

\subsection*{J.6 The Regge bridge: the substrate writes the volume term}

The centering constant $\mu$ of J.2 doubles as a prediction. Left uncentered, the closure term contributes its extensive part $-\beta\mu N_{\rm cells}$ to the action; against the quadratic volume pin $\varepsilon(N_{41}-\bar N)^2$, with $N_{\rm cells}=N_{41}/2$, that linear tilt displaces the equilibrium volume by
\[
\Delta N_{41}=-\frac{\beta\,\mu}{4\,\varepsilon},
\]
and $\mu$ is computed analytically by transfer-matrix thermodynamic integration from $\eta_*$ and the injectivity penalty, with no lattice input. At the demonstration point ($\beta=1$, $\mu=2.488$, $\varepsilon=0.01$) the prediction is $\Delta N_{41}=-62.2$. Three arms from a common clean base at $N_{41}=2000$ measure it:

\begin{center}
{\small
\begin{tabular}{lcc}
\hline
arm & mean $N_{41}$ & shortfall \\
\hline
bare ($\beta=0$) & $1993.2$ & $-6.8$ \\
centered ($\beta=1$) & $1969.8$ & $-30.2$ \\
uncentered ($\beta=1$) & $1907.0$ & $-93.0$ \\
\hline
\end{tabular}
}
\end{center}

The uncentered-minus-centered difference is $-62.8$ against the predicted $-62.2$. The subtraction isolates the $\mu$-term, since the two $\beta=1$ arms share every other piece of machinery; the centered arm's own $-30$ is the second-order non-neutrality of the thermodynamic-integration centering, real, cancelling in the difference, and gated by matched volume in production. The reading is the bridge of Section~23.4: a number computed from the ultraviolet specification appears, at the predicted magnitude, as a coefficient of the host action's volume sector in a dynamical measurement. The measurement is demonstration-volume and single-seed; the curvature-sector counterpart, the $c_0$ and field shifts of J.4 read as phase-boundary displacements, is pre-registered for production statistics.

\subsection*{J.7 The pinned-defect experiment}

Both arms resume from the thermalized theory-point checkpoint of J.5 and run $3000$ further sweeps with $100$ pins each. The real arm pins cells at maximal closure failure (all four faces $m=0$: six collisions, $K^2=48$); the placebo arm pins cells at the best-closed injective configuration ($m=\{0,1,2,3\}$, $K^2=40.67$). Pinned faces are excluded from the label heat bath, and the carrier cells are protected from removal by the geometry moves identically in both arms, so the real-minus-placebo difference isolates the failure content; pin survival is a hard gate, passed at $100/100$ in both arms. Freshly sampled unpinned cells far from every pin are measured inside both arms as the vacuum reference. Shell observables are accumulated by graph distance in the slice adjacency, excluding the pinned cells themselves (they are the source, the shells are the medium), for roughly $10^5$ defect-environment samples per arm.

\begin{center}
{\small
\begin{tabular}{llcccc}
\hline
arm & shell & mean $E$ & collision fraction & mean coordination & cells per shell \\
\hline
failure   & 1 & $1.6782\pm0.0010$ & $0.0690\pm0.0003$ & $5.2997\pm0.0007$ & $4$ \\
          & 2 & $1.7325\pm0.0008$ & $0.0790\pm0.0003$ & $5.5789\pm0.0009$ & $7.237\pm0.003$ \\
          & 3 & $1.7274\pm0.0006$ & $0.0787\pm0.0002$ & $5.6459\pm0.0007$ & $11.296\pm0.006$ \\
\hline
closed    & 1 & $1.7291\pm0.0011$ & $0.0794\pm0.0004$ & $5.2513\pm0.0008$ & $4$ \\
(placebo) & 2 & $1.7304\pm0.0008$ & $0.0792\pm0.0003$ & $5.5773\pm0.0009$ & $7.116\pm0.005$ \\
          & 3 & $1.7280\pm0.0007$ & $0.0788\pm0.0002$ & $5.6460\pm0.0008$ & $10.905\pm0.007$ \\
\hline
vacuum    & 1 & $1.7306\pm0.0011$ & $0.0797\pm0.0004$ & $5.3165\pm0.0018$ & $4$ \\
          & 2 & $1.7303\pm0.0008$ & $0.0795\pm0.0003$ & $5.6108\pm0.0014$ & $7.235\pm0.006$ \\
          & 3 & $1.7288\pm0.0007$ & $0.0789\pm0.0002$ & $5.7310\pm0.0012$ & $10.909\pm0.012$ \\
\hline
\end{tabular}
}
\end{center}

The reading follows the ladder fixed before the run. The capacity state responds at the first shell and only there: mean closure energy and collision fraction around the failure pins are displaced well outside the quoted errors from both the placebo and the vacuum, and both return to vacuum values by the second shell, the sub-lattice range the correlation length of J.4 requires. The local geometry responds through the third shell: at identical anchoring the failure pins carry higher first-shell coordination than the closed pins and systematically more cells per shell, a volume excess of $3.6\%$ at the third shell where the placebo matches the vacuum. Both responses separate real from placebo, so both are attributable to the failure content. The quoted errors are statistical accumulation errors, not yet thinned for autocorrelation, and the result is single-seed; those caveats set the replication burden for the numbers, and the pre-registered reading grades the mechanism as demonstrated at that level.

\subsection*{J.8 The conservation law and the transport instrument}

The dichotomy of Section~23.7 in operator form. If free capacity is a per-cell budget re-equilibrated by maximum entropy --- link shares $c_l\ge0$ with $\sum_{l\ni x}c_l=C-m_x$ --- stationarity gives one multiplier per cell, and the linearized constraint reads $(2z\,\mathbb{1}-L)\,\delta\mu=-m/\chi$ with $z=4$ and $L$ the slice Laplacian. The operator has no small-momentum pole; on a measured $169$-cell slice its Green function falls eight orders of magnitude within thirteen steps. If instead free capacity $f_x$ obeys a continuity law,
\[
f_x(t+dt)-f_x(t)=-\sum_{\text{faces }xy}J_{xy}-{\rm commit}_x+{\rm release}_x,
\qquad J_{xy}=-J_{yx},
\qquad \sum_x\,[\,f_x+m_x\,]={\rm const},
\]
with a defect drawing steady maintenance flux $Q_x\propto m_x$, then any local flux law linearizes to $J\propto-D\nabla f$ and the steady state solves $L\,\delta f=(Q/D)\,(\delta_{\rm source}-{\rm uniform\ return})$: the massless graph-Poisson equation, whose Green function on a three-dimensional slice falls as $1/r$. The deficit field is $\delta f\propto m/r$, and the relative deficit is the weak-field bridge variable of Section~11.

The instrument implements the conservation reading with the sink left to emerge. A field on the slice cells is initialized at the vacuum anchor $g_{\mathrm{share,eff}}=7.4198$ and transported by antisymmetric face fluxes; global conservation of $\sum_x(f_x+m_x)$ is asserted at machine precision every sweep, and the run aborts on any violation. Absorption is proportional to a cell's instantaneous label collisions, so a pinned defect drains capacity only through the closure failure it induces in itself and its neighborhood by the ordinary label dynamics: the maintenance flux each defect draws is measured, not set. Pins are placed at four commitment levels ($1$, $2$, $3$, and $6$ colliding pairs), giving the two headline readings: the emergent flux against commitment, which tests maintenance linearity dynamically, and the first-shell deficit per unit flux, whose constancy across the ladder is the junction statement that fixes the lattice analog of $G$. The two profile gates are the deficit against the graph-Coulomb Green function and a fitted screening mass consistent with zero, any conservation leakage appearing as Yukawa screening. Absorption reads the failure a cell carries above the vacuum's own standing level (the $46\%$ transient density of J.4), so the vacuum is absorption-free and the transported field massless by construction. The first full-volume run demonstrated why this matters: with absorption reading total failure instead, the vacuum itself absorbs, and the measured deficit range of four to five steps reproduced the predicted self-screening length $\xi=\sqrt{D/(\kappa\,\langle n_{\rm coll}\rangle_{\rm vac})}\simeq5$--$7$ --- the model's own Yukawa mass, a post-diction confirming that the transport sector behaves as derived. The ladder also carries a level-zero rung, a frozen best-closed pin, which measures the commitment-independent boundary-dressing flux directly; the linearity reading is the dressing-subtracted fit. The first implementation freezes the geometry while the labels stay fully dynamical, isolating the field sector; conservative transport through the geometry moves is specified as the extension. On small validation volumes the instrument behaves physically: deficit wells ordered by commitment, monotone recovery with distance, emergent flux rising along the ladder, conservation exact --- and the highest-commitment pin draws less flux than the linear extrapolation because it depletes its own neighborhood, a self-screening the model produces without being told. The production measurement at $N_{41}=20{,}000$, where shells to graph distance eight and beyond exist for the profile gates, is in progress. Its readings are pre-registered: linearity failure strikes the maintenance postulate, a nonzero screening mass strikes the conservation reading, and either would be reported as such.

\section*{Appendix K: Numerical Checks and Robustness}

Appendix~K is intentionally modest. It does not add new derivations, and these scripts and checks do not establish the ontology: they audit the numerical consequences of the stated derivations, and their role is reproducibility, not independent proof. It collects the main numerical cross-checks that make it easier to see that the same coefficient chain survives repeated contact with independent benchmark calculations, and it records, in Section~K.3, a complete self-contained script that recomputes the entire numerical spine of the manuscript from the stated definitions. The dynamical-lattice program is a separate layer with a different job: where this appendix recomputes the chain, Section~23 and Appendix~J couple it to a dynamical substrate and measure what it does.

The numbers recomputed here, and where each enters the main text, are: the admissibility entropy $g_{\mathrm{share,eff}}$ (Section~13, Appendices~B and~D.4); the substrate length $L_*$ (Sections~2 and~13, Appendix~D.4); the induced scale $G_*$ (Sections~10 and~13, Appendices~D.4 and~K); the galactic scale $a_0$ (Section~14, Appendix~C.7); the diamond-lattice Green constant $G_{\mathrm{tet}}(0)$ (Appendix~C.5); the edge-kernel chain $J_{\mathrm{bare}}\to\gamma$ and the return sum $\Sigma_{\mathrm{ret}}$ (Appendices~C.3--C.4); and the charged-lepton ladder with its support-exponent cross-check (Section~13, Appendix~I.1).

\subsection*{K.1 Cross-sector numerical checks}

The cross-check program includes:
\begin{itemize}
\item the one-bit fermionic defect check $\Delta S_f=\ln 2$;
\item the rooted-shell convergence check $\sigma_{\mathrm{ind}}^{(2)}\simeq \sigma_{\mathrm{ind}}^{(3)}$;
\item the UV closed-branch moments $\langle K^2\rangle_{\eta_*}$, $\mathrm{Var}_{\eta_*}(K^2)$, and $a_{\mathrm{UV}}$;
\item cross-sector consistency among the electron anchor, the substrate length $L_*$, the induced scale $G_*$, the Green-matched Newton closure, and the galactic scale $a_0$.
\end{itemize}
These checks do not replace the derivations, but they show that the same coefficient chain survives independent numerical scrutiny across the sectors where closure is claimed.

\subsection*{K.2 Reproducibility ledger for the substrate scale}

The substrate-length calculation is short enough to record as a numerical ledger. Enumerating the $1680$ oriented injective tetrahedral states with labels $m=-3,\ldots,3$, weighting them by $e^{-\eta K^2}$, and solving
\[
\langle K^2\rangle_\eta=\frac{3}{2\eta}
\]
gives
\[
\eta_*=0.02986684439352237,
\qquad
g_{\mathrm{share,eff}}=7.419800023570903.
\]
The one-pass seven-sector history support and its transverse export are then
\[
e^{7g_{\mathrm{share,eff}}}
=
3.602860521062804\times10^{22},
\]
\[
\frac23 e^{7g_{\mathrm{share,eff}}}
=
2.401907014041869\times10^{22}.
\]
Using $\lambda_e=\hbar/(m_ec)=3.861592671986303\times10^{-13}\ {\rm m}$ gives
\[
L_*=\frac{\lambda_e}{(2/3)e^{7g_{\mathrm{share,eff}}}}
=
1.607719470158885\times10^{-35}\ {\rm m},
\]
and hence
\[
G_*=\frac{c^3L_*^2}{\hbar}
=
6.603991270884347\times10^{-11}\ {\rm m^3\,kg^{-1}\,s^{-2}}.
\]

The source-side Green constant is obtained independently from the diamond-lattice integral in Appendix~C.5. Uniform-grid quadrature with endpoint extrapolation gives
\[
G_{\mathrm{tet}}(0)=0.448220394388\ldots ,
\]
confirming the exact Joyce value
\[
G_{\mathrm{tet}}(0)
=
\frac{3\Gamma(1/3)^6}{2^{14/3}\pi^4}.
\]
This is the value used in the source theorem. These numbers are not additional inputs; they are the numerical evaluation of the finite spectrum, the seven-sector history support, and the graph Green function already defined in the derivation.

\subsection*{K.3 Reproduction script}

The script below recomputes every load-bearing number in the manuscript from the stated definitions alone: the $1680$-state $K^2$ spectrum and its multiplicities, the closure solution $\eta_*$, the admissibility entropy $g_{\mathrm{share,eff}}$, the substrate length $L_*$ and induced scale $G_*$, the Joyce constant and the exact $n_{\rm hor}S_\infty^{\rm cell}=1/4$ identity, the edge-kernel chain, the projection coefficient $\epsilon$ and horizon target $\sigma_*$, the charged-lepton ladder with its support-exponent cross-check, and the slot mutual information. No value in the chain is entered by hand; each printed line carries the manuscript value it must reproduce. The script runs in seconds under Python~3 with \texttt{numpy} and \texttt{scipy}, and is maintained alongside the manuscript at \texttt{jaigp.org}.

{\footnotesize\begin{verbatim}
# ============================================================
# Reproduction script for the Entropic Scalar EFT manuscript.
# Recomputes the full numerical spine from stated definitions:
# the 1680-state K^2 spectrum, eta*, g_share,eff, L*, G*, the
# Joyce constant, the 1/4 horizon identity, the edge-kernel
# chain, epsilon, sigma*, the lepton ladder, and the slot MI.
# Dependencies: Python 3, numpy, scipy.
# Every printed value is compared against the manuscript.
# ============================================================
import itertools, math
from math import log, exp, pi, gamma, sqrt
from collections import Counter
import numpy as np
from scipy.optimize import brentq

# 1. Enumerate 1680 states: injective 4-of-7 labels m in -3..3, x2 parity
labels = range(-3,4)
states = list(itertools.permutations(labels,4))
print("P(7,4) =", len(states), " total with parity:", 2*len(states))

def K2(ms):
    S = sum(ms); Sig2 = sum(m*m for m in ms)
    return 48 - (S*S - Sig2)/3

spec = Counter()
for st in states:
    spec[round(K2(st)*3)] += 2   # parity doubling; key = 3*K2 to keep exact
print("Distinct K2 values and multiplicities (K2 as fraction /3):")
for k in sorted(spec): print(f"  K2={k}/3 = {k/3:.4f}  mult={spec[k]}")
print("Total multiplicity:", sum(spec.values()))

# 2. Solve closure condition <K2> = 3/(2 eta)
Ks = np.array([k/3 for k in sorted(spec)])
ns = np.array([spec[k] for k in sorted(spec)], dtype=float)
def avgK2(eta):
    w = ns*np.exp(-eta*Ks)
    return (w*Ks).sum()/w.sum()
f = lambda eta: avgK2(eta) - 3/(2*eta)
eta_star = brentq(f, 1e-4, 1.0, xtol=1e-16)
print("\neta* =", eta_star, " (paper: 0.02986684439352237)")

w = ns*np.exp(-eta_star*Ks); Z = w.sum(); p_class = w/Z
p_micro = np.exp(-eta_star*Ks)/Z
H = -(ns*p_micro*np.log(p_micro)).sum()
print("g_share,eff =", H, " (paper: 7.419800023570903)")
avg = avgK2(eta_star)
var = (p_class*(Ks-avg)**2).sum()
print("<K2> =", avg, " (paper 50.2229154254)   3/(2eta*) =", 3/(2*eta_star))
print("Var(K2) =", var, " (paper 15.6889750078)   a_UV = 1/Var =", 1/var)
print("C_cl = eta* <K2> =", eta_star*avg)

# 3. L*, G*
hbar=1.054571817e-34; c=299792458.0; me=9.1093837015e-31
lam_e = hbar/(me*c)
print("\nlambda_e =", lam_e, " (paper 3.861592671986303e-13)")
Lstar = 1.5*lam_e*exp(-7*H)
print("e^{7g} =", exp(7*H), " (paper 3.602860521062804e22)")
print("(2/3)e^{7g} =", (2/3)*exp(7*H), " (paper 2.401907014041869e22)")
print("L* =", Lstar, " (paper 1.607719470158885e-35)")
Gstar = c**3*Lstar**2/hbar
print("G* =", Gstar, " (paper 6.603991270884347e-11)")
G_codata = 6.67430e-11
LP = sqrt(hbar*G_codata/c**3)
print("L* vs Planck length:", (LP-Lstar)/LP*100, "% below  (paper 0.528%)")
print("G* vs CODATA:", (G_codata-Gstar)/G_codata*100, "% below  (paper 1.053%)")

# 4. Joyce constant
Gtet = 3*gamma(1/3)**6/(2**(14/3)*pi**4)
print("\nG_tet(0) =", Gtet, " (paper 0.4482203943883814)")
print("ln(7/6)/(pi*Gtet) =", log(7/6)/(pi*Gtet), " (paper 0.109472228)")
print("S_inf_cell = 3ln2/(32 pi Gtet) =", 3*log(2)/(32*pi*Gtet), " (paper 0.0461482516)")
print("n_hor*S_inf_cell =", (8*pi*Gtet/(3*log(2))) * (3*log(2)/(32*pi*Gtet)))

# 5. Edge kernel chain
Jbare = 2/3*eta_star
Jtree = Jbare/3
Sig = 7+2/9
cloop = 1/(1+Jtree*Sig)
Jren = Jtree*cloop
print("\nJ_bare =", Jbare, " (paper 0.0199112296)")
print("J_tree =", Jtree, " (paper 0.0066370765)")
print("c_loop =", cloop, " (paper 0.95426)")
print("J_ren =", Jren, " (paper 0.00633348)")
print("gamma coeff 4 J_ren/(3 pi^2) =", 4*Jren/(3*pi**2), " (paper 8.556e-4)")

# 6. epsilon and cluster numbers
epsv = H/(4*pi**2)
print("\nepsilon = g/(4pi^2) =", epsv, " 1-eps =", 1-epsv, " ceiling =", 2-epsv)
print("sigma* = pi/g =", pi/H, " (paper 0.42340665)")

# 7. lepton ladder
mmu = 720*2/7; mtau = 720**2*(2/7)**4
print("\nm_mu/m_e =", mmu, " PDG 206.7683 err%:", (206.7683-mmu)/206.7683*100)
print("m_tau/m_e =", mtau, " PDG 3477.23 err%:", (3477.23-mtau)/3477.23*100)
print("alpha_mu =", log(206.7683)/log(mmu), " alpha_tau =", log(3477.23)/log(mtau))

# 8. mixing-time ratio
tau_star = Lstar/c; tau_e = hbar/(me*c**2)
print("\ntau* =", tau_star, " tau_e =", tau_e, " ratio =", tau_e/tau_star)

# 9. slot mutual information (marginal of slot under admissibility)
joint = Counter()
for st in states:
    wgt = exp(-eta_star*K2(st))
    joint[(st[0],st[1])] += 2*wgt
tot = sum(joint.values())
pj = {k:v/tot for k,v in joint.items()}
p0 = Counter(); p1 = Counter()
for (a,b),v in pj.items(): p0[a]+=v; p1[b]+=v
I = sum(v*log(v/(p0[a]*p1[b])) for (a,b),v in pj.items())
H0 = -sum(v*log(v) for v in p0.values())
print("\nslot MI =", I, " (paper 0.1545)   H(slot) =", H0, " I/H =", I/H0, " (paper 0.079)")
\end{verbatim}}

\section*{Appendix L: Fork accounting for the scale chain}

Because the substrate length induces a value of $G_*$ near the observed Newton constant, this appendix records the construction choices behind that result and quantifies the discrete-choice fork space around it, so the scale chain can be weighed against the freedom available in the construction. It is referenced wherever $G_*$ is discussed in the main text.

\paragraph{Provenance.} By the author's recollection, a target entropy near $7.42$ was identified early in the framework's development by inverting the measured Newton constant, with the ensemble constructed to deliver it jointly with other structural constraints; the documentary record does not reach that genesis layer. What the record establishes is this: the realized ensemble was fixed, and motivated in print by information-independence arguments, roughly ten weeks before the induced-$G_*$ chain existed; the export factor's structural components predate the chain, its deployment is contemporaneous with the first landing, and its candidate menu was published with residuals disclosed; and both recorded routes to $G$ land at one-percent misses rather than exact values. The Newton constant was the sole observational target of that calibration; the galactic acceleration scale, the capacity ceiling, the charged-lepton structure, and the committed-phase abundance are downstream outputs of the calibrated number, compared against measurement only after fixation. The induced $G_*$ is accordingly treated as a calibration, not a prediction --- with the calibration understood as a discrete selection rather than a continuous adjustment: within the selected construction the entropy is derived with no adjustable parameter, is reproducible from the stated rules alone, and was never re-adjusted against any subsequent observable, which is why the downstream misses stand in the record as written. The audit below quantifies how much construction freedom the discrete selection could have consumed. The evidential weight of the scale chain rests on results obtained after the ensemble was frozen --- the cluster capacity bound and its twenty-four-point test (Section~17.5), the committed-phase abundance against the measured dark-to-baryonic ratio (Section~18.5), the redshift direction of the acceleration scale, and the sub-$g_c$ rotation-curve regime --- none of which entered the construction.

\paragraph{The audit.} The substrate length depends exponentially on the realized ensemble, $L_*\propto e^{-7g_{\mathrm{share,eff}}}$, so the calibration is meaningful only relative to the space of constructions that could have been written instead. This appendix quantifies that space. The grammar varies five discrete choices around the realized chain: the label count $M\in\{4,\dots,9\}$ with $j=(M-1)/2$; the multiplicity rule (ordered or unordered injective selections, with and without parity doubling, and the non-injective $M^4$); the closure coefficient $d\in\{1,2,3,4,6\}$ in $\langle K^2\rangle=d/(2\eta)$; the support exponent $n\in\{4,\,M,\,2M,\,M(M-1)/2\}$; and the export prefactor in $\{\tfrac12,\tfrac23,1,\tfrac32,2\}$. This yields $3300$ candidate constructions, each evaluated end to end: ensemble, closure, entropy, induced scale.

The results are as follows. The induced $G_*$ scatters across the grammar with a standard deviation of $137$ natural-log units. Within ten percent of the measured value lie four constructions, and all four are the realized one, differing only in the closure coefficient, which shifts $g_{\mathrm{share,eff}}$ by less than $0.1\%$ because the closure weighting is a small correction to $\ln\Omega_{\mathrm{tet}}$; the closure coefficient therefore carries almost no load. The nearest structurally distinct construction misses by eleven percent, and within a factor of two of the measured value lie only two construction families in the entire grammar. One further menu member belongs in the ten-percent set: replacing the export prefactor with the quantum-speed-limit coefficient $\pi/2$, the value a saturated orthogonality bound would select, lands at $+8.5\%$; the two-percent landing remains unique to the transverse-export reading. At the same time, the menu density implies that roughly one family is expected to land near any target by chance alone, so the numerical landing by itself carries little evidential weight. Any residual weight would reside in the a priori force of the discrete arguments that select the realized construction --- the seven-state face sector (Section~5), the parity doubling of the orientation classes, the ordered injective assignment over distinguishable faces, the seven-fold support exponent (Appendix~D.4), and the transverse export factor (Appendix~C.5). Given the provenance above, these arguments are historically retrodictive, and the framework's evidential case is not rested on them here; it is rested on the post-fixation tests listed in the provenance statement.

The same audit redistributes the evidential weight of the galactic scale. Because $a_0=(g_{\mathrm{share,eff}}/4\pi^2)\,cH_0$ depends on the ensemble only logarithmically, sixty-three percent of the entire grammar lands $a_0$ within twenty-five percent of the observed value: the galactic acceleration scale tests the coupling form, not the realized ensemble. The induced $G_*$, with its exponential sensitivity, is the sharp test of the ensemble itself, with the charged-lepton ladder its weaker corroborator. The audit script is maintained alongside the reproduction script of Appendix~K.3 at \texttt{jaigp.org}.

\section*{Appendix M: Why the saturated CMB carrier must be committed capacity, not a lagged response}

The microwave background requires a gravitating component that carries the growing mode of the baryon perturbations while rejecting their acoustic oscillation. This appendix records the quantitative exclusion of every relaxational realization of that component and the no-go that selects the constraint reading of Section~18.5.

\paragraph{Growth-limited kernels.} For development dynamics $dW/dt=(W_{\max}-W)/\tau$ with $\tau=\beta\,t_{\rm ff}(\rho_{\rm local})$, the measured cluster radial decline of the source weight fixes $\beta\simeq15.6$, while saturation of the cosmological bath by recombination requires $\beta\lesssim0.56$: the window is empty by a factor of order thirty. The radial ratio between $R_{500}$ and $R_{200}$, a $\beta$-independent prediction of the free-fall clock, agrees with the data at the several-percent level, so the exclusion is of the absolute clock, not of the radial structure.

\paragraph{Lagged response.} A first-order lag transmits a fraction $T=[1+(\omega\tau)^2]^{-1/2}$ of an oscillation at frequency $\omega$. The transmitted component acts as additional effective baryon loading $(1/\epsilon-1)\,T$ in the acoustic driving; Einstein--Boltzmann computation shows a $0.3\%$ temperature-spectrum tolerance bounds $T\le2\times10^{-3}$, requiring $\beta\gtrsim37$ against the development requirement $\beta\lesssim0.56$: empty by a factor of order sixty-six. More generally, by recombination a third-peak mode has completed only a few oscillations, so no causal filter of any order achieves the required rejection.

\paragraph{Sound-speed classification.} For $L=f(X)$, perturbations carry $c_s^2=f_X/(f_X+2Xf_{XX})$. The released branches give $c_s^2=\tfrac12$ (deep) and $1$ (Newtonian): free-streaming. Plateau approaches $f=f_{\max}-A/X^n$ give $c_s^2=-1/(2n{+}1)$: gradient-unstable. Both natural cap readings --- stationarity of the canonical momentum and of the energy density --- impose $f_X+2Xf_{XX}=0$, the pole of $c_s^2$: a rigid medium carrying no perturbation. The only locus with $c_s^2=0$ and healthy density response is the constraint surface of the mimetic class, $X$ pinned with a multiplier, which is the reading of Section~18.5.

\paragraph{No-go for relaxation-plus-cap carriers.} Within dynamics consisting of relaxation toward a demand together with a hard cap, with the cell weight as the only state variable, no coupling of the weight to the demand carries a secular component of the modulation while rejecting its oscillation: a pinned weight retains no memory of the modulation, and any instantaneous coupling of bounded periodic inputs is itself bounded and periodic. A secular component requires an additional conserved integrating variable. The conserved spatial density of committed cells is that variable, and it is native to the saturated phase rather than added to it.

\begin{thebibliography}{99}

\bibitem[McGaugh et al.(2016)]{McGaugh2016}
S.~S. McGaugh, F. Lelli, and J.~M. Schombert,
``Radial acceleration relation in rotationally supported galaxies,''
\emph{Physical Review Letters} \textbf{117}, 201101 (2016).

\bibitem[Clowe et al.(2006)]{Clowe2006}
D. Clowe et al.,
``A direct empirical proof of the existence of dark matter,''
\emph{The Astrophysical Journal Letters} \textbf{648}, L109--L113 (2006).

\bibitem[Bertotti et al.(2003)]{Bertotti2003}
B. Bertotti, L. Iess, and P. Tortora,
``A test of general relativity using radio links with the Cassini spacecraft,''
\emph{Nature} \textbf{425}, 374--376 (2003).

\bibitem[Tian et al.(2020)]{Tian2020}
Y. Tian et al.,
``The radial acceleration relation in CLASH galaxy clusters,''
arXiv:2001.08340 (2020).

\bibitem[Guttmann(2010)]{Guttmann2010}
A.~J. Guttmann,
``Lattice Green functions in all dimensions,''
\emph{Journal of Physics A: Mathematical and Theoretical} \textbf{43}, 305205 (2010).

\bibitem[Joyce(1994)]{Joyce1994}
G.~S. Joyce,
``On the cubic lattice Green functions,''
\emph{Proceedings of the Royal Society of London A} \textbf{445}, 463--477 (1994).

\bibitem[Arnowitt et al.(1962)]{ADM1962}
R. Arnowitt, S. Deser, and C.~W. Misner,
``The dynamics of general relativity,''
in \emph{Gravitation: An Introduction to Current Research},
ed. L. Witten (Wiley, New York, 1962).

\bibitem[Bekenstein(1973)]{Bekenstein1973}
J.~D. Bekenstein,
``Black holes and entropy,''
\emph{Physical Review D} \textbf{7}, 2333--2346 (1973).

\bibitem[Hawking(1975)]{Hawking1975}
S.~W. Hawking,
``Particle creation by black holes,''
\emph{Communications in Mathematical Physics} \textbf{43}, 199--220 (1975).

\bibitem[Gibbons and Hawking(1977)]{GibbonsHawking1977}
G.~W. Gibbons and S.~W. Hawking,
``Action integrals and partition functions in quantum gravity,''
\emph{Physical Review D} \textbf{15}, 2752--2756 (1977).

\bibitem[Regge and Wheeler(1957)]{ReggeWheeler1957}
T. Regge and J.~A. Wheeler,
``Stability of a Schwarzschild singularity,''
\emph{Physical Review} \textbf{108}, 1063--1069 (1957).

\bibitem[Zerilli(1970)]{Zerilli1970}
F.~J. Zerilli,
``Effective potential for even-parity Regge--Wheeler gravitational perturbation equations,''
\emph{Physical Review Letters} \textbf{24}, 737--738 (1970).

\bibitem[Sanderson et al.(2003)]{Sanderson2003}
A.~J.~R. Sanderson, T.~J. Ponman, A. Finoguenov, E.~J. Lloyd-Davies, and M. Markevitch,
``The Birmingham--CfA cluster scaling project --- I. Gas fraction and the $M$--$T_X$ relation,''
\emph{Monthly Notices of the Royal Astronomical Society} \textbf{340}, 989--1010 (2003).

\bibitem[Vikhlinin et al.(2006)]{Vikhlinin2006}
A. Vikhlinin, A. Kravtsov, W. Forman, C. Jones, M. Markevitch, S.~S. Murray, and L. Van Speybroeck,
``Chandra sample of nearby relaxed galaxy clusters: mass, gas fraction, and mass--temperature relation,''
\emph{The Astrophysical Journal} \textbf{640}, 691--709 (2006).

\bibitem[Eckert et al.(2022)]{Eckert2022}
D. Eckert, S. Ettori, E. Pointecouteau, R.~F.~J. van der Burg, and S.~I. Loubser,
``The gravitational field of X-COP galaxy clusters,''
\emph{Astronomy \& Astrophysics} \textbf{662}, A123 (2022).

\bibitem[Li et al.(2023)]{Li2023}
P. Li, Y. Tian, M.~P. J\'ulio, M.~S. Pawlowski, F. Lelli, S.~S. McGaugh, J.~M. Schombert, J.~I. Read, P.-C. Yu, and C.-M. Ko,
``Measuring galaxy cluster mass profiles into the low-acceleration regime with galaxy kinematics,''
\emph{Astronomy \& Astrophysics} \textbf{677}, A24 (2023).

\bibitem[Milgrom(1983)]{Milgrom1983}
M. Milgrom,
``A modification of the Newtonian dynamics as a possible alternative to the hidden mass hypothesis,''
\emph{Astrophysical Journal} \textbf{270}, 365 (1983).

\bibitem[Bekenstein(2004)]{Bekenstein2004}
J.~D. Bekenstein,
``Relativistic gravitation theory for the modified Newtonian dynamics paradigm,''
\emph{Physical Review D} \textbf{70}, 083509 (2004).

\bibitem[Famaey and McGaugh(2012)]{FamaeyMcGaugh2012}
B. Famaey and S.~S. McGaugh,
``Modified Newtonian dynamics (MOND): observational phenomenology and relativistic extensions,''
\emph{Living Reviews in Relativity} \textbf{15}, 10 (2012).

\bibitem[Jacobson(1995)]{Jacobson1995}
T. Jacobson,
``Thermodynamics of spacetime: the Einstein equation of state,''
\emph{Physical Review Letters} \textbf{75}, 1260 (1995).

\bibitem[Verlinde(2011)]{Verlinde2011}
E. Verlinde,
``On the origin of gravity and the laws of Newton,''
\emph{Journal of High Energy Physics} \textbf{2011}(4), 29 (2011).

\bibitem[Verlinde(2017)]{Verlinde2017}
E. Verlinde,
``Emergent gravity and the dark universe,''
\emph{SciPost Physics} \textbf{2}, 016 (2017).

\bibitem[Skordis and Z\l{}o\'snik(2021)]{Skordis2021}
C. Skordis and T. Z\l{}o\'snik,
``New relativistic theory for modified Newtonian dynamics,''
\emph{Physical Review Letters} \textbf{127}, 161302 (2021).

\bibitem[Collins et al.(2004)]{Collins2004}
J. Collins, A. Perez, D. Sudarsky, L. Urrutia, and H. Vucetich,
``Lorentz invariance and quantum gravity: an additional fine-tuning problem?,''
\emph{Physical Review Letters} \textbf{93}, 191301 (2004).

\bibitem[Bombelli et al.(1987)]{Bombelli1987}
L. Bombelli, J. Lee, D. Meyer, and R.~D. Sorkin,
``Space-time as a causal set,''
\emph{Physical Review Letters} \textbf{59}, 521 (1987).

\bibitem[Dowker et al.(2004)]{Dowker2004}
F. Dowker, J. Henson, and R.~D. Sorkin,
``Quantum gravity phenomenology, Lorentz invariance and discreteness,''
\emph{Modern Physics Letters A} \textbf{19}, 1829 (2004).

\bibitem[Chae(2023)]{Chae2023}
K.-H. Chae,
``Breakdown of the Newton--Einstein standard gravity at low acceleration in internal dynamics of wide binary stars,''
\emph{Astrophysical Journal} \textbf{952}, 128 (2023).

\bibitem[Banik et al.(2024)]{Banik2024}
I. Banik, C. Pittordis, W. Sutherland, B. Famaey, R. Ibata, S. Mieske, and H. Zhao,
``Strong constraints on the gravitational law from Gaia DR3 wide binaries,''
\emph{Monthly Notices of the Royal Astronomical Society} \textbf{527}, 4573 (2024).

\bibitem[Genzel et al.(2017)]{Genzel2017}
R. Genzel et al.,
``Strongly baryon-dominated disk galaxies at the peak of galaxy formation ten billion years ago,''
\emph{Nature} \textbf{543}, 397 (2017).

\bibitem[Prince and Dunkley(2019)]{PrinceDunkley2019}
H. Prince and J. Dunkley,
``Data compression in cosmology: A compressed likelihood for Planck data,''
\emph{Physical Review D} \textbf{100}, 083502 (2019).

\bibitem[Ciocan et al.(2026)]{Ciocan2026}
B.~I. Ciocan, N.~F. Bouch\'e, J. Fensch, D. Krajnovi\'c, J. Freundlich, H. Desmond, B. Famaey, and R. Techi,
``MUSE-DARK III: The evolution of the radial acceleration relation at intermediate redshifts,''
\emph{Astronomy \& Astrophysics} \textbf{709}, L16 (2026).

\bibitem[Varasteanu et al.(2025)]{Varasteanu2025}
A.~A. V\u{a}r\u{a}\c{s}teanu, M.~J. Jarvis, A.~A. Ponomareva, H. Desmond, I. Heywood, T. Yasin, N. Maddox, M. Glowacki, M. Maksymowicz-Maciata, P.~E.~M. Mancera Pi\~na, and H. Pan,
``MIGHTEE-HI: The radial acceleration relation with resolved stellar mass measurements,''
\emph{Monthly Notices of the Royal Astronomical Society} \textbf{541}, 2366 (2025).

\bibitem[Lang et al.(2017)]{Lang2017}
P. Lang et al.,
``Falling rotation curves of star-forming galaxies at $0.6<z<2.6$,''
\emph{Astrophysical Journal} \textbf{840}, 92 (2017).

\bibitem[Lelli et al.(2016)]{Lelli2016}
F. Lelli, S.~S. McGaugh, and J.~M. Schombert,
``SPARC: mass models for 175 disk galaxies with Spitzer photometry and accurate rotation curves,''
\emph{Astronomical Journal} \textbf{152}, 157 (2016).

\bibitem[Brouwer et al.(2021)]{Brouwer2021}
M.~M. Brouwer et al.,
``The weak lensing radial acceleration relation: measuring the dark matter law with KiDS-1000,''
\emph{Astronomy \& Astrophysics} \textbf{650}, A113 (2021).

\bibitem[Unruh(1976)]{Unruh1976}
W.~G. Unruh,
``Notes on black-hole evaporation,''
\emph{Physical Review D} \textbf{14}, 870 (1976).

\bibitem[Tully and Fisher(1977)]{TullyFisher1977}
R.~B. Tully and J.~R. Fisher,
``A new method of determining distances to galaxies,''
\emph{Astronomy \& Astrophysics} \textbf{54}, 661 (1977).

\bibitem[McGaugh et al.(2000)]{McGaugh2000}
S.~S. McGaugh, J.~M. Schombert, G.~D. Bothun, and W.~J.~G. de Blok,
``The baryonic Tully--Fisher relation,''
\emph{Astrophysical Journal Letters} \textbf{533}, L99 (2000).

\bibitem[Planck Collaboration(2020)]{Planck2020}
Planck Collaboration,
``Planck 2018 results. VI. Cosmological parameters,''
\emph{Astronomy \& Astrophysics} \textbf{641}, A6 (2020).

\bibitem[Riess et al.(2022)]{Riess2022}
A.~G. Riess et al.,
``A comprehensive measurement of the local value of the Hubble constant with 1 km s$^{-1}$ Mpc$^{-1}$ uncertainty from the Hubble Space Telescope and the SH0ES team,''
\emph{Astrophysical Journal Letters} \textbf{934}, L7 (2022).

\bibitem[Freedman et al.(2025)]{Freedman2025}
W.~L. Freedman, B.~F. Madore, I.~S. Jang, et al.,
``Status report on the Chicago--Carnegie Hubble Program (CCHP): measurement of the Hubble constant using the Hubble and James Webb Space Telescopes,''
\emph{The Astrophysical Journal} (2025), doi:10.3847/1538-4357/adce78.

\bibitem[Di Valentino et al.(2021)]{DiValentino2021}
E. Di Valentino et al.,
``In the realm of the Hubble tension --- a review of solutions,''
\emph{Classical and Quantum Gravity} \textbf{38}, 153001 (2021).

\bibitem[Will(2014)]{Will2014}
C.~M. Will,
``The confrontation between general relativity and experiment,''
\emph{Living Reviews in Relativity} \textbf{17}, 4 (2014).

\bibitem[Brans and Dicke(1961)]{BransDicke1961}
C. Brans and R.~H. Dicke,
``Mach's principle and a relativistic theory of gravitation,''
\emph{Physical Review} \textbf{124}, 925 (1961).

\bibitem[Adelberger et al.(2003)]{Adelberger2003}
E.~G. Adelberger, B.~R. Heckel, and A.~E. Nelson,
``Tests of the gravitational inverse-square law,''
\emph{Annual Review of Nuclear and Particle Science} \textbf{53}, 77 (2003).

\bibitem[Padmanabhan(2010)]{Padmanabhan2010}
T. Padmanabhan,
``Thermodynamical aspects of gravity: new insights,''
\emph{Reports on Progress in Physics} \textbf{73}, 046901 (2010).

\bibitem[Sanders and McGaugh(2002)]{SandersMcGaugh2002}
R.~H. Sanders and S.~S. McGaugh,
``Modified Newtonian dynamics as an alternative to dark matter,''
\emph{Annual Review of Astronomy and Astrophysics} \textbf{40}, 263 (2002).

\bibitem[Oriti(2016)]{Oriti2016}
D. Oriti,
``Group field theory as the second quantization of loop quantum gravity,''
\emph{Classical and Quantum Gravity} \textbf{33}, 085005 (2016).

\bibitem[Gielen et al.(2013)]{Gielen2013}
S. Gielen, D. Oriti, and L. Sindoni,
``Cosmology from group field theory formalism for quantum gravity,''
\emph{Physical Review Letters} \textbf{111}, 031301 (2013).

\bibitem[Barbieri(1998)]{Barbieri1998}
A. Barbieri,
``Quantum tetrahedra and simplicial spin networks,''
\emph{Nuclear Physics B} \textbf{518}, 714 (1998).

\bibitem[Baez and Barrett(1999)]{BaezBarrett1999}
J.~C. Baez and J.~W. Barrett,
``The quantum tetrahedron in 3 and 4 dimensions,''
\emph{Advances in Theoretical and Mathematical Physics} \textbf{3}, 815 (1999).

\bibitem[Rovelli and Smolin(1995)]{RovelliSmolin1995}
C. Rovelli and L. Smolin,
``Discreteness of area and volume in quantum gravity,''
\emph{Nuclear Physics B} \textbf{442}, 593 (1995).

\bibitem[Gleason(1957)]{Gleason1957}
A.~M. Gleason,
``Measures on the closed subspaces of a Hilbert space,''
\emph{Journal of Mathematics and Mechanics} \textbf{6}, 885 (1957).

\bibitem[Deutsch(1999)]{Deutsch1999}
D. Deutsch,
``Quantum theory of probability and decisions,''
\emph{Proceedings of the Royal Society of London A} \textbf{455}, 3129 (1999).

\bibitem[Zurek(2005)]{Zurek2005}
W.~H. Zurek,
``Probabilities from entanglement, Born's rule $p_k=|\psi_k|^2$ from envariance,''
\emph{Physical Review A} \textbf{71}, 052105 (2005).

\bibitem[Popescu and Rohrlich(1994)]{PopescuRohrlich1994}
S. Popescu and D. Rohrlich,
``Quantum nonlocality as an axiom,''
\emph{Foundations of Physics} \textbf{24}, 379 (1994).

\bibitem[Eckert et al.(2019)]{Eckert2019}
D. Eckert et al.,
``Non-thermal pressure support in X-COP galaxy clusters,''
\emph{Astronomy \& Astrophysics} \textbf{621}, A40 (2019).

\bibitem[Dupourqu\'e et al.(2023)]{Dupourque2023}
S. Dupourqu\'e, N. Clerc, E. Pointecouteau, D. Eckert, S. Ettori, and F. Vazza,
``Investigating the turbulent hot gas in X-COP galaxy clusters,''
\emph{Astronomy \& Astrophysics} \textbf{673}, A91 (2023).

\bibitem[XRISM Collaboration(2025)]{XRISM2025}
XRISM Collaboration,
``XRISM reveals low nonthermal pressure in the core of the hot, relaxed galaxy cluster Abell 2029,''
\emph{Astrophysical Journal Letters} \textbf{982}, L5 (2025).

\bibitem[Chamseddine \& Mukhanov(2013)]{ChamseddineMukhanov2013}
A.~H. Chamseddine and V. Mukhanov,
``Mimetic dark matter,''
\emph{Journal of High Energy Physics} \textbf{1311}, 135 (2013).

\bibitem[Skordis \& Z\l o\'snik(2021)]{SkordisZlosnik2021}
C. Skordis and T. Z\l o\'snik,
``New relativistic theory for modified Newtonian dynamics,''
\emph{Physical Review Letters} \textbf{127}, 161302 (2021).

\bibitem[Lewis et al.(2000)]{Lewis2000}
A. Lewis, A. Challinor, and A. Lasenby,
``Efficient computation of cosmic microwave background anisotropies in closed Friedmann--Robertson--Walker models,''
\emph{Astrophysical Journal} \textbf{538}, 473 (2000).

\bibitem[More et al.(2015)]{More2015}
S. More, B. Diemer, and A. Kravtsov,
``The splashback radius as a physical halo boundary and the growth of halo mass,''
\emph{Astrophysical Journal} \textbf{810}, 36 (2015).

\bibitem[Regge(1961)]{Regge1961}
T. Regge,
``General relativity without coordinates,''
\emph{Il Nuovo Cimento} \textbf{19}, 558--571 (1961).

\bibitem[Ambj{\o}rn et al.(2004)]{AmbjornJurkiewiczLoll2004}
J. Ambj{\o}rn, J. Jurkiewicz, and R. Loll,
``Emergence of a 4D world from causal quantum gravity,''
\emph{Physical Review Letters} \textbf{93}, 131301 (2004).

\bibitem[Ambj{\o}rn et al.(2012)]{AmbjornGJL2012}
J. Ambj{\o}rn, A. G{\"o}rlich, J. Jurkiewicz, and R. Loll,
``Nonperturbative quantum gravity,''
\emph{Physics Reports} \textbf{519}, 127--210 (2012).

\bibitem[Engle et al.(2008)]{EPRL2008}
J. Engle, E. Livine, R. Pereira, and C. Rovelli,
``LQG vertex with finite Immirzi parameter,''
\emph{Nuclear Physics B} \textbf{799}, 136--149 (2008).

\bibitem[Gozzini(2021)]{Gozzini2021}
F. Gozzini,
``A high-performance code for EPRL spin foam amplitudes,''
\emph{Classical and Quantum Gravity} \textbf{38}, 225010 (2021).

\end{thebibliography}

\end{document}
